% !TEX program = latexmk
\documentclass[uplatex]{jsarticle}
\usepackage{amsthm,amssymb,amsmath,url,mathrsfs,bm}
\usepackage{xparse}
\usepackage{ascmac}
\usepackage[dvipdfmx]{graphicx}
\usepackage[dvipdfmx]{xcolor}
\usepackage{tikz}
\usetikzlibrary{arrows, arrows.meta, positioning, calc, fit, shapes, backgrounds}
\usepackage{float}
\usepackage{array}
\usepackage{tabularx}
\usepackage{longtable}
\usepackage{multirow}
\usepackage[dvipdfmx,unicode=true,pdfencoding=auto]{hyperref}
\usepackage{pxjahyper}

\pdfstringdefDisableCommands{%
  \def\ref#1{ref}%
  \def\label#1{}%
}

\newcounter{thm}[section]
\renewcommand{\thethm}{\thesection.\arabic{thm}}
\newcommand{\boxedthmtitle}[3]{%
  \IfNoValueTF{#3}{#1~#2}{#1~#2（#3）}%
}
\makeatletter
\newcommand{\beginboxedthm}[2]{%
  \refstepcounter{thm}%
  \IfNoValueTF{#2}{\def\@currentlabelname{}}{\def\@currentlabelname{#2}}%
  \par\medskip\noindent
  \begin{itembox}[l]{\boxedthmtitle{#1}{\thethm}{#2}}%
}
\makeatother
\newcommand{\finishboxedthm}{%
  \end{itembox}%
  \par\medskip
}
\NewDocumentEnvironment{thm}{o}{\beginboxedthm{定理}{#1}}{\finishboxedthm}
\NewDocumentEnvironment{prop}{o}{\beginboxedthm{命題}{#1}}{\finishboxedthm}
\NewDocumentEnvironment{dfn}{o}{\beginboxedthm{定義}{#1}}{\finishboxedthm}
\NewDocumentEnvironment{lem}{o}{\beginboxedthm{補題}{#1}}{\finishboxedthm}
\NewDocumentEnvironment{exa}{o}{\beginboxedthm{例}{#1}}{\finishboxedthm}
\NewDocumentEnvironment{quiz}{o}{\beginboxedthm{問}{#1}}{\finishboxedthm}
\NewDocumentEnvironment{nota}{o}{\beginboxedthm{記号}{#1}}{\finishboxedthm}
\NewDocumentEnvironment{cor}{o}{\beginboxedthm{系}{#1}}{\finishboxedthm}
\NewDocumentEnvironment{rem}{o}{\beginboxedthm{注意}{#1}}{\finishboxedthm}
\renewcommand{\labelenumi}{(\arabic{enumi})}
\setcounter{tocdepth}{3}

\title{Designing Secure Software ノート\\[0.5em]\large 開発者のためのセキュアなソフトウェア設計}
\author{\thanks{この資料は、Loren Kohnfelder, \textit{Designing Secure Software: A Guide for Developers}（No Starch Press, 2021）の学習目的のために作成されたものです。許可なく他の場所での再配布・転載などはしないでください。}}
\date{\today}
\begin{document}

%% =======================================================
%%  part: front
%% =======================================================
\maketitle
\tableofcontents
\newpage

%% =======================================================
%%  このノートについて（原著の Introduction のうち構成の部分）
%% =======================================================
\section*{このノートについて}
\addcontentsline{toc}{section}{このノートについて}

このノートは、Loren Kohnfelder, \textit{Designing Secure Software: A Guide for Developers} \cite{Kohnfelder2021} の内容を、「定義」「例」「注意」に整理し直したものである。
原著の第1章から第14章までの本文と付録A--Cを含め、章立て（第 $N$ 章 $\to$ \verb|\section|）も原著に合わせている。
原著の図は画像としては収録せず、本文から内容が決まる図に限り、地の文・表・図式（TikZ）に書き直した。
原著の献辞、序文（Foreword）、まえがき（Preface）、謝辞、想定読者の節は収録していない。

以下の「本書の構成と各章の概要」は原著の Introduction の内容であり、文中の「筆者」は原著者を指す。

\subsection*{本書の構成と各章の概要}
% src: What Topics Does the Book Cover?

本書は 14 の章からなり、概念・設計・実装の 3 部と結論で構成されている。

\subsubsection*{第1部：概念}
% src: Part 1: Concepts

第1章から第5章では、本書全体の土台となる概念を説明する。
\begin{description}
  \item[第1章 基礎] 情報セキュリティとプライバシーの基本を概観する。
  \item[第2章 脅威] 脅威モデリングを紹介し、資産をどう守るかという観点から、中心となる概念である攻撃対象領域と信頼境界を具体的に説明する。
\end{description}
続く 3 つの章では、セキュアなソフトウェアを作るのに役立つ道具を紹介する。
\begin{description}
  \item[第3章 緩和策] 洗い出した脅威を防御的に緩和するための、よく使われる戦略を論じる。
  \item[第4章 パターン] 効果的なセキュリティ設計パターンをいくつか示し、避けるべきアンチパターンも指摘する。
  \item[第5章 暗号] 暗号を道具箱として捉え、標準的な暗号ライブラリでよくあるリスクを緩和する方法を説明する。背後にある数学（実務ではめったに必要にならない）には立ち入らない。
\end{description}

\subsubsection*{第2部：設計}
% src: Part 2: Design

この部は、おそらく本書で最も独自性があり、読者にとって最も重要な部分である。
\begin{description}
  \item[第6章 セキュアな設計、第7章 セキュリティ設計レビュー] セキュアなソフトウェア設計の指針と、それを実現する実践的な技法を示す。前者は設計者の視点から、後者はレビュー担当者の視点から扱う。
\end{description}
あわせて、ソフトウェアの設計に最初からセキュリティを組み込む（bake in）ことがなぜ重要なのかも説明する。
両章では、第1部で紹介した考え方を土台に、それを取り入れてセキュアな設計を作るための具体的な方法論を示す。
レビューの方法論は筆者が業界で得た経験をそのまま下敷きにしたもので、自分の仕事の進め方に合わせて調整できる段階的なプロセスを含んでいる。
これらの章を読む際には、考え方を実践に移した例として、付録Aのサンプル設計書にも目を通すとよい。

\subsubsection*{第3部：実装}
% src: Part 3: Implementation

第8章から第13章では実装段階のセキュリティを扱い、デプロイ、運用、製品寿命の終了（end-of-life）にも触れる。
この部では、セキュアな設計ができた後、新たな脆弱性を持ち込まずにソフトウェアを開発する方法を説明する。
コードの断片を交えながら、脆弱性がどのようにコードに入り込むのか、それをどう防ぐのかを示す。
\begin{description}
  \item[第8章 セキュアプログラミング] プログラマーが直面するセキュリティ上の課題と、現実の脆弱性がコードの中でどのような姿をしているかを紹介する。
  \item[第9章 低レベルのコーディング上の欠陥] コンピューターの算術に潜む落とし穴と、C 言語流に動的メモリ割り当てを明示的に管理することがセキュリティをどう損ないうるかを扱う。
  \item[第10章 信頼できない入力、第11章 Web セキュリティ] 長年よく知られているのにいっこうになくならない、ありふれたバグの多く（インジェクション、パストラバーサル、XSS、CSRF の脆弱性など）を扱う。
  \item[第12章 セキュリティテスト] コードがセキュアであることをテストで確かめるという、あまりにも活用されていない実践を扱う。
  \item[第13章 セキュアな開発のベストプラクティス] セキュアな実装の指針の締めくくりとして、一般的なベストプラクティスをいくつか扱い、よくある落とし穴について警告する。
\end{description}

この部のコード例は、たいてい、まず避けるべき脆弱性を示し、続いてコードをセキュアにするパッチ済みの版を示す形になっている（それぞれ「脆弱なコード」（vulnerable code）、「修正済みのコード」（fixed code）と表示される）。
したがって、ここに載せたコードは、製品のソフトウェアにそのまま写して使うためのものではない。
修正済みのコードでさえ、ほかの問題が絡めば別の文脈では脆弱になりうる。本書のコードは、どれも、どんな用途でもセキュアだと保証されたものではないと考えてほしい。

\subsubsection*{結論}
% src: Conclusion

最終章の第14章「今後の展望」は短い。筆者の水晶玉は曇っているからである。
この章では本書の要点をまとめ、未来をのぞき見ようと試みる。そして、本書がよりセキュアなソフトウェアにどう貢献できるかという展望を手始めに、ソフトウェアセキュリティを一段引き上げる助けになりそうなアイデアを、思索的なものも含めて示す。

\subsubsection*{付録}
% src: Appendices

\begin{description}
  \item[付録A] セキュリティを意識した設計が実際にどのようなものかを示すサンプル設計書である。
  \item[付録B] 本書に登場するソフトウェアセキュリティの用語をまとめた用語集である。
  \item[付録C] 意欲的な読者が調べて楽しめるような、答えの決まっていない演習と問いを収めている。
\end{description}

\newpage

%% =======================================================
%%  part: ch1
%% =======================================================
%% =======================================================
%%  第1章 基礎
%% =======================================================
\section{基礎}
% src: 1: Foundations

この章では、\cite{Kohnfelder2021} の第1章に沿って、セキュリティ、信頼、情報セキュリティの古典的原則、およびプライバシーの基礎を扱う。

\begin{quote}
「誠実さは土台であり、たいていは揺るぎない土台である。たとえ自分の言ったことで面倒なことになっても、そこに立っていられるのだから。」
\hfill ---Charlamagne tha God
\end{quote}

ソフトウェアセキュリティは、論理的な営みであると同時に、直観的な意思決定に基づく技芸でもある。
現代のデジタルシステムを理解している必要があるうえに、そうしたシステムとやり取りし、その影響を受ける人間への感受性も欠かせない。
これを手強いと感じたなら、本書が説明しようとしている根本的な課題をすでによくつかんでいる。
この見方に立つと、ソフトウェアセキュリティの問題がなぜこれほど長くこの分野を悩ませてきたのかも見えてくる。これまでの着実な進歩が問題の一部を少しずつ削り取ったにすぎないとしても、なぜそこまでの労力が必要だったのかもわかる。
しかし、この状況には非常に良い知らせも含まれている。
プロセスのあらゆる段階でセキュリティへの意識を高め、より良いセキュリティに関わっていけば、私たち全員が実際に違いを生み出せるのである。

まず、セキュリティとは正確には何なのかを考えるところから始める。
セキュリティには主観的な面があるので、その基礎について明晰に考えることが決定的に重要である。
本書で示すのは、著者自身の経験に基づいて、現時点で最良と考える考え方を著者なりに理解したものである。
信頼はセキュリティ全体の土台である。
誰も真空の中で仕事をしているわけではないし、現代のデジタルシステムはあまりにも複雑で、シリコンから一人で作り上げることなどできない。そのため、自分で作らないもの（ハードウェア、ファームウェア、オペレーティングシステム、コンパイラに始まるすべて）は、提供してくれる他者を信頼するしかない。
この土台の上で、次にセキュリティの 6 つの古典的原則を示す。古典的な情報セキュリティの構成要素と、それを実施するための「ゴールドスタンダード」である。
最後に情報プライバシーの節で、デジタル製品やサービスが現代生活の最もセンシティブな領域にまで入り込んでいくなかで考慮すべき、人間的・社会的に重要な要素を補う。

セキュリティ、信頼、機密性といった言葉の意味については、読者も十分な直観を持っているはずである。しかし本書ではこれらの言葉が特定の技術的な意味を帯びるので、その意味を注意深く読み解いておく価値がある。
そのため、この章は丁寧に読むことが勧められている。
上級の読者には、原著者から挑戦として、自分でもっと良い説明を書いてみてほしいという提案がある。
誰にとっても学びの多い演習になるはずである。

%----------------------------------------------------------------------
\subsection{セキュリティを理解する}
% src: Understanding Security

すべての生物は、危険を避けて進み、攻撃から身を守り、見つけられる限りの安全な場所を目指す本能を備えている。

人間に生まれつき備わった物理的なセキュリティの感覚は、うまく働けば実に見事なものである。そのことをよく理解しておくことが大切である。
対照的に、仮想世界では頼りにできる本物の手がかりはほとんどなく、偽の手がかりはいくらでもでっち上げられる。
技術的な観点からセキュリティに取り組む前に、人間に何ができるのかを示す現実世界の話を一つ見ておこう（後で見るように、デジタルの領域ではまったく新しいスキルが必要になる）。

\begin{exa}[自動車販売員の実話]\label{exa:ch1-car-salesman}
以下は、ある自動車販売員の実話である。
顧客の試乗が終わり、販売員と顧客は販売店の敷地に戻ってきた。
販売員は車を降り、顧客と話を続けながら車の前方へ回り込んだ。
「彼の目を見たとき」と販売員は振り返る。「そのとき『しまった。こいつはこの車を盗むつもりだ』と思ったんです。」
事態は急展開した。
客から泥棒に早変わりした男は車のギアを入れて急発進し、販売員は命がけで車のボンネットにしがみついた。
犯人は販売員を振り落とそうと車を乱暴に走らせたが、振り落とせなかった
（幸い、販売員に大きなけがはなく、犯人はまもなく逮捕され、有罪となり、損害賠償を命じられた）。
\end{exa}

二人の男の目が合ったとき、微妙なリスク計算が行われていた。
ほんの一瞬のうちに、販売員は顧客の表情や身ぶりから得られる複雑な視覚的な手がかりを処理し、敵対的な行動をとろうとする明確な意図を読み取ったのである。

\begin{dfn}[スピアフィッシング]\label{dfn:ch1-spear-phishing}
不特定多数ではなく、特定の標的をだますために作られた詐欺メールによる攻撃を、\underline{スピアフィッシング}（spear phishing）攻撃という。
\end{dfn}

では、同じ販売員がスピアフィッシング攻撃の標的になったと想像してみよう。
デジタルの領域には、攻撃者と対面したときに感じ取ったような手がかりがないので、彼ははるかに簡単にだまされてしまうはずである。

情報セキュリティ、コンピュータ、ネットワーク、ソフトウェアに関しては、デジタルシステムを守れる見込みを少しでも持ちたければ、直面するリスクを分析的に考えて評価しなければならない。
しかも、ビットやコードを直接見たり、嗅いだり、聞いたりできないなかで、それをやらなければならない。
オンラインでデータを調べるときは、いつもソフトウェアを使って情報を人間が読めるフォントで表示しているのであり、たいていは自分と実際のビットとの間に大量のコードが挟まっている。
それどころか、合わせ鏡の迷宮（hall of mirrors）になっているおそれさえある。
だから、自分のツールを信頼し、自分が調べているつもりのデータを本当に調べているのだと信頼するしかない。

ソフトウェアセキュリティの中心は、さまざまな脅威からデジタル資産を守ることである。その取り組みを主に方向づけるのが一連の基本的なセキュリティ原則であり、これがこの章の残りの話題となる。
こうした第一原理からシステムを分析すれば、脆弱性がどのようにソフトウェアに入り込むのか、問題をどう事前に回避し緩和すればよいのかが見えてくる。
これらの基礎的な原則は、後の章で扱うほかの設計手法とともに、ソフトウェアだけでなく、自転車の錠前、銀行の金庫室、刑務所の設計と運用にも当てはまる。

\begin{dfn}[情報セキュリティとソフトウェアセキュリティ]\label{dfn:ch1-information-software-security}
\underline{情報セキュリティ}（information security）という用語は、特にデータの保護と、アクセスをどのように許可するかを指す。
\underline{ソフトウェアセキュリティ}（software security）はより広い用語であり、情報セキュリティを確実に実施することも含めて、信頼に値するソフトウェアシステムの設計・実装・運用に焦点を当てる。
\end{dfn}

%----------------------------------------------------------------------
\subsection{信頼}
% src: Trust

デジタルの領域でも信頼は同じくらい決定的に重要であるが、あまりにも当然のものとみなされがちである。
ソフトウェアセキュリティは、突き詰めれば信頼に依存している。
システムのすべての部分を制御することも、ソフトウェアをすべて自分で書くことも、依存するコンポーネントの供給元をすべて審査することもできないからである。
現代のデジタルシステムは非常に複雑で、巨大テクノロジー企業でさえ技術スタック全体をゼロから構築することはできない。
シリコンからオペレーティングシステム、ネットワーク、周辺機器、そしてそれらすべてを動かす多数のソフトウェア層に至るまで、私たちが日常的に頼っているシステムは、途方もない規模と複雑さを持つ驚くべき技術的成果である。
すべてを自分で作れる者はいないので、組織は、機能や価格で選んだハードウェアやソフトウェア製品に頼ることが多い。
しかし、依存関係の一つひとつに信頼についての判断も伴っていることを忘れてはならない。

すべてを調査し検証する時間や資源は誰にもないとしても、セキュリティを確保するには、こうした信頼関係を詳しく吟味しなければならない。
信頼が足りなさすぎると、現実の脅威がまず起こりそうにないところまで守ろうとして、余計な作業を大量に抱え込むことになる。
逆に、あまりに無造作に信頼すると、後で不意打ちを食らいかねない。
率直に言えば、ある主体を完全に信頼するということは、その主体が何の報いも受けずにあなたを裏切れるということである。
信頼される側の動機や意図しだいでは、ごまかし、嘘、不公正、怠慢、無能、ミス、そのほかあらゆる形で信頼を裏切られるかもしれない。

不完全な情報しかないなかで、重大な判断を素早く下さなければならない。信頼が最も力を発揮するのは、まさにこうした場面である。
しかし、人間に生まれつき備わった信頼の感覚は微妙な感覚入力に頼っており、それはデジタルの領域にはまったく当てはまらない。
以下ではまず信頼という概念そのものを取り上げ、私たちが経験する信頼とは何かを分析してから、ソフトウェアに関わる信頼へと話を進める。
読みながら共通する筋道を探し、ソフトウェアについての考え方を自分の信頼の直観と結びつけてみるとよい。
すでに持っている信頼のスキルを生かすのは強力な手法である。時間をかけるうちに、どれだけ技術的な分析を重ねるよりも効果的な、ソフトウェアセキュリティについての勘が身につく。

\subsubsection{信頼を感じる}
% src: Feeling Trust

信頼を理解する最善の方法は、信頼に身を委ねるとどんな感じがするのかを、実際に経験しながら注意深く観察することである。

\begin{exa}[目を閉じて道路を渡る思考実験]\label{exa:ch1-crossing-road}
信頼とは何かを実感できる思考実験（あるいは、本当に信頼している人と実際に試してみる演習）を一つ挙げる。
友人と一緒に交通量の多い大通りを歩いていて、ほんの数フィート先を車が次々と通り過ぎていくところを想像してほしい。
前方に横断歩道が見えてきたところで、あなたは友人にこう説明する。安全に道路を渡れるよう誘導してほしい、安全に渡れるかどうかは君に任せる、自分は目を閉じて素直についていく、と。
手をつないで横断歩道まで進むと、友人はあなたをそっと道路のほうへ向かせ、体に触れて待つよう合図する。
猛スピードで走り過ぎる車の音を聞きながらも、友人（今やあなたの保護者である）が安全に渡れるようになるまで待ってくれていることは、よくわかっている。
それでも、おそらく心臓の鼓動がはっきりと速まり、迫りくる危険の音がしないかと耳をそばだてている自分に気付くはずである。

やがて友人は、はっきりとあなたを前へ導き、縁石から一歩下りるよう促す。
目を閉じたまま道路に踏み出すと決めたなら、そのとき感じているのは純粋な信頼である（あるいは、いくらかの不信かもしれない）。
心ははっきりとしたリスクを鋭く感じ取り、感覚は安全を自分で確かめようと張り詰め、体の奥深くの何かが「やめておけ」と警告している。
あなた自身の内なるセキュリティ監視システムは証拠が足りないと判断し、動く前に目を開けてほしいと思っている。
もし友人が何かの理由で状況を見誤ったら、あるいはもっと悪いことに、命に関わる悪質ないたずらを仕掛けていたらどうなるのか
（これが先に述べた、信頼に対する二つの脅威、すなわち無能と悪意である）。
こうした本能を押しとどめて道路を渡らせているのは、あなたが友人に寄せた信頼なのである。
\end{exa}

自分がデジタルな信頼の判断を下していることをもっと意識しよう。そして、ほかの人にもそれが見えるように、それがセキュリティにどれほど重要な影響を与えるかがわかるように手助けしよう。
理想的には、重要なサービスのためにコンポーネントやベンダーを選ぶときに、今の演習のような信頼の判断を導くのとまったく同じ直観を生かせるようになるとよい。

\subsubsection{ビットは見えない}
% src: You Cannot See Bits

ここまでの話はすべて、「データを直接見ている」と思っているときでも、実際にははるか遠くにある表現を見ているにすぎない、という事実を強調するためのものである。
実際に見ているのは画面上のピクセルであり、それが物理的にどこにあるのか正確には知らない特定のバイトの内容を表していると信じているだけである。しかも、データを画面上で人間が読める形に変換するために、おそらく何百万もの命令が実行されている。
デジタル技術は非常に抽象的で、恐ろしく高速で、直接見えないところに隠れているので、信頼の問題がとりわけ厄介になる。
また、現代のネットワークでは、遠く離れた場所から匿名で接続できる。
データを調べるときには、メモリ上の実際のデータと、私たちがデータの値として読み取る文字を形作るピクセルとの間に、いつも大量のソフトウェアとハードウェアがあることを思い出してほしい。
その中のどれかが悪意をもって実際のデータを偽って表示していたら、どうやってそれに気付けるだろうか。
デジタル情報について実際のありのままの姿（ground truth）を直接観察することは、どんな意味でも極めて難しい。

\begin{exa}[ブラウザの鍵アイコン]\label{exa:ch1-lock-icon}
Web ブラウザのアドレスバーにある、Web サイトへの安全な接続を示す鍵アイコンを考えてみる。
この特徴的なピクセルが表示されるかどうかによって、ユーザーにはただ 1 ビット、つまり「ここは安全か、安全でないか」だけが伝わる。
その裏側には、第11章で詳しく述べるように大量のデータとかなりの量の計算があり、そのすべてが「はい／いいえ」の二値のセキュリティ表示に集約されている。
熟練した開発者であっても、たった 1 つの事例が妥当かどうかを自分で確かめようとすれば、途方もない作業になるはずである。
だから私たちはソフトウェアを信頼するしかないし、信頼するだけの十分な理由もある。
ここで言いたいのは、その信頼がどれほど深く広く行き渡っているかを認識し、当たり前のものと思わないようにしよう、ということである。
\end{exa}

\subsubsection{能力と不完全さ}
% src: Competence and Imperfection

ほとんどの攻撃は、ソフトウェアの欠陥や設定ミスを突くところから始まる。そうした欠陥やミスは、プログラマーや IT スタッフが正直に誠実に努力した結果として生じたものである。
彼らもたまたま人間であり、人間である以上、不完全なのである。
ライセンスでは事実上あらゆる責任が日常的に否認されているので、どのソフトウェアも「買い手注意」（caveat emptor）の原則のもとで使われている。
よく言われるように「すべてのソフトウェアにはバグがある」のだとすれば、そのバグの一部は悪用可能であり、いずれ悪者がそのいくつかを見つけ、悪用する機会を得ることになる。
悪意のあるソフトウェアを誤って信頼したせいで直接攻撃を許し、被害者になるケースは比較的まれである。

幸い、オペレーティングシステムやプログラミング言語についての大きな信頼の判断は、たいていは容易である。
多くの大企業には、高品質なハードウェアやソフトウェア製品を提供しサポートしてきた豊富な実績があり、そうした企業を信頼するのは十分に合理的である。
実績の少ない相手を信頼するのは、もっとリスクが高いかもしれない。
そうした組織にも熟練した意欲ある人々が大勢いて熱心に働いている可能性は高いが、業界の透明性が低いため、その製品のセキュリティは判断しにくい。
オープンソースには透明性があるが、それが役に立つかどうかは、プロジェクトのオーナーがどこまで監督しているかにかかっている。貢献者がバグのあるコードや明らかに悪意のあるコードを紛れ込ませるのを防げるかどうかが、そこで決まるからである。
驚くべきことに、より高いレベルのセキュリティや、攻撃を受けた場合の補償を約束して差別化を図ろうとするソフトウェア企業すら存在しない。そのため、消費者には選ぶ余地さえない。
法律、規制、ビジネス上の契約は、いずれも信頼の判断にまつわる不確実性を緩和する追加の手段となる。

信頼の判断を真剣に受け止めつつ、誰も 100 パーセント正しく判断することはできないと認めよう。
悪い知らせは、信頼の判断は未来の予測なので、常に不完全だということである。
米国証券取引委員会（SEC）が警告するように、「過去の実績は将来の結果を保証しない」。
良い知らせは、人間は信頼を見極める能力を高度に進化させてきたということである（ただし、その能力が最もよく働くのは対面のときであり、デジタルメディアを介するとまったくうまく働かない）。正確な情報を持ち、意識して行動すれば、大多数の場合に信頼の判断を正しく下せる。

\subsubsection{信頼はスペクトルである}
% src: Trust Is a Spectrum

信頼には常に程度があり、信頼の評価には常にいくらかの不確実性が伴う。
スペクトルの極端な一端は、たとえば大手術を受ける場合である。このとき私たちは文字どおり自分の命を医療専門家に預け、自分の身体の制御だけでなく、意識そのものや手術を見届ける能力まで進んで手放すこともある。
最悪の場合、彼らが失敗して私たちが命を落とせば、（遺産に関する法的権利は別として）文字どおり何の救済手段もない。
日常的な信頼はもっと限定的である。
クレジットカードには限度額があり、支払われない場合に銀行が被りうる損失に上限を設けている。
車にはバレーキー（valet key）があり、トランクを開けられないようにできる。

信頼はスペクトルなので、「信頼せよ、されど検証せよ」（trust but verify）という方針は、完全な信頼と完全な不信の間を埋める便利な道具になる。
ソフトウェアでは、認可と入念な監査を組み合わせることでこれを実現できる。
たいていは、自動監査（大量の、ほとんど同じことの繰り返しのアクティビティログを正確にチェックする）と手動監査（抜き取り検査を行い、例外的なケースに対処し、最終判断を下す人間をループの中に置く）を組み合わせることになる。
監査については、この章の後半で詳しく扱う。

\subsubsection{信頼の判断}
% src: Trust Decisions

ソフトウェアでは、信頼するかしないかの二択になる。
アプリケーションにさまざまな権限を課すシステムもあるが、その場合でも、個々の権限は許可するか許可しないかのどちらかである。
迷ったときは、候補となる解決策のうち少なくとも 1 つが相応の信頼を得られる限り、信頼しない側に倒しておけば安全である。
評価が厳しすぎてどの製品も信頼を得られないなら、そのコンポーネントを自分で作ることを検討するしかなくなる。

信頼の判断を下すことは、放っておけば事実上無限に広がってしまう決定木の枝を切り落とすことだと考えるとよい。
あるサービスやコンピュータが安全だと信頼できれば、それ以上深く分析する手間が省ける。
逆に、信頼することに慎重であれば、すべてのサブコンポーネントを含め、システムのより多くの部分を自分で構築して保護しなければならない。

\begin{exa}[クラウドストレージについての信頼の判断]\label{exa:ch1-cloud-storage-trust}
原著の図1-1は、信頼の判断を下す例を示している。
データの保存先として完全に信頼できるクラウドストレージサービスが見当たらない場合の代替案の一つは、データを保存する前にローカルで暗号化しておき（そうすればベンダーから漏れても害はない）、2 つ以上のサービスを互いに独立に冗長に使う（そうすればすべてのサービスで同時にデータが失われる確率は極めて小さくなる）ことである。
\end{exa}
% 図 1-1: 信頼の判断の例（原文では "graphic" のみで、キャプション・内容は欠落している）。

明示的に信頼していない入力（事実上すべての入力、とくに公衆インターネットや任意のクライアントから来るものはすべてここに含めるべきである）は、疑いの目で扱い、最大限の注意を払う（詳しくは第4章の「信頼への慎重さ」を参照）。
信頼している入力であっても、完全に信頼できると仮定するのは危険なことがある。
システム全体を壊れにくくし、無害なバグが起きたときにエラーが広がるのを防ぐためだけでもよいので、手軽にできるところでは機会をとらえて安全性のチェックを追加することを検討しよう。

\subsubsection{暗黙のうちに信頼されるコンポーネント}
% src: Implicitly Trusted Components

どのソフトウェアプロジェクトも、暗黙のうちに信頼している途方もない技術スタックに依存している。
ハードウェア、オペレーティングシステム、開発ツール、ライブラリ、そのほかの依存関係がそうであり、これらを審査するのは現実的でないので、ベンダーの評判を頼りに信頼している。
それでもなお、自分が何を暗黙のうちに信頼しているのかをある程度把握しておき、とくに暗黙の信頼の範囲を大きく広げる前には、その判断をよく検討すべきである。

暗黙の信頼を管理する簡単な手法はないが、役に立つ考え方が一つある。
信頼する相手の数をできるだけ少なくすることである。
たとえば、Microsoft（あるいは Apple など）のオペレーティングシステムを使うとすでに決めているなら、露出を最小限にするために、コンパイラ、ライブラリ、アプリケーション、そのほかの製品やサービスも同じ会社のものを使う方向に寄せるとよい。
大まかに言えば、信頼する企業が増えるほど、そのどれかに裏切られる機会も増える、という理屈である。
加えて、同じ会社の製品群は一緒に使ったときの互換性が高く、よくテストされていることが多い、という実際的な利点もある。

\subsubsection{信頼に値する存在であること}
% src: Being Trustworthy

最後に、信頼の判断の裏返し、すなわち自分が製品やサービスを提供する側として信頼を得ることも忘れてはならない。
どんなソフトウェア製品も、自分が信頼に値することをエンドユーザーに納得してもらわなければならない。
たいていは、しっかりとしたプロフェッショナルなイメージを見せるだけで十分であるが、製品が重要な機能を担っているなら、顧客に信頼の確かな根拠を示すことが決定的に重要である。

自分の仕事への信頼を生み出す基本的な方法として、次のようなものが挙げられている。
\begin{itemize}
\item 透明性は信頼を生む。オープンに作業すれば、顧客は製品を評価できる。
\item 第三者に関わってもらうと、その独立性によって信頼が築かれる（たとえば、監査人を雇う）。
\item 自分の製品が、ほかの製品と統合される第三者の立場になることもある。互いに距離を置いた関係にある二者が結託するのは難しいので、信頼が増す。
\item 問題が実際に起きたときは、フィードバックに耳を傾け、断固として行動し、調査の結果と再発防止のためにとった措置をすべて公表する。
\item 特定の機能や設計要素によって、信頼を目に見える形にできる。たとえば、分散した拠点で保存・検証されたバックアップの数をリアルタイムで示すアーカイブソリューションがそうである。
\item 行動は信頼を生むが、中身のない主張は、目の肥えた顧客の信頼をかえって損なう。信頼に値することを示す具体的な証拠を、できれば顧客が自分で検証できる形で示す。オープンソースのコードの品質を実際に審査する人はほとんどいないとしても、審査\emph{できる}と知っていること（そして、おそらく誰かほかの人が審査していると想定できること）には、ほぼ同じくらいの説得力がある。
\end{itemize}

%----------------------------------------------------------------------
\subsection{古典的原則}
% src: Classic Principles

情報セキュリティの指導原則が生まれたのは、コンピュータが施錠され空調の効いた二重床の特別な部屋から出て、ネットワークにつながり始めた 1970 年代である。
こうした伝統的なモデルは、現代の情報セキュリティにとっての「ニュートン力学」にあたる。
つまり、多くの用途では簡潔で優れた指針になるが、それですべてが片付くわけではない。
情報プライバシーは、伝統的な情報セキュリティの原則ではカバーされない、現代のデータ保護と管理責任（stewardship）にまつわるより微妙な考慮事項の一例である。

基礎的な原則は、3 つずつの 2 組にうまくまとまる。
最初の 3 つの原則は C-I-A と呼ばれ、データアクセスの要件を定める。
残りの 3 つは、アクセスをどう制御し監視するかに関わる。
こちらをゴールドスタンダードと呼ぶ。
2 組の原則は互いに依存しており、全体がそろって初めてデータ資産を保護できる。

許可されていないデータアクセスを防いだとしても、その先には、どの人、どのコンポーネントやシステムにアクセスを委ねるべきかという問題が残る。
これは信頼に関するより難しい問題である。どんなデジタルシステムを守るにも避けては通れないが、突き詰めれば情報セキュリティの範囲を超えている。

\subsubsection{情報セキュリティの C-I-A}
% src: Information Security’s C-I-A

ソフトウェアセキュリティは伝統的に、情報セキュリティの 3 つの基本原則、すなわち機密性・完全性・可用性の上に築かれる。
本書ではこれらをまとめて C-I-A と呼ぶ。
いずれもデータ保護の基本を軸に定式化されており、3 本の柱それぞれの意味は直観的にわかる。

\begin{dfn}[C-I-A]\label{dfn:ch1-cia}
情報セキュリティの次の 3 つの基本原則をまとめて \underline{C-I-A} という。
\begin{itemize}
\item \underline{機密性}（Confidentiality）：許可されたデータアクセスだけを許す。情報を漏らさない。
\item \underline{完全性}（Integrity）：データを正確に保つ。許可されていない変更や削除を許さない。
\item \underline{可用性}（Availability）：データを利用できる状態に保つ。著しい遅延や許可されていない停止を許さない。
\end{itemize}
\end{dfn}

この短い定義はそれぞれ、目標と、その目標を損なわせないための防御とを述べている。
設計をレビューするときは、セキュリティを損なう方法を考え、そこから防御策へとさかのぼって考えると役に立つことが多い。

C-I-A の 3 つの構成要素はいずれも理想を表しており、完璧を求めないことが決定的に重要である。

% 注: 付録Bの用語集の語義を混ぜた定義を置いていたが、本文にない記述なので削除した。
\begin{exa}[暗号化された通信からの推測]\label{exa:ch1-encrypted-traffic}
しっかりと暗号化されたネットワークトラフィックであっても、分析すれば、執念深い盗聴者は 2 つのエンドポイント間の通信について、交換されたデータ量のような何らかの情報を推測できるかもしれない。
厳密に言えば、こうしたデータのやり取りによって、エンドポイント間のやり取りの機密性は弱まっている。
しかし実用上は、極端な手段をとらない限り防ぎようがなく、たいていはリスクが十分に小さいので無視しても差し支えない。
ネットワークトラフィックから情報を推測することは\underline{サイドチャネル攻撃}（side channel attack）の一例であり、それが問題になるかどうかは、それがもたらす脅威を評価して判断する。
そのトラフィックはどんな活動に対応しているのか、敵対者はその知識をどう利用しうるのか、といった点である。
次の章では、こうした脅威の評価を詳しく説明する。
\end{exa}

C-I-A のどの構成要素にも、認可という考え方が組み込まれていることに注意してほしい。
C-I-A が求めているのは、正しい開示、正しいデータの変更、正しい可用性の制御だけである。
何が「正しい」のかは重要な詳細であり、別途定めるべき認可ポリシーであるが、これらの基本的なデータ保護の基礎概念そのものには含まれない。
その部分は「ゴールドスタンダード」で論じる。

% src: Confidentiality
\underline{機密性：}
機密性を保つとは、私的な情報を許可された方法でのみ開示することである。
単純に聞こえるが、実際にはいくつもの複雑な問題がある。

まず、どの情報を私的なものとみなすかを注意深く特定することが大切である。
設計文書でこの区別を明確にしておくべきである。
何がセンシティブかは自明に思えることもあるが、実際には人によって驚くほど意見が分かれ、明示的な仕様がなければ誤解が生じかねない。
最も安全な前提は、外部から集めた情報をすべて、デフォルトで私的なものとして扱うことである。その扱いは、どのように、またなぜ緩めてよいのかを説明した明示的なポリシーで別途宣言されるまで続ける。

データを私的なものとして扱うべき理由のうち、見落とされがちなものとして次のものがある。
\begin{itemize}
\item エンドユーザーは、たとえ明かされても害がないデータであっても、特に説明がない限り、自分のデータは私的に扱われると当然のように期待しているかもしれない。
\item 人々は、別の用途のためのテキストフィールドにセンシティブな情報を入力するかもしれない。
\item 情報の収集、取り扱い、保存には、多くの人が知らない法律や規制が適用される場合がある（たとえば、ヨーロッパの人々があなたの Web サイトを閲覧するなら、EU の GDPR 規制の対象となりうる）。
\end{itemize}

私的な情報を扱うときは、何が適切なアクセスなのかを決めておく。
情報をいつ、どのように開示するかの設計は突き詰めれば信頼の判断であり、規則を明記するだけでなく、その規則の背後にある主観的な選択を説明しておくことにも価値がある。
これについては、第4章でパターンを論じる際にさらに述べる。

機密性の侵害には程度の幅がある。
完全な開示では、攻撃者はメタデータを含むデータセット全体を手に入れる。
幅の下端には、内部エラーメッセージのように、実質的に何の影響もない軽微な情報の開示があるかもしれない。

\begin{exa}[連番の顧客番号による部分的な開示]\label{exa:ch1-sequential-customer-numbers}
部分的な開示の例として、新規顧客に連番を割り当てる慣行を考えてみる。
抜け目のない競合他社は、ときどき新規顧客として登録して新しい顧客番号を手に入れ、番号の差を順に計算すれば、各期間に獲得された顧客数を知ることができる。
\end{exa}

保護されたデータの詳細が少しでも漏れれば、程度の差はあれ機密性の侵害である。

軽微な開示が持ちうる価値は、とかく過小評価されやすい。
攻撃者は、開発者が当初意図したのとはまったく違う方法でデータを利用するかもしれない。また、小さな情報の断片を組み合わせると、個々の断片からは得られない強力な知見が得られることがある。

\begin{dfn}[匿名化の解除（再識別）]\label{dfn:ch1-deanonymization}
いくつもの小さな情報を組み合わせて個人を特定する過程を、\underline{匿名化の解除}（deanonymization）または\underline{再識別}（reidentification）という。
\end{dfn}

\begin{exa}[郵便番号と Netflix のデータセット]\label{exa:ch1-netflix-reidentification}
ある人の郵便番号（ZIP コード）を知っても大したことはわからないかもしれない。しかし、その人のおおよその年齢と、医師（MD）であることもわかっていれば、人口のまばらな地域ではこれらの情報を組み合わせて本人を特定できるかもしれない。
また、研究者は、Netflix が公開した、匿名化されているはずのデータセットを分析して、多数のユーザーアカウントを IMDb のアカウントと照合することに成功した。
お気に入りの映画が、個人を一意に識別する効果的な手がかりになったのである。
\end{exa}

% src: Integrity
\underline{完全性：}
情報セキュリティの文脈でいう完全性とは、単にデータが真正かつ正確であり、許可されていない改ざんや削除から守られていることである。

\begin{dfn}[来歴]\label{dfn:ch1-provenance}
データの元の出所と、データに加えられた許可済みの変更の正確な記録を、データの\underline{来歴}（provenance）という。
\end{dfn}

許可されていない変更から守ることに加えて、データの来歴を正確に記録しておけば、完全性をより強力に保証する重要な手段になりうる。

多くの改ざん攻撃に対する古典的な防御の一つは、重要なデータのバージョンを保存し、その来歴を記録しておくことである。
簡単に言えば、良いバックアップをとっておくことである。
増分バックアップは、簡単かつ効率的に導入でき、どのデータがいつ変わったかを正確に示す一連のスナップショットが得られるので、優れた緩和策になりうる。
しかし、完全性が必要なのはデータの保護に限らない。コンポーネント、サーバーログ、ソフトウェアのソースコードとバージョン、さらに問題が起きたときに改ざんの元の出所を突き止めるのに必要なフォレンジック情報についても、完全性を確保しなければならないことが多い。
管理者アクセスを限定するアクセス制御に加えて、セキュアなダイジェスト（チェックサムに似たもの）とデジタル署名は強力な完全性チェックの手段である。
これらについては第5章で説明する。

改ざんにはさまざまな手口があり、保存されたデータを書き換えるとは限らないことを覚えておいてほしい。
たとえば Web アプリケーションでは、クライアント側で、クライアントとサーバー間の通信経路上で、許可された当事者をだまして変更させることで、ページ上のスクリプトを書き換えることで、あるいはそのほか多くの方法で改ざんが起こりうる。

% src: Availability
\underline{可用性：}
可用性に対する攻撃は、インターネットにつながった世界の悲しい現実であり、最も防御しにくい攻撃の一つになりうる。
最も単純な場合、攻撃者は桁外れに大量のトラフィックをサーバーに送りつけ、正当な利用に見えるリクエストでサーバーを圧倒するだけかもしれない。
可用性という用語は、情報が一時的に利用できない状態を指す。
恒久的に失われたデータも利用できないことに変わりはないが、それは一般に、本質的には完全性の侵害とみなされる。

% 注: 付録Bの用語集の語義を混ぜた定義を置いていたが、本文にない記述なので削除した。
匿名の\underline{サービス拒否}（denial-of-service, DoS）攻撃は、身代金目的で行われることが多く、あらゆるインターネットサービスを脅かす難しい課題である。
これに対する最善の防御は、重い負荷に耐えるインフラを持つ大規模なサービス上でホスティングし、問題が起きたらインフラを素早く移せる柔軟性を保っておくことである。
被害者の多くがこうしたインシデントを内々に解決するので、DoS 攻撃が実際にどれほど頻繁に起き、どれほどのコストを生んでいるのかは誰にもわからない。
しかし、そうしたインシデントに備えて、事前に詳細な計画を立てておくべきであることは疑いない。

可用性に対する脅威はほかにも多くの種類がある。
Web サーバーにとっては、バグを引き起こしてクラッシュや無限ループを招く不正な形式のリクエストが、サービスに壊滅的な打撃を与えうる。
さらに、アプリケーションのストレージ、計算、通信の容量を過負荷にしたり、キャッシュの効果をなくすようなパターンを使ったりする攻撃もあり、いずれも深刻な問題を引き起こす。
ソフトウェア、設定、データを許可なく破壊されることも、たとえバックアップがあっても復旧に時間がかかりうるので、可用性に悪影響を与えうる。

\subsubsection{ゴールドスタンダード}
% src: The Gold Standard

C-I-A がセキュアなシステムの目標だとすれば、ゴールドスタンダードはその目標を達成する手段を表している。
Aurum は金を表すラテン語で、それゆえ金の元素記号は「Au」であるが、セキュリティを実施するための 3 つの重要な原則は、たまたまいずれも英語でこの同じ 2 文字から始まる。

\begin{dfn}[ゴールドスタンダード]\label{dfn:ch1-gold-standard}
セキュリティを実施する次の 3 つの基本的な仕組みをまとめて、\underline{ゴールドスタンダード}（Gold Standard）という。
\begin{itemize}
\item \underline{認証}（Authentication）：プリンシパルの身元を高い確度で判定すること。
\item \underline{認可}（Authorization）：認証されたプリンシパルによる行為だけを確実に許すこと。
\item \underline{監査}（Auditing）：プリンシパルの行為についての信頼できる記録を、検査のために残しておくこと。
\end{itemize}
\end{dfn}

\begin{rem}[authN と authZ]\label{rem:ch1-authn-authz}
専門用語への注意：authentication と authorization は長くてよく似た単語なので、両者をはっきり区別できる便利な略語として、authentication を \underline{authN}、authorization を \underline{authZ} と書くことがある。
\end{rem}

\begin{dfn}[プリンシパル]\label{dfn:ch1-principal}
確実に認証された任意の主体を\underline{プリンシパル}（principal）という。
人、企業や組織、政府機関、アプリケーション、サービス、デバイス、そのほか行為する力を持つあらゆるエージェントがこれにあたる。
\end{dfn}

認証とは、プリンシパルのクレデンシャルが妥当であることを確実に確かめる過程である。
システムでは一般に、登録ユーザーが自分のユーザーアカウントに対応するパスワードを知っていることを証明して認証するが、認証はもっと広い概念である。
クレデンシャルには、プリンシパルが知っているもの（パスワード）、持っているもの（スマートカード）、あるいはプリンシパル自身の特徴（生体データ）がありうる。
これについては次の項で詳しく述べる。

認証されたプリンシパルによるデータアクセスには認可の判断が適用され、定められた規則に従って行為が許可または拒否される。
たとえば、アクセス制御を設定できるファイルシステムでは、特定のユーザーに対して特定のファイルを読み取り専用にできる。
銀行のシステムでは、窓口係は一定額までの取引は記録できるが、それより大きな取引にはマネージャーの承認が必要になることがある。

プリンシパルが何をしたかを、行為の失敗した試みも含めて正確に記録するセキュアなログをサービスが残していれば、管理者は事後に監査を行い、システムがどう動作したかを検査して、すべての行為が適切だったことを確かめられる。
正確な監査ログは、実際に起きた出来事を信頼できる形で伝えてくれるので、強固なセキュリティの重要な構成要素である。
詳細なログは何が起きたかの記録となり、異常な出来事や疑わしい出来事が起きたときに、正確に何があったのかを明らかにしてくれる。
たとえば、重要なファイルが消えていることに気付いたら、理想的には、誰がいつそのファイルを削除したかの詳細がログからわかり、さらなる調査の出発点が得られるはずである。

ゴールドスタンダードは、C-I-A を保護するための実施の仕組みとして働く。
機密性と完全性は、許可されていない開示や改ざんからの保護として定義した。
可用性もまた、許可された管理者による制御の対象である。
認可の判断を本当の意味で実施するには、システムを使うプリンシパルが適切に認証されていることが欠かせない。
最後に監査が全体を締めくくる。誰がいつ何をしたかの信頼できるログを残し、不正がないかを定期的にレビューし、行為した当事者に責任を負わせるのである。

セキュアな設計では、常に認証と認可をきれいに分けるべきである。
両者を混ぜると混乱を招くし、これらの段階がきれいに分かれているほうが監査証跡もわかりやすくなるからである。

\begin{exa}[金庫室に通した理由]\label{exa:ch1-vault-guard}
分離が重要である理由を、次の 2 つの現実世界の例が示している。
\begin{itemize}
\item 「なぜあの男を金庫室に入れたんだ？」「さっぱりわかりませんが、ちゃんとした人に見えたんです！」
\item 「なぜあの男を金庫室に入れたんだ？」「彼の身分証は『Sam Smith』のもので有効でしたし、Joan からの書面のメモを持っていました。」
\end{itemize}
2 番目の答えは 1 番目よりはるかに中身がある。
1 番目の答えは、警備員が間抜けだという証明になるだけで、ほかには何の役にも立たない。
金庫室が破られた場合、2 番目の答えなら調べるべき明確な手がかりがある。
Joan には金庫室へのアクセスを許可し、メモを書く権限があったのか。
警備員が身分証のコピーをとっていれば、その情報は Sam Smith を特定して見つけ出す助けになる。
対照的に、Joan のメモに単に「持参人を金庫室に入れよ」としか書かれていなかったら（認証なしの認可）、セキュリティが破られた後、調査員には何が起きたのか、侵入者が誰なのか、ほとんど見当がつかなかったはずである。
\end{exa}

% src: Authentication
\underline{認証：}
認証の過程では、プリンシパルが本当に名乗っているとおりの者であることを示すクレデンシャルに基づいて、プリンシパルの身元の主張を検査する。
サービスによっては、チャレンジに対するデジタル署名のような、より強い形のクレデンシャルを使うこともある。
デジタル署名は、プリンシパルがその身元に対応する秘密鍵を持っていることを証明するものであり、ブラウザが HTTPS で Web サーバーを認証するときにもこの方法が使われている。
デジタル署名がより強い認証になるのは、プリンシパルが秘密を明かさずに、それを知っていることを証明できるからである。

\begin{dfn}[認証の証拠の種類]\label{dfn:ch1-authentication-factors}
認証に適した証拠は、次の種類に分けられる。
\begin{itemize}
\item \underline{知っているもの}（something you know）：パスワードなど。
\item \underline{持っているもの}（something you have）：セキュアトークンなど。アナログの世界では、偽造できない何らかの証明書、パスポート、署名入りの文書など。
\item \underline{その人自身であるもの}（something you are）：すなわち生体情報（指紋、虹彩パターンなど）。
\item \underline{いる場所}（somewhere you are）：検証された位置。たとえば、セキュアな施設内のプライベートネットワークに接続していること。
\end{itemize}
\end{dfn}

こうした方法の多くは、かなり当てにならない。
知っているものは漏れうるし、持っているものは盗まれたりコピーされたりしうるし、位置はさまざまな方法で操作されうるし、その人自身であるものでさえ偽装されうる（しかも侵害されたら、自分自身であるものを後から変えることはできない）。
さらに今日のネットワーク化された世界では、認証はほとんど常にネットワーク越しに行われるので、対面で認証するよりも難しくなる。
たとえば Web では、ブラウザが信頼の仲介役となってローカルで認証を行い、成功したときだけ暗号学的なクレデンシャルをサーバーに渡す。
こうした懸念を緩和するために、システムでは一般に複数の認証要素を使う。認証要素を頻繁に監査することも重要な安全策である。
弱い認証要素でも、2 つあれば 1 つよりはましである（ただし、大してましではない）。

しかし組織は、誰かにクレデンシャルを割り当てる前に、厄介な問題に向き合わなければならない。その人が入社するとき、アカウントを登録するとき、あるいはパスワードを忘れてアクセスを取り戻そうとヘルプデスクに電話してきたときに、その人の本当の身元をどうやって判定するかという問題である。

\begin{exa}[Google の入社時の身元確認]\label{exa:ch1-google-onboarding}
原著者が Google に入社したとき、新入社員全員がある月曜日の朝に集められ、数人の IT 管理担当者と向かい合った。
担当者たちは、パスポートやそのほかの身分証を新入社員名簿と照合した。
それが済んで初めて、バッジと会社支給のノートパソコンが渡され、ログインパスワードを設定するよう求められた。

IT チームは、提示されたクレデンシャル（身分証）が、本人が名乗っているとおりの人物であることを正しく示しているかを確かめることで、新入社員の身元を確認したのである。
この身元確認のセキュリティは、提出した政府発行の身分証と裏付け書類（たとえば出生証明書）の完全性にかかっていた。
それらの書類は、どれほど正確に発行されたのか。
偽造したり不正に入手したりするのは、どれほど難しいのか。
理想を言えば、出生時の登録から始まる結びつきの連鎖が生涯を通じて途切れずに保たれ、私たち一人ひとりを真正かつ一意に識別できるとよい。
\end{exa}

人々を確実に識別するのが難しいのは、主に、最も効果的な手法が権威主義のにおいを放ち、社会的に受け入れられないからである。
そのため、日常生活ではある程度のプライバシーと自由を守るために、あえて弱い方法を選んでいる。
人の本当の身元をどう判定するかという問題は本書の範囲外である。本書は、このより難しい ID 管理（identity management）の問題ではなく、ゴールドスタンダードに焦点を当てる。

できる限り、既存の信頼できる認証サービスに頼り、不必要に車輪の再発明をしないこと。
単純なパスワード認証でさえ安全に行うのはかなり難しく、忘れたパスワードを安全に扱うのはさらに難しい。
一般的に言えば、認証の過程ではクレデンシャルを検査し、合格か不合格かのどちらかだけを返すべきである。
部分的に成功したことを示すのは避けるべきである。
試行錯誤でクレデンシャルを絞り込もうとする攻撃者に手を貸しかねないからである。
総当たりの推測という脅威を緩和するには、認証そのものを計算の重い処理にしたり、過程に段階的に長くなる遅延を入れたりするのが一般的な戦略である（第4章の「予測可能性の回避」も参照）。

ユーザーを認証した後、システムはその身元をプリンシパルに安全に結びつける方法を用意しなければならない。
通常は、認証モジュールがプリンシパルにトークンを発行し、プリンシパルは以降のリクエストで、完全な認証の代わりにそのトークンを使う。
考え方としては、プリンシパルが Web ブラウザなどのエージェントを介して、自分が名乗っているとおりの者であることを手短に保証するものとして認証トークンを提示し、以降のリクエストのための「セキュアなコンテキスト」を作る、というものである。
このコンテキストによって、保存されたトークンが結びつけられ、認証されたプリンシパルに代わって以降のリクエストとともに提示されるようになる。
Web サイトでは閲覧セッションに対応づけたセキュアなクッキーでこれを行うことが多いが、ほかの種類のプリンシパルやインターフェースにはさまざまな手法がある。

認証された身元の安全な結びつきは、根本的に異なる 2 つの方法で破られうる。
わかりやすいのは、攻撃者が被害者の身元を乗っ取る場合である。
もう一つは、認証されたプリンシパル自身がぐるになって、自分の身元を譲り渡したり、他人に押しつけたりしようとする場合である。
後者の例が、有料ストリーミングのサブスクリプションの共有である。
Web ではこの結びつきが緩く、プリンシパルの協力に頼っているため、これを防ぐあまり良い手段がない。

% src: Authorization
\underline{認可：}
重要な行為を許可するか拒否するかは、認証で確立されたプリンシパルの身元に基づいて判断すべきである。
システムは、ビジネスロジック、アクセス制御リスト、そのほかの形式的なアクセスポリシーによって認可を実装する。

匿名の認可（すなわち認証なしの認可）が役に立つ状況も、まれにはある。
現実世界の例としては、混雑した駅の公共ロッカーの鍵を持っている場合が挙げられる。
時間に基づくアクセス制限（たとえば、データベースへのアクセスを営業時間内に限る）も、よくある例である。

ある資源に対する認可は、単一のガードで実施すべきである。
コードベースのあちこちに認可コードが散らばっていると、保守も監査も悪夢になる。
そうではなく、一様にアクセスを許可する共通のフレームワークに認可を任せるべきである。
きれいな設計構造にしておけば、開発者が認可を正しく実装しやすくなる。
できる限り、その場しのぎのわかりにくいロジックではなく、数多くある標準的な認可モデルのどれかを使うべきである。

\begin{dfn}[ロールベースアクセス制御]\label{dfn:ch1-rbac}
認証されたプリンシパルにロール（役割）を割り当て、ロールに基づいてアクセスを許可する方式を、\underline{ロールベースアクセス制御}（role-based access control, RBAC）という。
RBAC は認証と認可の橋渡しをし、一様なフレームワークによってアクセス制御を単純にする。
\end{dfn}

\begin{exa}[銀行のロール]\label{exa:ch1-bank-roles}
銀行のロールとしては、たとえば窓口係、マネージャー、融資担当者、警備員、財務監査人、IT 管理者が考えられる。
システムは、一人ひとりのアクセス権限を個別に選ぶのではなく、各人の身元に応じて 1 つ以上のロールを割り当て、ロールに対応する権限を自動的かつ一様に与える。
より高度なモデルでは、1 人が複数のロールを持ち、アクセスのたびにどのロールを適用するかを明示的に選ぶこともある。
\end{exa}

認可の仕組みは、オペレーティングシステムが伝統的に提供してきた単純な読み取り／書き込みのアクセス制御よりも、はるかにきめ細かくできる。
より堅牢な認可の仕組みを設計すれば、有用な機能を失わずにアクセスを制限し、セキュリティを強化できる。
こうしたより高度な認可モデルには、\underline{属性ベースアクセス制御}（attribute-based access control, ABAC）や\underline{ポリシーベースアクセス制御}（policy-based access control, PBAC）などがあり、ほかにも数多くある。

\begin{exa}[銀行の窓口係のきめ細かな認可]\label{exa:ch1-bank-teller}
きめ細かな認可によってポリシーをどう引き締められるかを、銀行の窓口係の簡単な例で見てみる。
\begin{itemize}
\item \underline{レート制限}（Rate-limited）：窓口係は 1 時間に 20 件まで取引を行えるが、それを超えると疑わしいとみなされる。
\item \underline{時間帯}（Time of day）：窓口係の取引は、勤務時間である営業時間内に行われなければならない。
\item \underline{自己取引の禁止}（No self-service）：窓口係は自分の個人口座で取引をしてはならない。
\item \underline{複数のプリンシパル}（Multiple principals）：10,000 ドルを超える窓口係の取引には、別途マネージャーの承認が必要である（これにより、1 人の悪人が一度に大金を動かすリスクがなくなる）。
\end{itemize}
\end{exa}

最後に、パスワードのような特定のデータについては、読み取り専用のアクセスでさえ権限が大きすぎる場合がある。

\begin{exa}[ハッシュによるパスワードの照合]\label{exa:ch1-password-hash-check}
システムでは通常、ハッシュどうしを比べてログインパスワードを確認し、平文のパスワードが漏れる可能性をまったく生じさせないようにしている。
ユーザー名とパスワードはフロントエンドサーバーに送られる。フロントエンドサーバーはパスワードをハッシュして認証サービスに渡し、平文パスワードの痕跡をすぐに破棄する。
認証サービスは、クレデンシャルのデータベースから平文パスワードを読むことはできないが、ハッシュは読むことができ、それをフロントエンドサーバーから渡されたハッシュと比べる。
こうして認証サービスはクレデンシャルを確認するが、パスワードそのものには一切アクセスしないので、たとえ侵害されてもパスワードを漏らすことはありえない。
\end{exa}

インターフェースの設計でこうした代替手段がとれるようになっていなければ、データ漏えいの可能性を緩和するこうした機会を逃すことになる。
これについては、第4章で最小情報（Least Information）のパターンを論じる際にさらに掘り下げる。

% src: Auditing
\underline{監査：}
組織がシステムの活動を監査できるようにするには、セキュリティを保つうえで重要なすべての出来事について、システムが信頼できるログを生成しなければならない。
認証と認可のイベント、システムの起動と停止、ソフトウェアの更新、管理者によるアクセスなどがこれにあたる。
完全に信頼できる記録とみなすには、監査ログが改ざんされにくくなっている必要があり、理想的には管理者でさえ手を加えにくくしておくべきである。
監査はゴールドスタンダードの重要な柱である。
インシデントは実際に起こるし、認証や認可のポリシーに欠陥がある場合もあるからである。
監査はまた、必要な監視の目を提供することで、信頼されたプリンシパルが害をなす内部犯行への緩和策にもなる。

適切に行えば、監査ログは日常的な監視に欠かせない。システムの活動レベルを測り、エラーや疑わしい活動を検知し、インシデントが起きた後には、攻撃が実際にいつどのように起きたのかを突き止めて被害の範囲を見積もるのに使える。
デジタルシステムを完全に守るとは、ポリシーを正しく実施することだけではなく、情報資産の責任ある管理者であることだ、ということを忘れてはならない。
監査によって、信頼されたプリンシパルが自分に与えられた広い権限の範囲内で適切に行動したことを確かめられる。

\begin{exa}[Twitter のパスワードのログ出力]\label{exa:ch1-twitter-password-log}
2018 年 5 月、Twitter は恥ずべきバグを公表した。
あるコード変更のせいで、生のログインパスワードがうっかり内部ログに出力されていたことがわかったのである。
これが実際に悪用された可能性は低いが、顧客の信頼を傷つけたことは確かであり、決して起きてはならないことであった。
\end{exa}

ログは多くの技術スタッフが日常的に目にする可能性があるので、開示のリスクをできるだけ小さくするために、運用上の詳細は記録しても、実際の私的な情報は一切保存すべきではない。
この要件の詳しい扱いについては、まさにこの問題に対処するロギングツールを詳述した付録Aのサンプル設計文書を参照してほしい。

システムはまた、悪事を隠すためにログを改ざんすることを誰にも許してはならない。
攻撃者がログを変更できれば、自分の活動の痕跡をきれいに消し去ってしまうだけである。
とくにセンシティブでリスクの高いログについては、内部犯行の実行者が自分の痕跡を隠せないように、管理上・運用上の統制が異なる独立したシステムで監査ログを管理すべきである。
これを完全に行うのは難しいが、独立した監視の目があるというだけで、たいていはどんな不正に対しても強い抑止力になる。
ちょうど、ささやかなフェンスと目立つ監視カメラが不法侵入を効果的に思いとどまらせるのと同じである。

さらに、システムを回避しようとする試みはどれも非常に疑わしく見えるし、少しでも下手な動きをすれば違反者は深刻な結果を招くことになる。
いったん捕まれば、彼らが罪を言い逃れるのは難しいであろう。

\begin{dfn}[否認不可能性]\label{dfn:ch1-non-repudiability}
ログに記録された行為について、行為者が後からそれを否認できない性質を\underline{否認不可能性}（non-repudiability）という。
否認不可能性は監査ログの重要な性質である。
\end{dfn}

名前のわかっている管理者がある時刻にあるコマンドを実行し、その直後にシステムがクラッシュしたことがログに示されていれば、他人に責任をなすりつけるのは難しい。
対照的に、組織が複数の管理者に同じアカウントを共有させていたら（ひどい考えである）、誰が実際に何をしたのかを確実に知る方法がなく、全員にもっともらしい言い逃れ（plausible deniability）の余地を与えてしまう。

結局のところ、監査ログが役に立つのは、ログを監視し、異常な出来事を注意深く分析し、必要に応じて適切な措置をとってフォローアップする場合だけである。
そのためには、ゴルディロックスの原則（Goldilocks principle）に従って、ちょうどよい量の詳細をログに記録することが大切である。
ログが多すぎると監視すべきデータ量が膨れ上がり、ノイズが多すぎたり整理されていなかったりするログからは有用な情報を拾い出しにくくなる。
一方、詳細が足りないまばらなログでは重要な情報が抜け落ちるかもしれないので、適切なバランスを見つけることは常に課題であり続ける。

%----------------------------------------------------------------------
\subsection{プライバシー}
% src: Privacy

情報セキュリティの基礎である C-I-A とゴールドスタンダードに加えて、原著者が紹介したいもう一つの基本的な話題が、関連分野である情報プライバシーである。
セキュリティとプライバシーの境界をはっきり定めるのは難しく、両者は密接に関係していると同時に、かなり異なってもいる。
本書では、両者を一つにまとめようとするのではなく、ソフトウェアを作る過程にセキュリティとプライバシーの両方を組み込むために、両者が重なり合う部分に焦点を当てる。

人々のデジタルな情報プライバシーを尊重するには、次のような人的要素も考慮に入れて、機密性の原則を広げなければならない。
\begin{itemize}
\item 情報の収集と利用に対する顧客の期待
\item 情報の適切な利用と開示についての明確なポリシー
\item さまざまな種類の情報の取り扱いに関する法的・規制上の問題
\item 個人情報の処理に関する政治的・文化的・心理的側面
\end{itemize}

ソフトウェアが現代生活にますます浸透するにつれて、人々はソフトウェアをより私的な形で使い、生活のセンシティブな領域にも取り入れるようになった。その結果、多くの複雑な問題が生じている。
過去の事故や濫用によってリスクが目に見えるようになり、社会が政治的・法的な手段で新たな課題に取り組むなかで、私的な情報を適切に扱うことは難しくなっている。

ソフトウェアセキュリティの文脈では、これは次のことを意味する。
\begin{itemize}
\item あらゆるデータの収集と共有が、顧客や利害関係者にどんな結果をもたらすかを考える。
\item 潜在的な問題をすべて指摘し、必要に応じて専門家の助言を得る。
\item 私的な情報の利用について明確なポリシーとガイドラインを定め、それに従う。
\item ポリシーとガイダンスを、ソフトウェアで実施されるチェック・アンド・バランスに落とし込む。
\item データの取得・利用・共有・削除について正確な記録を残す。
\item コンプライアンスのために、データアクセスの認可と例外的なアクセスを監査する。
\end{itemize}

システムを適切に制御し適切なアクセスを提供するというセキュリティの仕事は比較的はっきりしているが、プライバシーの仕事はそれほど明確に定まっていないことが多い。
また、社会がより多くのデータを収集する未来へと深く踏み込んでいくなかで、プライバシーに関する期待や規範はまだ模索されている途中である。
こうした課題を考えると、データの利用について透明性を最大限に高めること（ポリシーを誰にでも理解できるほど簡潔にしておくことを含む）と、データ、とくに個人を特定できる情報の収集を最小限にすることを検討するのが賢明である。

情報は特定の目的のためだけに収集し、役に立つ間だけ保持すること。
許可された用途が設計で想定されていないなら、そもそも収集しないこと。
「いつか」使うためにむやみにデータを集めるのは危険であり、良い考えであることはまずない。
あるデータの最後の許可された用途が必要なくなったら、安全に削除するのが最善の保護である。
とくにセンシティブなデータや、プライバシーを最大限に守りたい場合には、この方針をさらに強め、開示されたときの潜在的なリスクが保持し続ける潜在的な価値を上回った時点でデータを削除すべきである。

\begin{exa}[長年分のメールの保持]\label{exa:ch1-email-retention}
何年分ものメールを保持しておくと、ときどき何かの役に立つかもしれないが、明確なビジネス上の必要があることはおそらくない。
しかも社内メールは、漏えいしたり、たとえば召喚状（subpoena）によって開示を強制されたりすれば、負債になりうる。
そうしたデータを「念のため」いつまでも抱え込むよりも、削除するのがたいていは最善のポリシーである。
\end{exa}

情報プライバシーを完全に扱うことは本書の範囲外であるが、人々についてのデータを収集するシステムを設計するときには、プライバシーとセキュリティは切り離せない二つの側面である。
そして人々は、何らかの形でほとんどすべてのデジタルシステムとやり取りしている。
強力なプライバシー保護はセキュリティが堅固であって初めて可能になる。だからここで訴えたいのは、プライバシーへの配慮を意識し、設計段階からソフトウェアに組み込んでほしいということである。

複雑ではあるものの、プライバシーについて広く知られたベストプラクティスが一つある。
プライバシーに関する期待を明確に伝えることである。
セキュリティとは対照的に、プライバシーポリシーには、情報サービスが顧客データをどこまで活用したいか、したくないかについて大きな裁量の余地がありうる。
「あなたのデータを再利用し、販売します」はプライバシーのスペクトルの一方の極端であるが、「日によってはあなたのデータを保護しないかもしれません」というのは、セキュリティに対する姿勢としては成り立たない。
プライバシーの失敗が起きるのは、ユーザーの期待が実際のプライバシーポリシーとずれているとき、あるいは明確なポリシーがあるのにそれが何らかの形で破られたときである。
前者の問題は、データの取り扱いをユーザーに前もって説明していないことから生じる。
後者は、ポリシーがあいまいなとき、責任あるスタッフにポリシーが無視されたとき、あるいはセキュリティが破綻してポリシーが損なわれたときに起こる。

\newpage

%% =======================================================
%%  part: ch2
%% =======================================================
%% =======================================================
%%  part: ch2
%% =======================================================
\section{脅威}
% src: 2: Threats

この章では、\cite{Kohnfelder2021} の第2章に沿って、脅威と脅威モデリングを扱う。

\begin{quote}
「脅威は、たいてい脅威の対象そのものよりも恐ろしい。」 ---Saul Alinsky
\end{quote}

脅威はどこにでもあるが、うまく管理すれば、脅威と共存していくことはできる。
ソフトウェアも同じである。ただしソフトウェアの場合、脅威に備えさせてくれる何百万年もの進化の恩恵がない。
だからこそ、ソフトウェアセキュリティの考え方（security mindset）を身につける必要がある。そのためには、作り手の視点から攻撃者の視点へと頭を切り替えなければならない。
システムに対する潜在的な脅威を理解することは、ソフトウェアの設計に堅固な防御と緩和策を組み込むうえで欠かせない出発点である。
しかし、そもそも脅威に気づくためには、典型的なユースケースや、ソフトウェアを意図どおりに使う場面を考えるのをやめなければならない。
代わりに、ソフトウェアをありのままの姿で見る必要がある。つまり、コードとコンポーネントの寄せ集めであり、その中をデータが流れ、あちこちに保存されている、という姿である。

\begin{exa}[ゼムクリップ]\label{exa:ch2-paperclip}
ゼムクリップは紙の束をまとめておくために巧みに設計されているが、うまく曲げれば、硬い針金に簡単に作り変えられる。
セキュリティの考え方を持つ人は、この針金を錠前の鍵穴に差し込んでタンブラーを動かせば、鍵がなくても開けられることに気づく。
\end{exa}

ここで強調しておきたいのは、脅威には害が生じるあらゆる経路が含まれるということである。
意図的に行われる敵対的な攻撃は議論の重要な焦点だが、だからといって、ソフトウェアのバグ、人為的ミス、事故、ハードウェア故障などによるほかの脅威を除外してよいわけではない。

\underline{脅威モデリング}（threat modeling）は、ソフトウェア開発プロセス全体を通じて、セキュリティに影響するあらゆる判断の指針となる視点を与えてくれる。
以下では、脅威モデリングを行うための数多くの具体的な方法論のどれかを取り上げるのではなく、概念と原則に焦点を当てる。
2000年代初頭に Microsoft で初めて実践された初期の脅威モデリングは効果を上げたが、大規模なトレーニングと相当な労力の投資を必要とした。
幸い、脅威モデリングのやり方はいくらでもある。概念さえ理解すれば、使える時間と労力に合わせてプロセスを調整しながら、意味のある結果を出すのは難しくない。

大規模なソフトウェアシステムのすべての脅威を列挙し、脆弱な箇所をすべて特定しようとするのは、気の遠くなる作業である。
しかし、賢いセキュリティ作業は完璧を目指すのではなく、\emph{少しずつ}水準を引き上げることを目標とする。
最初の取り組みでは潜在的な問題のごく一部しか見つからず、そのうちの一部しか緩和できないかもしれない。それでも大きな改善である。
ひょっとすると、その取り組みのおかげで重大なセキュリティインシデントを避けられるかもしれない。それは本当の成果である。
残念ながら、阻止した攻撃のことを知る機会はほとんどなく、手応えが得られないことを物足りなく感じるかもしれない。
セキュリティの考え方という筋肉は、鍛えれば鍛えるほど脅威を見抜くのがうまくなる。

最後に、脅威モデリングによって、セキュリティの範囲を超えて対象システムへの理解が一段と深まることがある、という点も押さえておきたい。
ソフトウェアを新しい角度から調べるうちに、セキュリティとは関係のないさまざまな改善、効率化、単純化、新機能につながる洞察が得られることがある。

%----------------------------------------------------------------------
\subsection{敵対者の視点}\label{subsec:ch2-adversarial-perspective}
% src: The Adversarial Perspective

\begin{quote}
「エクスプロイトは、現実世界で私たちが目にするもののうち、最も『魔法の呪文』に近い。正しい呪文を組み立てれば、デバイスを遠隔操作できる。」 ---Halvar Flake
\end{quote}

究極の脅威は人間の加害者である。セキュリティインシデントがひとりでに起こるわけではない。
ソフトウェアセキュリティを本腰を入れて分析するなら、潜在的な攻撃にきちんと備えるために、仮想的な敵対者が何を試みるかを考えることになる。
攻撃者は雑多な集団で、\underline{スクリプトキディ}（script kiddies：自動化されたマルウェアを使う、技術的スキルを持たない犯罪者）から高度な国家レベルの攻撃者まで、あらゆる層がいる。
敵対者の視点で考えられるなら、それに越したことはない。しかし、相手の動きをすべて正確に予測できると思い込んだり、狡猾な敵を出し抜く名探偵のように相手の頭の中に入り込もうとして時間をかけすぎたりしてはならない。
攻撃者の考え方を理解するのは有益だが、セキュアなソフトウェアを作るという目的にとっては、攻撃者が探索・侵入・データの持ち出しに使う実際の手法の細部は重要ではない。

システムの明らかな標的が何かを考え（敵対者にとって価値のあるものが自分にとってはそれほど価値がないこともあれば、その逆もある）、それらの資産を堅固に保護すべきである。ただし、仮想的な攻撃者の心を読もうとして時間を無駄にしてはならない。
攻撃者は余計な労力を費やすより、目的を達成するために最も弱い環に狙いを定めることが多い（あるいは、目的もなくあちこち探っていることもある。その場合、行動に方向性がなく恣意的に見えるので、防御は非常に難しい）。
バグは弱さの表れなので確実に注目を集める。明らかなバグに出くわした攻撃者は、本当に何かを壊せるかどうか、工夫を凝らしたさまざまな変種を試す。
システム内部の詳細を明かしてしまうエラーや副作用（例えば詳細なスタックダンプ）は、攻撃者が飛びついて利用する格好の材料である。

攻撃者は弱点をいったん見つけると、そこにさらに労力を注ぐ可能性が高い。小さな欠陥でも、集中的に攻撃されると広がって大きな結果を招くことがあるからである（第8章で詳しく見る）。
単独では問題にならない2つの小さな欠陥を組み合わせて重大な攻撃を生み出せることも多いので、どの脆弱性も真剣に受け止めるのが賢明である。
また、攻撃者は脅威モデリングのことを確実に知っている。ただし、内部情報なしで作業している（少なくとも、ある程度の侵入に成功するまでは）。

敵対者が何に時間を費やすかを本当に予期することはできないが、仮想的な攻撃者の動機を、熱心な攻撃を受ける可能性の目安として考えることには意味がある。
要するに、ある有名な犯罪者が銀行を襲った理由として語った「そこに金があるからだ」という言葉に尽きる。
つまり、システムを攻撃して得られる見込みの利益が大きいほど、潜在的な攻撃者はより高いスキルと多くの資源を投入してくると予想できる。
推測にすぎないとはいえ、相対的な指針としては役に立つ分析である。強力な企業や、政府・軍・金融機関は大きな標的である。あなたの猫の写真はそうではない。

結局のところ、あらゆる種類の暴力と同じく、攻撃して害を与えるほうが防御するよりもはるかに容易である。
攻撃者は侵入地点を選べるし、その気になれば好きなだけ多くのエクスプロイトを試せる。一度成功すればよいからである。
こうしたことすべてが、セキュリティ作業を優先すべき理由をさらに増やしている。防御側は、使える優位性はすべて使う必要がある。

%----------------------------------------------------------------------
\subsection{4つの質問}\label{subsec:ch2-four-questions}
% src: The Four Questions

Microsoft で長年にわたり脅威モデリングの旗振り役を務めた Adam Shostack は、その方法論を次の4つの質問に集約している。

\begin{dfn}[4つの質問]\label{dfn:ch2-four-questions}
脅威モデリングの方法論を集約した次の4つの質問を、\underline{4つの質問}（Four Questions）という。
\begin{enumerate}
\item 私たちは何に取り組んでいるのか？（What are we working on?）
\item 何がうまくいかなくなりうるか？（What can go wrong?）
\item それに対して私たちは何をするのか？（What are we going to do about it?）
\item 私たちはうまくやれたか？（Did we do a good job?）
\end{enumerate}
\end{dfn}

最初の質問の目的は、プロジェクトの文脈と範囲をはっきりさせることである。
この質問に答えるには、プロジェクトの要件と設計、コンポーネントとその相互作用を記述し、運用上の問題やユースケースも検討する。
続いてこの方法の核心として、2番目の質問では潜在的な問題を予期しようとし、3番目の質問では特定した問題への緩和策を探る（緩和策は第3章で詳しく見るが、まず緩和策と脅威の関係を調べる）。
最後の4番目の質問では、プロセス全体---ソフトウェアが何をするのか、どのようにうまくいかなくなりうるのか、脅威をどれだけうまく緩和できたか---を振り返り、リスクがどれだけ下がったかを評価して、システムが十分にセキュアになることを確かめる。
未解決の問題が残っていれば、もう一度質問を繰り返して、残った隙間を埋める。

脅威モデリングにはこのほかにも多くの要素があるが、4つの質問から取り組むだけでも驚くほど遠くまで進める。
これらの概念を身につけ、本書のほかの考え方や手法と組み合わせれば、自分が構築・運用するシステムのセキュリティ水準を大きく引き上げられる。

%----------------------------------------------------------------------
\subsection{脅威モデリング}\label{subsec:ch2-threat-modeling}
% src: Threat Modeling

\begin{quote}
「いったい何がまずいことになるっていうんだ？」（What could possibly go wrong?）
\end{quote}

この質問は、皮肉な冗談として口にされることが多い。
しかし、皮肉抜きで問えば、脅威モデリングの出発点を簡潔に言い表している。
この最初の質問に答えるには、脅威を特定し評価する必要がある。そのうえで脅威に優先順位を付け、重要な脅威のリスクを下げる緩和策に取り組める。

この文を分解してみよう。基本的な脅威モデリングのプロセスは、おおよそ次の手順で表せる。

\begin{dfn}[脅威モデリングの基本プロセス]\label{dfn:ch2-threat-modeling-process}
\underline{脅威モデリング}（threat modeling）の基本的なプロセスは次の手順からなる。
\begin{enumerate}
\item 範囲内のすべてを漏れなく検討するために、システムのモデルに基づいて作業する。
\item システム内で保護が必要な\underline{資産}を特定する。
\item システムのモデルをコンポーネントごとにくまなく調べて潜在的な脅威を探し、\underline{攻撃対象領域}（attack surface：攻撃が始まりうる場所）、資産（価値あるデータと資源）、\underline{信頼境界}（trust boundary：システムのうち信頼度の高い部分と低い部分をつなぐインターフェース）、そしてさまざまな種類の脅威を特定する。
\item 特定した潜在的な脅威を、最も具体的なものから仮想的なものまで分析する。
\item 脅威を、重大なものから順に順位付けする。
\item 最も重大な脅威について、リスクを下げる緩和策を提案する。
\item 効果が最も大きく、最も容易な緩和策から順に追加し、収穫逓減が始まるまで続ける。
\item 最も重大な脅威に対する緩和策から順に、その有効性をテストする。
\end{enumerate}
\end{dfn}

複雑なシステムでは、潜在的な脅威をすべて挙げた一覧は膨大になり、完全な分析はほぼ確実に実行不可能である（何かを行う考えうる方法をすべて列挙しようとしても、想像力を働かせれば決して終わらないのと同じである。そして攻撃者はしばしば想像力を働かせる）。
実際には、最初の脅威モデリングでは、価値の高い資産に対する最大かつ最も起こりやすい脅威だけに集中すべきである。
それらの脅威を理解して第一線の緩和策を講じたら、すでに特定してある残りの小さな脅威を繰り返し検討していくことで、残存リスクを評価できる。
そこから先は、必要に応じて追加の脅威モデリングを1回以上行い、そのたびに網を広げて、ほかの資産、より深い分析、起こりにくい脅威や軽微な脅威をさらに含めていく。
最も重要な脅威を十分に深く理解し、必要な緩和策を計画し、残る既知のリスクを許容できると判断した時点で、このプロセスは終わる。

\begin{exa}[日常生活における脅威モデリング]\label{exa:ch2-private-message}
人は日常生活の中で、常識的な用心と呼ばれる、脅威モデリングに似たことを直観的に行っている。
人前でプライベートなメッセージを送るとき、ほとんどの人は携帯電話に声で吹き込むのではなく、文字で入力する。
脅威モデリングの言葉でいえば、メッセージの内容が情報資産であり、開示が脅威である。
他人に聞こえる範囲で話すことが攻撃対象領域であり、声を出さない別の入力方法を使うことは良い緩和策である。
詮索好きな他人がのぞき込んでいるなら、もう一方の手で携帯電話を覆って画面を隠すといった緩和策を追加することもできる。
\end{exa}

現実世界ではこうしたことをごく自然に、いつも行っている。しかし、慣れ親しんだ物理的な直観が通用しない複雑なソフトウェアシステムに同じ手法を当てはめるには、はるかに多くの規律が必要である。

%----------------------------------------------------------------------
\subsubsection{モデルに基づいて作業する}\label{subsec:ch2-work-from-model}
% src: Work from a Model

脅威を漏れなく特定するには、厳密なアプローチが必要である。
脅威モデリングでは伝統的に、システムを\underline{データフロー図}（data flow diagram, DFD）や \underline{UML}（Unified Modeling Language）で記述したものを使うが、モデルは何を使ってもよい。
DFD、UML、設計文書、あるいは非公式な「ホワイトボードセッション」のどれであれ、システムの高レベルの記述として何を選んでも構わない。分析に必要な詳細を捉えられるだけの粒度がある限り、システムを抽象化したものを見る、というのが基本的な考え方である。

形式的なアプローチほど厳密で正確な結果が得られる傾向があるが、その分、時間と労力が余計にかかる。
セキュリティコミュニティは長年にわたって、それぞれ異なるトレードオフを持つ代替の方法論を数多く考案してきた。その大きな理由は、本格的な脅威モデリングの方法（DFD のような形式的なモデルを含むもの）が非常にコストと労力のかかるものだからである。
今日では、このプロセスを支援する専用のソフトウェアも使える。
最良のツールは作業の大部分を自動化してくれるが、結果の解釈とリスク評価には常に人間の判断が必要である。
本書では、システムを十分に理解して4つの質問に漏れなく答えられさえすれば、特別な図やツールなしで自分で脅威モデリングを行えるよう、知っておくべきことをすべて伝える。
そこから先は、望むならより高度な形式へ進めばよい。

どのモデルに基づいて作業するにしても、対象システムを適切な解像度で漏れなくカバーすること。
分析の詳細度は、\underline{ゴルディロックスの原則}（Goldilocks principle）に従って選ぶ。細かくしすぎると作業が終わらなくなり、高レベルすぎると重要な細部を見落とす。
ほとんど成果のないままプロセスがすぐに終わってしまうなら、粒度が粗すぎる確かな兆候である。同様に、何時間作業してもほとんど進まないなら、モデルの粒度が細かすぎる可能性がある。

\begin{exa}[一般的な Web サーバーのモデルの粒度]\label{exa:ch2-web-server-granularity}
一般的な Web サーバーにとって、適切な粒度とはどの程度かを考えよう。
左に「インターネット」、中央に「フロントエンドサーバー」、右に3番目のコンポーネント「データベース」があり、それらが線でつながっただけのブロック図のモデルを渡されたとする。
これでは役に立たない。これまでに考案されたほぼすべての Web アプリケーションがこのモデルに当てはまってしまうからである。
資産はおそらくすべてデータベースの中にあるが、具体的には何なのか。
システムとインターネットの間には信頼境界があるはずだが、信頼境界はそれだけなのか。
明らかに、このモデルはレベルが高すぎる。
反対の極端な例は、すべてのライブラリ、フレームワークのすべての依存関係、分析したいアプリケーションのレベルよりはるかに下にあるコンポーネントどうしの関係まで、細かく分解して示したモデルであろう。

ゴルディロックス版は、この両極端の間のどこかに収まる。
データベースに格納されたデータ（資産）は、それぞれをひとまとまりとして扱えるカテゴリーにまとめる。例えば、顧客データ、在庫データ、システムログである。
サーバーコンポーネントは、複数のプロセス（それぞれがどの権限で動くかを含む）、場合によってはホストマシン上の内部キャッシュ、そしてインターネットやデータベースとの通信に使う通信路とネットワークの説明が見える程度の粒度に分解する。
\end{exa}

%----------------------------------------------------------------------
\subsubsection{資産を特定する}\label{subsec:ch2-identify-assets}
% src: Identify Assets

モデルを順序立てて調べながら、資産と、資産に対する潜在的な脅威を特定する。

\begin{dfn}[資産]\label{dfn:ch2-asset}
システム内で保護しなければならない実体を\underline{資産}（asset）という。
資産の多くはデータだが、ハードウェア、通信帯域、計算能力、さらには電力のような物理的資源も含まれうる。
\end{dfn}

脅威モデリングの初心者は当然すべてを守りたがるし、完璧な世界ならそれが理想である。
しかし実際には、資産に優先順位を付ける必要がある。
例えば、任意の Web アプリケーションを考えよう。インターネット上の誰もが、こちらではまったく制御できないブラウザやほかのソフトウェアを使ってアクセスできるので、クライアント側を完全に保護することはできない。
また、内部のシステムログは常に非公開にしておくべきだが、ログに外部の者にとって価値のない無害な詳細しか含まれていないなら、ログの保護に多くの労力を注ぐのは理にかなわない。
そうしたリスクを完全に無視せよという意味ではない。重要度の低い緩和策のせいで、ほかの場所で必要な労力が奪われないようにせよということである。
例えば、機密性のないログを、管理者だけが内容を読めるようにパーミッションを設定して保護するのは文字どおり1分で済むので、費やす価値のある労力である。

一方、金融取引を表すデータは、本物のお金として扱い、それに見合った優先順位を付けるとよい。
個人情報も、機密性がますます高まっている資産のカテゴリーである。ある人の所在地やその他の識別情報を知られると、その人のプライバシーが損なわれたり、身の危険にさらされたりしかねないからである。

\begin{rem}[複雑なリスク評価計算を避ける]\label{rem:ch2-avoid-risk-calculations}
著者は一般に、複雑なリスク評価計算はしないよう勧めている。
例えば、リスクを順位付けするために金額を割り当てようとするのは避けるべきである。
そのためには、知りようのない多くの事柄について、何とかして確率をひねり出さなければならない。
何人の攻撃者があなたを標的にし、どれほど熱心に、何をしようとするのか。
どれくらいの頻度で、どの程度成功するのか。
そもそも顧客データベースにはいくらの価値があるのか（会社にとっての価値と、攻撃者が売却できる金額は異なることが多く、ユーザーが自分のデータに認める価値もまた異なりうることに注意）。
仮想的なセキュリティインシデントによって、何時間分の作業やどれだけの費用が発生するのか。
\end{rem}

\begin{dfn}[T シャツサイズによる資産の順位付け]\label{dfn:ch2-tshirt-sizes}
資産を「T シャツのサイズ」で順位付けするという、簡単なのに驚くほど効果的な方法がある（著者が有用と考える単純化であり、業界標準の慣行ではない）。
\begin{itemize}
\item 最大限に保護しなければならない主要な資産には「Large」を割り当てる。
\item 価値はあるがそれほど重大ではない資産には「Medium」を割り当てる。
\item 重要性の低い軽微な資産には「Small」を割り当てる（通常は一覧にさえ載せない）。
\item 価値の高いシステムには、金融機関の銀行口座残高や、通信のセキュリティの基盤となる秘密の暗号鍵のように、並外れた水準の保護に値する「Extra-Large」の資産があることもある。
\end{itemize}
この単純な方式では、保護と緩和の取り組みをまず Large の資産に集中させ、次に Medium の資産に\underline{機会主義的に}（opportunistically）取り組む。
機会主義的な保護とは、少ない労力で済み、欠点もほとんどない作業のことである。
Small の資産をごく機会主義的に保護できるとしても、そこに時間を割く前に、Large の資産をすべて守ること。
\end{dfn}

第13章では脆弱性の順位付けを詳しく論じており、その多くは脅威の評価にも当てはまる。

\begin{exa}[PayPal による銀行口座の所有確認]\label{exa:ch2-paypal-deposit}
本物のお金が資産を保護するための資源として使われる、珍しいがわかりやすい例を考えよう。
銀行口座を PayPal に接続するとき、Web サイトはその口座が本当にあなたのものであることを確認しなければならない。
この段階であなたはすでにアカウントを持っており、PayPal は検証済みのメールアドレスを知っている。しかし今度は、あなたがある銀行口座の正当な所有者であることを確かめる必要がある。
PayPal はこの課題に巧みな解決策を考え出したが、少しお金がかかる。
同社は、新しいユーザーが所有していると主張する銀行口座に、ランダムな金額を入金するのである（入金額を \$0.01 から \$0.99 の間とすると、平均コストは顧客1人あたり \$0.50 である）。
銀行間送金では、事前の承認なしにどの口座にも入金できる。起こりうる最悪の事態といっても、文字どおり、誰かの口座に謎の寄付が入るだけだからである。
入金後、PayPal は入金額を答えるよう求める。これは口座の所有者にしかできないので、正しく答えられれば所有の証明として扱う。
PayPal はこのプロセスで文字どおりお金を失うが、スタッフに給料を払って銀行口座の所有を確認させるほうが遅く、コストもかかるので、非常に理にかなったやり方である。
\end{exa}

\begin{exa}[PayPal の口座認証の脅威モデリング]\label{exa:ch2-paypal-threat-model}
いま述べた銀行口座の認証プロセスを脅威モデリングしてみよう（この議論のために実際のプロセスを単純化したものであり、著者は実際のプロセスについての詳しい情報を持っていない）。
例えば、偽の PayPal アカウントを100個開設し、それぞれについて入金額をランダムに推測すれば、一度は認証される見込みがそれなりにある。
そうなれば、その口座を乗っ取ったことになる。
PayPal はこのような攻撃をどう緩和できるだろうか。
ほかにどのような攻撃と緩和策が考えられるだろうか。

分析の手がかりとなる観点をいくつか挙げる。
大量の推測という脅威に対しては、いくつもの制限を設けて、敵対者により多くの労力を強いることができる。
同じアカウントからの銀行口座の設定は1日1回までとし、同じコンピュータ（IP アドレス、ユーザーエージェント、その他のフィンガープリントで識別する）からの新規アカウント作成も制限する。
このような制限を\underline{レート制限}（rate limiting）といい、理想的には、試行を繰り返すほど課す遅延を大きくすべきである（例えば、2回目の失敗の後は、攻撃者は再試行まで1週間待たなければならない、など）。

このプロセスには微妙な点がある。ユーザーの利便性とセキュリティのバランスを取らなければならないからである。
ユーザーが銀行情報を入力して検証を要求するだけの場合、入力ミスがあるとプロセスをやり直さなければならないことがある。そのため、正直な顧客のことを考えると、レート制限はかなり緩くせざるをえない。
そこで、入力がうまくできないだけの顧客を失わずにレート制限を厳しく保つために、銀行の詳細情報を入力する際のミスを減らす方法を検討すべきである。
一つの方法は、銀行情報を2回入力してもらい、両者が一致した場合にだけ先に進むことである。
これでうまくいくが、手間が増えるので、面倒くさがりの人はあきらめるかもしれない。そうなれば良い顧客を失うことになる。
法的な問題を別にすれば、無効にした小切手の写真をアップロードしてもらうほうがよいかもしれない。
システムが印刷された銀行情報を読み取って表示し、顧客に確認してもらえば、ミスの可能性は事実上なくなる。

しかし、このシステムを何年も使った後で、攻撃が何度も成功していたことが発覚したらどうなるか。
もしかすると、辛抱強い泥棒たちがレート制限が明けるのを待って推測を続け、99分の1という推測の確率では止めきれなかったことが判明するかもしれない。
ほかの条件がすべて同じなら、PayPal は「タダのお金」の入金額を最大 \$3 や \$5、あるいはそれ以上に引き上げることもできる。しかし、ある時点で（正確な損益分岐点はおそらくアクチュアリーが教えてくれるだろう）、入金にかかる金銭的コストが新規顧客獲得の価値を上回ってしまう。

その場合、会社はまったく別のアプローチを検討しなければならない。
一つのアイデアを挙げる（ほかのアイデアもぜひ考えてみてほしい）。新規顧客の設定を、生身のカスタマーサポート担当者がビデオ通話で行うのである。
本物の人間と向き合わなければならないというだけで、多くの攻撃者はそもそも怖気づくであろう。
担当者は銀行の取引明細書などの証拠を見せるよう求め、その場で承認できる（これは単純化した例であり、実際のビジネス上の提案ではない）。
\end{exa}

優先して保護する資産には、おそらくデータが含まれる。例を挙げれば、顧客の資源、個人情報、業務文書、運用ログ、ソフトウェアの内部などである。
データ資産の保護に優先順位を付けるときは、情報セキュリティ（第1章で論じた C-I-A の三要素）を含む多くの要因を考慮する。データの漏えい、改変、破壊がもたらす害は大きく異なりうるからである。
情報の漏えいは、部分的な開示（例えばクレジットカード番号の下4桁）も含めて評価が難しい。攻撃者がその情報で何ができるかを考えなければならないからである。
攻撃者が複数の断片的な情報をつなぎ合わせて、完全なデータセットに近いものを作れる場合、分析はさらに難しくなる。

資産をひとまとめにすれば分析はかなり単純になるが、その過程で解像度を失わないよう注意すること。
例えば、いくつかのデータベースを一緒に管理し、同じようにアクセスを許可し、似た発生源からのデータに使い、同じ場所に保存しているなら、まとめて一つとして扱うのは十分に理にかなっている。
しかし、これらの要因のどれかが大きく異なるなら、別々に扱う十分な理由がある。
そうした違いは、リスク分析でも緩和策の検討でも考慮すること。

最後に、資産の価値は常に、関係するすべての当事者の視点から考えること。
例えば、ソーシャルメディアサービスは、社内の計画、広告データ、顧客データなど、あらゆる種類のデータを管理している。
これらの資産の価値は、あなたがその会社の CEO なのか、広告主なのか、顧客なのか、あるいは金銭的利益や政治的な目的を追い求める攻撃者なのかによって、それぞれ異なる。
実際、顧客の間でさえ、通信におけるプライバシーをどれほど重視するか、自分のデータにどれほどの価値を置くかは、人によって大きく異なるであろう。
優れたデータ管理（data stewardship）の原則に従えば、顧客やパートナーのデータは、会社自身の独自データ以上に手厚く保護すべきだとさえ言える（実際に会社の幹部がこれを方針として掲げたという話を著者は聞いたことがある）。

\begin{exa}[Facebook の Beacon]\label{exa:ch2-facebook-beacon}
すべての企業がこのアプローチを取るわけではない。
Facebook の Beacon 機能は、ユーザーの購入の詳細を自動的にそのユーザーのニュースフィードに投稿したが、たちまち顧客の怒りが噴出し、いくつかの訴訟も起こされたため、すぐに停止された。
Beacon が Facebook を危険にさらしたことはなかった（ブランドの評判を傷つけたことを除けば）が、顧客には現実の危険をもたらした。
顧客にとっての情報漏えいの結果を脅威モデリングしていれば、クリスマスや誕生日のプレゼント、さらに悪いことには婚約指輪の購入が意図せず知られてしまうことが問題になりそうだと、すぐにわかったはずである。
\end{exa}

%----------------------------------------------------------------------
\subsubsection{攻撃対象領域を特定する}\label{subsec:ch2-identify-attack-surfaces}
% src: Identify Attack Surfaces

\begin{dfn}[攻撃対象領域]\label{dfn:ch2-attack-surface}
攻撃者が最初に侵入する地点を\underline{攻撃対象領域}（attack surface）という。
\end{dfn}

攻撃対象領域は攻撃者が最初に侵入する地点なので、特に注意を払うこと。
攻撃対象領域を小さくできる機会があれば、どんなものでも大きな成果とみなすべきである。そうすれば、潜在的な問題の発生源そのものを完全に断てるからである。
多くの攻撃はシステム全体に広がっていく可能性があるので、早い段階で食い止めることは優れた防御になりうる。
セキュリティの厳しい政府の建物が、唯一の一般用入口を入ってすぐの場所に金属探知機付きの検問所を設けているのはこのためである。

ソフトウェアの設計は通常、物理的な建物の設計よりもはるかに複雑なので、攻撃対象領域全体を特定するのはそれほど簡単ではない。
システムを信頼できるセキュアな環境に組み込めない限り、ある程度の攻撃対象領域ができるのは避けられない。
インターネットは常に巨大な露出点になる。文字どおり、どこの誰でもインターネットを通じて匿名で接続できるからである。
イントラネット（プライベートネットワーク）は信頼できるとみなしたくなるかもしれないが、物理的セキュリティと IT セキュリティの両方で非常に高い水準を満たしていない限り、おそらくそうすべきではない。
少なくとも、リスクが低めの攻撃対象領域として扱うこと。
デバイスやキオスクアプリケーションの場合は、画面やユーザーインターフェースのボタンを含む筐体の外側を攻撃対象領域とみなすこと。

\begin{rem}[デジタル領域の外にある攻撃対象領域]\label{rem:ch2-non-digital-attack-surface}
攻撃対象領域はデジタル領域の外にも存在することに注意する。
キオスクを例にとると、公共の場所にある画面からは「\underline{ショルダーサーフィン}」（shoulder surfing：肩越しののぞき見）によって情報が漏れうる。
攻撃者はさらに巧妙な\underline{サイドチャネル攻撃}（side-channel attack）を仕掛け、電磁放射、熱、消費電力、キーボードの音などを監視して、システムの内部状態を推測することもできる。
\end{rem}

%----------------------------------------------------------------------
\subsubsection{信頼境界を特定する}\label{subsec:ch2-identify-trust-boundaries}
% src: Identify Trust Boundaries

次に、システムの信頼境界を特定する。
信頼と権限はほぼ常に対になっているので、そのほうがわかりやすければ、\underline{権限境界}（privilege boundary）という観点で考えてもよい。
人間の世界で信頼境界に相当するものとしては、管理職と従業員の間のやり取りや、誰を中に入れるかを自分で選ぶ自宅のドアが挙げられる。

\begin{exa}[カーネルとユーザーランドの境界]\label{exa:ch2-kernel-userland}
信頼境界の古典的な例として、オペレーティングシステムのカーネル／ユーザーランドのインターフェースを考えよう。
このアーキテクチャは、メインフレームコンピュータが希少で、多くのユーザーで共有されることが多かった時代に普及した。
システムはカーネルを起動し、カーネルは別々のユーザーランドのプロセスインスタンス（別々のユーザーアカウントに対応する）にあるアプリケーションを隔離して、互いに干渉したりシステム全体をクラッシュさせたりしないようにした。
ユーザーランドのコードがカーネルを呼び出すたびに、実行は信頼境界を越える。
\end{exa}

信頼境界が重要なのは、より高い権限での実行に移る場面が、より大きな問題を引き起こす機会になるからである。

\begin{rem}[信頼と権限]\label{rem:ch2-trust-vs-privilege}
本書では、権限の高低と信頼の高低について述べるが、両者は非常に密接に関連していて、きれいに切り分けるのが難しいため、混乱を招きやすい。
信頼と権限の本質的な性質から、両者はほぼ例外なく相関する。信頼が高いところでは通常、権限も高く、その逆も同様である。
本書の範囲を超える話だが、これらの表現（信頼と権限）は互換的に使われることが多い。他人を正そうとせず、自分にとって最も筋が通るように寛大に解釈するのが、たいていは最善である。
\end{rem}

\begin{exa}[SSH デーモンの信頼境界]\label{exa:ch2-sshd}
SSH セキュアシェルデーモン（sshd(8)）は、信頼境界を伴うセキュアな設計の優れた例である。
SSH プロトコルを使うと、許可されたユーザーはホストにリモートでログインし、インターネット越しのセキュアなネットワーク通信路を介してシェルを実行できる。
しかし、プロトコルを開始するための接続を常時待ち受ける SSH デーモンは、信頼境界をまたぐため、非常に注意深い設計を必要とする。
待ち受けプロセスには通常スーパーユーザー権限が必要である。許可されたユーザーが正当な認証情報を提示したとき、どのユーザーのプロセスでも作成できなければならないからである。
しかし同時に、公衆インターネット上で待ち受けなければならず、世界中からの攻撃にさらされる。

SSH のログイン要求を受け付けるために、デーモンはまず、盗聴も改ざんも受け付けないセキュアな通信路を確立し、次に機密性の高い認証情報を処理・検証しなければならない。
そこまで済んで初めて、ホストコンピュータ上に適切な権限でシェルプロセスを生成できる。
このプロセス全体には大量のコードが関わる。そのコードは（どのユーザーアカウントのプロセスでも作成できるように）最高レベルの権限で動作し、完璧に動作しなければシステムが深刻に侵害される危険がある。
入ってくる要求はインターネット上のどこからでも来る可能性があり、最初は攻撃と見分けがつかない。これほど魅力的で、しかも失敗したときの代償が大きい標的はなかなか想像できない。

攻撃対象領域の大きさと、脆弱性があった場合の深刻さを考えると、デーモンプロセスについては、大がかりな取り組みでリスクを緩和することが正当化される。
図 \ref{fig:ch2-sshd} は、この重要な信頼境界を保護するためにデーモンがどのように設計されているかを単純化して示したものである。

上から順に見ていくと、接続が入ってくるたびに低権限の子プロセスがフォークされ、その子プロセスがソケットを待ち受けて親（スーパーユーザー）プロセスと通信する。
子プロセスは、プロトコルの複雑なセキュアな通信路の暗号化も確立し、ログインの認証情報を受け付けて特権を持つ親プロセスに渡す。親プロセスは、入ってきた要求を信頼してシェルを与えるかどうかを判断する。
要求ごとに新しい子プロセスをフォークすることは、信頼境界における戦略的な保護になる。作業をできるだけ隔離でき、意図しない副作用がメインのデーモンプロセス内に積み重なるリスクも最小限に抑えられるからである。
ユーザーがログインに成功すると、デーモンは認証されたユーザーアカウントの権限で新しいシェルプロセスを作成する。
ログインの試行が認証に失敗すると、その要求を処理した子プロセスは終了するので、後でシステムに悪影響を及ぼすことはない。
\end{exa}

% 図 2-1: How the design of the SSH daemon protects critical trust boundaries
% 注: 原文の図は画像のみ。上の本文（子プロセスのフォーク、親との通信、認証成功時のシェル生成）から決まる範囲だけを描いた。
\begin{figure}[H]
\centering
\begin{tikzpicture}[
  box/.style={draw, rounded corners, align=center, minimum width=3.4cm, minimum height=1.0cm, font=\small},
  tr/.style={-{Stealth[length=2.5mm]}, thick},
  lb/.style={font=\scriptsize, fill=white, inner sep=1pt, align=center}
]
\node[box] (net) at (0,0) {インターネットからの\\接続要求};
\node[box] (child) at (0,-2.2) {子プロセス（低権限）\\セキュアな通信路の確立・\\認証情報の受け付け};
\node[box] (parent) at (6,-2.2) {親プロセス\\（スーパーユーザー）\\信頼するかを判断};
\node[box] (shell) at (6,-4.6) {シェルプロセス\\（認証されたユーザーの権限）};
\node[box] (term) at (0,-4.6) {子プロセスの終了};
\draw[tr] (net) -- node[lb, right] {接続ごとにフォーク} (child);
\draw[tr] (child) -- node[lb, above] {認証情報} (parent);
\draw[tr] (parent) -- node[lb, right] {認証成功} (shell);
\draw[tr] (child) -- node[lb, right] {認証失敗} (term);
\draw[dashed] (-2.6,-1.1) -- (8.6,-1.1);
\node[font=\scriptsize, anchor=west] at (2.3,-0.85) {信頼境界};
\end{tikzpicture}
\caption{SSH デーモンの設計が重要な信頼境界を保護する仕組み（単純化）}\label{fig:ch2-sshd}
\end{figure}

資産の場合と同じく、信頼レベルについても、どこでひとまとめにし、どこで分けるかを決めることになる。
オペレーティングシステムでは、もちろんスーパーユーザーの信頼レベルが最も高く、ほかの管理ユーザーの中にも、まったく同等の権限を持つとみなすべきほどスーパーユーザーに近い者がいるかもしれない。
信頼の序列でその次に来るのは、通常、許可されたユーザーである。
特別な権限を持つ、より信頼されたグループを形成するユーザーもいるかもしれないが、通常は、その中で誰をどれだけ多めに、あるいは少なめに信頼するかまで決める必要はない。
ゲストアカウントは通常、信頼が最も低い。ゲストの資源を保護することよりも、ゲストからシステムを保護することに重点を置くべきであろう。

\begin{exa}[Web サービスにおける信頼レベル]\label{exa:ch2-web-trust-levels}
Web サービスは悪意あるクライアントユーザーに耐えなければならないので、Web のフロントエンドシステムは、入ってくるトラフィックを検証し、整った形式のサービス要求だけを転送することがある。つまり、事実上インターネットとの信頼境界の上に立つ役割を担う。
Web サーバーは、ファイアウォールの背後にある、より信頼されたデータベースやマイクロサービスに接続することが多い。
お金が関わる場合（例えばクレジットカード処理サービス）には、支払いは専用の高信頼のシステムで扱うべきであり、理想的にはそのシステムをデータセンター内の囲い込まれた領域に隔離すべきである。
認証されたユーザーは、自分のアカウントデータへのアクセスについては信頼してよいが、それ以外についてはまったく信頼できないものとして扱うべきである。通常、ログインは誰でも作れるからである。
匿名の公開 Web アクセスはさらに信頼レベルが低い。静的な公開コンテンツなら、プライベートなデータサービスにまったく接続されていないマシンから配信することもできる。
\end{exa}

信頼境界を越える遷移は、常に明確に定義されたインターフェースとプロトコルを通じて行うこと。
これは、国境や入国港にある、武装した警備員が配置された検問所のようなものと考えればよい。
国境管理官がパスポートの提示を求め（認証の一形態）、所持品を検査する（入力検証の一形態）のと同じように、信頼境界は潜在的な攻撃を緩和する絶好の機会として扱うべきである。

最大のリスクは、当然ながら、SSH の待ち受けプロセスの例のような、低い信頼から高い信頼への遷移に潜んでいることが多い。
しかし、だからといって高い信頼から低い信頼への遷移を無視してよいわけではない。
システムが信頼度の低いコンポーネントにデータを渡すときは、その都度、情報を開示していないか、開示しているならそれが問題にならないかを検討する価値がある。
例えば、低権限のプロセスでさえ、自分が動作しているコンピュータのホスト名を読み取れる。攻撃者が足がかりを得てシステム上でコードを実行できた場合にヒントを与えかねないので、機密情報を使ってマシンに名前を付けてはならない。
さらに、高信頼のサービスが低信頼の要求に代わって作業するときには、ユーザーランドの要求者がカーネルに過大な負荷をかけることに成功すれば、サービス拒否攻撃を受けるリスクが生じる。

%----------------------------------------------------------------------
\subsubsection{脅威を特定する}\label{subsec:ch2-identify-threats}
% src: Identify Threats

ここから、脅威モデリングの核心となる作業、つまり潜在的な脅威の特定に取りかかる。
モデルに基づいて、システムの各部分をじっくり調べる。
脅威は資産の周りや信頼境界に集まりやすいが、どこに潜んでいてもおかしくない。

まず大まかに調べ（例えばシステムを1万フィート上空から眺めるような視点）、後で実りが多そうな部分や興味深い部分に戻って、より念入りに調べる（1000フィートからの視点）ことを勧める。
先入観を持たず、エクスプロイトの具体的な方法がまだわからなくても、可能性は必ず含めること。

資産に対する直接的な脅威や、信頼境界での脅威（攻撃者が信頼されたコンポーネントを簡単にだまして言いなりにさせうる）を特定するのは容易なはずである。
具体的な状況でのそうした脅威の例は、本書全体を通じて数多く示す。
しかし、間接的な脅威が見つかることもある。直ちに害を与えられる資産も、越えるべき信頼境界も見当たらないような脅威である。
そのような脅威を、一連の出来事の一部としてどう機能しうるかを考えないまま、すぐに無視してはならない。ビリヤードのクッションショットや、道を形作る飛び石のようなものと考えるとよい。
害を与えるために、攻撃者は複数の間接的な脅威を組み合わせなければならないかもしれない。あるいは、バグや設計の悪い機能と組み合わさることで、間接的な脅威が攻撃者に足がかりを与える糸口になるかもしれない。
どれほど有望に見えるか、危険にさらされる資産がどれほど重要かによっては、小さな脅威でさえ緩和する価値があるかもしれない。

%----------------------------------------------------------------------
\subsubsection{銀行の金庫の例}\label{subsec:ch2-bank-vault}
% src: A Bank Vault Example
% 注: 原典ではこの項と次の2項は「Identify Threats」の下位見出しと思われるが、分量を考慮して \subsubsection とした。

ここまでの概念はまだかなり抽象的に感じられるかもしれないので、架空の銀行の金庫を脅威モデリングして、具体的な文脈の中で見てみよう。
このウォークスルーを読む間は概念に集中してほしい。注意深く読んでいれば、著者が挙げる論点（意図的に網羅していない）をさらに発展させられるはずである。

% 注: この例は 1 ページを超えて itembox（ページをまたげない）からはみ出すため、exa 環境に入れず地の文にした。
\underline{モデルと資産：}
地元の銀行の支店を思い浮かべてほしい。
古い建物で、正面の重厚な無垢のオーク材の両開きの扉を、堂々としたローマ風の円柱が縁取っているとしよう。
人件費や資材が安かった時代に建てられたので、厚い鉄筋コンクリートの壁は破りようがないように見える。
この例では、銀行の建物の中心にあるセキュアな金庫に保管された大量の金だけに焦点を当てる。これが保護したい主要な資産である。
モデルには建物の建築図面を使い、建物全体の配置がひと目でわかる、1インチ10フィートの縮尺の平面図をもとに作業する。

\underline{信頼境界と攻撃対象領域：}
主要な信頼境界は明らかに金庫の扉だが、カウンターの奥にある従業員専用区域の施錠された扉がもう一つ、そして顧客ロビーと外部を隔てる銀行の正面扉が3つ目の信頼境界である。
簡単のため、裏口はモデルから省く。裏口は常に非常に厳重に施錠されており、警備員がいるときにまれに開けられるだけだからである。
すると、重要な攻撃対象領域として残るのは、正面扉と、誰でも簡単に立ち入れる顧客ロビーの区域だけになる。

\underline{脅威の特定（1回目）：}
これで、潜在的な脅威を見つけるという本番の作業の舞台が整った。
金を盗まれることが最大の脅威なのは明らかだが、それではあまりに漠然としていて、どう防げばよいかの手がかりがほとんど得られないので、さらに具体的な脅威を探す。
金を盗むには、攻撃者は金庫に不正にアクセスする必要がある。
そのためには、金庫がある従業員専用区域に不正に立ち入る必要がある。
今のところ、こうした抽象的な脅威がどのように起こりうるかはわからないが、分解すればより具体的にできる。
潜在的な脅威をいくつか挙げる。
\begin{itemize}
\item 金庫の暗証番号をひそかに盗み見る。
\item 金庫の暗証番号を推測する。
\item 化粧とかつらで銀行の頭取になりすます。
\end{itemize}
確かに、これらの作り話の脅威はかなりばかげている。しかし、モデルから脅威を導き出したことと、抽象的な脅威から具体的な脅威へと移っていったことに注目してほしい。

\underline{脅威の特定（2回目）：}
より詳細な2回目の作業では、完全な建築図面、電気と配管の配置、金庫の設計仕様を含むモデルを使う。
詳しい情報があれば、具体的な攻撃は簡単に思い描ける。
先ほど挙げた最初の脅威、つまり攻撃者が金庫の暗証番号を知るという脅威を取り上げよう。
この脅威はいくつかの形で起こりうる。そのうちの3つを見てみよう。
\begin{itemize}
\item 目ざとい強盗がロビーをうろついて、金庫が開けられるところを観察する。
\item 金庫の暗証番号が付箋に書かれていて、カウンターにいる客から見える。
\item 通りの向かいにいる共犯者が、望遠鏡で金庫の暗証番号ダイヤルをのぞく。
\end{itemize}
当然ながら、金庫の暗証番号を知っただけでは、侵入者は金を手に入れられない。
部外者に暗証番号を知られるのは重大な脅威だが、完全な攻撃の一段階にすぎない。完全な攻撃には、従業員専用区域への侵入、金庫への侵入、そして金を持っての逃走が含まれていなければならない。

\underline{緩和策：}
ここまで来れば、列挙した脅威に優先順位を付け、緩和策を提案できる。
特定した潜在的な攻撃のそれぞれに対する、素直な緩和策を挙げる。
\begin{itemize}
\item ロビーをうろつく者：金庫の前に不透明な衝立を置く。
\item 付箋による漏えい：保護されていない書面の写しを禁止する方針を設ける。
\item 望遠鏡によるスパイ：不透明なすりガラスの窓を設置する。
\end{itemize}
これらは考えうる多くの防御的な緩和策のほんの一部である。
建物の設計時にこの種の攻撃が考慮されていれば、配置を工夫することで、こうした脅威の一部をそもそもなくせたかもしれない（例えば、外部のどの窓からも金庫の区域が直接見通せないようにしておけば、後から不透明なガラスを取り付ける必要はなかった）。

本物の銀行のセキュリティや金融リスク管理はもちろんはるかに複雑だが、この単純化した例から、脅威モデリングのプロセスがどう進むのか、そして分析がどう前に進んでいくのかがわかる。
金庫の中の金はこれ以上ないほど単純な資産である。ここで、複雑なソフトウェアシステムのモデルを具体的にどう調べれば、そのシステムが直面する脅威が見えてくるのか、という疑問がわいてくるはずである。

%----------------------------------------------------------------------
\subsubsection{STRIDE による脅威の分類}\label{subsec:ch2-stride}
% src: Categorizing Threats with STRIDE

1990年代後半、パーソナルコンピューティングの世界は Microsoft Windows が支配していた。
PC が企業でも家庭でも欠かせない道具になるにつれ、同社の売上は際限なく伸びると多くの人が考えていた。
しかし Microsoft は、ネットワークがどう機能すべきかを理解し始めたばかりであった。
インターネット（当時はまだ頭文字を大文字の I で綴るのが普通だった）と World Wide Web という新しいものが急速に人気を集めており、Microsoft の Internet Explorer Web ブラウザは、先駆者である Netscape Navigator から積極的に市場シェアを奪っていた。
こうして同社は、セキュリティという新たな問題に直面した。世界中のコンピュータをつなげばどんな厄介事が持ち上がるのか、誰にもわからなかった。

Microsoft のテスターのチームは創意工夫を凝らしてセキュリティ上の欠陥を探していたが、世間の人々のほうがはるかに速く欠陥を見つけているようであった。
ネットワーク越しに顧客を危険にさらす脆弱性にパッチを出すという後手の対応を数年続けた後、同社は先手を打つためにタスクフォースを結成した。
この取り組みの一環として、著者は Praerit Garg と共同で、開発者が自分の製品のセキュリティ上の欠陥を見つけられるようにするための簡単な方法論を説明した論文を執筆した。
STRIDE という脅威の分類体系に基づく脅威モデリングは、同社の全製品グループにわたる大規模な教育活動の原動力となった。
20年以上たった今も、業界中の研究者が、脅威を列挙するために STRIDE や、それから独立に派生した多くの手法を使い続けている。

\begin{dfn}[STRIDE]\label{dfn:ch2-stride}
検討すべき具体的な種類の脅威のチェックリストを示すことで脅威の特定プロセスに焦点を与える、ソフトウェアに対する脅威の分類体系（taxonomy）を \underline{STRIDE} という。
STRIDE は次の問いからなる。何が\underline{なりすまし}（Spoofing, S）され、\underline{改ざん}（Tampering, T）され、\underline{否認}（Repudiation, R）されうるか。
どんな情報が開示（\underline{情報漏えい}、Information disclosure, I）されうるか。
\underline{サービス拒否}（Denial of service, D）や\underline{権限昇格}（Elevation of privilege, E）はどのようにして起こりうるか。
\end{dfn}

これらのカテゴリーは、分析の焦点を絞れるほど具体的でありながら、特定の設計に関係する細部を頭の中で肉付けし、そこから掘り下げていけるほど一般的でもある。

\begin{rem}[STRIDE は方法論ではない]\label{rem:ch2-stride-not-methodology}
セキュリティコミュニティの人々は STRIDE を脅威モデリングの方法論と呼ぶことが多いが、これは（少なくとも、この頭字語を考案した著者の考えでは）用語の誤用である。
STRIDE は単なるソフトウェアに対する脅威の分類体系である。
この頭字語は、脅威のカテゴリーを見落としていないことを確かめるための、簡単で覚えやすい記憶術になる。
STRIDE は完全な脅威モデリングの方法論ではない。完全な方法論には、この章ですでに見てきたほかの多くの構成要素が含まれていなければならない。
\end{rem}

STRIDE の使い方を見るために、なりすましから始めよう。
モデルをコンポーネントごとに調べ、セキュアな動作がユーザー（あるいはマシン、コードのデジタル署名など）の身元にどのように依存しているかを考える。
ここで身元をなりすませたら、攻撃者はどんな利点を得られるか。
こう考えていくと、たぐってみるべき糸口がたくさん見つかるはずである。
モデルの文脈の中で各コンポーネントに脅威の観点から向き合うと、本来どう動作すべきかという考えを脇に置きやすくなり、どう悪用されうるかが見えてくる。

\begin{rem}[ホワイトボードによるウォームアップ]\label{rem:ch2-stride-whiteboard}
著者が何度もうまく使ってきた優れた手法がある。脅威モデリングのセッションを、6つの脅威の名前をホワイトボードに書き出すところから始めるのである。
弾みをつけるために、細部を掘り下げる前に、これらの抽象的な脅威をいくつかブレインストーミングする。
「ブレインストーミング」という言葉にはさまざまな意味がありうるが、ここでは、考えすぎず、アイデアの良し悪しもまだ判断せずに（駄目なものは後で飛ばせばよい）、素早く広い範囲をカバーすることを指す。
このウォームアップによって、何に注意すべきかを意識する準備ができ、必要な考え方に頭を切り替えやすくもなる。
これらの脅威のカテゴリーをよく知っていても、すべてを一通り検討する価値はある。また、あまりなじみがなく技術的な2、3のカテゴリーは、注意深く説明しておく価値がある。
\end{rem}

表 \ref{tab:ch2-stride} に、6つのセキュリティ目標と、それぞれに対応する脅威のカテゴリー、各カテゴリーの脅威の例をいくつか挙げる。
セキュリティ目標と脅威のカテゴリーは表裏一体であり、防御（目標）と攻撃（脅威）のどちらか一方から考えたほうが取り組みやすい場合もある。

\begin{table}[H]
\centering
\caption{STRIDE の脅威カテゴリーのまとめ}\label{tab:ch2-stride}
\small
\begin{tabularx}{\linewidth}{|l|l|X|}
\hline
目標 & STRIDE の脅威 & 例 \\
\hline
真正性（Authenticity） & なりすまし（Spoofing） & フィッシング、盗まれたパスワード、なりすまし（impersonation）、メッセージのリプレイ、BGP ハイジャック \\
\hline
完全性（Integrity） & 改ざん（Tampering） & 不正なデータの改変と削除、Superfish の広告挿入 \\
\hline
否認不可能性（Non-repudiability） & 否認（Repudiation） & もっともらしい否認、不十分なログ記録、ログの破壊 \\
\hline
機密性（Confidentiality） & 情報漏えい（Information disclosure） & 漏えい、サイドチャネル、弱い暗号化、キャッシュに残されたデータ、Spectre/Meltdown \\
\hline
可用性（Availability） & サービス拒否（Denial of service） & 大量の同時要求による Web サーバーのパンク、ランサムウェア、MemCrashed \\
\hline
認可（Authorization） & 権限昇格（Elevation of privilege） & SQL インジェクション、xkcd の「Little Bobby Tables」 \\
\hline
\end{tabularx}
\end{table}

STRIDE の顔ぶれのうち半分は、第1章で学んだ情報セキュリティの基本を直接脅かすものである。
情報漏えいは機密性の敵、改ざんは完全性の敵であり、サービス拒否は可用性を損なう。
STRIDE の残りの半分はゴールドスタンダードを標的にする。
なりすましは、偽の身元を名乗ることで真正性を破る。
権限昇格は、適切な認可を破る。
残る否認は監査に対する脅威である。ただ、これはすぐにはぴんと来ないかもしれないので、詳しく見ておく価値がある。

ゴールドスタンダードに従えば、システム内で行われた重要な操作を正確に記録しておき、その操作を監査すべきである。
\underline{否認}（repudiation）は、誰かがある操作を行ったことを、もっともらしく否定するときに起こる。

\begin{exa}[否認の実際]\label{exa:ch2-repudiation-in-practice}
ソフトウェアセキュリティに携わってきた長年の間、著者は誰かが何かを面と向かって否認する場面を見たことがない（著者の前で「やっただろう！」「やってない！」と怒鳴り合った人は一人もいない）。
実際に起こるのは、例えば、データベースが突然消え、何もログに記録されていなかったために誰にも理由がわからず、失われたデータが跡形もなく消えている、といったことである。
組織は侵入があったのではないかと疑うかもしれない。
あるいは不正を働く内部者の仕業かもしれないし、管理者の痛恨のミスかもしれない。
しかし証拠がなければ、誰にもわからない。
これは大きな問題である。インシデントの後で何が起こったのかを説明できなければ、再発を防ぐのは非常に難しいからである。
物理的な世界では、こうした完全犯罪はまれである。銀行強盗のような行為にはその場にいることが必要で、そこには本質的にさまざまな痕跡が残るからである。
ソフトウェアは違う。証拠を確実に収集し、イベントをログに記録する手段を用意しておかない限り、証拠となる指紋も泥だらけの靴跡も残らない。
\end{exa}

否認の脅威を緩和するには、通常、正確な監査証跡があることを管理者とユーザーが知っていて、自分の行為に責任を負うと理解しているシステムを運用する。
管理者のパスワードを全員で共有する付箋に書いておくべきでない理由が、ここにもう一つある。
そうしていると、問題が起きたときに、全員が「誰かほかの人がやったに違いない」ともっともらしく主張できてしまう。
これは全員を完全に信頼している場合にも当てはまる。事故は起こるものであり、問題が生じたときに使える証拠が多いほど、復旧も是正も容易になるからである。

%----------------------------------------------------------------------
\subsubsection{映画の中の STRIDE}\label{subsec:ch2-stride-movies}
% src: STRIDE at the Movies

息抜きとして（そしてこれらの概念を定着させるために）、映画『オーシャンズ11』（Ocean's Eleven）の筋書きに STRIDE の脅威を当てはめてみよう。
この古典的な強盗映画は、攻撃者と防御者の両方の視点から、STRIDE のすべてのカテゴリーを含む脅威モデリングの概念をうまく示している。
簡潔さと焦点を優先して筋書きを単純化したこと、そしてネタバレについてはご容赦いただきたい。

\begin{exa}[『オーシャンズ11』と STRIDE]\label{exa:ch2-oceans-eleven}
ダニー・オーシャンは仮釈放の条件を破り（権限昇格）、昔の犯罪仲間に会いに飛行機で飛び、ラスベガスへ向かう。
彼は裕福なカジノの内部関係者に大胆な強盗計画を持ちかける。その関係者からカジノの運営の詳細を教わり（情報漏えい）、ダニーは元受刑者の仲間を集める。
一味は、練習用に作った実物大の金庫のレプリカを使って作戦を練る。
決行の夜、ダニーはカジノに姿を現し、予想どおり警備員に捕まって、完璧なアリバイを手にする（罪の否認）。
やがて彼は通気ダクトから抜け出し、さまざまな策略を駆使して、仲間とともに金庫からお金の半分を抜き取り（金庫の完全性に対する改ざん）、リモコン操作のバンで運び出す。

一味は、金庫に残った数百万ドルを爆破すると脅し（非常に高くつくサービス拒否）、バンの中のお金を自分たちのものにするよう交渉する。
カジノのオーナーは拒否して SWAT チームを呼び、続く混乱の中で一味は金庫の中身を破壊して逃げる。
騒ぎが収まった後、カジノのオーナーは金庫を確かめて全損を嘆くが、ふと細かな点がおかしいことに気づく。
オーナーは、まるで一度も抜け出していなかったかのように留置場に戻っているダニーを問い詰める。そこで観客は、SWAT チームが実は一味だったこと（警察を装うなりすまし）、偽の銃撃戦の後、戦術装備のバッグにお金を隠して歩いて出ていったことを知る。
練習用の金庫の模型は、本物の金庫が破られたように見せかけるための映像を提供していた（場所のなりすまし）。本物の金庫が実際に破られたのは、カジノが偽の SWAT チームに完全なアクセスを許したとき（一味にとっての権限昇格）であった。
ダニーはヒロインを取り戻し、全員がお金を持ってまんまと逃げおおせる。犯人たちにとってはハッピーエンドだが、カジノが脅威モデリングのコンサルタントを雇っていたら、まったく違う結末になっていたかもしれない。
\end{exa}

%----------------------------------------------------------------------
\subsubsection{脅威を緩和する}\label{subsec:ch2-mitigate-threats}
% src: Mitigate Threats

この段階で、潜在的な脅威がひととおり手元に集まっているはずである。
次は、効果的な防御につなげるために、それらの脅威を評価し、優先順位を付ける必要がある。
脅威はせいぜい将来の出来事についての経験に基づく推測なので、どんな評価にもある程度の主観が入る。

脅威を理解するとは、正確にはどういうことだろうか。
この問いに簡単な答えはないが、知っていることを磨き上げること、そしてすべてわかったつもりになる罠に陥らないよう健全な懐疑心を保つことが含まれる。
実際には、ほとんどが抽象的な脅威を素早く洗い出して集め、それぞれをもう少し突ついて詳しく知る、ということになる。
すると、かなりはっきりした攻撃が1つか2つ、あるいは攻撃を構成しうる要素の一部が見えてくるかもしれない。
収穫逓減の壁に突き当たるまで、検討を深めていく。

\begin{dfn}[脅威への対処の4つの方法]\label{dfn:ch2-four-ways-to-handle-threats}
特定した脅威には、次の4つの方法のいずれかで対処できる。
\begin{enumerate}
\item \underline{緩和}（mitigate）：再設計するか防御を追加することで、リスクが現実になる頻度を減らすか、害の程度を許容できる水準まで下げる。
\item \underline{除去}（remove）：脅かされている資産が必要でなければ、資産そのものを除去する。除去できない場合は、資産の露出を減らすか、脅威を増大させるオプション機能を制限することを目指す。
\item \underline{移転}（transfer）：通常は対価と引き換えに、責任を第三者に肩代わりさせることでリスクを移転する（例えば保険はリスク移転の一般的な形である。また、機密性の高いデータの処理を、機密性を保護する義務を負うサービスに外部委託することもできる）。
\item \underline{受容}（accept）：リスクを十分に理解したうえで、負うのが妥当なものとして受け入れる。
\end{enumerate}
\end{dfn}

重大な脅威については常に緩和を試みるべきだが、結果は一長一短になることが多いと心得ておくこと。
実際には、さまざまな理由から、最善の解決策がいつも実行できるとは限らない。大きな変更はコストがかかりすぎるかもしれないし、自分では制御できない外部の依存関係を使わざるをえないかもしれない。
ほかのコードが脆弱な機能に依存していて、修正すると何かが壊れるかもしれない。
このような場合、緩和とは、脅威を減らすことなら何でもすることを意味する。
防御側にとっては、どんな優位性でも、たとえ小さなものでも役に立つ。

\begin{dfn}[部分的な緩和]\label{dfn:ch2-partial-mitigation}
部分的な緩和の方法には次のものがある。
\begin{itemize}
\item \emph{害が起こる可能性を低くする}：例えば、攻撃が10\% の確率でしか成功しないようにする。
\item \emph{害の深刻さを小さくする}：例えば、データのごく一部しか破壊できないようにする。
\item \emph{害を元に戻せるようにする}：例えば、失われたデータをバックアップから簡単に復元できるようにしておく。
\item \emph{害が起こったことを明らかにする}：改変された製品を簡単に見分けられる開封明示包装（tamper-evident packaging）を使って、消費者を保護する（ソフトウェアでは、きちんとしたログ記録がこの役割を果たす）。
\end{itemize}
\end{dfn}

本書の残りの大部分は緩和、つまり脅威を最小限に抑えるようにソフトウェアを設計する方法と、さまざまな種類の緩和策を考え出すのに役立つ戦略やセキュアなソフトウェアのパターンを扱う。

%----------------------------------------------------------------------
\subsection{プライバシーへの配慮}\label{subsec:ch2-privacy}
% src: Privacy Considerations

プライバシーへの脅威はセキュリティへの脅威と同じくらい現実的なものであり、システムに対する脅威を漏れなく評価するには、別途検討する必要がある。プライバシーへの脅威は、情報漏えいのリスクに人間的な要素を加えるからである。
規制や法律上の考慮事項がありうるだけでなく、個人情報の取り扱いには倫理的な問題が絡むこともあり、利害関係者の期待に応えることが重要である。

どんな種類であれ個人データを収集しているなら、基本姿勢としてプライバシーを真剣に受け止めるべきである。
自分は人々の私的な情報を預かる管理人（steward）だと考えること。
ユーザーが抱きうる幅広いプライバシーの懸念を注意深く考慮するなど、常にユーザーの視点を忘れないよう努め、迷ったら慎重なほうを選ぶこと。
ソフトウェアの作り手は、システム構築の論理に没頭していると、個人データがどれほど機密性の高いものになりうるかを軽く見がちである。
コードの中ではデータベーススキーマのフィールドの一つにすぎないように見えるものが、漏えいすれば実在の人物に現実の影響を及ぼす情報であることもある。
現代の生活のデジタル化が進み、モバイルコンピューティングがどこにでも広がるにつれて、プライバシーはますますコードに左右されるようになり、それも想像しにくい新しい形をとるかもしれない。
要するに、今のうちから細心の注意を払い、時代のはるか先を行っておくのが賢明である。

\begin{rem}[プライバシーへの脅威を最小化するための一般的な考慮事項]\label{rem:ch2-privacy-considerations}
プライバシーへの脅威を最小限に抑えるための、ごく一般的な考慮事項には次のものがある。
\begin{itemize}
\item 抽象的に考えるのではなく、実際のユースケースのシナリオをモデル化してプライバシーを評価する。
\item どのようなプライバシーポリシーや法的要件が適用されるかを調べ、その条件に厳格に従う。
\item データの収集を必要なものだけに制限する。
\item 気味悪がられるおそれに敏感になる。
\item 使用目的がはっきりしない私的な情報は決して収集・保存しない。
\item 収集済みの情報が使われなくなったり役に立たなくなったりしたら、積極的に削除する。
\item 第三者との情報共有を最小限にする（共有する場合は、きちんと文書化すべきである）。
\item 機密情報の開示を最小限にする。理想的には、知る必要がある場合（need-to-know）に限るべきである。
\item 透明性を保ち、データ保護の取り組みをエンドユーザーが理解できるよう手助けする。
\end{itemize}
\end{rem}

%----------------------------------------------------------------------
\subsection{どこにでもある脅威モデリング}\label{subsec:ch2-threat-modeling-everywhere}
% src: Threat Modeling Everywhere

ここで述べた脅威モデリングのプロセスは、私たちが世の中をどう渡り歩いているかを形式化したものである。私たちは、リスクを機会と天秤にかけることでリスクを管理している。
危険な環境では、あらゆる生物がこれと同じ基本原則に基づいて判断を下している。
いったん探し始めれば、脅威モデリングの実例はどこにでも見つかる。

\begin{exa}[子どもの安全対策]\label{exa:ch2-child-proofing}
幼い子どもを連れた友人が訪ねてくるとき、著者はいつも数分かけて特別な準備をする。
活発な3歳児のアレックスは好奇心旺盛なので、家の中を「チャイルドプルーフ（子どもの安全対策）」して回るのである。
これはまさに脅威モデリングである。脅威をカテゴリーごとに想像し---何がアレックスを傷つけうるか、何が壊されそうか、何を幼い子の目に触れないようにしておくほうがよいか---次にそのパターンに当てはまる資産を探すのだから。
典型的な脅威としては、アレックスが壁のコンセントに差し込みかねない鋭いペーパーナイフ、簡単に壊してしまいそうな壊れやすいアンティークの花瓶、あるいは子どもにふさわしくない画像を含む写真集などがある。
攻撃対象領域は、活発な幼児の手が届くあらゆる場所である。
緩和策は一般に、露出点や脆弱な点を取り除く、減らす、あるいはなくすことである。例えば、壊れやすい花瓶をドライフラワーだけを挿したプラスチックの花瓶に替えたり、暖炉の棚の上に移したりできる。
子どもがいる人なら、子どもが何をしでかすかを予想するのがどれほど難しいかを知っている。
例えば、手が届かないと思っていた棚に、アレックスが本を積み上げてよじ登り、届いてしまうことまで予想できていただろうか。
これがソフトウェアの外での脅威モデリングの姿であり、先手を打った緩和が十分に労力に見合う理由もよくわかる。
\end{exa}

\begin{exa}[日常生活で見かける脅威モデリング]\label{exa:ch2-everyday-threat-modeling}
日常生活で目にしたことがあるかもしれない、脅威モデリングのほかの例をいくつか挙げる。
\begin{itemize}
\item 店は、万引きした商品を返品して店のクレジットを得たり、新しい服を一度着てから返品して返金を受けたりするような悪用を緩和するために、返品ポリシーを工夫している。
\item Web サイトの利用規約は、ユーザーがサイトを悪意をもって悪用するさまざまな手口を防ごうとしている。
\item 交通安全法、制限速度、運転免許、自動車保険の加入義務は、いずれも運転をより安全にするための緩和の仕組みである。
\item 図書館は、蔵書の盗難、抱え込み、損傷を緩和するように貸出ポリシーを定めている。
\end{itemize}
\end{exa}

あなた自身もこうした手法を使っている場面をたくさん思いつくであろう。
世界についての物理的な直観を頼りにできる場合、ほとんどの人にとって脅威モデリングは驚くほど簡単である。
ソフトウェアの脅威モデリングも、ほかの場面ですでに十分に磨かれている技能と同じように働くのだと気づけば、生まれ持った能力をソフトウェアのセキュリティ分析に生かし始め、すぐに技能を次の段階へ引き上げられる。

\newpage

%% =======================================================
%%  part: ch3
%% =======================================================
%% =======================================================
%%  part: ch3
%% =======================================================
\section{緩和策}
% src: 3: Mitigation

この章では、\cite{Kohnfelder2021} の第3章に沿って、脅威に対する緩和策を扱う。

\begin{quote}
「技とたゆまぬ努力をもってすれば、緩和できないものはない。」\\
\hfill --- Gaius Plinius Caecilius Secundus（Pliny the Younger）
\end{quote}

この章で取り上げるのは、第2章の 4つの質問（Four Questions）のうち 3 番目の「それに対して何をするのか？」（What are we going to do about it?）である。
脅威を予測し、潜在的な脆弱性に備えて守りを固めることで、セキュリティの考え方は実際に効果のある行動になる。

\begin{dfn}[緩和策]\label{dfn:ch3-mitigation}
問題の深刻度・範囲・影響を小さくするために先回りして講じる対応を、\underline{緩和策}（mitigation）という。
\end{dfn}

前章で見たように、緩和は誰もが日常的に行っていることである。

\begin{exa}[身の回りの緩和策]\label{exa:ch3-everyday-mitigations}
乳児に食事をさせるときに必ず出る食べこぼしを受け止めるよだれかけ、シートベルト、制限速度、火災報知器、食品衛生の習慣、公衆衛生上の対策、産業安全の規制などは、緩和策のほんの一例である。
いずれも、リスクを前にして予想される害を避けたり小さくしたりするために、前もって手を打っている点で共通している。
ソフトウェアをよりセキュアにするための取り組みの多くも、まさにこれである。
\end{exa}

\begin{rem}[緩和策はリスクを減らすが除去はしない]\label{rem:ch3-mitigation-not-elimination}
緩和策はリスクを\emph{軽減する}ものであって、\emph{除去する}ものではない。この点を心に留めておくことが重要である。
念のために言えば、何らかの方法でリスクそのものを取り除ける場合 --- たとえば、安全でないとわかっている古い機能を削除する場合 --- は、ぜひそうすべきである。ただし、著者はそれを緩和策とは呼ばない。
緩和策が目指すのは、攻撃を起こりにくくすること、難しくすること、そして実際に起きたときの害を小さくすることである。
エクスプロイトを検知しやすくする措置も、迅速な対応と修復につながるなら緩和策であり、開封されたことがわかる包装（tamper-evident packaging）と同じ発想である。
小さな努力の一つひとつがシステム全体のセキュリティを着実に引き上げ、ささやかな改善でも積み重なれば保護は大きく強化される。
\end{rem}

この章では、まず緩和策を概念的に論じ、そこから一般的な手法をいくつか紹介する。
中心となるのは、脅威モデリングで得られる視点に基づく\emph{構造的な}緩和策で、ほぼどのようなシステム設計をセキュアにする場合にも役立つ。
以降の章では、ここでの考え方を土台に、特定の技術や脅威に踏み込んだ、より詳しい方法を示す。

章の後半では、ソフトウェア設計で繰り返し直面するセキュリティ上の課題、すなわちアクセスポリシーとアクセス制御の整備、インターフェースの設計、通信とストレージの保護について指針を示す。
これらの議論は全体として、よくあるセキュリティ上のニーズに対処するための手引き（playbook）となり、本書の残りの部分で肉付けしていく。

%----------------------------------------------------------------------
\subsection{脅威への対処}\label{subsec:ch3-addressing-threats}
% src: Addressing Threats

脅威モデリングによって何がうまくいかなくなりうるかが明らかになり、セキュリティ上の注意を本当に重要な箇所に向けられるようになる。
とはいえ、脆弱性はいつでも除去できると考えるのは楽観的すぎる。
リスクが生じるポイント --- 重大なイベントや判断の分かれ目 --- こそ、緩和策を講じる絶好の機会である。

前章で学んだように、最大の脅威から先に対処し、できる限り抑え込むべきである。
そのうえで、最大のリスクがどこに残っているかを見極めてそこに狙いを定める、という作業を繰り返す。

\begin{exa}[機微な個人情報を扱うシステム]\label{exa:ch3-sensitive-personal-data}
機微な個人情報を処理するシステムでは、許可のない開示という脅威がどうしても大きくなる。
この大きなリスクに対しては、次のうちいずれか、あるいはすべてを検討するとよい。
\begin{itemize}
\item データへのアクセスを最小限にする。
\item 収集する情報の量を減らす。
\item 不要になった古いデータを積極的に削除する。
\item 侵害が起きたときに早く気付けるよう監査する。
\item 攻撃者がデータを持ち出す（exfiltrate）能力を弱める措置をとる。
\end{itemize}
\end{exa}

優先度の最も高いリスクに手を打ったら、次は設計にあまりオーバーヘッドや複雑さを加えずに済むところで、機会をとらえて小さなリスクも緩和していく。

\begin{exa}[ソルト付きハッシュによるパスワード照合]\label{exa:ch3-salted-hash}
賢い緩和策の好例として、ログインのたびに送られてきたパスワードを、実際のパスワードの平文ではなくソルト付きハッシュ（salted hash）と照合するというベストプラクティスがある。
パスワードが漏れると認証の根幹が揺らぐので、パスワードの保護は極めて重要である。
ハッシュどうしを比べる手間は元の値どうしを比べる場合とほとんど変わらないが、平文のパスワードを保存しなくて済むという大きな利点がある。
つまり、攻撃者が何らかの方法でシステムに侵入できたとしても、実際のパスワードを知ることはできない。
\end{exa}

この例は害の低減（harm reduction）という考え方をよく表しているが、パスワード照合にかなり特化した話である。
次に、もっと幅広く使える緩和戦略を考える。

%----------------------------------------------------------------------
\subsection{構造的な緩和戦略}\label{subsec:ch3-structural-mitigations}
% src: Structural Mitigation Strategies

緩和策は、多くの場合、常識的な話に行き着く。つまり、リスクを減らす機会があればそこで減らす、ということである。
脅威モデリングを使うと、潜在的な脆弱性を、攻撃対象領域（attack surface）、信頼境界（trust boundary）、資産（保護が必要な標的）という観点から捉えやすくなる。
\underline{構造的な緩和策}（structural mitigation）は、一般にモデルのこうした要素そのものに対して講じるものだが、具体的な実現方法は設計の詳細によって変わる。
以下の各項では、モデルの抽象度で働くため幅広く適用できるはずの手法を紹介する。

\subsubsection{攻撃対象領域の最小化}\label{subsec:ch3-minimize-attack-surfaces}
% src: Minimize Attack Surfaces

システムの攻撃対象領域を特定すれば、エクスプロイトがどこから来る可能性が最も高いかがわかる。したがって、システムの「外殻」を固めるための取り組みは、どんなものでも大きな効果がある。
攻撃対象領域の削減は、入口ごとにその下流でどれだけのコードやデータが関わるか、という観点で考えるとよい。
同じ機能を実行するインターフェースが複数あるシステムでは、それらを一つにまとめると、脆弱性を心配しなければならないコードが減るので有益な場合がある。
広く使われているこの手法の例をいくつか挙げる。

\begin{exa}[攻撃対象領域の削減]\label{exa:ch3-attack-surface-reduction}
\begin{itemize}
\item クライアント／サーバーシステムでは、機能をクライアント側に移すことで、サーバーの攻撃対象領域を減らせる。
サーバーへのリクエストを必要とする操作は、どれも攻撃対象領域を増やす。不正な形式のリクエストや偽造されたクレデンシャルに悪用されるおそれがあるからである。
これに対して、必要な情報と計算能力がクライアント側にそろっていれば、サーバーの負荷と攻撃対象領域の両方が減る。
\item 誰でも匿名で呼び出せる公開 API から認証付きの API へ機能を移すと、攻撃対象領域を効果的に減らせる。
アカウント作成という手間が加わるので攻撃のペースが落ち、攻撃者の追跡やレート制限もしやすくなる。
\item カーネルのサービスを使うライブラリやドライバは、カーネルへのインターフェースとカーネル内のコードを最小限にすることで、攻撃対象領域を減らせる。
そうすれば攻撃の対象となるカーネル遷移が減るうえ、たとえ攻撃が成功しても、ユーザーランドのコードではそれほど大きな害を及ぼせない。
\item デプロイと運用の場面にも、攻撃対象領域を減らす機会がたくさんある。
企業ネットワークでは、ファイアウォールの内側に移せるものはすべて移すのが手っ取り早い。
ネットワーク越しのリモート管理を有効にする設定はよい例である。便利な機能ではあるが、めったに使わないなら無効にしておき、必要なときは有線でアクセスすることを検討するとよい。
\end{itemize}
\end{exa}

これらは、攻撃対象領域の削減が効果を発揮する典型的な場面の一部にすぎない。
個々のシステムでは、はるかに創意工夫に富んだ、そのシステムならではの機会が見つかるかもしれない。
外部からのアクセスを減らし、機能とインターフェースを最小限にし、必要もなく露出しているサービスを保護する方法を、常に考え続けるべきである。
ある機能が実際にどこでどのように使われているかをよく理解しているほど、こうした緩和の機会を多く見つけられる。

\subsubsection{脆弱性の窓を狭める}\label{subsec:ch3-narrow-windows}
% src: Narrow Windows of Vulnerability

この緩和手法は攻撃対象領域の削減に似ているが、比喩としての表面積ではなく、脆弱性が生じうる実質的な\emph{時間}の幅を狭める。
これも常識に基づくもので、ハンターが撃つ直前にだけ安全装置を外し、撃った後すぐに掛け直すのも同じ理由からである。

この緩和策は、通常、信頼度の低いデータやリクエストが信頼度の高いコードとやり取りする信頼境界に適用する。
信頼度の高いコードをうまく隔離するには、そのコードが行わなければならない処理をできるだけ少なくする。
たとえば、可能であれば信頼度の高いコードを呼び出す前にエラーチェックを済ませておき、そのコードが自分の仕事だけをしてすぐに終われるようにする。

\begin{exa}[Code Access Security における権限の表明]\label{exa:ch3-cas}
\underline{Code Access Security}（CAS）は、今ではめったに使われないセキュリティモデルだが、コードの実効権限をきめ細かく制御できるので、この緩和策を説明するのにうってつけである。
（正直に言っておくと、著者は CAS を主要なセキュリティ機能として前面に押し出した .NET Framework バージョン 1.0 で、セキュリティ担当のプログラムマネージャーを務めていた。）

CAS のランタイムは、信頼の度合いに応じて、コードの単位ごとに異なる権限を付与する。
次の疑似コードは、汎用的な権限についてよく使われるイディオムである。ここでいう権限は、特定のファイルへのアクセスやクリップボードへのアクセスの許可などにあたる。
要するに CAS は、信頼度の高いコードが、自分を呼び出したコードの低い権限を引き継ぐようにする。ただし、必要なときには自分の高い権限を一時的に表明（assert）できる。
権限の表明は次のように動作する。
\begin{verbatim}
Worker(parameters) {
  // When invoked from a low-trust caller, privileges are reduced.
  DoSetup();
  permission.Assert();
  // Following assertion, the designated permission has been granted.
  DoWorkRequiringPrivilege();
  CodeAccessPermission.RevertAssert();
  // Reverting the assertion undoes its effect.
  DoCleanup();
}
\end{verbatim}
（コメントの意味は順に、「信頼度の低い呼び出し元から呼ばれると権限は縮小される」「表明の後は、指定した権限が付与されている」「表明を取り消すと、その効果は元に戻る」である。）

この例のコードは強力な権限を持っているが、もっと信頼度の低いコードから呼び出されることもある。
信頼度の低いコードから呼び出された場合、このコードは最初、呼び出し元の縮小された権限で動作する。
厳密に言えば、実効権限は、そのコード自身、その呼び出し元、呼び出し元の呼び出し元、……とスタックの最上位まで遡ったすべてのコードに付与された権限の共通部分（つまり最小値）である。
\texttt{Worker} メソッドの処理の一部には、呼び出し元が持っていないかもしれない高い権限が必要である。そこでセットアップを済ませた後、\texttt{DoWorkRequiringPrivilege}（これもその権限を持っていなければならない）を呼び出す前に、必要な権限を表明する。
その部分の仕事が終わったらすぐに \texttt{RevertAssert} を呼んで特別な権限を手放し、特別な権限のいらない残りの処理を済ませてから戻る。
CAS モデルでは、時間の窓を最小にするという考え方のもとで、こうした権限の表明を必要なときにだけ使い、不要になったらすぐに取り消せるようになっている。
\end{exa}

\begin{exa}[オンラインバンキングにおける端末ごとの権限レベル]\label{exa:ch3-mobile-banking}
脆弱性の窓を狭めるという考え方を、少し違う形で当てはめてみる。
オンラインバンキングは便利で速く、モバイル端末があればどこからでも銀行取引ができる。
しかし、銀行のクレデンシャルを携帯電話に保存しておくのは危険である。端末をなくしたときに誰かに口座を空にされてはたまらないし、モバイル端末は紛失がずっと起こりやすい。
著者が銀行業界全体で実装されてほしいと考える優れた緩和策は、端末ごとに、自分が安心できる権限レベルを設定できるようにすることである。
慎重な顧客なら、モバイルアプリでできることを残高照会と、1 日あたりの控えめな取引金額の上限までに制限するかもしれない。
そうすれば、顧客は安心して電話で銀行取引ができる。
ほかにも、利用できる時間帯、位置情報、国内通貨のみ、といった制限が役に立つだろう。
どの緩和策も、どんな種類の侵害が起きたとしても最悪の事態を限定できるという点で役に立つ。
\end{exa}

\subsubsection{データの露出の最小化}\label{subsec:ch3-minimize-data-exposure}
% src: Minimize Data Exposure

データ開示のリスクに対するもう一つの構造的な緩和策は、機微なデータがメモリ上に存在する期間を短くすることである。
前の手法とよく似ているが、ここで短くするのは、コードが高い権限で動作している時間ではなく、機微なデータにアクセスでき、露出するおそれのある時間である。
思い出してほしいが、プロセス内のアクセスは制御が難しいので、データはメモリ上にあるというだけで危険にさらされている。
このように重大なものがかかっている場合は、「メーターが回っている」と考えるとよい。
特に重要な情報 --- 暗号の秘密鍵のようなデータや、パスワードのような認証クレデンシャル --- は、不要になった時点でメモリ上のコピーをすぐに上書きしておく価値があるかもしれない。
そうすれば、何らかの経路で漏えいが起こりうる時間が短くなる。

\begin{exa}[Heartbleed とメモリ上のデータの保持期間]\label{exa:ch3-heartbleed}
第9章で見るように、Heartbleed 脆弱性は Web の大部分でセキュリティを根底から覆し、メモリ上に残っていたあらゆる種類の機微なデータを露出させた。
そうしたデータを保持する時間を短くしていれば、エクスプロイトを予見していなかったとしても、おそらく有効な緩和策（いわば「止血」）になっていただろう。
\end{exa}

この手法は、データストレージの設計にも応用できる。

\begin{exa}[アカウント削除後のデータの扱い]\label{exa:ch3-account-deletion}
ユーザーがシステム上で自分のアカウントを削除すると、通常はそのユーザーのデータも破棄される。ただし、誤って、あるいは悪意によってアカウントが閉鎖された場合に備えて、手動で復元できる仕組みを用意していることが多い。
これを実装する簡単な方法は、閉鎖されたアカウントに削除予定の印を付けてデータはたとえば 30 日間（手動復元の期間が過ぎるまで）残しておき、その後にシステムがすべてを最終的に削除する、というものである。
この方式では、ユーザーが破棄を指示したアカウントのデータにうっかりアクセスしないよう、多くのコードが、そのアカウントが削除予定かどうかを確認しなければならない。
たとえば一括メールのジョブがこの確認を忘れると、アカウントを閉鎖したユーザーに何かの通知を誤って送ってしまうかもしれない。ユーザーから見れば、自分の意思を踏みにじられたように感じるだろう。

この緩和策の考え方からは、もっとよい選択肢が見えてくる。ユーザーがアカウントを削除したら、システムはその内容をオフラインのバックアップに移し、データ本体はすぐに削除すればよい。
手動復元が必要になるまれなケースでもバックアップのデータを使って対応でき、しかもバグのせいでこの種の誤りが起こる余地はなくなる。
\end{exa}

\begin{rem}[データの積極的な消去が適切な場面]\label{rem:ch3-wiping-extreme}
一般に、データのコピーを積極的に消去するのは極端な措置であり、ふさわしいのは特に機微なデータや、アカウントの閉鎖のような重要な操作の場合に限られる。
これを自動的に行えるようにする言語やライブラリもある。性能が問題になる場合を除けば、メモリを再利用する前に中身を消去する簡単なラッパー関数を用意すれば済む。
\end{rem}

%----------------------------------------------------------------------
\subsection{アクセスポリシーとアクセス制御}\label{subsec:ch3-access-policy}
% src: Access Policy and Access Controls

標準的なオペレーティングシステムの権限機能が提供するのは、ごく初歩的なファイルアクセス制御だけである。
これは、プロセスを所有するユーザーとグループに基づいて、個々のファイルに対する読み取り（機密性）や書き込み（完全性）のアクセスを、許可するかしないかの二択で決めるものである。
この機能があると、資産やリソースの保護を設計するときにも同じ限られた枠組みで考えてしまいがちである。しかし、本当に適切なアクセスポリシーは、もっときめ細かく、ほかの多くの要因によって決まるかもしれない。

まず、従来のアクセス制御が現代の多くのシステムにどれほど合っていないかを考えてみる。
Web サービスやマイクロサービスは、プリンシパルに代わって動作するように設計されており、そのプリンシパルはたいていプロセスの所有者とは一致しない。
この場合、一つのプロセスが認証済みのリクエストをすべて処理するので、そのプロセスには常に全クライアントのデータにアクセスする権限が必要になる。
つまり、脆弱性が一つでもあれば、すべてのクライアントが危険にさらされるおそれがある。

効果的な\underline{アクセスポリシー}（access policy）を定めることは重要な緩和策である。本来許可すべきアクセスと、システムがたまたま備えているアクセス制御とのずれを埋められるからである。
利用できるオペレーティングシステムのアクセス制御から考え始めるのではなく、システムを通じて行動するさまざまなプリンシパルが何を必要としているかを考え抜き、どのようなアクセスが適切なのかを正確に記述した理想的なアクセスポリシーを定めるべきである。
きめ細かいアクセスポリシーなら、選択肢は豊富にある。
たとえば、1 分・1 時間・1 日あたりのアクセス回数の上限、扱えるデータ量の上限、勤務時間に合わせた時間帯の制限、同僚の利用状況や過去の頻度に応じて変わるアクセス上限などは、すぐに思いつく仕組みである。

安全なアクセス制限を決めるのは手間のかかる作業だが、アプリケーションのセキュリティ要件を理解する助けになるので、やるだけの価値はある。
ポリシーをコードで完全には実装しきれなくても、少なくとも効果的な監査の指針にはなる。
適切な制御の仕組みがそろっていれば、まずは緩い制限から始めて実際の利用状況を見極め、システムが実際にどうアクセスされているかがわかってくるにつれて、ポリシーを徐々に絞り込んでいける。

\begin{exa}[カスタマーサービス担当者のアクセス上限]\label{exa:ch3-customer-service}
カスタマーサービス担当者のチームが使う、架空のシステムを考える。
担当者は、問い合わせてくる可能性のある顧客なら誰の記録にでもアクセスできなければならないが、1 日に対応する顧客の数は限られている。
妥当なアクセスポリシーとしては、各担当者が 1 シフトでアクセスできる顧客記録を 100 件までに制限することが考えられる。
いつでもすべての記録にアクセスできるなら、不正をはたらく担当者は全顧客データのコピーを持ち出せてしまう。制限付きのポリシーなら、1 日あたりの最悪の被害を大きく抑えられる。
\end{exa}

きめ細かいアクセスポリシーができたら、次は上限をどこに設定するかという課題に直面する。
極端なエッジケースでも正当な利用を妨げないようにしなければならない場合、これは難しい問題になる。
たとえばカスタマーサービスの例では、ふだんは 50 件の記録が必要になる日すらめったにないにもかかわらず、季節的な繁忙期に対応できるよう、担当者がアクセスできる顧客記録を 1 シフトあたり最大 100 人分とするかもしれない。
なぜだろうか。
ポリシーの設定を 1 年を通じてこまめに調整するのは現実的でないし、上限が仕事の妨げになることが決してないよう、余裕を持たせておきたいからである。
また、いつ予想外に業務が急増するかわからないので、特定の日付に基づいてもっと細かく具体的なポリシーを定めても、うまくいかないかもしれない。

では、通常時と、システムとして許容すべきまれな最大需要時との差を縮める方法はないだろうか。
この厄介な状況に対処する優れた方法の一つが、非常時に使う\underline{自己申告による例外}（self-declared exception）をポリシーに設けておくことである。
この仕組みがあれば、担当者は理由を示すことで、短い期間だけ自分の上限を引き上げられる。

\begin{exa}[アクセス上限の「逃がし弁」]\label{exa:ch3-relief-valve}
このような「逃がし弁」（relief valve）を用意しておけば、基本のアクセスポリシーは厳しく絞っておける。
必要なときには、担当者はアクセス上限に達した時点で、たとえば「今日は問い合わせが多いので、残業して片付けます」といった簡単な届けを出せば、追加のアクセス許可を得られる。
こうした届けは監査でき、届けが日常的に出されるようになれば、経営側は需要が正当に増えたことをその理由とともに把握したうえで、ポリシーの上限を引き上げられる。
\end{exa}

このような柔軟な手法を使えば、恣意的になりがちな固定の厳格な制限ではなく、融通のきく上限を持ったアクセスポリシーを作れる。

%----------------------------------------------------------------------
\subsection{インターフェース}\label{subsec:ch3-interfaces}
% src: Interfaces

ソフトウェアの設計は、システムの機能ごとの部分に対応するコンポーネントで構成される。
こうした設計は、部分どうしのつながりを線で表したブロック図として描ける。
このつながりが\underline{インターフェース}（interface）であり、セキュリティ分析の主な焦点となる。データと制御の流れがわかるだけでなく、緩和策を追加できるはっきりしたチョークポイントにもなるからである。
特に信頼境界がある場合、セキュリティ上で最も注目すべきなのは、信頼度の低いコンポーネントから高いコンポーネントへ向かうデータと制御の流れであり、防御策が必要になることが多いのもそこである。

大規模なシステムには、通常、ネットワーク間、プロセス間、プロセス内のインターフェースがある。
最も強い隔離を実現するのはネットワークインターフェースである。端点どうしのやり取りは、ほぼ確実にすべて回線を通じて行われるからである。ほかの種類のインターフェースでは、事情はもっと複雑になる。
オペレーティングシステムはプロセス境界で強い隔離を実現しているので、プロセス間通信のインターフェースはネットワークインターフェースとほぼ同じくらい信頼できる。
どちらの場合も、一般にこれらのチャネルを迂回して別の経路でやり取りすることはできない。
攻撃対象領域がきれいに限定されているので、重要な信頼境界のほとんどはここにある。
そのため、プロセス間通信とネットワークインターフェースが脅威モデリングの主な焦点となる。

インターフェースはプロセス内にもあり、そこでのやり取りには比較的制約が少ない。
よく書かれたソフトウェアなら、プロセス内にも意味のあるセキュリティ境界を作れる。ただし、それが有効なのは、すべてのコードが互いにうまく協調し、決められた範囲からはみ出さない場合に限られる。
攻撃者から見れば、プロセス内の境界のほうがはるかに突破しやすい。
それでも、一つの脆弱性を通じて攻撃者が得られる制御は限られている場合もあるので、どんな保護でも、ないよりはましである。
たとえるなら、数秒しか行動する時間のない強盗を考えてみるとよい。ささやかな予防策でも、損失を防ぐには十分かもしれない。

大規模なソフトウェア設計はどれも、高い権限でアクセスする領域ができるだけ小さくなるようにコンポーネントを構成し、さらに機微な情報の流れを制限してセキュリティリスクを減らす、という繊細な課題に向き合うことになる。
情報にアクセスできるコンポーネントを、よく隔離された最小限の集合に設計で絞り込むほど、攻撃者は機微なデータにアクセスしにくくなる。
これに対して弱い設計では、あらゆる種類のデータがあちこちに流れ、コンポーネント内のどこかに脆弱性が一つあるだけで露出が大きくなる。
インターフェースのアーキテクチャは、システムが資産をうまく保護できるかどうかを左右する大きな要因である。

%----------------------------------------------------------------------
\subsection{通信}\label{subsec:ch3-communication}
% src: Communication

ネットワークにつながったシステムはすっかり当たり前になり、どのネットワークにもつながっていないスタンドアロンのコンピュータはまれな例外になった。
クラウドコンピューティングのモデルとモバイル接続が組み合わさって、ネットワークにはどこからでもアクセスできるようになっている。
その結果、インターネット接続であれ、プライベートネットワークであれ、Bluetooth や USB などによる周辺機器との接続であれ、今使われているほぼすべてのソフトウェアシステムにとって通信は欠かせないものになっている。

こうした通信を守るには、チャネルを物理的に保護して盗聴やのぞき見を防ぐか、そうでなければデータを暗号化して完全性と機密性を確保しなければならない。
物理的なセキュリティに頼るやり方は、一般に脆い。攻撃者に物理的な防護を迂回されると、たいていはデータの流れ全体にアクセスされてしまい、しかもそうした侵入は検知しにくいからである。
現代のプロセッサは十分に速く、暗号化による計算のオーバーヘッドはたいていごくわずかなので、通信を暗号化しない正当な理由はめったにない。
基本的な暗号化は第5章で、Web のための HTTPS は特に第11章で扱う。

\begin{rem}[暗号化は通信の事実を隠せない]\label{rem:ch3-fact-of-communication}
とはいえ、どんなに優れた暗号化も万能薬ではない。
残る脅威の一つとして、暗号化では\emph{通信しているという事実}（fact of communication）を隠せない。
つまり、攻撃者がチャネルを流れる生のデータを読めるなら、内容は解読できなくても、回線上でデータが送受信されていることはわかり、データ量もおおよそ見積もれる。
さらに、攻撃者が通信チャネルを改ざんできるなら、暗号化されたデータの送信を妨害できるかもしれない。
\end{rem}

%----------------------------------------------------------------------
\subsection{ストレージ}\label{subsec:ch3-storage}
% src: Storage

データストレージのセキュリティは、通信のセキュリティとよく似ている。データを保存するということは、データを未来に送り、未来のある時点で何かの目的のために取り出すということだからである。
そう考えると、通信中のデータが回線上で脆弱であるのとまったく同じように、保存されたデータも記憶媒体の上で保存中（at rest）に脆弱である。
保存中のデータを改ざんや開示から守るには、物理的なセキュリティか暗号化のどちらかが必要である。
同じように、可用性は、バックアップコピーがあるか、物理的な保護がうまくいくかにかかっている。

ストレージはシステム設計のあらゆるところに登場するので、データセキュリティの細部は運用側に任せて後回しにしがちである。しかしそうすると、設計の段階でデータ損失を先回りして緩和するよい機会を逃してしまう。
たとえば、データバックアップの要件はソフトウェア設計の重要な一部である。何が求められるかは決して自明ではなく、トレードオフも多いからである。
障害が起きてもデータが失われないように設計した冗長なストレージシステムを計画することもできるが、費用がかさむことがあり、性能面のコストもかかる。
バックアップは、データセット全体のコピーでもよいし、差分（incremental）、つまり積み重ねれば正確なコピーを再構築できるトランザクションの記録でもよい。
いずれにしても、バックアップは独立した場所に確実に保存し、許容できる遅延の範囲内で、決まった頻度で取るべきである。
クラウドのアーキテクチャなら、ほぼリアルタイムでデータを冗長に複製でき、継続的なバックアップとしてはおそらく最良の解決策になるが、それには費用がかかる。

バックアップコピーを含め、保存中のデータはすべて許可のないアクセスにさらされる危険があるので、物理的に保護するか暗号化して守らなければならない。
バックアップコピーを多く作るほど、コピーが多いせいで漏えいが起こるリスクは大きくなる。
両極端の状況を考えてみると、このことがよくわかる。

\begin{exa}[家族写真のバックアップ]\label{exa:ch3-family-photos}
写真は大切な思い出であり、どの家族にとってもかけがえのない歴史の一部なので、バックアップを何部も残しておくのは賢明である。
コピーが一つもない状態で元のファイルが失われたり、傷ついたり、壊れたりすれば、取り返しのつかない損失になりかねない。
これに備えて、家族写真のコピーをできるだけ多くの親戚に送って保管してもらう、という方法もある。
しかし、これにも欠点がある。親戚の誰かが（マルウェアに感染したり、ノートパソコンを盗まれたりして）データを盗まれる可能性が高まるからである。
こちらも深刻な事態になりうる。写真は私的な思い出であり、それがすべて Web 上に公開されて拡散されればプライバシーの侵害になる。さらに、見知らぬ人が子どもの身元を特定でき、それが搾取につながるようなことがあれば、もっと大きな脅威にもなりうる。

これは、データ損失のリスクと漏えいのリスクを天秤にかけざるをえない、根本的なトレードオフである。
両方を同時に最小にすることはできないが、いくつかの方法で、ある程度はバランスをとれる。

二つの脅威の間をとる妥協案として、写真を暗号化してから親戚に送る方法がある。（もちろん、そうすると親戚は写真を見られない。）
しかしその場合は、親戚に預けないことにした鍵を自分で保管する責任を負うことになり、その鍵をなくせば暗号化したコピーは役に立たなくなる。
\end{exa}

写真の保存の話からは、データのバックアップに関するもう一つの重要な側面、つまり記憶媒体の寿命と陳腐化の問題も見えてくる。
物理的な媒体（ハードディスクや DVD など）は時間とともに必ず劣化し、新しいハードウェアが登場するにつれて古い媒体はサポートされなくなっていく。
（著者自身、ずっと以前に、今では骨董品のようなコンピュータでしか読めない何十枚ものフロッピーディスクから 1 本の USB メモリにデータを移したことがあり、そのデータは今ではクラウドにコピーしてある。）
媒体や機器がまだ動いたとしても、新しいソフトウェアは古いデータ形式のサポートを打ち切りがちである。
そのため、データ形式の選択は重要であり、広く使われているオープンな標準形式を強く勧める。独自形式は、公式にサポートが終了すると、リバースエンジニアリングしなければ読めなくなるからである。
さらに長い期間で見れば、ソフトウェアの標準が移り変わり、アプリケーションが古い形式のサポートを廃止していくのに合わせて、ファイル形式を変換する必要が出てくるかもしれない。

この章で挙げた例は説明のために単純化したものである。特定した脅威を緩和するための手法を数多く紹介してきたが、これらは可能性のほんの氷山の一角にすぎない。
個々の緩和策は、アプリケーションごとのニーズに合わせて調整し、できれば設計に欠かせない一部として組み込むべきである。
言うのは簡単だが、実際には効果的な緩和策を講じるのは難しい。システムごとの文脈で多種多様な脅威を考慮しなければならず、しかもできることには限りがあるからである。
次の章では、セキュアな設計の一部としてこうした緩和策を作り上げるのに役立つ、有用なセキュリティ特性を持った主要なパターンと、注意すべきアンチパターンを紹介する。

\newpage

%% =======================================================
%%  part: ch4
%% =======================================================
%% =======================================================
%%  part: ch4
%% =======================================================
\section{パターン}
% src: 4: Patterns

この章では、\cite{Kohnfelder2021} の第4章に沿って、セキュアな設計を促すパターンとアンチパターンを扱う。

\begin{quote}
「芸術とは、感性によって形づくられたパターンである。」 ---Herbert Read
\end{quote}

建築家は昔から、新しい建物を構想する際に\underline{デザインパターン}（design pattern）を用いてきた。このやり方は、ソフトウェア設計の指針としてもそのまま役に立つ。
この章では、セキュアな設計を促すパターンのうち、とくに有用なものを数多く紹介する。
なかには古くからの知恵に由来するパターンもあり、肝心なのは、それをソフトウェアにどう当てはめ、それによってセキュリティがどう高まるのかを知ることである。

これらのパターンは、さまざまなセキュリティ脆弱性を緩和したり回避したりするもので、潜在的な脅威に対処するための重要な道具箱となる。
多くは単純だが、なかには理解しにくく、例で説明するのがいちばんよいものもある。
単純なパターンを軽く見てはならない。広く適用でき、最も効果の高いものに数えられるからである。
また、してはいけないことを記述する\underline{アンチパターン}（anti-pattern）として捉えたほうがわかりやすい概念もある。
ここでは、共通の特徴に基づいてパターンをグループに分けて示す。各グループは道具箱の仕切りだと考えるとよい。

% 図 4-1: Groupings of secure software patterns this chapter covers.
%   （画像のみ。グループと各パターンの名前は以下の節構成と同じだが、図の配置は本文から決まらないため描かない。）

これらのパターンをいつ、どこで使うかには判断が必要である。
設計上の決定は、必要性と単純さを基準に下すようにする。
パターンは強力だが、やりすぎてはならない。ドアに7つのデッドボルトとチェーンが要らないのと同じで、一つの問題を解決するのに考えうるデザインパターンを全部当てはめる必要はない。
複数のパターンが当てはまる場合は、最良のものを1つか2つ選ぶ。セキュリティ上の要求がきわめて重要な場合は、もう少し増やしてもよい。
使いすぎは逆効果になりうる。複雑さとオーバーヘッドが増えるせいで収穫逓減が起き、その損失がセキュリティ上の上積みをすぐに上回ってしまうからである。

\subsection{設計の性質}
% src: Design Attributes

最初のグループのパターンは、セキュアな設計がどのようなものかを大局的に述べる。一言でいえば、単純で透明であることである。
これらは「単純にしておけ（keep it simple）」と「隠すべきものは何もないはずだ（you should have nothing to hide）」という格言に由来する。
基本的で、当たり前にさえ思えるパターンだが、広く適用でき、非常に強力である。

\subsubsection{設計の簡潔さ}
% src: Economy of Design

\begin{dfn}[設計の簡潔さ]\label{dfn:ch4-economy-of-design}
\underline{設計の簡潔さ}（Economy of Design）とは、「設計はできる限り単純であるべきである」というパターンである。
\end{dfn}

設計の簡潔さがセキュリティの水準を引き上げるのは、設計が単純なほどバグが少なく、したがって見つからずに残る脆弱性も少ないと考えられるからである。
開発者は「どんなソフトウェアにもバグがある」と言うが、単純なプログラムであれば、バグがまったくないことも十分ありうる。
セキュリティ機構については、候補となる設計のうち最も単純なものを選び、きわめて重要なセキュリティ機能を担う複雑な設計には用心すること。

\begin{exa}[LEGO ブロック]\label{exa:ch4-lego}
LEGO ブロックはこのパターンの好例である。
標準となる部品をいったん設計・製造してしまえば、それを使って数え切れないほど多様で独創的な作品を組み立てられる。
汎用性の低い部品を数多くそろえた同じようなシステムでは、組み立てがもっと難しくなる。どの作品にもより多くの種類の部品の在庫が必要になり、ほかにも技術的な課題が出てくるはずである。
\end{exa}

\begin{exa}[大規模 Web サービスのシステムアーキテクチャ]\label{exa:ch4-web-service-architecture}
巨大なデータセンターで動くように作られた大規模 Web サービスのシステムアーキテクチャには、設計の簡潔さの例が数多く見られる。
大規模でも信頼性を保てるよう、こうした設計では機能を小さな自己完結したコンポーネントに分解し、それらを組み合わせて複雑な処理を行う。
多くの場合、基本的なフロントエンドが HTTPS リクエストを終端し、受け取ったデータを解析・検証して内部のデータ構造に変換する。
そのデータ構造は多数のサブサービスに送られて処理され、サブサービスはさらにマイクロサービスを使ってさまざまな機能を実行する。

Web 検索のようなアプリケーションでは、応答の各部分を別々のマシンが並列に独立して作り、さらに別のマシンがそれらを組み合わせて完全な応答に仕上げることがある。
クエリの解析、スペル修正、テキスト検索、画像検索、結果のランキング、ページのレイアウトといった処理を、それぞれ担当する小さなサービスを多数作るほうが、すべてを一つの巨大なプログラムで行うよりもはるかに簡単である。
\end{exa}

\begin{rem}[設計の簡潔さは絶対的な命令ではない]\label{rem:ch4-economy-not-absolute}
設計の簡潔さは、何でも常に単純にしなければならないという絶対的な命令ではない。
むしろ、単純さには大きな利点があることを強調し、複雑さは大きな価値をもたらす場合にだけ受け入れるべきだと説くものである。
\end{rem}

\begin{exa}[*nix と Windows のアクセス制御リスト]\label{exa:ch4-acl}
基本的な *nix と Windows とで、\underline{アクセス制御リスト}（access control list, ACL）の設計がどう違うかを考える。
前者は単純で、ユーザーごと、ユーザーグループごと、または全員に対して、読み取り・書き込み・実行の権限を指定する。
後者ははるかに込み入っていて、許可と拒否のアクセス制御エントリをどちらも任意の数だけ含められ、継承の機能もある。しかも注目すべきことに、評価の結果はリスト内のエントリの順序によって変わる。
（ここでの説明は設計上の論点を示すために単純化したもので、完全な説明ではない。）

このパターンからわかるとおり、単純な *nix のパーミッションのほうが正しく強制しやすい。それだけでなく、システムの利用者にとっても ACL の仕組みを正しく理解しやすく、そのぶん正しく使いやすい。
とはいえ、あるアプリケーションにとって Windows の ACL がちょうど適切な保護を提供し、しかも正確に設定できるのであれば、それは優れた解決策になりうる。

設計の簡潔さのパターンは、単純な選択肢が無条件に優れているとも、複雑な選択肢が必ず問題をはらむとも言っていない。
この例でいえば、*nix の ACL が本質的に優れているわけではないし、Windows の ACL に必ずバグがあるわけでもない。
ただ、Windows の ACL は開発者にとっても利用者にとっても習得が難しく、複雑な機能を使うと人を混乱させ、意図しない結果を招きやすい。
ここで重要な設計上の選択は、どちらの ACL の設計が利用者のニーズに最もよく合うかである（原著者はこの点について立場を表明しない）。
*nix の ACL は単純すぎて現実の要求を満たせないのかもしれないし、逆に Windows の ACL は典型的な使い方には機能が多すぎて扱いにくいのかもしれない。
これは各自が自分の目的に照らして答えるべき難しい問いだが、このデザインパターンはそれを考える手がかりを与えてくれる。
\end{exa}

\subsubsection{透明な設計}
% src: Transparent Design

\begin{dfn}[透明な設計]\label{dfn:ch4-transparent-design}
\underline{透明な設計}（Transparent Design）とは、「強固な保護は決して秘密に依存してはならない」というパターンである。
これに対応するアンチパターンを\underline{隠蔽によるセキュリティ}（Security by Obscurity）という。
\end{dfn}

\begin{exa}[デス・スター]\label{exa:ch4-death-star}
透明な設計のパターンに従わずに失敗した設計として最も有名なのは、おそらく『スター・ウォーズ』のデス・スターである。その排熱口は、戦闘基地の中枢をまっすぐ狙い撃てる通り道になっていた。
ダース・ベイダーがモッティ提督に向けたのと同じくらいの厳しさで、この原則について設計者たちの責任を追及していたら、物語はまったく違う展開になっていたかもしれない。
\end{exa}

よくできたシステムであれば、設計を明かすことは、その堅牢さを示して攻撃者に攻撃を思いとどまらせるはずであり、攻撃者の仕事を楽にするものであってはならない。
このパターンに対応するアンチパターン、すなわち隠蔽によるセキュリティのほうが、よく知られているかもしれない。

このパターンがとくに戒めているのは、設計が秘密であることに頼ることである。
設計を公開しなければならないとか、秘密の情報を持つこと自体が悪いとかいう意味ではない。
設計を完全に公開するとその設計が弱くなるのであれば、秘密にしておくことに頼るのではなく、設計を直すべきだということである。
暗号鍵やユーザーのアイデンティティのように、漏えいすると実際にセキュリティが損なわれる正当な秘密情報には、この話はまったく当てはまらない。
このパターンの名前が「絶対的な透明性（Absolute Transparency）」ではなく「透明な設計」なのはそのためである。
暗号化方式の設計（鍵長、メッセージ形式、暗号アルゴリズムなど）をすべて開示しても、セキュリティは少しも弱まらないはずである。

\begin{rem}[隠蔽によるセキュリティは危険信号である]\label{rem:ch4-obscurity-red-flag}
このアンチパターンは大きな危険信号とみなすべきである。
例えば、すばらしすぎて詳細を公開できない画期的な暗号アルゴリズムを発明した、と主張する自称「専門家」を信用してはならない。そうした主張は例外なくまがい物である。
\end{rem}

隠蔽によるセキュリティの問題は、一時的に敵を食い止めるのには役立つかもしれないが、きわめて脆いことである。
例えば、ある設計が時代遅れの暗号アルゴリズムを使っていたとする。そのソフトウェアがいまだに、例えば DES（1970年代のレガシーな対称暗号アルゴリズム）を使っていると悪者に知られれば、1日もかからずに簡単に破られてしまう。
そうではなく、設計の詳細が公開されていてもいなくても隠すべきものが何もないよう、しっかりしたセキュリティの土台を築くのに必要な作業をすることである。

\subsection{露出の最小化}
% src: Exposure Minimization

最大のグループのパターンは、用心することを求める。「安全側に倒す（err on the safe side）」と考えるとよい。
これらのパターンは、そうしない重要な理由がない限り安全策をとるという、リスクとリターンについての基本的な戦略を表している。

\subsubsection{最小権限}
% src: Least Privilege

\begin{dfn}[最小権限]\label{dfn:ch4-least-privilege}
\underline{最小権限}（Least Privilege）とは、「仕事にちょうど足りるだけの権限を使うのが常に最も安全である」というパターンである。
タスクを行う際に誤りを犯すリスクを減らすことを目的とする。
\end{dfn}

銃は弾を抜いた状態でだけ扱う。電動のこぎりの刃を交換するときはプラグを抜く。
こうしたありふれた安全上の習慣は、最小権限のパターンの例である。
重要なシステムの管理者が、職場でログインしたまま気ままにインターネットを閲覧すべきでないのも、このパターンによる。悪意ある Web サイトを訪れて侵害されれば、その攻撃で深刻な被害が生じかねないからである。

\begin{exa}[sudo コマンド]\label{exa:ch4-sudo}
*nix の \texttt{sudo} コマンドは、まさにこの目的のためにある。
高い権限を持つユーザーアカウント（\underline{sudoers} と呼ばれる）は、その強大な力をうっかり使ったり、侵害されて使われたりしないよう注意する必要がある。
この保護のために、スーパーユーザーのコマンドを実行するには先頭に \texttt{sudo} を付けなければならず、その際にパスワードを求められることもある。
この仕組みのもとでは、ほとんどのコマンド（\texttt{sudo} が不要なもの）はユーザー自身のアカウントにしか影響せず、システム全体に影響を与えることはできない。

これは、火災報知器のスイッチに付いている「緊急時はガラスを割ってください（IN CASE OF EMERGENCY BREAK GLASS）」というカバーが誤作動を防ぐのに似ている。スイッチを入れる前に、明示的な一手間（\texttt{sudo} の接頭辞に当たる）を強いるのである。
ガラスのカバーがあれば、誰も火災報知器をうっかり引いてしまったとは言えない。同様に、有能な管理者がまったくの偶然で \texttt{sudo} とシステムを壊すコマンドを打ち込むことはありえない。
\end{exa}

このパターンが重要なのは単純な理由による。脆弱性が悪用されたとき、攻撃者が足がかりにできる権限は少ないほどよいからである。
スーパーユーザー権限のような万能の権限は、どうしても必要なときに、できるだけ短い時間だけ使うこと。
スーパーマンでさえ最小権限を実践していた。仕事があるときだけユニフォームを着て、世界を救ったらすぐにクラーク・ケントの姿に戻っていたのである。

実際には、昇格した権限を必要な分だけ選んで控えめに使うのは手間がかかる。
電動工具の手入れをするのにいちいちプラグを抜くのが手間なのと同じで、権限を慎重に使うには規律が要るが、正しく行うほうが常に安全である。
エクスプロイトが起きたとき、権限を慎重に使っていたかどうかが、小規模な侵入で済むかシステム全体が侵害されるかの分かれ目になる。
最小権限を実践すれば、バグや人的ミスによる被害も緩和できる。

\begin{rem}[最小権限とバランス]\label{rem:ch4-least-privilege-balance}
どんな経験則もそうだが、このパターンも、複雑になりすぎないようバランス感覚を持って使うこと。
最小権限は、システムが常に文字どおり最小限の権限だけを付与すべきだという意味ではない（例えば、ファイル X に書き込むコードに、その1つのファイルへの書き込み権限だけを与える、といったことまでは求めない）。

この優れたパターンを、なぜいつも限界まで適用しないのかと思うかもしれない。
全体のバランス感覚を保つことや、どの緩和策にも収穫逓減があると認識することに加えて、ここで大きく効いてくるのは、権限を制御する仕組みの粒度と、権限を上げ下げするたびにかかるコストである。
例えば *nix のプロセスでは、権限はユーザー ID とグループ ID のアクセス制御リストに基づいて付与される。
実効 ID と実 ID を切り替えられる柔軟性（\texttt{sudo} が行っているのはこれである）を除けば、プロセスを fork せずに不要な権限を一時的に手放す簡単な方法はない。
コードは、可能なところでは低い周辺的な権限（ambient privilege）で動作し、必要な部分でだけ高い権限を使い、判断の自然な区切りで切り替えるべきである。
\end{rem}

\subsubsection{最小情報}
% src: Least Information

\begin{dfn}[最小情報]\label{dfn:ch4-least-information}
\underline{最小情報}（Least Information）とは、「仕事に必要な最小限の個人情報だけを収集し、アクセスするのが常に最も安全である」というパターンである。
最小権限（定義 \ref{dfn:ch4-least-privilege}）のデータプライバシー版であり、意図しない開示のリスクを最小限に抑えるのに役立つ。
\end{dfn}

サブルーチンを呼び出すとき、サービスを要求するとき、要求に応答するときには、必要以上の個人情報を渡さないようにし、機会があるごとに不要な情報の流れを減らすこと。
このパターンを実際に実装するのは難しいことがある。ソフトウェアは用途に合わせて最適化されていない標準的なコンテナでデータを受け渡しがちなので、本当は必要のない余分なデータが紛れ込むことが多いからである。
実のところ、このパターンがほかの文献で取り上げられているのを目にすることはまずないだろう。

ソフトウェアがこのパターンを満たせないことは、あまりにも多い。インターフェースの設計は時とともに多くの目的に応えるよう発展していくものであり、一貫性のために同じパラメータやデータ構造を使い回すのが便利だからである。
その結果、厳密には必要のないデータが、一見無害な余分な荷物として一緒に送られる。
もちろん、システム中を流れるこの不要なデータが新たな攻撃の機会を生むと、問題になる。

\begin{exa}[CRM システム]\label{exa:ch4-crm}
企業内のさまざまな従業員が使う、大規模な\underline{顧客関係管理}（customer relationship management, CRM）システムを考える。
従業員たちは、営業、生産、出荷、サポート、保守、研究開発、経理など、多種多様な目的でこのシステムを使う。
役割に応じて、各従業員がアクセスを許される情報の範囲は異なる。
最小情報を実践するには、この企業のアプリケーションは、特定のタスクを行うのに必要な最小限のデータだけを要求すべきである。

電話に応対するカスタマーサポート担当者を考えよう。システムが発信者番号通知（Caller ID）を使って顧客レコードを検索するなら、担当者は顧客の電話番号を知る必要はなく、購入履歴さえわかればよい。
これと対照的なのが、全項目を含む顧客レコードの検索を、許可するか不許可にするかしか選べない素朴な設計である。
理想を言えば、担当者がもっと広いアクセス権限を持っていても、個々のタスクでは必要最小限の情報だけを要求してそれで作業できるようにすべきであり、そうすれば開示のリスクを最小限に抑えられる。
\end{exa}

実装のレベルでは、ローカルにキャッシュした情報を不要になった時点で消去することや、ユーザーが明示的に詳細の表示を求めるまでは利用可能なデータの一部だけを画面に表示することなども、最小情報の設計に含まれる。
パスワードを \texttt{********} と表示するおなじみの慣行は、このパターンによって\underline{ショルダーサーフィン}（shoulder surfing、肩越しの盗み見）のリスクを緩和している。

このパターンは、設計の段階で適用することがとくに重要である。インターフェースの両側を同時に変更しなければならないので、後から実装するのはきわめて難しくなりうるからである。
必要とするデータの組がそれぞれ異なる特定のタスクに合わせて、独立したコンポーネントを設計すれば、うまくいく可能性が高まる。
機密性の高いデータを扱う API は、情報の露出を最小限にするため、呼び出し側が必要なデータの部分集合を指定できるようにしておくべきである（表 \ref{tab:ch4-least-information-api}）。

\begin{table}[H]
\centering
\caption{最小情報に準拠しない API と準拠する API の例}\label{tab:ch4-least-information-api}
% 表 4-1: Examples of Least Information Compliant and Non-Compliant APIs
\small
\begin{tabular}{|p{0.44\textwidth}|p{0.44\textwidth}|}
\hline
最小情報に準拠しない API & 最小情報に準拠する API \\
\hline
\texttt{RequestCustomerData(id='12345')}
&
\texttt{RequestCustomerData(id='12345', items=['name', 'zip'])} \\
\hline
\texttt{\{'id': '12345', 'name': 'Jane Doe', 'phone': '888-555-1212', 'zip': '01010', . . .\}}
&
\texttt{\{'name': 'Jane Doe', 'zip': '01010'\}} \\
\hline
\end{tabular}
\end{table}

左側の \texttt{RequestCustomerData} API は最小情報のパターンを無視している。呼び出し側は、ID を指定して完全なデータレコードを要求するしかないからである。
呼び出し側は電話番号を必要としないので、要求する必要もない。受け取った電話番号を無視したとしても、電話番号を狙う攻撃者にとっての攻撃対象領域は広がったままである。
右側は同じ API の別バージョンで、呼び出し側が必要な項目を指定でき、その項目だけが返されるので、個人情報の流れが最小限になる。

セキュアバイデフォルトのパターン（定義 \ref{dfn:ch4-secure-by-default}）もあわせて考えると、呼び出し側が情報の流れを抑えるために必要なものを正確に要求できることを前提に、\texttt{items} パラメータのデフォルトは最小限の項目の組にしておくべきである。

\subsubsection{セキュアバイデフォルト}
% src: Secure by Default

\begin{dfn}[セキュアバイデフォルト]\label{dfn:ch4-secure-by-default}
\underline{セキュアバイデフォルト}（Secure by Default）とは、「ソフトウェアは常に『箱から出した状態で』セキュアであるべきである」というパターンである。
初期状態も含めてソフトウェアをセキュアバイデフォルトに設計し、運用者が何もしなくてもリスクが生じないようにする。
これはシステム全体の設定だけでなく、コンポーネントの設定オプションや API のパラメータにも当てはまる。
\end{dfn}

デフォルトパスワードが設定されたデータベースやルーターは、このパターンに違反していることで悪名高く、この設計上の欠陥は今日に至るまで驚くほど広く残っている。

セキュリティに真剣に取り組むなら、後でセキュアにするつもりで安全でない状態に設定することは決してしてはならない。脆弱な期間が生まれるうえ、後でセキュアにすることはあまりにも頻繁に忘れられるからである。
例えば、デフォルトパスワードのある機器を使わなければならない場合は、ネットワークに配置する前に、まずファイアウォールの内側のプライベートネットワークで安全に設定すること。
この分野の先駆者であるカリフォルニア州は、このパターンを法律で義務付けている。同州の Senate Bill No.~327（2018年）は、ネットワークに接続される機器のデフォルトパスワードを禁止している。

セキュアバイデフォルトが当てはまるのはデフォルトパスワードに限らない。セキュリティに悪影響を与えうる設定や構成はすべて対象になる。
パーミッションは、デフォルトでは制限の強い設定にしておくべきである。緩い設定が必要な場合は、安全であるときに限って、ユーザーが明示的に変更しなければならないようにする。
潜在的に危険なオプションは、すべてデフォルトで無効にしておくこと。
逆に、セキュリティ保護を提供する機能はデフォルトで有効にし、最初から機能するようにしておくこと。
そしてもちろん、ソフトウェアを常に最新の状態に保つことも重要である。古いバージョン（既知の脆弱性を含んでいるかもしれない）で運用を始め、いつか更新されるだろうと期待するようなことをしてはならない。

理想を言えば、安全でないオプションはそもそも必要ないはずである。
設定可能なオプションの提案は慎重に検討すること。安全でないオプションを用意するのは簡単かもしれないが、それは後々、ほかの人にとっての落とし穴になる。
また、オプションが一つ増えるごとに設定の組み合わせの数も増え、そのすべての組み合わせが本当に有用で安全であることを確かめる作業は、オプションが増えるほど難しくなることも忘れてはならない。
どうしても安全でない設定を用意しなければならない場合は、そのリスクを管理者に積極的に説明するよう心がけること。

もっとも、セキュアバイデフォルトが当てはまる範囲は、設定オプションよりもはるかに広い。
API のパラメータが指定されなかった場合のデフォルトは、安全な選択であるべきである。

\begin{exa}[HTTPS のデフォルト]\label{exa:ch4-https-default}
アドレスバーにプロトコルなしで入力された URL を受け付けるブラウザは、そのサイトが HTTPS を使うものと仮定し、それで接続できなかった場合にだけ HTTP にフォールバックすべきである。
新しい HTTPS 接続を確立するためにネゴシエートする2つのピアは、デフォルトで、より安全な暗号スイートを優先して受け入れるべきである。
\end{exa}

\subsubsection{拒否リストより許可リスト}
% src: Allowlists over Blocklists

\begin{dfn}[許可リストと拒否リスト]\label{dfn:ch4-allowlist-blocklist}
\underline{許可リスト}（allowlist）は安全なものを列挙したもので、本質的に有限である。
これに対して\underline{拒否リスト}（blocklist）は安全でないものをすべて列挙しようとするもので、そうすることで、安全であってほしいと願う無限の集合を暗黙のうちに許可してしまう。
\underline{拒否リストより許可リスト}（Allowlists over Blocklists）とは、セキュリティ機構を設計する際に拒否リストより許可リストを選ぶべきであるというパターンである。
どちらのやり方のほうが危険かは明らかである。
\end{dfn}

まず、許可リストと拒否リストという選択肢が何を意味するのか、そしてなぜ常に許可リストを選ぶべきなのかをはっきり理解するために、ソフトウェア以外の例を挙げる。

\begin{exa}[COVID-19 下の海岸閉鎖命令]\label{exa:ch4-beach-closure}
COVID-19 による外出禁止の緊急命令が出た最初の数か月間、原著者の州の知事は、次のような但し書き付きで海岸の閉鎖を命じた（ここでは単純化して示す）。
\begin{itemize}
\item 「何人も、いかなる海岸においても、座り、立ち、横たわり、くつろぎ、日光浴をし、またはぶらついてはならない……」
\item ……ただし、「社会的距離の要件が守られている限り、海岸を走り、ジョギングし、または歩く」場合を除く（サーフィンをするために海岸を横切ることも許可される）。
\end{itemize}
第1の条項は許可しない活動を列挙しているので拒否リストであり、第2の例外条項は列挙した活動を許可しているので許可リストである。
法律上の事情から、このような文言にする十分な理由があるのかもしれないが、純粋に論理の観点から見ると改善の余地が大いにある、と原著者は考えている。

まず拒否リストを考える。海岸で人がしかねない危険な活動のなかには、第1の条項が禁止し損ねているものがほかにもあるはずである。
命令の意図が人々を動き続けさせることにあったのなら、抜け落ちているものはたくさんある。例えば、ひざまずくこと、ヨガ、生きた彫像のパフォーマンスなどである。
拒否リストの問題は、抜け落ちがそのまま欠陥になることである。したがって、考えうる悪いケースをすべて列挙し尽くせない限り、安全でないシステムになる。

次に、海岸で許可される活動の許可リストを考える。
こちらも不完全ではあるが（スキップしても問題ないことに誰が異を唱えるだろうか）、大きなセキュリティ上の問題は引き起こさない。
海岸でスキップしていた人のごく一部が不当に罰せられるかもしれないが、その害は小さい。もっと重要なのは、列挙が不完全でも、危険な活動を許してしまう穴が開くことはないという点である。
最初に漏れていた安全な活動は、必要に応じて簡単に許可リストに追加できる。
\end{exa}

より一般的に、左端の「不許可」から右端の「許可」へと少しずつ移り変わっていく連続体を考えてみよう。
その途中のどこかに境界線がある。
目標は、線の右側にある良いものを許可し、左側にある悪いものを不許可にすることである。
許可リストは右端から線を引き、リストが大きくなるにつれて線を少しずつ左へ動かして、スペクトルのより広い範囲を含めていく。
許可リストから何か良いものが漏れていても、本当の境界であるとらえがたい線より安全側にいることに変わりはない。
安全な行為をすべて許可するぴったりの点（それ以上何かをリストに加えると行き過ぎになる点）には決してたどり着けないかもしれないが、この方法なら安全側にとどまるのは簡単である。
これを拒否リストのやり方と比べてみよう。本当の境界より左側にあるものを残らず列挙しない限り、許可すべきでないものを許可してしまうことになる。
最も安全な拒否リストはほとんどすべてを含むものになるが、それではおそらく制限が厳しすぎるので、どちらにしてもうまくいかない。

許可リストの使い方はしばしばあまりにも当たり前なので、それがパターンだと意識されないことも多い。
例えば銀行なら、高額取引を承認する権限を、信頼できる少数の管理職に与えるのが妥当だろう。
権限のない従業員全員の拒否リストを管理して、それ以外の従業員には暗黙のうちにその権限を与える、などということは誰も考えもしない。

\begin{rem}[拒否リストによる入力検証]\label{rem:ch4-blocklist-input-validation}
ところが、ずさんなプログラマーは、値に無効な文字のリストのどれも含まれていないことを確かめる、というやり方で入力検証をしようとし、その過程で NUL（ASCII 00）や DEL（ASCII 127）のような文字をいとも簡単に忘れてしまう。
\end{rem}

\begin{exa}[アンチウイルス]\label{exa:ch4-antivirus}
皮肉なことに、おそらく最も売れている消費者向けのソフトウェアセキュリティ製品であるアンチウイルスは、既知のマルウェアをすべてブロックしようとする。
現代のアンチウイルス製品は、ダイジェストを既知のマルウェアのデータベースと照合するだけだった昔のものよりはるかに洗練されているが、それでもどれも、ある程度は拒否リストに基づいて動作しているように見える。
（ほとんどの商用アンチウイルスソフトウェアはプロプライエタリなので、推測するしかない。隠蔽によるセキュリティの好例である。）
アンチウイルスが拒否リストの手法から抜け出せないのは、もっともなことである。マルウェアの見本を集める方法はわかっているが、世界中の良いソフトウェアを一つ残らずリリース前に何らかの方法で許可リストに入れるという構想は、とても実現しそうにないからである。
原著者の主眼は特定の製品やその価値の評価にあるのではなく、拒否リストによって保護するという設計上の選択と、それがなぜ必然的にリスクを伴うかにある。
\end{exa}

\subsubsection{予測可能性の回避}
% src: Avoid Predictability

\begin{dfn}[予測可能性の回避]\label{dfn:ch4-avoid-predictability}
\underline{予測可能性の回避}（Avoid Predictability）とは、「予測可能なデータ（や振る舞い）は秘密にしておけない。攻撃者が推測で知ることができるからである」という認識に基づくパターンである。
\end{dfn}

ソフトウェア設計において、データが予測可能であることは情報の漏えいにつながりうるため、深刻な欠陥のもとになることがある。

\begin{exa}[連番の顧客アカウント ID]\label{exa:ch4-sequential-account-id}
新しい顧客アカウント ID を割り当てるという単純な例を考える。
新しい顧客が Web サイトに登録すると、システムはそのアカウントを指し示す一意の ID を必要とする。
すぐに思いつく簡単な方法の一つは、最初のアカウントを 1、2番目のアカウントを 2、というように番号を振ることである。
これでも動作はするが、攻撃者の視点から見ると、何が明らかになってしまうだろうか。

新しいアカウント ID を見れば、攻撃者はそれまでに作られたユーザーアカウントの数を簡単に知ることができる。
例えば攻撃者が定期的に使い捨てのアカウントを新しく作れば、その時点でその Web サイトに顧客アカウントがいくつあるかを正確に把握できる。これは、ほとんどの企業が競合他社に知られたくない情報である。
システムの細部によっては、ほかにも多くの落とし穴がありうる。
このまずい設計がもたらすもう一つの結果は、次に作られる新しいアカウントに割り当てられる ID を、攻撃者が簡単に推測できることである。この知識があれば、攻撃者は自分がその新しいアカウントであると名乗って登録システムを混乱させ、新しいアカウントの設定を妨害できるかもしれない。
\end{exa}

予測可能性の問題はさまざまな形で現れ、設計によって漏れる情報の種類も変わる。
例えば、アカウント保有者の名前や郵便番号の一部を含むアカウント ID は、アカウント所有者のアイデンティティの手がかりを不必要に漏らしてしまう。
もちろん、同じ問題は Web ページやイベントなどの ID にも当てはまる。
こうした問題に対する最も単純な緩和策は、ID の目的が一意のハンドルであるなら、ID をハンドル以外の何物でもないものにすることである。ユーザー数のカウントや、ユーザーのメールアドレスや、その他の識別情報をもとに作ってはならない。

こうした問題を避ける簡単な方法は、セキュアにランダムな ID を使うことである。
真にランダムな値は推測できないので、情報を漏らさない。
（厳密に言えば、ID の長さからは ID の最大数がわかってしまうが、通常これは機密情報ではない。）
標準的なシステム機能として提供される乱数生成器には、擬似乱数生成器とセキュアな乱数生成器の2種類がある。
予測可能であっても無害だと確信できない限り、低速ではあってもセキュアな方を使うべきである。
セキュアな乱数生成器について詳しくは第5章を参照。

\subsubsection{安全な失敗}
% src: Fail Securely

\begin{dfn}[安全な失敗]\label{dfn:ch4-fail-securely}
\underline{安全な失敗}（Fail Securely）とは、「問題が起きた場合には、必ずセキュアな状態に落ち着くようにせよ」というパターンである。
\end{dfn}

\begin{exa}[電気ヒューズ]\label{exa:ch4-fuse}
物理的な世界では、このパターンは常識そのものである。
昔ながらの電気ヒューズがよい例である。電流が流れすぎると、熱で金属が溶けて回路が切れる。
物理法則により、過剰な電流が流れ続けるような壊れ方はありえない。
\end{exa}

このパターンはいちばん当たり前のものに思えるかもしれないが、ソフトウェアの性質上（物理法則は味方してくれない）、簡単に軽んじられてしまう。

最初はほとんど取るに足りないように見えるコーディングのタスクも、エラー処理のために複雑になることが多い。
通常のプログラムの流れは単純でも、接続が切れたとき、メモリの確保に失敗したとき、入力が無効なとき、その他さまざまな問題が起きたときには、可能であれば処理を続け、無理ならうまく撤退する必要がある。
コードを書いていると、本来のタスクよりもこうした雑事の処理に多くの時間を取られているように感じることがある。そのためエラー処理のコードを重要でないと軽く片付けてしまいがちで、これが脆弱性のよくある原因になっている。
攻撃者は、悪用できる脆弱性があることを期待して、できることなら意図的にこうしたエラーケースを引き起こそうとする。
落とし穴は無数にあるが、よくある罠のいくつかは取り上げておく価値がある。

エラーケースを徹底的にテストするのは面倒なことが多く、複数のエラーが重なって新しいコードパスが生じうる場合はなおさらである。そのため、エラーケースは攻撃の温床になりうる。
個々のエラーが安全に処理されるか、さもなければ要求全体が拒否されるか、そのどちらかになるようにすること。

\begin{exa}[写真共有サービスへの画像アップロード]\label{exa:ch4-image-upload}
誰かが写真共有サービスに画像をアップロードしたら、ただちにその画像が正しい形式かどうかを確認し（不正な形式の画像は悪意を持って使われることが多いため）、正しくなければ、それ以上使われないよう、すぐにそのデータをストレージから削除すること。
\end{exa}

\subsection{強固な強制}
% src: Strong Enforcement

このグループのパターンは、規則を徹底して強制することで、コードが確実に正しく振る舞うようにする方法を扱う。
抜け穴はあらゆる法律や規制の悩みの種であり、これらのパターンは、システムの裏をかく手段を作り込まない方法を示す。
コードを書いてから、それが禁じられたことをするはずはないと推論するよりも、禁じられたことが起こりえないように構造的に設計するほうがよい。

\subsubsection{完全仲介}
% src: Complete Mediation

\begin{dfn}[完全仲介]\label{dfn:ch4-complete-mediation}
\underline{完全仲介}（Complete Mediation）とは、「すべてのアクセス経路を保護し、例外なく同じアクセス制御を強制せよ」というパターンである。
つまり、保護された資産へのアクセスは、すべて一貫した方法でセキュアに検査するということである。
あるリソースにアクセスする方法が複数ある場合、どの方法にも同じ認可チェックを課さなければならず、フリーパスで通れたり、より緩いポリシーが適用されたりする近道があってはならない。
\end{dfn}

完全仲介は、当たり前の考え方に付けられた、わかりにくい名前である。

\begin{exa}[投資会社の納税者番号]\label{exa:ch4-tax-id}
ある金融投資会社の情報システムのポリシーで、一般従業員は管理職の承認なしに顧客の納税者番号（tax ID）を調べられないと定められているとする。そのため、システムは一般従業員にその項目を省いた縮小版の顧客レコードを提供している。
管理職は完全なレコードにアクセスでき、管理職でない従業員に正当な必要がまれに生じた場合は、管理職に調べてもらうよう頼める。

従業員はさまざまな形で顧客を支援しており、その一つとして、何らかの理由で顧客のもとに郵送の納税申告用書類が届かなかった場合に、書類を再発行している。
顧客の本人確認をした後、従業員は書類の複製（PDF）を要求し、それを印刷して顧客に郵送する。

このシステムの問題は、従業員がアクセスできないはずの顧客の納税者番号が、その納税書類に記載されていることである。これは完全仲介の失敗である。
不正を働こうとする従業員は、再発行を装って任意の顧客の納税書類を要求し、その納税者番号を知ることができる。こうして、従業員への開示を防ぐはずのポリシーが骨抜きになる。
\end{exa}

このパターンに従う最善の方法は、可能な限り、あるセキュリティ判断を行う箇所を一か所にまとめることである。
その箇所はしばしば\underline{ガード}（guard）、くだけた言い方では\underline{ボトルネック}（bottleneck）と呼ばれる。
ある資産へのアクセスは、すべて一つの門を通らなければならない、という考え方である。
それが無理で、複数の経路にガードを置く必要がある場合は、同じアクセスに対するチェックはすべて機能的に等価であるべきで、理想的には同一のコードで実装すべきである。

実際には、このパターンを一貫して守るのは難しいことがある。
ガードの置き方によって、準拠の度合いには次のような違いがある。

\begin{dfn}[完全仲介への準拠の度合い]\label{dfn:ch4-mediation-compliance}
\begin{itemize}
\item \underline{高い準拠}（high compliance）：リソースへのアクセスが一つの共通ルーチンを介してのみ許される（ボトルネック型のガード）。
\item \underline{中程度の準拠}（medium compliance）：リソースへのアクセスがさまざまな箇所で行われるが、どの箇所も同一の認可チェックで守られている（共通の複数ガード）。
\item \underline{低い準拠}（low compliance）：リソースへのアクセスがさまざまな箇所で行われ、それぞれがまちまちで一貫しない認可チェックで守られている（不完全な仲介）。
\end{itemize}
\end{dfn}

あるリソースについて、認可ポリシーを単純にし、認可チェックを一つのボトルネックのコードパスに集中させる設計が、なぜこのパターンを正しく実現する最善の方法なのか。それを反例で示そう。
最近ある Reddit ユーザーが、これを誤るのがいかに簡単かを示す事例を報告した。

\begin{exa}[iMessage 経由のスクリーンタイム回避]\label{exa:ch4-screen-time}
\begin{quote}
8歳の妹が iOS 12.4.6 の iPhone 6 で、スクリーンタイムの制限時間を過ぎても YouTube を使っているのを見かけた。
調べてみると、妹はメッセージにあるスクリーンタイムのバグを見つけていた。それを使うと、iMessage App Store で入手できるアプリが使えてしまうのである。
\end{quote}
Apple は iMessage に独自のアプリを組み込めるように設計したため、YouTube アプリを起動する方法が複数できたが、この動画視聴への別経路にはスクリーンタイムのチェックを実装していなかった。これは完全仲介の典型的な失敗である。
\end{exa}

同じリソースにアクセスする複数の経路を、それぞれ微妙に異なる動作をしかねない独自のコードで用意するのは避けること。少しでも食い違いがあれば、ある経路のガードがほかの経路より弱くなっているおそれがあるからである。
ガードが複数あると、本質的に同じチェックを何度も実装しなければならず、変更のたびに何か所も直す必要があるので、保守も難しくなる。
重複したガードを使えば、誤りを犯す機会が増え、徹底的にテストする手間も増える。

\subsubsection{最小共通メカニズム}
% src: Least Common Mechanism

\begin{dfn}[最小共通メカニズム]\label{dfn:ch4-least-common-mechanism}
\underline{最小共通メカニズム}（Least Common Mechanism）とは、「共有されるメカニズムを最小化することによって、独立したプロセス間の分離を維持せよ」というパターンである。
\end{dfn}

このパターンが何を意味し、どう役に立つのかは、例で考えるとわかりやすい。

\begin{exa}[マルチユーザー OS のカーネル]\label{exa:ch4-kernel}
マルチユーザーのオペレーティングシステムのカーネルは、異なるユーザーコンテキストで動くプロセスのためにシステムリソースを管理する。
カーネルの設計は、プロセスが明示的にリソースや通信チャネルを共有しない限り、プロセス同士の分離を根本から保証している。
内部では、カーネルはすべてのユーザープロセスからの要求に応えるために、さまざまなデータ構造を保持している。
このパターンが指摘するのは、こうしたデータ構造という共通のメカニズムが、うっかりプロセス間を\underline{橋渡し}（bridge）してしまうおそれがあり、そのような機会はできるだけ減らすのが最善だということである。
例えば、ある機能をユーザーランドのコードで実装できるなら、その機能はプロセス境界によって必然的にそのプロセスの中に閉じ込められるので、何らかの形でユーザープロセス間を橋渡ししてしまう可能性は低くなる。
ここでいう「橋渡し」とは、とくに、情報を漏らすことや、あるプロセスが認可なしに別のプロセスに影響を与えられるようにすることを指す。
\end{exa}

それでもまだ抽象的に感じるなら、ソフトウェア以外の次のたとえを考えてみよう。

\begin{exa}[会計士の机]\label{exa:ch4-accountant}
申告期限の前日に、納税申告書を確認してもらうために会計士のもとを訪れたとする。
会計士の机の上には、書類やフォルダーの山がミニチュアの摩天楼のように積み上がっている。
雑然とした山をかき回した末に、会計士はあなたの書類を取り出して打ち合わせを始める。
待っている間、他人の名前や納税者番号が書かれた納税書類や銀行の明細書が丸見えになっている。
会計士は、あなたの税金についての走り書きのメモを、うっかり他人のファイルに書き込んでしまうかもしれない。
これこそ、会計士が仕事机を共通のメカニズムとして共有しているせいで生じる独立した当事者間の橋渡しであり、最小共通メカニズムが避けようとするものである。

翌年、別の会計士を雇って会いに行くと、その会計士はキャビネットからあなたのファイルを取り出す。
会計士が机の上でファイルを開くと、机は整頓されていて、ほかの顧客の書類はどこにも見当たらない。
これが最小共通メカニズムの正しいやり方であり、書類の取り違えや、詮索好きな顧客が他人の書類を見てしまうリスクが最小限に抑えられる。
\end{exa}

ソフトウェアの分野では、独立したプロセスや異なるユーザーとやり取りするサービスを設計する際に、このパターンを適用すること。
全員のデータが入った一枚岩のデータベースの代わりに、ユーザーごとに別々のデータベースを用意したり、コンテキストに応じてアクセスの範囲を限定したりできないか検討しよう。
すべてのデータを一か所に置くもっともな理由がある場合もあるが、このパターンに従わないことを選ぶなら、それによって増えるリスクに注意を払い、必要な分離を明示的に強制すること。

\begin{exa}[Web のクッキー]\label{exa:ch4-cookies-least-common}
Web のクッキーは、各クライアントが自分のクッキーデータを独立に保存するので、このパターンを使った好例である。
\end{exa}

\subsection{冗長性}
% src: Redundancy

\underline{冗長性}（redundancy）は工学における安全性の中心的な戦略であり、自動車のスペアタイヤのように、多くの常識的な慣行に表れている。
このグループのパターンは、冗長性を使ってソフトウェアをよりセキュアにする方法を示す。

\subsubsection{多層防御}
% src: Defense in Depth

\begin{dfn}[多層防御]\label{dfn:ch4-defense-in-depth}
\underline{多層防御}（Defense in Depth）とは、「独立した保護の層を組み合わせると、全体としてより強固な防御になる」というパターンである。
\end{dfn}

この強力な手法は、どうしてもバグを抱えてしまうソフトウェアシステムを、個々の構成要素よりもセキュアにするための、最も重要なパターンの一つである。

\begin{exa}[穴の開いた合板と暗室]\label{exa:ch4-plywood}
窓を合板でふさいで部屋を暗室にしたいとする。
合板はたくさんあるが、誰かがどの板にも小さな穴をでたらめにいくつか開けてしまっている。
板を1枚打ち付けただけでは、無数のピンホールのせいで暗くならない。
その上に2枚目の板を打ち付ければ、たまたま2つの穴が重なりでもしない限り、部屋は真っ暗になる。
金属探知機とボディチェックの両方を備えたセキュリティ検査所も、このパターンの例である。
\end{exa}

ソフトウェア設計では、とくに重要なセキュリティ判断を守るために、2つ以上の独立した保護機構を重ねることで多層防御を取り入れる。
穴の開いた合板と同じく、それぞれの実装には欠陥があるかもしれないが、ある攻撃が両方を突破する可能性はごくわずかである。合板の2つの穴がたまたま重なって光が漏れるようなものである。
独立したチェックを2つ設ければ労力は2倍、時間も2倍かかるので、この手法は控えめに使うべきである。

\begin{exa}[サンドボックス]\label{exa:ch4-sandbox}
労力とオーバーヘッドを利益とうまく釣り合わせたこの手法の好例が、\underline{サンドボックス}（sandbox）の実装である。サンドボックスとは、信頼できない任意のコードを安全に実行できるコンテナである。
（現代の Web ブラウザは Web Assembly をセキュアなサンドボックス内で実行する。）
信頼できないコードをシステム内で実行すると、何かがうまくいかなかった場合に壊滅的な結果を招きかねないので、保護の層を複数重ねるオーバーヘッドも正当化される（図 \ref{fig:ch4-sandbox}）。

サンドボックスで実行するコードは、まず解析器（防御の第1層）でスキャンされ、一連の規則に照らして検査される。
違反があれば、システムはそのコードを丸ごと拒否する。
例えば、カーネルの呼び出しを禁止する規則や、特定の特権マシン命令の使用を禁止する規則などがありうる。
スキャナを通過したコードだけが、インタプリタにロードされる。インタプリタはコードを実行しながら、同じ種類の過剰な特権操作を防ぐための制限を数多く強制する。
攻撃者がこのシステムを破るには、まずスキャナの規則チェックをすり抜け、さらにインタプリタをだまして禁止された操作を実行させなければならない。

この例がとくに効果的なのは、コードのスキャンと解釈が根本的に異なる処理なので、同じ欠陥が両方の層に現れる可能性が低いからである。それぞれが独立に開発されていればなおさらである。
スキャナがある攻撃手法を見逃す確率が百万分の一で、インタプリタも同じだとすれば、両者を組み合わせたシステム全体が実際に破られる確率はおよそ一兆分の一になる。
これがこのパターンの力である。
\end{exa}

% 図 4-2: An example of a sandbox as the Defense in Depth pattern
%   （画像のみ。下図は上の本文の記述（スキャナ→インタプリタの2層、違反時は拒否）から決まる範囲で描いた。）
\begin{figure}[H]
\centering
\begin{tikzpicture}[
  bx/.style={draw, rounded corners, align=center, minimum width=2.6cm, minimum height=1.0cm, font=\small},
  ar/.style={-{Stealth[length=2.5mm]}, thick},
  lb/.style={font=\scriptsize, fill=white, inner sep=1pt}
]
\node[bx] (C) at (0,0) {信頼できない\\コード};
\node[bx] (S) at (4,0) {解析器（スキャナ）\\防御の第1層};
\node[bx] (I) at (8,0) {インタプリタ\\防御の第2層};
\node[bx] (R) at (4,-2.2) {拒否};
\draw[ar] (C) -- (S);
\draw[ar] (S) -- node[lb, above] {規則を通過} (I);
\draw[ar] (S) -- node[lb, right] {規則に違反} (R);
\end{tikzpicture}
\caption{多層防御のパターンとしてのサンドボックスの例}\label{fig:ch4-sandbox}
\end{figure}

\subsubsection{権限の分離}
% src: Separation of Privilege

\begin{dfn}[権限の分離]\label{dfn:ch4-separation-of-privilege}
\underline{権限の分離}（Separation of Privilege）とは、「2者は1者よりも信頼できる」というパターンである。
\underline{職務の分離}（Separation of Duty）とも呼ばれる。
異なる2人に別々の鍵を託すなら、錠が2つあるほうが1つより強い、という議論の余地のない事実を指している。
\end{dfn}

その2人が共謀する可能性もあるが、そうしたことはめったに起きない。しかもそのリスクを最小限に抑えるよい方法があり、いずれにしても1人の個人に全面的に頼るよりはるかによい。

\begin{exa}[貸金庫]\label{exa:ch4-safe-deposit-box}
貸金庫は、すべての貸金庫が収められた金庫室のセキュリティを銀行が管理し、利用者はそれぞれ自分の箱を開ける別々の鍵を持つように設計されている。
銀行員は、錠をドリルで壊すといった力ずくの手段に訴えない限りどの箱も開けられず、一方でどの顧客も金庫室を開ける番号を知らない。
銀行が顧客にアクセスを認め、そのうえで顧客が自分の鍵を使ったときにだけ、箱を開けることができる。
\end{exa}

このパターンは、保護されたリソースに対して、別々の、重なり合う責任がある場合に適用すること。

\begin{exa}[データセンターのセキュリティ]\label{exa:ch4-datacenter}
データセンターのセキュリティは典型的な事例である。
データセンターには、スーパーユーザーのアクセス権を持ってマシンを運用する責任を負うシステム管理者（大規模な運用ならそのチーム）がいる。
それとは別に、警備員が施設への物理的なアクセスを管理している。
これらの別々の職務は、それぞれの資格情報やアクセスキーの管理も含めて、組織内で異なる役員の配下にある従業員に割り当てるべきである。そうすれば共謀が起こりにくくなり、一人の上司が手順に反する異例の行動を命じることも防げる。
具体的には、リモートで作業する管理者はデータセンター内のマシンに物理的にアクセスできるべきではなく、データセンター内に物理的にいる人々は、マシンにログインするためのアクセスコードやストレージユニットを復号するのに必要な鍵を一切知るべきではない。
物理的なアクセスと管理者のアクセスの両方を手に入れてセキュリティを完全に破るには、それぞれの管理領域から1人ずつ、計2人が共謀しなければならない。
大きな組織では、データセンター内で管理されるさまざまなデータセットについて別々のグループが責任を負い、分離の度合いをさらに高めることもある。
\end{exa}

このパターンのもう一つの使い方は、一つの責任を複数の職務に分け、一人の行為者の誤りや悪意によって深刻な結果が生じるのを防ぐことである。ふつうは最も重要な機能のためにとっておかれる使い方である。
データのバックアップコピーが漏えいする可能性に備えた追加の保護として、別々の人に託した異なる鍵で2回暗号化しておけば、後でそのコピーを使うには両方の当事者の協力が必要になる。
極端な例として、核ミサイルの発射には、10フィート離れた錠に差し込んだ2つの鍵を同時に回す必要があり、一人で行動する人間が発射させることは決してできないようになっている。

監査ログは権限の分離によって保護すること。イベントの記録とレビューに責任を負うチームを、イベントを起こすチームとは別にするのである。
つまり、管理者はユーザーの活動を監査できるが、管理者を監査するには別のグループが必要だということである。
そうしなければ、悪意ある行為者が自分の不正な活動の記録を妨げたり、監査ログを改ざんして痕跡を隠したりできてしまう。

スーパーユーザー権限を持つ管理者はすべてを制御できるので、単一のコンピュータの中で権限の分離を実現することはできないが、それでもうまく近似する方法はたくさんある。
複数の独立したコンポーネントからなる設計を実装すれば、最終的には管理者がそれを破れるとしても、緩和策として意味がある。破壊工作が複雑になって攻撃に時間がかかるようになり、攻撃者が途中で誤りを犯して捕まる可能性が高まるからである。
管理者に対する強固な権限の分離は、管理者に、別の管理下にある専用の ssh ゲートウェイを経由して作業させ、そのゲートウェイでセッションを完全かつ詳細に記録し、場合によってはほかの制限も課す、という設計で実現できる。

内部者の脅威をなくすのは難しく、場合によっては不可能だが、だからといって緩和策が時間の無駄だということにはならない。
誰かが見ているとわかっているだけで、それ自体が大きな抑止力になる。
こうした予防措置は、単に相手を疑っているから講じるものではない。誠実なスタッフなら、説明責任を明確にし、自分自身の誤りによるリスクを減らしてくれる権限の分離はどんなものでも歓迎すべきである。
不正な内部者に、痕跡をきれいに隠す手間をかけさせれば、その行動を遅らせ、現行犯で捕まる確率を高められる。
幸い、人間には同僚と向き合うための十分に発達した信頼の仕組みが備わっており、そのおかげで内部者の裏切りは実際にはきわめてまれである。

\subsection{信頼と責任}
% src: Trust and Responsibility

信頼と責任は、協力を成り立たせる接着剤である。
ソフトウェアシステムは互いにますます接続され、依存し合うようになっているので、このグループのパターンは重要な道しるべになる。

\subsubsection{信頼への慎重さ}
% src: Reluctance to Trust

\begin{dfn}[信頼への慎重さ]\label{dfn:ch4-reluctance-to-trust}
\underline{信頼への慎重さ}（Reluctance to Trust）とは、「信頼は常に、確かな証拠に基づいた明示的な選択であるべきである」というパターンである。
\end{dfn}

このパターンは、信頼が貴重なものであると認めたうえで、懐疑的であることを促す。
ソフトウェアが登場する前から、犯罪者は人を信頼するという人間の自然な傾向につけ込んできた。作業員に変装して侵入したり、いんちきの万能薬を売りつけたり、ほかにもありとあらゆる詐欺を働いたりしてきたのである。
信頼への慎重さが教えるのは、制服を着ているからといって本物だと思い込まないこと、そして FBI を名乗る電話の相手がペテン師かもしれないと疑うことである。
ソフトウェアでいえば、インストールの前にコードの真正性を確認することや、認可の前に強力な認証を求めることが、このパターンに当たる。

\begin{exa}[HTTP クッキー]\label{exa:ch4-http-cookies}
第11章で詳しく説明するように、HTTP クッキーの使い方はこのパターンの好例である。
Web サーバーはクライアントへの応答でクッキーを設定し、クライアントがその後の要求でそれらのクッキーを送り返してくることを期待する。
しかしクライアントには実際にはそうする義務はないので、サーバーは常にクッキーを割り引いて受け取るべきである。クライアントが常に律儀にこの作業をこなすと全面的に信頼するのは、大きなリスクである。
\end{exa}

信頼への慎重さは、悪意がない場合でも重要である。
例えば、きわめて重要なシステムでは、全体が損なわれないよう、すべてのコンポーネントが同じ高い品質とセキュリティの基準を満たしていることを確かめなければならない。
きわめて重要な機能に匿名の開発者のコード（マルウェアを含んでいるかもしれないし、単にバグがあるだけかもしれない）を使うといった、信頼についてのまずい判断は、たちまちセキュリティを損なう。
このパターンは単純明快で筋が通っているが、人は生来信じやすく、信頼を控えるのは疑り深すぎるように感じられることもあるので、実践するのは難しい場合がある。

\subsubsection{セキュリティ責任の受容}
% src: Accept Security Responsibility

\begin{dfn}[セキュリティ責任の受容]\label{dfn:ch4-accept-security-responsibility}
\underline{セキュリティ責任の受容}（Accept Security Responsibility）とは、「すべてのソフトウェア専門家はセキュリティに対して責任を負う明確な義務があり、その姿勢を自らが作るソフトウェアに反映させるべきである」というパターンである。
\end{dfn}

例えば設計者は、システムに組み込む外部コンポーネントを審査する際に、セキュリティ要件も審査項目に含めるべきである。
また、2つのシステムの間のインターフェースでは、それぞれの側が果たすべき責任を明示的に引き受けると同時に、呼び出し側が果たすことを当てにしている保証を確認しておくべきである。

避けたいアンチパターンは、いつか問題にぶつかったときに、2人の開発者が互いに「セキュリティはそっちで面倒を見ていると思っていたから、こっちではやらなくていいと思っていた」と言い合う事態である。
大きなシステムでは、双方が互いに責任を押し付け合う状況に簡単に陥る。

\begin{exa}[フロントエンドとビジネスロジックの責任の空白]\label{exa:ch4-responsibility-gap}
コンポーネント A が信頼できない入力を受け取り（例えば、匿名のインターネット要求を受け取る Web フロントエンドサーバー）、場合によっては何らかの処理や形式の変換を施したうえで、コンポーネント B のビジネスロジックに渡す状況を考える。
コンポーネント A は、B が適切な検証とエラーチェックで信頼できない入力を安全に扱うものと仮定して、セキュリティの責任をまったく負わずに、すべての入力をそのまま通過させることもできる。
一方コンポーネント B から見れば、フロントエンドがすべての要求を検証し、安全な要求だけを B に渡してくるのだから、B はセキュリティの心配をまったくしなくてよいと思い込みやすい。
この状況の正しい扱い方は、明示的に取り決めることである。誰が要求を検証するのか、そして下流に対してどのような保証を（もしあるなら）提供するのかを決めておく。
安全を最大限に重視するなら、両方のコンポーネントがそれぞれ独立に入力を検証する多層防御（定義 \ref{dfn:ch4-defense-in-depth}）を検討すること。
\end{exa}

責任の空白がソフトウェアの設計者と利用者の間に生じる、あまりにもよくある別のケースを考えよう。
セキュアバイデフォルトのパターンの議論で挙げた設定の例、とくに安全でないオプションを用意する場合を思い出してほしい。
ある設定オプションが安全性に劣るとわかっているなら、設計者はそのオプションを本当に用意する必要があるのかを慎重に検討すべきである。
つまり、簡単だから、あるいは「いつか誰かが欲しがるかもしれない」からといって、ユーザーにオプションを与えてはならない。
それは、いつか誰かが知らずにはまる罠を仕掛けるのと同じである。
リスクを伴う設定にもっともな理由がある場合は、まず、問題を安全に解決できるよう設計を変える方法を検討すること。
それができず、要件そのものが本質的に安全でない場合は、設計者がユーザーに助言し、結果を知らないままユーザーがそのオプションを設定してしまうことのないよう守るべきである。
リスクを文書化し、脆弱性を補う緩和策を提案することも重要だが、それだけでなく、ユーザーには明確なフィードバックを返すべきである。理想を言えば、責任を放り出すだけの「よろしいですか？（詳しくはこちら：[リンク]）」というダイアログよりもましなものであるべきだ。

\begin{rem}[「よろしいですか？」ダイアログの何がいけないのか]\label{rem:ch4-are-you-sure}
% src（コラム）: What's Wrong with the "Are you Sure" Dialog?
原著者は個人的に、「よろしいですか？」ダイアログやそのたぐいは、ほとんど常に設計の失敗であり、しかもセキュリティを損なうことも多いと考えている。
原著者は、そうしたダイアログが問題に対する最善の解決策になっている例に、まだ一度も出会ったことがない。
セキュリティ上の影響がある場合、この慣行はセキュリティ責任の受容のパターンに反する。設計者がユーザーに責任を押し付けているうえ、ユーザーは「よろしい」かどうか確信が持てないかもしれないのに、ほかに選択肢がないからである。
念のため言っておくと、ここでの指摘は、\texttt{rm} コマンドの確認プロンプトや、誤って実行しないことが重要なその他の操作についての、ふつうの確認を対象にしたものではない。

こうしたダイアログは\underline{ダイアログ疲れ}（dialog fatigue）という現象の犠牲になりやすい。何かを片付けようとしている人は反射的にダイアログを閉じてしまい、ほぼ例外なく、ダイアログを助けではなく邪魔だと受け止める。
セキュリティを強く意識している原著者でさえ、こうしたダイアログが出ると「では、ほかにどうすればやりたいことができるのか」と思うという。
選択肢は、やりたいことを諦めるか、かなりのリスクを自分で負って進めるかのどちらかである。しかも、そのリスクが具体的に何なのかは推測するしかない。「詳しくはこちら」の説明があっても、よい解決策が示されていたためしがないからである。
この段階で「よろしいですか？」が伝えているのは、後悔するかもしれないことをしようとしている、ということだけである。具体的に何が起こりうるのかは説明されず、たぶん後戻りはできないということがほのめかされているにすぎない。

原著者は、こうしたダイアログに「いいえ、確かではないがとにかく進める」という新しい第3の選択肢を加え、ソフトウェアがユーザーの期待に応えられなかったのだから、それを重大なエラーとして記録することを提案している。
セキュリティがきわめて重要な状況では、この種の責任の押し付けを洗い出して精査し、いずれ解決すべき重大なバグとして扱うこと。
どう取り除くかは具体的な状況によるが、責任を引き受けるための一般的なやり方がいくつかある。
\begin{itemize}
\item これから具体的に何が起こるのか、そしてその理由をはっきり示す。
\item 文言は簡潔にしつつ、詳しい説明とよい文書へのリンク、またはそれに相当する参照先を示す。
\item あいまいな文言（「本当にこれを実行しますか？」）を避け、操作の対象が具体的に何なのかを示す（重要な情報をダイアログボックスで隠してはならない）。
\item 二重否定や紛らわしい言い回し（「いいえ」と答えるとその操作が選ばれてしまう「本当に戻りますか？」など）は決して使わない。
\item 可能なら、元に戻すオプションを用意する。最近よく見かけるよいパターンは、大きな操作の後に、元に戻す機能をさりげなく提示するものである。
\item 元に戻す方法がない場合は、リンク先の文書で回避策を示すか、確信が持てないなら事前にデータをバックアップするよう勧める。
\end{itemize}
こうした\underline{ホブソンの選択}（Hobson's choice、選択の余地のない選択）をできるだけ減らし、理想的には、責任を引き受けるだけのノウハウを持つ専門の管理者が使う場面に限るよう努めよう。
\end{rem}

\subsection{アンチパターン}
% src: Anti-Patterns

\begin{quote}
「他人の災難の中に、自分が避けるべき災いを見ることを学べ。」 ---Publilius Syrus
\end{quote}

熟練者の仕事ぶりを見て覚えるのがいちばんの技能もあるが、もう一つの重要な学び方は、他人が過去に犯した誤りを避けることである。
化学の初学者は、酸を薄めるときは必ず水の入った容器に酸を加え、決して逆にしてはならないと教わる。大量の酸があるところに水を加えると、最初の一滴の水が激しく反応して大量の熱を出し、その熱で水が一瞬で沸騰して、水と酸が爆発的に飛び散るおそれがあるからである。
この教訓を、身をもって真似して学びたい人はいない。同じ精神で、ここではセキュリティのために避けたほうがよいアンチパターンをいくつか紹介する。

以下の短い項で、ソフトウェアセキュリティのアンチパターンをいくつか挙げる。
これらは一般にセキュリティ上のリスクを伴うので避けるのが最善だが、それ自体が実際の脆弱性というわけではない。
前の節で扱った名前付きのパターンは一般に通用する用語であるのに対し、ここで挙げるものには定着した名前がないものもあるので、便宜上、内容を表す呼び名を付けている。

\subsubsection{混乱した使節}
% src: Confused Deputy

\begin{dfn}[混乱した使節]\label{dfn:ch4-confused-deputy}
\underline{混乱した使節}（Confused Deputy）の問題は、多くのソフトウェア脆弱性の核心にある根本的なセキュリティ上の課題である。
\end{dfn}

すべてのアンチパターンの母とさえ言えるかもしれない。
その名前と意味を説明するには、短い物語から始めるのがよい。

\begin{exa}[ノーマン・ベイツの逮捕]\label{exa:ch4-norman-bates}
判事が令状を出し、保安官代理（deputy）にノーマン・ベイツを逮捕するよう命じたとする。
保安官代理はノーマンの住所を調べ、そこに住んでいる男を逮捕する。
男は何かの間違いだと言い張るが、保安官代理はそんな言い訳は聞き飽きている。
この物語のどんでん返しは（『サイコ』とは何の関係もない）、ノーマンが捕まることを見越して、何年も前から偽の住所を使っていたことである。
このごまかしに惑わされた保安官代理は、逮捕の権限を誤って行使してしまった。ノーマンは保安官代理を手玉に取り、保安官代理に正当に与えられた権限を、まんまと自分の邪悪な目的のために使わせたのである。
（\underline{スワッティング}（swatting）、すなわち緊急事態を虚偽通報して無実の被害者のもとに警察を差し向ける卑劣な犯罪は、混乱した使節の問題の完璧な例であるが、原著者はそうした悲しい話を詳しく語りたくはなかったという。）
\end{exa}

この問題のよくある例としては、ユーザーランドのコードから呼び出されたときのカーネルや、インターネットから呼び出されたときの Web サーバーがある。
呼び出される側が\underline{使節}（deputy）である。高い権限を持つコードが、低い権限の呼び出し元に代わって何かを行うために呼び出されるからである。
このリスクは信頼境界を越えることから直接生じる。脅威モデリングで信頼境界にとりわけ強い関心が向けられるのはそのためである。
後の章では、使節を混乱させるさまざまな方法を扱う。いくつか挙げるだけでも、バッファオーバーフロー、不十分な入力検証、クロスサイトリクエストフォージェリ（CSRF）攻撃などがある。
人間の保安官代理は本能や過去の経験、その他の手がかり（常識も含む）に頼れるが、ソフトウェアはそうはいかない。必要な予防措置をすべて見越して設計・実装されていない限り、意図していなかったことをするよう簡単にだまされてしまう。

% src: Intention and Malice
\underline{意図と悪意：}
第1章のおさらいをすると、ソフトウェアが信頼に足るものであるための要件は2つある。誠実であり、かつ品質の高い製品を届ける能力があると信頼できる人々によって作られていることである。
2つの条件の違いは意図にある。
ノーマン・ベイツの逮捕で問題だったのは、保安官代理が不正を働いたことではなく、方針どおりに被逮捕者の身元をきちんと確認しなかったことである。
もちろん、コードが命令に背いたり怠けたりすることはないが、まずく書かれたコードは、意図されたのとは別の動作を簡単にしてしまう。
だまされやすい多くのコンピュータ利用者、ときには技術に精通したソフトウェア専門家でさえ、悪意あるソフトウェアを信頼するようだまされることはある。しかし多くの攻撃は、正当に信頼されているが、たまたま欠陥を抱えたソフトウェアの中の混乱した使節を悪用することで成り立っている。

混乱した使節の脆弱性は、元の要求のコンテキストがコードの途中で失われたとき、例えば要求者のアイデンティティがもう参照できなくなったときに生じることが多い。
この種の混乱は、高い権限の呼び出しと低い権限の呼び出しの両方から使われる共通のコードで、とくに起こりやすい。
図 \ref{fig:ch4-confused-deputy} は、そのような呼び出しの様子を示している。

% 図 4-3: An example of the Confused Deputy anti-pattern
%   （画像のみ。下図は直後の本文（中央に Deputy、右に High、Low からの呼び出しは信頼境界を越える、Deputy は Utility を使う）から描いた。
%    Low を左に置いたのは「右の High」との対比による読み取りである。）
\begin{figure}[H]
\centering
\begin{tikzpicture}[
  bx/.style={draw, rounded corners, align=center, minimum width=2.2cm, minimum height=1.0cm, font=\small},
  ar/.style={-{Stealth[length=2.5mm]}, thick},
  lb/.style={font=\scriptsize, fill=white, inner sep=1pt}
]
\node[bx] (L) at (0,0) {Low\\（低い権限）};
\node[bx] (D) at (4.5,0) {Deputy};
\node[bx] (H) at (9,0) {High\\（高い権限）};
\node[bx] (U) at (4.5,-2.3) {Utility};
\draw[ar] (L) -- (D);
\draw[ar] (H) -- (D);
\draw[ar] (D) -- (U);
\draw[dashed, thick] (2.25,1.3) -- (2.25,-1.3) node[below, font=\scriptsize] {信頼境界};
\end{tikzpicture}
\caption{混乱した使節のアンチパターンの例}\label{fig:ch4-confused-deputy}
\end{figure}

中央の Deputy のコードは、低い権限のコードと高い権限のコードの両方のために仕事をする。
右の High から呼び出されたとき、Deputy は信頼できる呼び出し元のために、危険を伴いうる操作を行うこともある。
Low からの呼び出しは信頼境界を越えてくるので、Deputy は低い権限の呼び出し元にふさわしい安全な操作だけを行うべきである。
内部の実装では、Deputy は仕事をするのにサブコンポーネントの Utility を使う。
Utility のコードには高い権限の呼び出し元と低い権限の呼び出し元という区別がないため、低い権限の呼び出し元にはできないはずの危険を伴う操作を、Deputy に頼まれるまま誤って行ってしまうおそれがある。

% src: Trustworthy Deputy
\underline{信頼できる使節：}
信頼できる使節になる方法を、どこに危険があるかを考えるところから分析してみよう。
混乱の可能性が始まるのは信頼境界であることを思い出してほしい。混乱した使節を攻撃する目的は、使節の高い権限を利用することだからである。
使節が要求とその要求者を把握しており、適切な認可チェックが行われている限り、すべてはうまくいくはずである。

先ほどの Deputy のコードの例を思い出そう。そこでは、Low から呼び出されたときに信頼境界を考慮していなかった下位の Utility のコードで問題が起きた。
ある意味で、Deputy は知らないうちに Utility を混乱した使節にしてしまったのである。
Utility がもともと低い権限の呼び出し元から身を守ることを想定していなかったのなら、Deputy がだまされないよう Utility を徹底的に保護するか、Utility が低い権限の呼び出しを見分けられるよう修正する必要があるかもしれない。

混乱した使節のもう一つのよくある失敗は、要求に応じて行う動作の中で起きる。
\underline{データ隠蔽}（data hiding）は、実装が使う仕組みを抽象化の背後に隠す基本的なデザインパターンであり、使節はその仕組みを直接操作できるが、要求者にはできない。
例えば、使節が要求の副作用として情報をログに記録することがあっても、要求者はそのログにアクセスできない。
使節にログを書かせるということは、要求者が使節の権限を利用しているということなので、意図しない副作用に注意することが重要である。
要求者が不正な形式の文字列を使節に渡すことができ、その文字列がログに流れ込んでデータを壊し、判読できなくしてしまうなら、それは実質的にログを消去する混乱した使節攻撃である。
この場合の防御は、要求者からの文字列がログに流れ込みうることに気づき、それがどんな影響を及ぼしうるかを考えたうえで、例えば入力検証を求めることから始まる。

第3章で触れた\underline{コードアクセスセキュリティ}（Code Access Security）モデルは、混乱した使節の脆弱性が生じるのを防ぐことを特に狙って設計されたものである。
低い権限のコードが高い権限の使節のコードを呼び出すと、実効的な権限はそれに合わせて引き下げられる。
使節がより大きな権限を必要とするときは、それを明示的に宣言しなければならず、そうすることで、自分が低い権限のコードの求めに応じて動作していることを認めることになる。

まとめると、信頼境界では混乱した使節にならないよう、信頼度の低いデータや低い権限からの呼び出しを注意深く扱うこと。
必要に応じて認可を完全にチェックできるよう、タスクを実行する過程全体を通じて、要求に関連付けられたコンテキストを保持しておくこと。
副作用によって要求者が自分の権限を超えられないよう注意すること。

\subsubsection{信頼の逆流}
% src: Backflow of Trust

\begin{dfn}[信頼の逆流]\label{dfn:ch4-backflow-of-trust}
信頼度の低いコンポーネントが信頼度の高いコンポーネントを制御しているときに必ず存在するアンチパターンを、\underline{信頼の逆流}（Backflow of Trust）という。
\end{dfn}

\begin{exa}[個人のコンピュータからのリモート管理]\label{exa:ch4-admin-home-computer}
システム管理者が個人のコンピュータを使って企業のシステムをリモートで管理するのは、信頼の逆流の例である。
その人物自身は正当に権限を与えられ信頼されているが、自宅のコンピュータは企業の管理体制の外にあり、管理者権限を使うセッションの場にすべきではない。
\end{exa}

本質的には、これは起こるべくして起こるのを待っている、構造的な権限昇格だと考えることができる。

まともな人なら実生活でこのアンチパターンに陥ることはないだろうが、情報システムでは驚くほど見落とされやすい。
ここで重要なのは、コンポーネントにどれだけの信頼を与えているかではなく、コンポーネントがどれだけの信頼に値するかである。このことを忘れてはならない。
脅威モデリングで信頼境界を明示的に検討すれば、この種の潜在的な問題を浮かび上がらせることができる。

\subsubsection{サードパーティのフック}
% src: Third-Party Hooks

\begin{dfn}[サードパーティのフック]\label{dfn:ch4-third-party-hooks}
システム内のコンポーネントにあるフックが第三者に不当なアクセスを与えることを、\underline{サードパーティのフック}（Third-Party Hooks）という。
信頼の逆流のアンチパターンの一種である。
\end{dfn}

\begin{exa}[AI 予測コンポーネントのトンネル]\label{exa:ch4-ai-component}
何らかの専門的な処理を行うプロプライエタリなコンポーネントを組み込んだ、きわめて重要な業務システムを考える。
そのコンポーネントは、例えば高度な AI を使って将来の業務動向を予測するもので、機密の売上指標を取り込んで毎日予測を更新しているとしよう。
この AI コンポーネントは最先端のものなので、開発元の会社が毎日手入れをしなければならない。
ターンキーシステムのように手間なく使えるようにするには、管理インターフェースにアクセスするための、ファイアウォールを貫く直通のトンネルが必要になる。

これもまたゆがんだ信頼関係である。この第三者は、管理者の監督がまったく及ばないところから、企業システムの中枢に直接アクセスできるからである。
AI の提供元が不誠実だったり侵害されたりしていれば、社内データを簡単に持ち出したり、もっと悪いことをしたりできるが、こちらにはそれを知るすべがない。
\end{exa}

\begin{rem}[限定的なフック]\label{rem:ch4-limited-hooks}
限定的な種類のフックであればこの問題はなく、受け入れられる場合もある。
例えば、フックが自動更新の仕組みを実装していて、ソフトウェアの新しいバージョンをダウンロードしてインストールすることしかできないのであれば、相応の信頼があれば問題ないかもしれない。
\end{rem}

\subsubsection{パッチ適用不能なコンポーネント}
% src: Unpatchable Components

広く使われているコンポーネントであれば、誰かが脆弱性を見つけるかどうかではなく、いつ見つけるかの問題であることが、ほぼ例外なく言える。
そうした脆弱性が公に知られたら、そのコンポーネントがどの攻撃対象領域からも完全に切り離されていない限り、速やかにパッチを当てる必要がある。

\begin{dfn}[パッチ適用不能なコンポーネント]\label{dfn:ch4-unpatchable-components}
システム内のパッチを当てられないコンポーネントを\underline{パッチ適用不能なコンポーネント}（Unpatchable Components）という。
そのようなコンポーネントは、いずれ恒久的な負債になる。
\end{dfn}

ソフトウェアがプリインストールされたハードウェアコンポーネントはパッチを適用できないことが多いが、発行元がサポートを打ち切ったり廃業したりしたソフトウェアも、実質的に同じである。
実際には、実質的にパッチを適用できないソフトウェアはほかにもたくさんある。
\begin{itemize}
\item バイナリ形式でしか提供されていない、サポート対象外のソフトウェア
\item 時代遅れのコンパイラやその他の依存関係でビルドされたコード
\item 経営判断によって引退させられたコード
\item 訴訟に巻き込まれたコード
\item ランサムウェアによる侵害で失われたコード
\item そして驚いたことに、COBOL のように、古すぎて今では経験のあるプログラマーが不足している言語で書かれたコード
\end{itemize}
主要なオペレーティングシステムの提供元は、通常、一定期間はサポートとアップグレードを提供するが、その期間が過ぎるとそのソフトウェアは実質的にパッチ適用不能になる。
更新可能なソフトウェアであっても、提供元が適時にリリースを出せなければ、実質的には変わらないかもしれない。
必要なときにすぐ更新できる確信が持てないものを使って、運を天に任せるようなことをしてはならない。

\newpage

%% =======================================================
%%  part: ch5
%% =======================================================
%% =======================================================
%%  part: ch5
%% =======================================================
\section{暗号}
% src: 5: Cryptography

この章では、\cite{Kohnfelder2021} の第5章に沿って、暗号の基本的な道具立てとその正しい使い方を扱う。

\begin{quote}
「暗号はたいてい、破られるのではなく、迂回される。」\\
\hfill --- Adi Shamir
\end{quote}

高校時代、原著者は自動車教習の授業で落第しかけた。
公立学校に運転を教える予算があり、ガソリンにまだ鉛が入っていた（この素晴らしいアイデアを脅威モデリングした者は誰もいなかった）ほど昔の話である。
最初のうち、運転はまるでうまくいかなかった。
とくに覚えているのは、マニュアル車の Volkswagen Beetle のハンドルを初めて握った日のことと、助手席に座った体育教師の無表情な顔にありありと浮かんだ不安である。
下り坂でクラッチを踏み込むと、思ったように減速するどころか逆に車が加速することを、すぐに思い知った。
ところがその失敗を境に何かがかみ合い、突然運転できるようになった。
この思いがけない展開に、教師は驚きと安堵を隠さなかった。
振り返ってみると、マニュアル車を運転する手応えのおかげで車とより直接につながっている感覚が得られ、初めて直感で運転できるようになったのだろう、と原著者は考えている。

自動車教習が生徒に車の安全な運転の仕方を教えるのと同じように、この章では暗号の仕組みの詳細には立ち入らず、正しい使い方を論じながら暗号の基本的な道具立てを紹介する。
数学が得意でない人にも暗号をわかりやすくするため、この章では数学を避ける。
ただし一箇所だけ、あまりに巧妙なので入れずにはいられなかった例外がある。

この話題の扱い方としては型破りだが、重要なやり方でもある。
暗号の道具が十分に活用されていないのは、まさに暗号が参入障壁の高い専門家の領域とみなされるようになったからである。
現代のライブラリには暗号機能が用意されているが、効果を発揮させるには、開発者が使い方（と正しい使い方）を知っていなければならない。
原著者は、この章が暗号の使いどころについての勘を養う踏み台になることを願っている。
具体的な用途に応じて、必要なだけ自分で調べて補ってほしい。

%----------------------------------------------------------------------
\subsection{暗号の道具}
% src: Crypto Tools

現代の暗号の多くは根本的に純粋数学に基づいているので、正しく使えば確かに機能する。
これはアルゴリズムが絶対に破られないと証明されているという意味ではなく、破るには数学上の大きなブレークスルーが必要になるという意味である。

暗号はさまざまなセキュリティの道具を提供してくれるが、効果を上げるには考えて使わなければならない。
本書で繰り返し勧めているように、完全な解決策を備えた高品質なコードのライブラリに頼るべきである。
適切な抽象度のインターフェースを持つライブラリを選び、そのライブラリが何をしているのかを完全に理解しておくことが重要である。

暗号の歴史や背後にある数学は興味深いが、セキュアなソフトウェアを作るという目的にとっては、現代の道具箱はほんの一握りの基本的な道具でできている。
基本的な暗号のセキュリティ機能と、それぞれが何をするのか、それぞれの安全性が何に依存するのかを次に挙げる。

\begin{dfn}[基本的な暗号の道具]\label{dfn:ch5-crypto-tools}
\begin{itemize}
\item \underline{乱数}（random number）は、パディングやナンスとして役に立つが、\emph{予測不可能である場合に限る}。
\item \underline{メッセージダイジェスト}（message digest、\underline{ハッシュ関数}（hash function））は、データの指紋として働くが、\emph{衝突に対して耐性がある場合に限る}。
\item \underline{対称暗号}（symmetric encryption）は、当事者が共有する秘密鍵に基づいてデータを秘匿する。
\item \underline{非対称暗号}（asymmetric encryption）は、受信者が知っている秘密に基づいてデータを秘匿する。
\item \underline{デジタル署名}（digital signature）は、署名者だけが知っている秘密に基づいてデータを認証する。
\item \underline{デジタル証明書}（digital certificate）は、ルート証明書への信頼に基づいて署名者を認証する。
\end{itemize}
\end{dfn}

この章の残りでは、これらの道具と用途をさらに詳しく扱う。

%----------------------------------------------------------------------
\subsection{乱数}\label{subsec:ch5-random}
% src: Random Numbers

人間の頭にとって、ランダム性という概念をつかむのは難しい。
セキュリティの目的では、乱数の最も重要な性質である\emph{予測不可能性}に焦点を当てればよい。
後で見るように、乱数が決定的に重要になるのは、予測しやすいパスワードが弱いのと同じ理屈で、攻撃者に正しく推測されては困る場面である。
乱数は認証、ハッシュ、暗号化、鍵生成に使われ、そのどれもが予測不可能性を前提にしている。
以下の項では、ソフトウェアで使える 2 種類の乱数について、予測しやすさがどう違うのか、どちらをいつ使うべきかを述べる。

\subsubsection{疑似乱数}
% src: Pseudo-Random Numbers

\begin{dfn}[疑似乱数生成器]\label{dfn:ch5-prng}
決定論的な計算によって、無限に続く乱数列のように見えるものを生成するものを、\underline{疑似乱数生成器}（pseudo-random number generator, PRNG）という。
\end{dfn}

PRNG の出力は人間のパターン検出能力をたやすく上回るが、解析や敵対的なソフトウェアにかかれば、PRNG のふるまいは簡単にまねられてしまう。
PRNG は予測可能なので、セキュリティの文脈で使うには不適格である。

一方、疑似乱数の計算は非常に高速なので、セキュリティと関係のない幅広い用途には理想的である。
例えば、Monte Carlo シミュレーションを実行するときや、A/B テストのために Web ページのデザインの変種をランダムに割り当てるときは、PRNG を使えばよい。
万一誰かにアルゴリズムを予測されても、実害は何もないからである。

\begin{exa}[円周率の数字列]\label{exa:ch5-pi-digits}
疑似乱数がなぜ真にランダムではないのかを理解するために、次の数字列を考える。
\begin{verbatim}
94657640789512694683983525957098258226205224894077267194782684826
\end{verbatim}
この列はランダムだろうか。
1 と 3 が比較的少なく、2 がやけに多いが、真にランダムな数でも一様分布からこの程度ずれることは十分にありうる。
しかし、この列がどれほどランダムに見えても、種明かしを知っていれば次の数字を予測するのは簡単である。
そして、透明な設計（Transparent Design）のパターンが警告するように、自分のやり方を秘密にしておけると思い込むのは危険である。
実際、この数字列を Web で検索すればすぐに、円周率の小数点以下 200 桁目からの数字で、次の数字は 0147 だとわかる。

円周率は無理数なので、その数字は統計的に正規な分布をしており、日常的な意味では完全にランダムである。
その一方で、円周率は簡単に計算できるよく知られた数なので、この列は完全に予測可能であり、セキュリティの目的には使えない。
\end{exa}

\subsubsection{暗号論的に安全な疑似乱数}
% src: Cryptographically Secure Pseudo-Random Numbers

\begin{dfn}[暗号論的に安全な疑似乱数生成器]\label{dfn:ch5-csprng}
セキュリティのために乱数ビットが必要なとき、PRNG の欠点を補うために現代のオペレーティングシステムが提供している関数を、\underline{暗号論的に安全な疑似乱数生成器}（cryptographically secure pseudo-random number generator, CSPRNG）という。
CSRNG や CRNG と書かれることもあるが、肝心なのは「C」の部分で、暗号（crypto）に使っても安全だという意味である。
\end{dfn}

名前に「疑似」が入っているのは、完全なランダム性には届かないかもしれないと認めているからだが、専門家は、実用上はどんな目的にも安全に使えるほど予測不可能だと判断している。

セキュリティがかかっている場面では、この種の乱数生成器を使うべきである。
言い換えると、ランダムであるはずの数の値を予測できたと仮定したときにセキュリティが弱まるのであれば、CSPRNG を使うべきである。
本書で取り上げる乱数のセキュリティ上の用途は、すべてこれに当てはまる。

\begin{dfn}[ハードウェア乱数生成器]\label{dfn:ch5-hrng}
真にランダムなデータは、定義上、アルゴリズムで生成されるのではなく、予測不可能な物理過程から得られる。
そのような装置を\underline{ハードウェア乱数生成器}（hardware random number generator, HRNG）といい、\underline{エントロピー源}（entropy source）ともいう。
\end{dfn}

\begin{exa}[エントロピー源]\label{exa:ch5-entropy-sources}
放射性崩壊が起きるタイミングはランダムなので、Geiger カウンターは HRNG になりうる。
HRNG は現代の多くのプロセッサーに組み込まれているほか、後付けのハードウェアとして購入することもできる。
ソフトウェアもエントロピーを提供できる。たいていは、ディスクアクセス、キーボードやマウスの入力イベント、外部の相手との複雑なやり取りに左右されるネットワーク送信といった事象のタイミングからエントロピーを取り出す。
\end{exa}

\begin{exa}[溶岩ランプによるエントロピー]\label{exa:ch5-lava-lamps}
ある大手インターネット企業は、溶岩ランプ（lava lamp）をずらりと並べて、色鮮やかなランダム入力を生成している。
しかし、この手法の脅威モデルを考えてみよう。
この企業は溶岩ランプを本社オフィスの、それも受付エリアに展示しているので、攻撃者になりうる者がランプの状態を観察し、エントロピー源について見当をつけられるかもしれない。
実際には、溶岩ランプは舞台裏にある（おそらく）もっと普通のエントロピー源にエントロピーを上乗せしているだけなので、この展示が原因で企業のシステムがたやすく侵害されるリスクは緩和されている。
\end{exa}

\begin{rem}[CSPRNG の速度]\label{rem:ch5-csprng-speed}
エントロピー源がランダム性を生み出すには時間がかかるので、あまりに多くのビットをあまりに速く要求すると、CSPRNG は極端に遅くなる。
これが安全なランダム性の代償であり、確実に高速な代替手段として PRNG に重要な存在意義がある理由でもある。
高速な HRNG がない限り CSPRNG は控えめに使い、スループットが問題になる場面では、それがボトルネックにならないことをテストで確かめるべきである。
\end{rem}

%----------------------------------------------------------------------
\subsection{メッセージ認証コード}
% src: Message Authentication Codes

\begin{dfn}[メッセージダイジェスト]\label{dfn:ch5-message-digest}
一方向関数（one-way function）を使ってメッセージから計算される固定長の値を、\underline{メッセージダイジェスト}（message digest）または\underline{ハッシュ}（hash）という。
\end{dfn}

つまり、一意なメッセージにはそれぞれ固有のダイジェストがあり、どんな改ざんをしてもダイジェストの値が変わる。
一方向であることは重要である。ダイジェストの計算を逆にたどれないので、攻撃者は、たまたま同じダイジェストになる別のメッセージを見つけられないからである。
ダイジェストが一致すれば、メッセージの内容が改ざんされていないとわかる。

\begin{dfn}[衝突、衝突攻撃、原像攻撃]\label{dfn:ch5-collision}
2 つの異なるメッセージが同じダイジェストになることを、\underline{衝突}（collision）という。
攻撃者が同じダイジェスト値になる 2 つの異なる入力を見つけたとき、\underline{衝突攻撃}（collision attack）が成功したという。
ダイジェスト関数に対する最も壊滅的な攻撃は\underline{原像攻撃}（preimage attack）であり、特定のダイジェスト値が与えられたときに、攻撃者がその値になる入力を見つけられるというものである。
\end{dfn}

ダイジェストは大きなデータの塊を固定長の値に写すので、ありうるメッセージの数はダイジェスト値の数より多く、衝突は避けられない。
良いダイジェスト関数の条件は、衝突を見つけるのが極めて難しいことである。

暗号論的に安全なダイジェストアルゴリズムは強力な一方向関数で、衝突が起きる見込みは極めて小さいので、衝突は決して起きないと仮定してよい。
ダイジェストの力を生かすには、この仮定が欠かせない。
この仮定があるからこそ、2 つのダイジェストが等しいかどうかを比べることが、実質的にメッセージ全体を比べることになるのである。
2 つの指紋（fingerprint。ダイジェストのくだけた呼び名でもある）を比べて、同じ指で付けられたものかどうかを判断するようなものだと考えればよい。

もし誰もがあらゆる用途に同じダイジェスト関数を使っていたら、攻撃者はその関数を集中的に研究・解析し、いずれ衝突やその他の弱点を見つけるかもしれない。
これを防ぐ方法の一つが、鍵付きハッシュ関数を使うことである。

\begin{dfn}[鍵付きハッシュ関数とメッセージ認証コード]\label{dfn:ch5-mac}
ダイジェストの計算を変化させる秘密鍵のパラメータを追加で受け取るハッシュ関数を、\underline{鍵付きハッシュ関数}（keyed hash function）という。
256 ビットの鍵を取る鍵付きハッシュ関数は、実質的に $2^{256}$ 個の異なる関数からなる族である。
ハッシュ関数の鍵が秘密である限り攻撃者には偽造できないので、これらの関数は\underline{メッセージ認証コード}（message authentication code, MAC）とも呼ばれる。
\end{dfn}

つまり、一意な鍵を使えば、自分専用にカスタマイズしたダイジェスト関数が手に入る。

\subsubsection{MAC による改ざんの防止}
% src: Using MACs to Prevent Tampering

MAC は、攻撃者によるデータの改ざんを防ぐためによく使われる。

\begin{exa}[Alice と Bob の MAC 通信]\label{exa:ch5-alice-bob-mac}
Alice が公開チャネルを通じて Bob にメッセージを送りたいとする。
2 人はある秘密鍵を内密に共有している。
盗聴は気にしないのでデータを暗号化する必要はないが、偽のメッセージを見逃すと困る。
邪悪な Mallory は回線上の通信を改ざんできるが、鍵は知らないとする。
Alice は鍵を使って MAC を計算し、メッセージごとに添えて送る。
Bob は通信を受け取ると、受信したメッセージの MAC を計算し、Alice が添えて送った MAC と比べる。一致しなければ、偽物として無視する。
\end{exa}

この仕組みは、賢い Mallory に対してどれほど安全だろうか。
まず、すぐに思いつく攻撃を考える。
\begin{itemize}
\item Mallory がメッセージを改ざんすると、MAC がメッセージのダイジェストと一致しなくなる（そして Bob は無視する）。
\item Mallory が MAC を改ざんすると、今度はそれがメッセージのダイジェストと一致しなくなる（そして Bob は無視する）。
\item Mallory がまったく新しいメッセージをでっち上げても、MAC を計算する手段がない（そして Bob は無視する）。
\end{itemize}
しかし、防がなければならないケースがもう一つある。
Mallory に残されたもう一つの抜け道と、その防ぎ方がわかるだろうか。

\subsubsection{リプレイ攻撃}
% src: Replay Attacks

先に述べた MAC 通信の方式にはまだ問題が残っている。これは、決意の固い攻撃者を相手に暗号の道具を使いこなすのがいかに難しいかを示している。

\begin{exa}[ウィジェットの注文のリプレイ]\label{exa:ch5-widget-replay}
Alice は毎日 Bob に、翌日に何個のウィジェットを配送してほしいかを伝える注文を送るとする。
Mallory はこの通信を観察し、Alice が送るメッセージと MAC の組を集めている。
Alice は 1 日目に 3 個、翌日に 5 個のウィジェットを注文する。
3 日目に、Alice は 10 個のウィジェットを注文する。
ここで Mallory は、Alice のメッセージを改ざんする方法を思いつく。
Mallory は Alice のメッセージを横取りし、1 日目のメッセージ（3 個のウィジェットを指定したもの）のコピーにすり替える。
このコピーには、Alice が親切にもすでに計算してくれていて、Mallory が以前に記録しておいた正しい MAC も付いている。
\end{exa}

\begin{dfn}[リプレイ攻撃]\label{dfn:ch5-replay-attack}
例 \ref{exa:ch5-widget-replay} のように、以前の正当なメッセージを再送して行う攻撃を、\underline{リプレイ攻撃}（replay attack）という。
\end{dfn}

安全な通信プロトコルはリプレイ攻撃に対処しなければならない。
問題は暗号が弱いことではなく、暗号の使い方が正しくなかったことである。
この例の根本的な問題は、3 個のウィジェットを注文する正当なメッセージがどれもまったく同じだという点にあり、これは本質的には予測可能性の問題である。

\subsubsection{安全な MAC 通信}
% src: Secure MAC Communications

Alice と Bob のプロトコルを修正してリプレイ攻撃を防ぐ方法はいくつかあるが、どれも、メッセージが常に一意で予測不可能になるようにすることが鍵になる。

\begin{exa}[タイムスタンプによる修正]\label{exa:ch5-timestamp}
単純な修正として、Alice はメッセージにタイムスタンプを入れ、Bob は古いタイムスタンプの付いたメッセージを無視する、と取り決める方法がある。
こうすれば、Mallory が月曜日の 3 個のウィジェットの注文を水曜日にリプレイしても、Bob はタイムスタンプを比べて気付き、不正を見破れる。
ただし、メッセージが頻繁に送られる場合や、ネットワークの遅延が大きい場合には、タイムスタンプはうまく機能しないかもしれない。
\end{exa}

\begin{dfn}[ナンス]\label{dfn:ch5-nonce}
一度だけ使うための乱数を\underline{ナンス}（nonce）という。
\end{dfn}

\begin{exa}[ナンスによる修正]\label{exa:ch5-nonce-fix}
リプレイ攻撃の脅威に対するもっとよい解決策は、Alice がメッセージを送る前に、毎回 Bob が Alice にナンスを送ることである。
Alice は、メッセージと一緒に、Bob のナンスと、メッセージとナンスを組み合わせたものの MAC を送り返す。
ナンスはやり取りのたびに変わるので、リプレイ攻撃は封じられる。
Mallory は Bob が送るナンスを横取りして書き換えることはできるが、違うナンスが返ってくれば Bob は気付く。
\end{exa}

この単純な例には、メッセージがウィジェットの個数だけの短いものだという問題もある。
リプレイ攻撃の危険とは別に、非常に短いメッセージは総当たり攻撃に弱い。
鍵付きハッシュ関数の計算時間はたいていメッセージのデータ長に比例するので、数ビットしかなければ計算はすぐに終わる。
Mallory がハッシュ関数の鍵を次々と速く試せるほど、正当なメッセージの MAC に一致する正しい鍵を当てやすくなる。
鍵がわかってしまえば、Mallory は Alice になりすましてメッセージを送れる。

短いメッセージのこの弱点は、メッセージを適切な最小長になるまでランダムなビットでパディングすれば緩和できる。
長くなったメッセージの MAC を計算するには時間がかかるが、そのおかげで Mallory の総当たり攻撃は現実的に不可能なほど遅くなるので、むしろ好都合である。
実際、まさにこの理由から、ハッシュ関数は計算コストが高いほうが望ましい。
Mallory にできるだけ苦労させるには、パディングが（予測可能な疑似乱数ではなく）本当にランダムであることが重要になる。

%----------------------------------------------------------------------
\subsection{対称暗号}
% src: Symmetric Encryption

\begin{dfn}[平文、暗号文、対称暗号]\label{dfn:ch5-symmetric-encryption}
暗号化はどれも、元のメッセージである\underline{平文}（plaintext）を、\underline{暗号文}（ciphertext）と呼ばれる判読できない形に変換して、メッセージを秘匿する。
\underline{対称暗号}（symmetric encryption）のアルゴリズムは、秘密鍵を使って、メッセージの変換を通信者だけが使えるようにカスタマイズする。通信者は事前に鍵を取り決めておかなければならない。
復号アルゴリズムは、同じ秘密鍵を使って暗号文を平文に戻す。
秘密鍵を知っていれば暗号化も復号もできるので、この可逆な変換を\underline{対称暗号}（symmetric cryptography）と呼ぶ。
\end{dfn}

この節では、対称暗号のアルゴリズムをいくつか紹介してセキュリティ上の性質を示し、安全に使うために必要な注意点を説明する。

\subsubsection{ワンタイムパッド}
% src: One-Time Pad

暗号学者はずっと昔に理想的な暗号アルゴリズムを発見している。
後で見るように実際にはほとんど使われないが、きわめて単純なので、暗号化を論じる出発点としてはうってつけである。

\begin{dfn}[ワンタイムパッド]\label{dfn:ch5-one-time-pad}
通信者が暗号鍵として秘密のランダムなビット列を事前に取り決めておき、送信者はメッセージと鍵の排他的論理和をとって暗号文を作り、受信者は暗号文と同じ鍵ビットの排他的論理和をとって平文のメッセージを復元する暗号アルゴリズムを、\underline{ワンタイムパッド}（one-time pad）という。
\end{dfn}

排他的論理和（$\oplus$）では、鍵ビットが 0 なら対応するメッセージビットはそのままで、鍵ビットが 1 ならメッセージビットが反転することを思い出そう。
% 図 5-1: Alice and Bob using one-time pad encryption（ワンタイムパッドによる暗号化と復号の簡単な例。具体的なビット列は本文にないため図は描かない）
原著の図 5-1 は、ワンタイムパッドによる暗号化と復号の簡単な例を図示している。

次のメッセージは、秘密鍵のビット列の続きの部分を使って暗号化する。
鍵を使い切ったら、通信者は何らかの方法で新しい秘密鍵を取り決める必要がある。
この鍵が一度限りなのにはちゃんとした理由があり、それはすぐ後で説明する。
鍵がランダムだとすれば、メッセージの各ビットはランダムに反転したりしなかったりするので、攻撃者は鍵を知らない限り元のメッセージを見分けようがない。
ちょうど半分のビットをランダムに反転させれば、メッセージは完璧に偽装される。大部分のビットがそのままだったり反転していたりすれば、平文が部分的にばれてしまうからである。
このように解析による攻撃には強いかもしれないが、この方法がめったに使われない理由はすぐにわかる。鍵の長さでメッセージの長さが制限されてしまうのである。

\begin{exa}[ワンタイムパッドの鍵の再利用]\label{exa:ch5-otp-key-reuse}
ワンタイムパッドの鍵を再利用してはいけない理由を考える。
Alice と Bob が同じ秘密鍵 $K$ を使って 2 つの異なる平文メッセージ $M_1$ と $M_2$ を暗号化したとする。
Mallory は両方の暗号文 $(M_1 \oplus K)$ と $(M_2 \oplus K)$ を傍受する。
Mallory が 2 つの暗号文の排他的論理和をとると、鍵は打ち消し合って消える。
どんな数も自分自身との排他的論理和は 0 になる（1 は 0 に反転し、0 はそのまま）からである。
結果は、2 つのメッセージを弱く暗号化したものになる。
\[
(M_1 \oplus K) \oplus (M_2 \oplus K) = (M_1 \oplus M_2) \oplus (K \oplus K) = M_1 \oplus M_2
\]
これで平文が直接わかるわけではないが、情報は漏れ始める。
鍵ビットが取り除かれたので、解析すればメッセージ内のパターンについての手がかりが得られうる。
例えば、どちらかのメッセージに 0 のビットが連続する部分があれば、もう一方のメッセージの対応するビットがそのまま漏れてしまう。
\end{exa}

鍵を一度しか使えないという制限は、ほとんどの用途にとって致命的である。
Alice と Bob は、どれだけのデータを暗号化することになるかを前もって知らないかもしれず、そうなると鍵の長さを決めようがないからである。

\subsubsection{Advanced Encryption Standard}
% src: Advanced Encryption Standard

\begin{dfn}[ブロック暗号と AES]\label{dfn:ch5-aes}
長いメッセージはブロックの大きさの塊に分割し、短いメッセージはブロックの残りをランダムなビットでパディングして暗号化する方式を、\underline{ブロック暗号}（block cipher）という。
\underline{Advanced Encryption Standard}（AES）は、よく使われている現代の対称暗号のブロック暗号アルゴリズムである。
AES は、通常 256 ビットの秘密鍵を使って、128 ビットのデータのブロックを暗号化する。
Alice は、Bob が復号に使うのと同じ、取り決め済みの秘密鍵でデータを暗号化する。
\end{dfn}

ありうる弱点を考えてみよう。
Alice が時間をおいてまったく同じメッセージブロックを Bob に送ると、暗号文もまったく同じになり、賢い Mallory はその繰り返しに気付く。
Mallory がメッセージの意味を解読できなくても、これは緩和が必要な重大な情報漏えいである。
この通信はリプレイ攻撃にも弱い。Alice が同じ暗号文を再送して同じ平文のメッセージを伝えられるなら、Mallory にも同じことができるからである。

\begin{dfn}[ECB モード]\label{dfn:ch5-ecb}
同じメッセージを毎回同じように暗号化する方式を、\underline{ECB モード}（electronic code book mode）という。
\end{dfn}

ECB モードはリプレイ攻撃に弱いので、たいていは悪い選択である。
この問題を避けるには、後続のブロックにフィードバックなどの違いを持ち込む別のモードを使い、暗号文が先行するブロックの内容や列の中での位置によって変わるようにすればよい。
こうすれば、平文のブロックが同じでも、暗号文はまったく違うものになる。
ただし、データストリームをブロック単位で連鎖的に暗号化する方式には利点がある一方で、正しく暗号化・復号するには、通信者が順序の文脈を保持し続けなければならない。
そのため、どの暗号モードを選ぶかは、たいていアプリケーションごとの事情によって決まる。

\subsubsection{対称暗号の利用}
% src: Using Symmetric Cryptography

対称暗号は、正しく使えば高速かつ安全なので、現代の暗号化の主力になっている。
暗号化によって、安全でないチャネルを流れるデータや、ストレージに保存されたデータ（data at rest）を保護できる。
使い始める前に、いくつかの基本的な制約を考えておくことが重要である。

\begin{rem}[対称暗号の基本的な制約]\label{rem:ch5-symmetric-limitations}
\begin{itemize}
\item \underline{鍵の確立}（key establishment）：暗号アルゴリズムは秘密鍵が事前に取り決められていることを前提にしているが、その鍵をどうやって確立するかは定めていない。
\item \underline{鍵の秘匿}（key secrecy）：暗号化が有効かどうかは、必要なときには鍵を使えるようにしつつ、鍵の秘密を守れるかどうかに完全にかかっている。
\item \underline{鍵の大きさ}（key size）：秘密鍵は大きいほど強い（理論上はワンタイムパッドが理想である）が、大きな鍵は管理にコストがかかり、扱いにくくなる。
\end{itemize}
\end{rem}

対称暗号は本質的に共有の秘密鍵に頼っているので、Alice と Bob が直接会って確実に鍵を交換できない限り、準備するのが難しい。
この制約を克服するのが非対称暗号であり、インターネットでつながった世界のニーズに合う、驚くほど便利な新しい能力をもたらす。

%----------------------------------------------------------------------
\subsection{非対称暗号}
% src: Asymmetric Encryption

非対称暗号は直感にまったく反する暗号化の形であり、そこにこそ力がある。
対称暗号では Alice も Bob も同じ鍵でメッセージを暗号化・復号できるが、非対称暗号では、Bob は自分では復号できない秘密のメッセージを Alice に送れる。
つまり Bob にとって暗号化は一方向関数であり、その関数を逆にたどる（すなわちメッセージを復号する）ための秘密を知っているのは Alice だけである。

\begin{dfn}[公開鍵と秘密鍵]\label{dfn:ch5-key-pair}
非対称暗号では鍵の対を使う。暗号化のための\underline{公開鍵}（public key）と、復号のための\underline{秘密鍵}（private key）である。
\end{dfn}

ここでは、Bob（あるいは世界中の誰でも）が Alice に暗号化されたメッセージを送る方法を述べる。
双方向の会話では、Alice は Bob のまったく別の鍵の対を使い、同じ手順で返信する。
2 つの鍵による変換は互いに逆関数になっているが、一方の鍵を知っていても他方の鍵を割り出す手がかりにはならない。
そのため、一方の鍵を秘密にしておけば、その計算ができるのは自分だけになる。
この非対称性のおかげで、Alice は鍵の対を作って一方の鍵（公開鍵）を世界に公開し、対応する秘密鍵を持つ自分だけが復号できるメッセージを、誰でも暗号化できるようにできる。
これは革命的である。秘密を知っていることに基づく、Alice だけの能力が得られるからである。
この能力によって何が可能になるかを、以下で見ていく。

非対称暗号のアルゴリズムは数多くあるが、暗号の道具として使い方を理解するうえで、その数学的な詳細は重要ではない。重要なのはセキュリティ上の意味を理解することである。
ここでは、数学的に最も単純な元祖である RSA に焦点を当てる。

\subsubsection{RSA 暗号系}
% src: The RSA Cryptosystem

原著者は MIT で、幸運にも RSA 暗号系の発明者のうち 2 人と一緒に仕事をし、学士論文では非対称暗号でセキュリティをどう向上させられるかを探究した。
以下の簡略化した説明は RSA の元の論文に沿っているが、現代の実装は（ここで立ち入る必要のないさまざまな技術的理由から）もっと込み入っている。

\begin{dfn}[RSA 暗号系]\label{dfn:ch5-rsa}
\underline{RSA} の核となる考え方は、2 つの大きな素数を掛け合わせるのは簡単だが、その積から元の素数に分解するのは事実上不可能だということである。
\begin{enumerate}
\item まず、ランダムな大きな素数の対を選び、秘密にしておく。
\item 次に、その素数の対を掛け合わせる。結果を $N$ とする。
\item $N$ から、一意な鍵の対を計算できる。それぞれの鍵を $N$ と組み合わせると、互いに逆関数になっている 2 つの関数 $D$ と $E$ を計算できる。すなわち、$x<N$ を満たす任意の正の整数 $x$ に対して $D(E(x))=x$ であり、$E(D(x))=x$ でもある。
\item 最後に、鍵の対の一方を秘密鍵とし、もう一方を対応する公開鍵として $N$ とともに世界に公開する。
\end{enumerate}
秘密鍵と元の 2 つの素数を秘密にしておく限り、関数 $D$ を効率よく計算できるのは自分だけである。
\end{dfn}

\begin{exa}[RSA による暗号化と復号]\label{exa:ch5-rsa-encrypt}
Bob が Alice 宛てにメッセージを暗号化し、Alice がそれを復号する手順は次のとおりである。
ここで関数 $E_A$ と $D_A$ は、それぞれ Alice の公開鍵と秘密鍵に $N$ を組み合わせたものに基づく。
\begin{itemize}
\item Bob は Alice の公開鍵を使って、メッセージ $M$ から Alice 宛ての暗号文 $C$ を作る：$C = E_A(M)$。
\item Alice は自分の秘密鍵を使って、Bob の暗号文 $C$ からメッセージ $M$ を復号する：$M = D_A(C)$。
\end{itemize}
\end{exa}

公開鍵は秘密ではないので、攻撃者 Mallory も知っていると考える。
そのため、公開鍵暗号に特有の新たな懸念が生じる。

\begin{dfn}[選択平文攻撃]\label{dfn:ch5-chosen-plaintext}
盗聴者が予測しやすいメッセージを推測できる場合、公開鍵を使って、ありそうなメッセージを片っ端から自分で暗号化し、回線上を流れる暗号文と比べることができる。
一致する暗号文が送られてくるのを見れば、その元になった平文がわかる。
このような攻撃を\underline{選択平文攻撃}（chosen plaintext attack）という。
\end{dfn}

選択平文攻撃は、推測が現実的でなくなるだけの数のランダムなビットでメッセージをパディングすれば、簡単に防げる。

RSA は最初に公開された非対称暗号系ではなかったが、大きな反響を呼んだ。
これを破る（すなわち、ある人の公開鍵から秘密鍵を導き出す）には、大きな素数の積を因数分解するという、よく知られた難問を解かなければならないからである。

\begin{rem}[RSA の意義の当時と今]\label{rem:ch5-rsa-history}
原著者は RSA が公に登場した当時、ささやかながらその発明者たちと協力していたので、当時と現在のその意義について、興味深いかもしれない歴史的な注釈を加えることができる。
当時のコンピューターにとってこのアルゴリズムは計算量が大きすぎ、使うには高価な専用ハードウェアが必要であった。
そのため、使うのは大手の金融機関や軍の情報機関くらいだろうと考えていた。
計算能力が時間とともに指数関数的に伸びるという Moore の法則は知っていたが、40 年後に普通の人々が、必要な計算をこなせるプロセッサーを積んだネットワーク接続のスマートフォンを日常的に持ち歩くようになるとは、当時は誰も想像していなかった。
\end{rem}

今日、RSA は楕円曲線アルゴリズム（elliptic curve algorithm）などの新しい方式に置き換えられつつある。
これらのアルゴリズムは、別の数学に基づいて同様の能力を実現し、より少ない計算で強力な暗号化ができるという点で「費用対効果」が高い。
非対称暗号はたいてい対称暗号より計算コストが高いので、暗号化は通常、ランダムな秘密鍵を選んでその鍵を非対称暗号で暗号化し、メッセージ本体は対称暗号で暗号化するという形で行う。

%----------------------------------------------------------------------
\subsection{デジタル署名}
% src: Digital Signatures

\begin{dfn}[デジタル署名]\label{dfn:ch5-digital-signature}
公開鍵暗号を使って作られ、受信者に真正性（authenticity）を保証するものを\underline{デジタル署名}（digital signature）という。
メッセージを暗号化するかどうかとは関係なく、Alice の署名は、メッセージが本当に彼女から来たものであることを Bob に保証する。
また、Alice が送っていないと言い張った場合には、通信があったことの証拠にもなる。
\end{dfn}

第2章で述べたように、真正性と否認不可能性（non-repudiability）は、通信にとって機密性に次いで重要な 2 つのセキュリティ上の性質である。

% 図 5-2: A comparison of symmetric and asymmetric cryptography（左に対称暗号、右に非対称暗号の比較。図の詳細は本文から決まらないため描かない）
原著の図 5-2 は、左側の対称暗号と右側の非対称暗号の根本的な違いをまとめている。
対称暗号では両方の通信者が秘密鍵を知っているので、署名はできない。
非対称暗号の安全性は、一方の通信者だけが知っている秘密鍵にかかっているので、その通信者だけが秘密鍵を署名に使える。
しかも検証には公開鍵しか要らないので、その過程で秘密が明かされることはない。

\begin{exa}[デジタル署名の作成と検証]\label{exa:ch5-sign-verify}
Alice は、公開鍵暗号に使うのと同じ鍵の対を使ってデジタル署名を作る。
秘密鍵を知っているのは Alice だけなので、署名関数 $S_A$ を計算できるのも彼女だけである。
Bob は、あるいは公開鍵（と $N$）を持っている人なら誰でも、関数 $V_A$ を使って Alice の署名を検証できる。
すなわち、
\begin{itemize}
\item Alice はメッセージ $M$ に署名して、署名 $S = S_A(M)$ を作る。
\item Bob は、$M = V_A(S)$ が成り立つかを確かめて、メッセージ $M$ が Alice から来たものであることを検証する。
\end{itemize}
\end{exa}

デジタル署名の仕組みを完全に理解するには、さらにいくつか細かい点を説明しておく必要がある。
検証には公開鍵しか使わないので、Bob は Alice の秘密鍵を危険にさらすことなく、Alice がメッセージに署名したことを第三者に証明できる。
また、メッセージへの署名と暗号化は互いに独立しており、アプリケーションに応じてどちらか一方だけでも、両方でも行える。
本書では RSA の背後にある数学には踏み込まないが、署名関数と復号関数（どちらも秘密鍵が必要）は実は同じ計算であり、検証関数と暗号化関数（公開鍵を使う）も同じ計算であることは知っておいてほしい。
混乱を避けるため、目的に応じて別々の名前で呼ぶのが一番である。

デジタル署名は、デジタル証明書（次の節の主題）、電子メール、アプリケーションのコード、法的文書の署名に広く使われているほか、Bitcoin のような暗号通貨の保護にも使われている。
慣例として、1 回の署名操作で文書全体をカバーできるよう、便宜上メッセージのダイジェストに署名する。
ここまで来れば、ダイジェスト関数に対する原像攻撃が成功するとなぜ非常にまずいのかがわかるだろう。
同じメッセージダイジェストを持つ支払い契約を Mallory がでっち上げられれば、Bob の約束手形 $P$ はその契約に対する有効な署名としても通用してしまう。
% 注: 原文 "Bob's promissory note P also serves as a valid signature for it." のとおり訳した。
%   約束手形 P は原文のこの箇所で初出であり、文脈（署名の対象が P のダイジェストで、
%   その署名が偽の契約にも有効になる）を補う記述は原文にない。

%----------------------------------------------------------------------
\subsection{デジタル証明書}
% src: Digital Certificates

原著者が RSA アルゴリズムについて学び始めた頃、チームのメンバーと将来どんな応用がありうるかをブレインストーミングした。
公開鍵暗号の決定的な利点は、その便利さであった。
自分の公開鍵を世界に告知して誰でも使えるようにさえできれば、通信相手ごとに別々の鍵を管理しなくても、すべての通信に 1 つの鍵を使える。
しかし、公開鍵を告知するにはどうすればよいのだろうか。

原著者は学位論文の研究でその答えを考え出し、そのアイデアはその後広く実装されている。

\begin{dfn}[認証局とルート証明書]\label{dfn:ch5-certificate-authority}
デジタルの公開鍵暗号という新しい仕組みを普及させるために必要とされた、新しい種類の組織を\underline{認証局}（certificate authority, CA）という。
新しい CA は、まず自分の公開鍵を広く公開する。
やがて、オペレーティングシステムやブラウザーには、信頼できる CA の\underline{ルート証明書}（root certificate）一式がプリインストールされるようになる。ルート証明書には CA の公開鍵が含まれている。
\end{dfn}

\begin{dfn}[デジタル証明書]\label{dfn:ch5-digital-certificate}
CA は申請者から（通常は有料で）公開鍵を集め、申請者ごとに、その名前（例えば「Alice」）やその他の詳細を公開鍵とともに記載した\underline{デジタル証明書}（digital certificate）を発行する。
CA は、デジタル証明書が本物であることを保証するために、証明書のダイジェストに署名する。
\end{dfn}

理論上は、申請が本当に Alice から来たものかを審査で確かめることが CA のサービスの重要な部分であり、人々は、CA が審査を確実に行う場合に限って、その CA を信頼するはずである。
しかし実際には、とくにインターネット越しに身元を確認するのは非常に難しく、この点が問題になっている。

デジタル証明書を手に入れた Alice は、誰かと通信したいときにはいつでも、そのコピーを相手に送ればよい。
相手がその証明書を発行した CA を信頼していれば、相手はその CA の公開鍵を持っているので、「Alice」の公開鍵を記載したデジタル証明書の署名を検証できる。
デジタル証明書とは、要するに「Alice の公開鍵は X である」と述べた、CA の署名付きメッセージである。
これを検証できれば、受信者はすぐに Alice 宛てのメッセージを暗号化し始められる。通常は、まず自分のデジタル証明書を送り、Alice のメッセージが正しい相手に届いたことを本人に保証するところから始める。
デジタル署名も同じように機能し、同じデジタル証明書で裏付けられる。

\begin{rem}[デジタル証明書の実際の複雑さ]\label{rem:ch5-certificate-complications}
ここでのデジタル証明書の説明は簡略化したもので、信頼できる CA が名前と秘密鍵の結び付きをどう認証するかに焦点を当てている。
% 注: 原文は "the association of a name with a private key" であり、そのまま訳した
%   （前段の説明では証明書が結び付けるのは公開鍵である）。
実際にはもっと多くの事情が絡む。人の名前は一意とは限らず、名前は変わることがあり、州が違えば同じ名前の企業が存在することもある、といった具合である。
（第11章では、Web セキュリティの文脈でこうした厄介な問題のいくつかを掘り下げている。）
今日、デジタル証明書は、Web サーバーのドメイン名や電子メールアドレスを含むさまざまなアイデンティティに鍵を結び付けるために使われているほか、コード署名などのいくつかの特定の目的にも使われている。
\end{rem}

%----------------------------------------------------------------------
\subsection{鍵交換}
% src: Key Exchange

最初の鍵交換アルゴリズムは、RSA が発明される少し前に Whitfield Diffie と Martin Hellman によって開発された。
鍵交換がどれほど奇跡的かを理解するために、Alice と Bob が何らかの方法で通信チャネルを確立したものの、秘密鍵を事前に取り決めてもおらず、公開鍵の出所として信頼できる CA すらない状況を想像してほしい。
信じがたいことに、鍵交換を使えば、Mallory がすべてを観察している開かれたチャネル上で、2 人は秘密を確立できる。
そんなことが可能だというのはあまりに直感に反するので、ここでは自分で仕組みを確かめられるよう、原著者はあえて数学を示している。

幸い、数学は十分に単純で、小さな数なら計算も簡単である。
読者によってはなじみがないかもしれない記法は、接尾辞 $\pmod{p}$ だけである。これは整数 $p$ で割った余りをとることを意味する。
例えば、$2^7 \pmod{103}$ は $25$ である。$128 - 103 = 25$ だからである。

\begin{dfn}[Diffie--Hellman 鍵交換]\label{dfn:ch5-diffie-hellman}
\underline{Diffie--Hellman 鍵交換}（Diffie--Hellman key exchange）のアルゴリズムの基本は次のとおりである。
\begin{enumerate}
\item Alice と Bob は、素数 $p$ とランダムな数 $g$（$1<g<p$）を公開の場で取り決める。
\item Alice はランダムな自然数 $a$（$1<a<p$）を選び、$g^a \pmod{p}$ を Bob に送る。
\item Bob はランダムな自然数 $b$（$1<b<p$）を選び、$g^b \pmod{p}$ を Alice に送る。
\item Alice は、2 人の共有秘密 $S$ として $S = (g^b)^a \pmod{p}$ を計算する。
\item Bob は $S = (g^a)^b \pmod{p}$ を計算し、Alice と同じ共有秘密 $S$ を得る。
\end{enumerate}
\end{dfn}

% 図 5-3: Alice and Bob securely choosing a shared secret via key exchange
%   （小さな数による鍵交換のおもちゃの例。p, g, a, b の具体的な値は本文にないため図は描かない。
%    本文から読み取れるのは、探索空間が約 60 通りであることと、Alice の選んだ数 6 が Bob の結果 g^b に等しいことだけである）
\begin{exa}[小さな数による鍵交換]\label{exa:ch5-toy-key-exchange}
原著の図 5-3 は、小さな数を使ったおもちゃの例で、この方法が実際に機能することを示している。
この例は安全ではない。約 60 通りの可能性をしらみつぶしに調べるのは簡単だからである。
しかし同じ数学は大きな数でも成り立ち、数百桁の規模になれば、そうしたしらみつぶしの探索はとうてい不可能になる。

この例では数を小さく抑えるように選んだため、たまたま Alice が選んだ 6 が Bob の計算結果（$g^b$）と同じになっている。
実際にはこんなことはまず起きないが、もちろん起きてもアルゴリズムは問題なく機能し、この偶然に気付くのは Alice だけである。
\end{exa}

\begin{rem}[鍵交換における乱数]\label{rem:ch5-key-exchange-randomness}
Mallory に両者の選んだ数を推測される余地をなくすため、両者が実際に CSPRNG から安全な乱数を選ぶことが重要である。
例えば、Bob が $p$ と $g$ から自分の数を計算する公式を使っていたとすると、Mallory は多くの鍵交換を観察してその公式を推定し、最終的には同じ計算をまねて、鍵交換の秘密を破ってしまうかもしれない。
\end{rem}

鍵交換は、要するに種も仕掛けもない手品である。
Alice と Bob が舞台袖から歩いて入ってくると、舞台の真ん中には Mallory が立っている。
Alice が数を読み上げ、Bob が答える。2 往復のやり取りが済んでも、Mallory には依然として見当もつかない。
Alice と Bob が共有した秘密の数をそれぞれ大きなカードに書き、合図とともに掲げると、取り決めた秘密を表すまったく同じ数が現れる。

今日、鍵交換は、インターネット上の任意の 2 つのエンドポイントの間に安全な通信チャネルを確立するために欠かせない。
ほとんどのアプリケーションは性能の高い楕円曲線鍵交換のアルゴリズムを使っているが、考え方はほぼ同じである。
鍵交換は、インターネット上で（TLS プロトコルのような）安全な通信チャネルを設定するときにとくに役に立つ。
2 つのエンドポイントは、まず TCP チャネル（Mallory が観察しているかもしれない通信）を使い、次に鍵交換を行って、まだ素性を確かめていない相手方の通信者と秘密を取り決める。
共有秘密が得られれば、暗号化した通信によって安全な私的チャネルが実現する。
こうして、どんな通信者の組でも、事前に取り決めた秘密なしに安全なチャネルを立ち上げられる。

%----------------------------------------------------------------------
\subsection{暗号の利用}
% src: Using Crypto

この章では、暗号の道具箱に入っている道具を「自動車教習」のレベルで説明した。
暗号論的に安全な乱数は、予測不可能性を加えて、推測に基づく攻撃を阻む。
ダイジェストは、データの一意性を対応するトークンに凝縮して完全性を検査するための、安全な方法である。
暗号化には対称と非対称の両方の形があり、機密性を保護する。
デジタル署名は、メッセージを認証する方法である。
デジタル証明書は、CA への信頼を利用して、本物の公開鍵を簡単に共有できるようにする。
そして鍵交換は、離れた当事者どうしが公開のネットワーク接続を介して秘密鍵を安全に取り決められるようにするもので、これで暗号の道具箱がそろう。

% 図 5-4: Security versus the $5 wrench (courtesy of Randall Munroe, xkcd.com/538)
原著の図 5-4 の漫画（xkcd \#538「Security」、Randall Munroe 提供、xkcd.com/538。5 ドルのレンチに対するセキュリティ）は、この章の冒頭のエピグラフが指摘する点を描いている。すなわち、よくできた暗号は非常に強力なので、主な脅威は暗号を迂回されることだという点である。
この章から得られる最も重要な教訓はおそらく、うっかりそうした攻撃の隙を与えないよう、暗号を正しく使うことが決定的に重要だということである。

暗号は、ソフトウェアの設計で生じるセキュリティ上の課題や、脅威モデリングで見つかるセキュリティ上の課題の多くに役立つ。
システムがインターネット経由でパートナーのデータセンターにデータを送らなければならないなら、データを（機密性のために）暗号化し、（完全性のために）デジタル署名すべきである。あるいは、エンドポイントを認証する TLS の安全なチャネルを使うという簡単な方法をとってもよい。
安全なダイジェストを使えば、データの完全なコピーを保存しなくても、データが等しいかどうかを検査できる（MAC としての使い方を含む）。これは気の利いた方法である。
たいていは自分で作るのではなく既存の暗号サービスを使うことになるが、この章では、それらをいつどう使うべきか、そしてこの技術を安全に使ううえでの課題のいくつかについて、考え方を示した。

\begin{exa}[受信直後の機密データの暗号化]\label{exa:ch5-encrypt-on-arrival}
金融口座の残高やクレジットカード情報は、絶対に保護しなければならないデータの典型例である。
この種の機密データは大きな分散システムの中を流れていくので、施設への立ち入りが制限されていても、誰かがネットワークタップを物理的に差し込んで機密データを吸い上げられるようでは困る。
強力な緩和策の一つは、入ってくる機密データをすべて、フロントエンドの Web サーバーに最初に届いた時点ですぐに暗号化することである。
クレジットカード番号をすぐに公開鍵で暗号化しておけば、取引を処理する間、暗号化されたデータを中身の見えない塊（opaque blob）として受け渡せる。
最終的に、このデータは厳重に保護された金融処理用のマシンに届く。このマシンは秘密鍵を知っているので、データを復号し、銀行システムと取引を照合できる。
この方法なら、アプリケーションコードの大部分は、自分自身が情報を漏らすリスクを負わずに、機密データを後続の処理に安全に受け渡せる。
\end{exa}

\begin{exa}[暗号化データと鍵の分離保管]\label{exa:ch5-separate-key-storage}
よく使われる手法がもう一つある。対称暗号で暗号化したデータと秘密鍵を、別々の場所に保管することである。
例えば、バックアップ用の長期データ保管を第三者に外部委託したい企業を考える。
企業は暗号化したデータを保管用に引き渡し、バックアップから復元が必要になったときに使う鍵は、自社の金庫に保管しておく。
脅威の観点から見ると、データ保管サービスには（データを失う可能性があるので）完全性の保護を任せていることになるが、鍵が安全で暗号化が正しく行われている限り、機密性が損なわれるリスクはない。
\end{exa}

これらはよくある使い方のほんの一部であり、ほかにもこれらの道具の使い道はたくさん見つかるだろう。
（暗号通貨はとくに巧妙な応用の一つである。）
現代のオペレーティングシステムやライブラリには、現在通用しているいくつかのアルゴリズムの成熟した実装が用意されているので、実際の計算を自分で実装することなど考える必要すらない。

\begin{exa}[暗号化しても漏れるメタデータ]\label{exa:ch5-metadata-leak}
とはいえ、暗号化は万能薬ではない。攻撃者が暗号化されたデータの頻度や量などのメタデータを観察できれば、何らかの情報が漏れうる。
例えば、家の中で動きを検知すると画像を撮影する、クラウドベースの防犯カメラシステムを考える。
家族が留守のときは動きがないので、カメラからの送信もない。
画像が暗号化されていても、家庭のネットワークを監視できる攻撃者は、カメラの通信量が減ることから家族の日々の行動パターンを簡単に推測し、家が留守になる時間帯を突き止められる。
\end{exa}

\begin{rem}[暗号の安全性の前提]\label{rem:ch5-crypto-limits}
暗号の安全性は、現在知られている数学の限界と、デジタルハードウェア技術の最先端の水準を前提にしており、そのどちらも容赦なく進歩し続けている。
現代のアルゴリズムを揺るがすような、もっと効率のよい計算方法を見つけた数学者には、大きな名声が待っている。
さらに、量子物理学のような別の種類の計算技術が実用化される見込みも、もう一つの潜在的な脅威である。
どこかの強大な国家がすでにそうしたブレークスルーを達成していて、手の内を明かさないよう密かに使っている可能性すらある。
ほかのあらゆる緩和策と同じく、暗号にも本質的にトレードオフと未知のリスクがつきまとうが、それでも暗号は優れた道具箱であり、使いこなす価値のある道具の集まりである。
\end{rem}

\newpage

%% =======================================================
%%  part: ch6
%% =======================================================
%% =======================================================
%%  part: ch6
%% =======================================================
\section{セキュアな設計}
% src: 6: Secure design

この章では、\cite{Kohnfelder2021} の第6章に沿って、セキュアなソフトウェア設計を扱う。

\begin{quote}
「過剰、散らかり、混乱は情報の性質ではない。設計の失敗である。」\\
\hfill --- Edward Tufte
\end{quote}

セキュリティの原則、パターン、緩和策をしっかり理解していれば、ソフトウェア設計にセキュリティを組み込む作業はそれほど難しくない。
設計に対する脅威を見極めながら必要に応じてこれらの道具を使い、リスクがおのずと小さくなるような、よりよい設計案を探っていけばよい。

この章ではセキュアなソフトウェア設計を取り上げる。
セキュリティ設計レビューを扱う第7章とは対になる章である。
この二つの話題は、同じ活動を異なる視点から見たものといえる。
ソフトウェアの設計者は、この章で論じる概念を踏まえ、設計プロセス全体を通じてこれらの方法を適用すべきである。システムのセキュリティを、後でレビュー担当者に繕ってもらうために残しておいてはならない。
一方、レビュー担当者は、セキュリティ評価の層をもう一枚重ねるつもりで、脅威と緩和策という観点から設計を見るべきである。
セキュアな設計のプロセスは統合的（integrative）であり、セキュリティ設計レビューは分析的（analytic）である。両者を組み合わせて相乗効果を生かせば、セキュリティが組み込まれたよりよい設計ができる。

ソフトウェア設計は一種の技芸であり、この章ではそのうちセキュリティの側面だけを取り上げる。
正式なプロセスに従って設計する人も、すべて頭の中で済ませる人も、ここで示す考え方を取り入れるために仕事のやり方を変える必要はない。
脅威モデリングとセキュリティの視点が設計を\emph{主導する}必要はないが、設計の\emph{判断材料}にはなるべきである。

ここで述べるセキュアな設計の実践は大企業でよく見られるプロセスに沿っているが、どのような仕事の進め方にも合わせて調整できる。
小規模な組織ではずっとくだけたやり方になるだろうし、設計者とレビュー担当者が同じ人ということもある。
ここで示す手法は問題を一般的な形で扱っているので、どのようなやり方でソフトウェアを設計する場合にも簡単に応用できる。

% src(sidebar): A Sample Design Document that Integrates Security
\begin{rem}[セキュリティを統合した設計書のサンプル]\label{rem:ch6-sample-design-document}
設計は「ハウツー」の手順に落とし込めない創造的なプロセスである。そこで著者は、本書で示した概念をどう適用するかを見せるために、設計書の完全な例を用意したいと考えた。
付録A のサンプルは、最初からセキュリティを組み込むやり方を示している。
熟練の技を見せる完璧な設計例を意図したものではなく、最終的な仕上がりの感触がつかめる程度に中身の詰まった、作業途中の初稿である。
簡潔さを優先して、ここでの目的にとって重要でない部分は省いてある。また、現実の設計のほとんどがそうであるように、欠点や粗い箇所が残ったままの磨かれていない状態で示している。

このサンプル設計書が想定しているのは、監査をしやすくしつつ個人情報の開示を最小限に抑えるロギングツールであり、実際に使える有用なコンポーネントになることを狙っている。
この種のツールは、機微なデータを扱う大きなシステムの中で、実用的な緩和策になりうる。興味があれば、設計を肉付けして実際に作ってみてもよい。
いずれにせよ、この例にはぜひ目を通してほしい。この章の指針が設計書の中で実際にどう形になるかを見れば、セキュアな設計の進め方がよくわかるはずである。
\end{rem}

%----------------------------------------------------------------------
\subsection{設計へのセキュリティの統合}\label{subsec:ch6-integrating-security}
% src: Integrating Security in Design

\begin{quote}
「概念の完全性（conceptual integrity）こそがシステム設計で最も重要な考慮事項である、と私は主張したい。」\\
\hfill --- Fred Brooks（\textit{The Mythical Man-Month} より）
\end{quote}

設計段階は、セキュリティの原則とパターンをソフトウェアプロジェクトに組み込む絶好の機会である。
この初期の段階なら、実装に手間をかけて過去の決定に縛られてしまう前に、代替案を気軽に検討できる。

設計段階では、開発者は\underline{設計書}（design document）を作成し、ソフトウェアプロジェクトの重要な高水準の特徴を書き留めておくべきである。建物でいえば建築設計図にあたるものである。
設計の文書化には、ぜひ労力を割いてほしい。文書化は設計の厳密さを保つ助けになるうえ、自分の下した決定を他の人が理解できる貴重な成果物にもなる。脅威と緩和策のバランスや、それに伴うトレードオフについては、とくにそうである。

設計書は通常、機能記述（外から見てソフトウェアがどう動くか）と技術仕様（内側でどう動くか）から成る。
利害の競合するステークホルダーがいる場合、大きな取り組みを調整する場合、設計が正式な要求仕様や厳しい互換性要求に従わなければならない場合、難しいトレードオフに直面した場合などには、より正式な設計がとりわけ役に立つ。

ソフトウェア設計の案を見るときは、「セキュリティの帽子」（security hat）をかぶろう。
そうすれば、コーディングを始める前に、脅威モデリングをし、攻撃対象領域を特定し、データの流れを描き出す、といったことができる。
提案された設計のままではシステムをセキュアにすることが構造的に難しいなら、本質的によりセキュアな代替案を検討するのは今をおいてほかにない。
また、重要なセキュリティ上の緩和策は設計書の中で指摘しておき、実装者が前もってその必要性を理解できるようにすべきである。

経験を積んだ設計者は、最初からセキュリティを設計に組み込む。
それが難しそうなら、まず「機能的に完成した」設計の草稿を作り、セキュリティに焦点を当てて 2 回目の見直しをしてもかまわない。ただし、そのほうがはるかに手間がかかる。
大きな変更は、プロセスの早い段階で気付くほど簡単に済み、後からやり直す無駄な労力も避けられる。
新しいアーキテクチャを探ったり基本的な要件をいじってみたりするのは、まだ簡単にできるうちに、なるべく早く済ませるべきである。
Josh Bloch が皮肉を込めて言ったとおり、「1 週間のコーディングで 1 時間の思考を節約できることがよくある」（A week of coding can often save an hour of thought.）。

\subsubsection{設計上の前提を明示する}\label{subsec:ch6-explicit-assumptions}
% src: Making Design Assumptions Explicit

\begin{exa}[ハードウェアとソフトウェアのエンディアンの食い違い]\label{exa:ch6-endianness}
1980 年代半ば、著者は、当時としては強力なコンピュータをハードウェアとソフトウェアの両方とも一から設計・構築する会社で働いていた。
何年もの開発を経て、ついにオペレーティングシステムを試作ハードウェアに載せ、両チームの成果が一つになった。ところが……即座に失敗した。
ハードウェアチームの多くはビッグエンディアンのアーキテクチャを使う IBM の出身で、ソフトウェアチームの多くは伝統的にリトルエンディアンを使ってきた HP の出身だった。そのため「ビット 0」が、ハードウェアでは最上位ビットを、ソフトウェアでは最下位ビットを指していたのである。
何年にもわたる計画、会議、試作のあいだ、誰もが出身企業の文化のエンディアンを当然のものと思い込んでいた。
（そして当然ながら、原因が判明した後で必要な変更をする羽目になったのはソフトウェアチームだった。）
\end{exa}

書かれていない前提はセキュリティ設計レビューの効果を損ないかねないので、設計者は前提を文書化するよう努めるべきである（そしてレビュー担当者は、不明確な点があれば何でも尋ねるべきである）。
こうした前提を明示して記録しておくのによい場所は、設計書の本体の前に置く「背景」の節である。

前提の文書化は、深刻な誤解を先回りして予測し、誰かに「でも私は……と思っていた」と言わせないための作業だと考えるとよい。
文書化することが重要なのに、設計では書き漏らされやすい一般的な前提を挙げる。
\begin{itemize}
\item 設計の選択肢を狭める予算・リソース・時間の制約
\item システムが攻撃の標的になりそうかどうか
\item レガシーシステムとの互換性のような、譲れない要件
\item システムが達成すべきセキュリティの水準についての期待
\item データの機微性と、そのデータをセキュアに保護することの重要性
\item 将来システムに変更が必要になる見込み
\item システムが達成すべき具体的な性能や効率の基準
\end{itemize}

前提をはっきりさせることがセキュリティにとって重要なのは、誤解が根本原因となって、インターフェース設計が弱くなったりコンポーネント間のやり取りがちぐはぐになったりし、そこを攻撃者に悪用されることが多いからである。
加えて、前提がはっきりしていれば、設計のレビュー担当者もプロジェクトを明確かつ一貫した形で把握できる。

企業の中や関連するプロジェクトの集まりの中では、こうした前提の多くが一連の設計を通じて変わらないことがよくある。その場合は、共通の背景をまとめた共有文書にリストを載せておけばよい。
そうすれば個々の設計はこの共通の土台を参照し、前提が異なる例外だけを詳しく書けば済む。
たとえば請求システムでは、クレジットカード処理のコンポーネントに、企業のほかのアプリケーションより高いセキュリティ基準が課され、特定の金融規制に準拠しなければならないかもしれない。

\subsubsection{スコープの定義}\label{subsec:ch6-defining-scope}
% src: Defining the Scope

レビューのスコープ（範囲）が不確かなままでは、設計のセキュリティをきちんとレビューすることはできない。
スコープをはっきりさせることは、第2章の 4つの質問の一つ「何に取り組んでいるのか？」（What are we working on?）に答えるうえでも欠かせない。

\begin{exa}[新しい顧客請求システムのスコープ]\label{exa:ch6-billing-scope}
なぜそうなのかを、新しい顧客請求システムの設計を例に考えてみる。
請求対象の作業時間の報告を集める Web アプリは、この設計に含まれるのか、それとも別の設計なのか。
このシステムが依存している既存のデータベースはどうか。それらのシステムのセキュリティはスコープに入るのか入らないのか。
そして、社内会計システムへの報告に使う新しい Web ベースの API の設計は、レビューに含めるべきなのか。
\end{exa}

スコープをどう定め、どこまでを引き受けるかは、通常は設計者が戦略的に判断する。
ほかの人がスコープを定めた場合には、設計者は決められたスコープとその理由を理解しておかなければならない。
設計のスコープは、関係者全員にとって明確でありさえすれば、プロセス内で動くコード、ブロック図に描かれたシステムの特定のコンポーネント、ライブラリのコード、ソースリポジトリの一区画など、最も理にかなう形で定めてよい。
前の段落で触れた請求システムの設計なら、新しい API は同じ設計の拡張なので、おそらくスコープに含めるべきである。
逆に既存のデータベースは、まったく新しい使い方をするわけではなく、すでにセキュリティ面で十分な注意が払われているなら、おそらくスコープ外である。

設計のスコープが曖昧だと、レビュー担当者はセキュリティ上の重要な側面をスコープ外だと思い込み、設計者のほうはその問題に気付いていない、ということが起こりうる。
こうして誰も触れないまま、その側面が抜け落ちてしまうかもしれない。

\begin{exa}[誰もが他の誰かがやっていると思うバックアップ]\label{exa:ch6-backup-falls-through}
ほぼすべてのソフトウェア設計は、何らかの形でデータの保存を伴う。
そしてデータが使い捨てでない限り（使い捨てのデータはまれである）、よいバックアップを取り続けることは、さまざまな脅威（悪意によるものも偶発的なものも）による完全性の喪失に対する自明な緩和策である。
設計者はこうした自明な点を書き漏らしがちである。しかし設計のスコープがはっきり書かれていないと、本番システムのストレージはすべて誰か他の人が定期的にバックアップしているものと全員が思い込み、その作業が放置されるかもしれない --- 最初の障害が起きて、手痛い形で教訓を学ぶまでは。
\end{exa}

設計のエコシステムの一部をスコープから外したせいで、その部分が抜け落ちてしまうことのないようにすべきである。
レガシーシステムを引き継いだときは、まずシステムの最も機微な部分、セキュリティにとって最も根本的な部分、あるいは攻撃の標的として最も目立ちそうな部分に絞って理解に努めるとよい。
その後、独立したコンポーネントを構成するほかの部分のレビューを、すべてを網羅するまで慎重に進めていく。

設計の反復やスプリント、既存システムの大きな改訂は、再設計する箇所に合わせて狭いスコープを定めれば扱える。
新しい設計作業の境界を切り出してしまえば、スコープの外側については設計で定めた明確な前提条件があり、内側はすべて自由に一から作り直せる。
既存の設計文書があればこの作業はずっと楽になり確実にもなる。また、設計を更新したら、文書も正確に更新すべきである。

再設計が当初の範囲からはみ出していくのはよくあることで、たいていは悪いことでもない。そうなったら、必要に応じてスコープを調整すればよい。
たとえば、段階的な設計変更のために既存のインターフェースやデータ形式を修正しなければならないことがある。その変更でより機微なデータを扱うことになるなら、新しいセキュリティ上の前提に合わせて、インターフェースの反対側にも変更が必要になるかもしれない。

真空の中に存在するソフトウェア設計はほとんどない。設計は既存のシステム、プロセス、コンポーネントに依存している。
設計が依存先とうまくかみ合うことを確かめるのは、決定的に重要である。
とくに鍵になるのは、セキュリティ上の期待をそろえることである。セキュアでないコンポーネントを組み合わせて、セキュアなアプリケーションを作ることはできないからである。
そして、セキュアかセキュアでないかは二者択一ではなく連続的なものであり、前提と期待をそろえる必要があることに注意しなければならない。
同種のシステムや依存先のセキュリティ設計レビューの報告書を読み、それらに対するセキュリティ上の期待の裏付けをとるべきである。

\subsubsection{セキュリティ要件の設定}\label{subsec:ch6-security-requirements}
% src: Setting Security Requirements

セキュリティ要件は、主に 4つの質問の 2 番目「何がうまくいかなくなりうるか？」（What can go wrong?）から導かれる。
出発点としては C-I-A の三つ組が役に立つ。私的なデータを許可のない開示から守る必要性（機密性）、データをバックアップすることの重要性（完全性）、システムがどの程度頑健で信頼できなければならないか（可用性）を書き出すのである。
多くのソフトウェアシステムのセキュリティ要件は単純だが、それでも漏れをなくし優先順位を伝えるために、詳しく書いておく価値は十分にある。
自分にはごく当たり前のことでも、ほかの人にとってはそうでないかもしれない。目指すセキュリティ上の姿勢は、はっきり言葉にしておくのがよい。

一方の極端として注意すべきなのは、セキュリティが重要でない場合 --- 少なくとも、誰かが重要でないと考えている場合 --- である。
チームの別の誰かはセキュリティが当然重要だと考えているかもしれないので、これは明示しておくべき重要な前提である（このような期待の食い違いが最終的にどんな状況で表面化するかは、想像がつくだろう）。

\begin{exa}[セキュリティの重要度が低いアプリケーション]\label{exa:ch6-low-security}
人工的なダミーデータを処理する試作品を設計しているなら、セキュリティレビューは省略してよい。ただし、そのコードが後で転用されて個人情報の処理に使われることのないよう、そのことを文書化しておくべきである。
セキュリティの重要度が低いアプリケーションのもう一つの例として、複数の研究グループが共有する気象データの収集が挙げられる。気温などの大気の状態は誰でも自由に測定でき、開示されても害はない。
\end{exa}

もう一方の極端として、セキュリティが決定的に重要なソフトウェアには、特別な注意を払い、セキュリティ関連の要件を注意深く列挙する価値がある。
列挙した要件は、最高水準の品質を確保するための脅威モデリング、セキュリティレビュー、テストの焦点になる。
セキュリティ要件が設計にどう反映されるかの基本的な例は、サンプル設計書（付録A）を参照してほしい。
複雑な規制の対象となる大規模なシステムでは、高い水準のコンプライアンスを確保するためにセキュリティ要件が厳密に規定されていることもあるが、それは専門的な取り組みであり、ここでは扱わない。

決定的に重要なセキュリティ要件や、通常とは異なるセキュリティ要件を持つソフトウェア設計については、次の一般的な指針を検討するとよい。
\begin{itemize}
\item セキュリティ要件は、「どうやるか」を指示せず、最終的な目標として表現する。
\item すべてのステークホルダーのニーズを考慮する。とくに、ニーズが対立しうる場合には、うまくバランスをとる必要がある。
\item 決定的に重要な緩和策について、許容できるコストとトレードオフを認めておく。
\item 通常と異なる要件がある場合は、その目標とあわせて動機も説明する。
\item 完璧を求める命令ではなく、達成可能なセキュリティ目標を設定する。
\end{itemize}

重大なセキュリティ上のニーズを持つシステムの要件記述がどのようなものになりうるかを、次の極端な例で示す。

\begin{exa}[重大なセキュリティ上のニーズを持つシステムの要件記述]\label{exa:ch6-extreme-requirements}
\begin{description}
\item[国家安全保障局（National Security Agency）で、国家の最も機微な秘密を守る]
システム管理者は膨大な最高機密文書に並外れたアクセス権を持つ。これが国家安全保障にもたらす脅威を考えると、内部者による攻撃は可能な限り徹底して緩和しなければならない。
具体的には、広範なアクセス権限を持つ高位の職員になりすませる管理者がいれば、大量のファイルを持ち出したうえで、多くの異なるプリンシパルによる多数の独立したアクセスのように見せかけて痕跡を隠せてしまうおそれがある。
（Edward Snowden が NSA の内部文書を持ち出した手口についての非公式な説明によれば、彼はこの種の手法を使ったらしい。）

\item[大手金融機関の認証サーバー]
サーバーの暗号の秘密鍵が侵害されると、インターネットに面したすべてのシステムのセキュリティが完全に崩れる。
内部者による攻撃は起こりにくいが、運用担当者は\underline{もっともらしい否認}（plausible deniability）ができなければならない。
要件としては、物理的に警備された場所に置いた、改ざんの痕跡が残るハードウェアデバイスに鍵を保存することや、鍵の生成とローテーションを正式な儀式（ceremony）として行い、あらゆるアクセスに少なくとも 2 人の信頼できる人物が立ち会うことなどが考えられる。
（注：ここでは、信頼の分散と、物理的セキュリティと論理的セキュリティを重ね合わせることを最も直接的に示すために、あえて「どうやるか」を含めている。）

\item[高価な科学実験のデータの完全性]
この実験は 1 回しか行わない予定であり、実験に必要な資金は今後何年も再び得られそうにない。そのため、計測機器が収集する情報を失うわけにはいかない。
ストリーミングデータは即座に複製して異なる記憶媒体に冗長に保存しなければならない。同時に、追加のバックアップとして、二つの別々のネットワークを通じて、物理的に離れた遠隔のストレージシステムにも送らなければならない。
\end{description}
\end{exa}

\subsubsection{脅威モデリング}\label{subsec:ch6-threat-modeling}
% src: Threat Modeling

ソフトウェアアーキテクチャのセキュリティを高める最良の方法の一つは、設計プロセスに脅威モデリングを組み込むことである。
ソフトウェアの設計では、競合する要件や戦略を創造的にやりくりし、システムのいくつかの側面を少しずつ決め、ときには方針を転換しながら、全体像へと近づいていく。
このプロセスを脅威モデリングの観点から見ると設計のトレードオフが浮かび上がるので、設計者を正しい方向に導く大きな力になりうる。ただし、どうすればよりよい結果が得られるのかを正確に見極めるには、ある程度の試行錯誤が必要である。

まず、脅威モデリングをソフトウェア設計に組み込む力ずくの方法がある。
設計の候補をいくつもこしらえ、一つずつ脅威モデリングし、何らかの総合評価で点数を付けて、最良のものを選ぶというやり方である。
実際には、こうしたセキュリティ中心の評価は、ユーザビリティ、性能、開発コストといったほかの重要な要因を判断する材料にもなる。
しかし、複数の設計を作ってそれぞれを個別に脅威モデリングするのはあまりにも手間がかかる。そのため設計者は多くの場合、どのトレードオフに見込みがありそうかを直感で見抜き、各設計案を一から評価し直すのではなく、案どうしの違いを分析して比較する必要がある。

ソフトウェアシステム設計の初期段階では、信頼境界と攻撃対象領域に注意深く目を向けるべきである。セキュリティを確保しやすいアーキテクチャを築くうえで、この二つが決定的に重要だからである。
機微な情報のデータの流れは、できる限り、トポロジーの中で最も露出した部分から遠ざけておくべきである。

\begin{exa}[外回りの営業担当者向けアプリケーション]\label{exa:ch6-sales-app}
外出先で営業電話をかけるために、顧客の連絡先情報にオフラインでアクセスする必要がある営業担当者向けのアプリケーションを考える。
顧客データベース全体を各モバイル端末に入れると、情報が露出するリスクは非常に大きくなる。とはいえ、担当者が通信環境のよくない遠隔地に出張するのであれば、それも必要だという言い分は成り立つかもしれない。
脅威モデリングを行えばこのリスクが浮き彫りになり、代替案を検討するきっかけになるだろう。
たとえば、担当者の居場所の変化や出張予定に合わせて動的に更新される、地域ごとのデータベースの一部だけで足りるかもしれない。
あるいは、顧客の電話番号を渡す代わりに顧客ごとのコードを各営業担当者に渡し、担当者固有の PIN と組み合わせて転送サービス経由で電話をかけられるようにすれば、担当者が電話番号にアクセスする必要はまったくなくなる。
\end{exa}

設計者はまた、構築しているソフトウェアの\underline{本質的な脅威モデル}（essential threat model）を、代替設計を評価する基準線のようなものとして考慮すべきである。
本質的な脅威モデルとは、どのように作るかにかかわらず、理想化された設計そのものに内在するセキュリティリスクのモデルのことである。
たとえば、クライアント／サーバーシステムがクライアントから個人を特定できる情報（personally identifiable information, PII）を収集しているなら、その情報がクライアント側で、転送中に、あるいはデータを処理するサーバー側で露出するというセキュリティリスクは避けられない。
どんなに巧みな設計でもこれらのリスクを消すことはできない。ただし、多くの場合は適切な緩和策が必要になる。

内在するセキュリティリスクが高い場合、設計者はできる限り設計の変更を検討すべきである。
PII の例を続けると、すべてのユースケースで、その情報をすべて（あるいは一つでも）収集することが本当に必要だろうか。
必要でないなら、情報の一部を発生源で収集しないで済むような下位のケースをサポートすることには、十分に手間をかける価値があるかもしれない。

本質的な脅威モデルが設計の指針になるもう一つの形は、設計上の決定から生じる新たなリスクの源を浮き彫りにすることである。

\begin{exa}[機微なデータのキャッシュ層]\label{exa:ch6-cache-layer}
たとえば、応答時間を改善しようとして、機微なデータのためのキャッシュ層を追加すると決めたとする。
データを余分に保存すれば（そのデータは攻撃者が狙う資産になりうる）、必ず新しいリスクが加わる。キャッシュストアが攻撃対象領域の近くにある場合はなおさらである。
\end{exa}

この例からわかるように、設計を変更すれば --- よい方向にせよ悪い方向にせよ --- 必ず脅威モデルも変わる。セキュリティへの影響を理解していれば、設計者は代替案の長所を賢明に比較検討できる。

結局のところ、よいソフトウェア設計は主観的な判断に左右される。
そうした判断は、関係するさまざまな要因のバランスをとり、最良とまではいかなくても、少なくとも満足のいく結果を見つけ出す。
セキュリティは重要だが、それがすべてではないので、難しい決定は避けられない。
長年の経験から著者が気付いたのは、ときに怖いと感じても、セキュリティ上の懸念が最優先だと宣言するより、妥協についての議論に開かれた姿勢でいるほうがずっと生産的だということである。

セキュリティを最大限に高めるコストが低ければ、そうするよう推し進めるのは簡単だが、いつもそうとは限らない。
妥協が必要なときに心に留めておくとよい戦略を挙げる。
\begin{itemize}
\item 後からセキュリティ上の保護を追加しやすいよう、柔軟性を持たせて設計する（つまり、セキュアでない袋小路に自分を追い込まない）。
\item 特に懸念している攻撃があるなら、悪用の試みを監視しやすくなるよう、システムに計測の仕組みを組み込んでおく。
\item ユーザビリティとセキュリティが対立する場合は、ユーザーインターフェースの別案を探る。また、試作品を作り、現実的な状況でユーザビリティを測定する。ユーザビリティ上の懸念は思い込みにすぎず、実際には問題にならないこともある。
\item 特定の設計が招きうる大きな不利益を示す潜在的なシナリオ（脅威モデルから導いたもの）を使ってセキュリティリスクを説明し、緩和策を実装しないことのコストを示す。
\end{itemize}

%----------------------------------------------------------------------
\subsection{緩和策の組み込み}\label{subsec:ch6-building-in-mitigations}
% src: Building in Mitigations

ソフトウェアシステムのスコープとセキュリティ要件を定めて 4つの質問の最初の二つに答えたら、次は 3 番目の「それについて何をするのか？」を検討する番である。
この質問は、必要な保護と緩和策を設計に組み込むよう設計者を導く。
以下の各項では、ソフトウェア設計で最もよく繰り返し現れるテーマのうち、インターフェースとデータの二つについて、緩和策をどう組み込むかを検討する。
以下の議論と例は、設計に取り入れられる緩和策のほんの一端にすぎない。
これまでの三つの章の考え方はすべて、個々の設計のニーズに応じて適用できる。

\subsubsection{インターフェースの設計}\label{subsec:ch6-designing-interfaces}
% src: Designing Interfaces

インターフェースはシステムの境界を定め、設計やその構成要素の範囲を画定する。
インターフェースには、システムコール、ライブラリ、ネットワーク（クライアント／サーバー型でもピアツーピア型でも）、プロセス間およびプロセス内の API、共通のデータストア内の共有データ構造などがある。
セキュアな通信プロトコルのような複雑なインターフェースは、それだけで独立した設計を起こす価値があることが多い。

設計のスコープ内のインターフェースはすべて定義し、そのインターフェースを共有するコンポーネントがそれぞれどのようなセキュリティ上の責任を負うのかを、はっきり把握しておくべきである。
入力が確実に検証されるのか、それとも信頼できないデータとして扱うべきなのかを文書化する。
信頼境界がある場合は、境界を越える際の認証と認可をどう扱うかを説明する。

外部コンポーネント（設計のスコープ外にあるもの）とのインターフェースは、そのコンポーネントの既存の設計仕様に従うべきである。
そうした情報が手に入らない場合は、前提を文書化するか、不確実さを補う防御的な戦術を検討する。
たとえば、入力が検証されているかどうかを確かめられないなら、信頼できない入力を想定しておく。

セキュアなインターフェースを設計するには、まずインターフェースの働きをしっかり記述することから始める。そこには必要なセキュリティ特性（すなわち C-I-A、ゴールドスタンダード、プライバシーの要件）も含める。
インターフェースのセキュリティをレビューするとは、インターフェースが正しく機能し、潜在的な脅威に対して頑健であり続けることを検証することにほかならない。
設計者がセキュリティ要件をはっきりさせていなければ、セキュリティのレビュー担当者（と、後でそのインターフェースを使う開発者）は設計者の意図を推測するしかなく、要件を低く見積もるにせよ高く見積もるにせよ、混乱が生じる。

セキュリティを念頭に置いて設計されていない既存のコンポーネントや、要件に照らして十分にセキュアでない既存のコンポーネントを、使わざるをえないこともある。単に、コンポーネントがどれほどセキュアなのかわからないこともある。
選択の余地がないならこれを課題として明記し、可能であれば、そのコンポーネントのセキュリティ特性についてわかることを調べる（テスト用のモックアップを攻撃してみることも含まれるかもしれない）。
場合によっては、インターフェースをラップしてセキュリティ上の保護を追加するという手もある。

\begin{exa}[暗号化層によるストレージコンポーネントのラップ]\label{exa:ch6-wrap-storage}
データ漏えいに対して脆弱なストレージコンポーネントがあるなら、暗号化と復号を担うソフトウェア層を追加で設計し、そのコンポーネントには暗号化されたデータ --- 開示されても害のないもの --- だけを保存させることができる。
\end{exa}

\subsubsection{データの取り扱いの設計}\label{subsec:ch6-designing-data-handling}
% src: Designing Data Handling

データの取り扱いは事実上すべての設計の中心にあるので、その部分をセキュアにすることは重要な一歩である。
セキュアなデータの取り扱いは、データ保護の目標を書き出すところから始めるとよい。
データの特定の部分集合に追加の保護が必要なら、そのことを明示し、設計全体を通じて一貫して扱われるようにする。
たとえば、オンラインショッピングのアプリケーションでは、クレジットカード情報に追加の保護措置を施す。

機微なデータを持ち回る必要をできるだけ減らす。
これは設計レベルでリスクの露出を大幅に減らせる重要な機会であり（第4章の最小情報（Least Information）パターンを参照）、実装の段階になってからではたいてい手遅れである。
データを持ち回る必要を減らす方法の一つは、データを\underline{不透明な識別子}（opaque identifier）に関連付け、その識別子を、必要なときに実際のデータに変換できるハンドルとして使うことである。

\begin{exa}[不透明な識別子によるログ記録]\label{exa:ch6-opaque-identifier-logging}
付録A のサンプル設計のように、こうした識別子を使ってトランザクションをログに記録すれば、顧客の詳細情報をシステムログに残さずに済む。
ログエントリを調査する必要が生じるまれなケースでは、監査担当者がその詳細を調べればよい。
\end{exa}

公開情報や、そのほかの理由で機密性の要件の対象外となるデータを特定する。
このようなデータはデータの取り扱い要件の重要な例外となり、理にかなう範囲で保護を緩めることができる。
ただし、この方法を使う際には、データの意味は文脈によって変わるので、公開データでもほかの情報と組み合わされると十分に機微な情報になりうることを忘れてはならない。

\begin{exa}[公開情報と組み合わせによる機微性]\label{exa:ch6-public-context}
ほとんどの企業の所在地や最高経営責任者の名前は、通常は公開情報である。
しかし、名前の挙がった人物が正確にいつ社屋にいるのかは、非公開にしておくべきである。
\end{exa}

はっきりと別の決定をしていない限り、個人情報は常に機微なものとして扱い、そもそも具体的な用途があるときにだけ収集すべきである。
機微なデータを無期限に保存すると、そのデータを守る義務がいつまでも続く。
この義務を避ける最良の方法は、使われなくなった情報を可能なときに（たとえば何年も使われていない場合に）破棄することである。
設計では、不要になった私的なデータをいずれシステムから削除する必要が生じることを見越し、バックアップコピーも含めて、どのような条件で削除を実行するのかを定めておくべきである。

%----------------------------------------------------------------------
\subsection{設計へのプライバシーの統合}\label{subsec:ch6-integrating-privacy}
% src: Integrating Privacy into Design

私的な情報の保護に失敗した事件は、日常的にニュースの見出しを飾っている。
著者は、情報プライバシーへの配慮をソフトウェア設計に組み込むことが、企業が状況を改善するための重要な方法だと考えている。
プライバシーの問題はデータ保護が人に及ぼす影響にかかわるもので、法律や規制の問題だけでなく、顧客の期待や、許可のない開示がもたらしうる影響も含む。
これを正しく扱うには、特別な専門知識と主観的な判断が必要である。
しかし問題の一部は、データを使う認可を第三者に与えること --- それにはアクセスを許すことが伴う --- にかかっている。その範囲については、よいソフトウェア設計によって、過ちを最小限に抑える制御を設けることができる。

出発点として、設計者は適用されるすべてのプライバシーポリシーを熟知し、それらが設計とどう関係するかを理解しておくべきである。
疑問点は質問し、要件がはっきりするよう、できればプライバシーポリシーの責任者から書面で回答をもらう。
パートナー経由で取得したデータに適用されうる、第三者のプライバシーポリシー上の義務もこれに含まれる。
プライバシーポリシーはデータの収集・利用・保存・共有を規定しているので、設計の中でこれらの活動が行われるなら、ポリシーの規定がそのまま要件になる。
公開されているプライバシーポリシーに詳細があまり書かれていない場合は、必要な詳細を記した内部向けの版を作ることを検討する。

プライバシーの過失は、人やプロセスがポリシーの約束を誤解したり、そもそも約束を考慮し損ねたりしたときに起こりやすい。
データセキュリティの保護策を利用すれば、ポリシーの順守を確実にするための制限を設計に組み込める。
まずプライバシーポリシーが明確に約束していることを検討し、可能であれば設計でその約束を強制するようにする。

\begin{exa}[「データを共有しません」というポリシーとクラウドストレージ]\label{exa:ch6-no-sharing-policy}
ポリシーに「当社はお客様のデータを共有しません」と書かれているなら、共有が簡単にできるクラウドストレージサービスを使うことには慎重になるべきである。設定ミスでデータが露出しないことを保証するほかの措置が講じられていれば別である。
\end{exa}

監査は、機微なデータへの適切なアクセスを確実に記録するという目的だけをとっても、プライバシーを管理するうえで重要な道具である。
アクセスを注意深く監視していれば、問題のあるアクセスや利用を早期に検知して是正できる。
漏えいが起きた後で、問題のデータに誰がアクセスしたのかの記録がなければ、効果的に対応するのは非常に難しい。

できる限り、明示的なプライバシー保護を設計に盛り込む。
プライバシーへの準拠について自分で判断できない場合は、プライバシーポリシーの責任者に設計を承認してもらう。
ソフトウェア設計にプライバシーを組み込む一般的な手法には、次のようなものがある。
\begin{itemize}
\item 新しい種類のデータを収集する箇所を特定し、その収集がプライバシーポリシーに準拠していることを確認する。
\item 意図した目的でそのデータを使うことをポリシーが認めているかを確認する。
\item 設計上データを無制限に利用できてしまう可能性がある場合は、プライバシーポリシーの制約と準拠状況の監査方法をよく知っているスタッフだけにアクセスを限ることを検討する。
\item ポリシーでデータの保持期間が制限されているなら、期限どおりの削除を保証するシステムを設計する。
\item 設計の進化に伴ってデータベースのあるフィールドが使われなくなったら、開示のリスクを減らすためにそのフィールドを削除することを検討する。
\item データを共有する際に、受け取る側が経営陣の承認を得ていることを確かめる承認プロセスを組み込むことを検討する。
\end{itemize}

%----------------------------------------------------------------------
\subsection{ソフトウェアのライフサイクル全体の計画}\label{subsec:ch6-full-lifecycle}
% src: Planning for the Full Software Lifecycle

あまりにも多くのソフトウェア設計が、システムが永遠に使われ続けることを暗黙のうちに仮定し、どんなソフトウェアにも寿命があるという現実を無視している。
最初のリリースとデプロイから、更新と保守を経て、最終的な廃止に至るまで、システムの一生にかかわる多くの側面が重要なセキュリティ上の意味を持ち、しかも後になると見落とされやすい。
どれほど素晴らしいソフトウェア設計であっても、軌道に乗るにせよ立ち消えになるにせよ、環境の変化に伴って変わっていく。
こうした変化の影響は、設計プロセスの段階で予測して手を打っておくのが最善であり、少なくとも後の人のために書き残しておくべきである。
企業の中では、こうした問題の多くは共通のものであり、一般的な形で扱えばほとんどのシステムをカバーできるはずである。例外は、必要に応じて個々の設計で定めればよい。

新しい設計を作っている最中にシステムの寿命の終わりを想像するのは難しいが、その影響の大半は明らかなはずであり、どの設計も少なくともデータの長期的な処分については検討すべきである。
特定の法的・業務上の理由から一定期間データを保持しなければならないこともあるが、不要になったらバックアップコピーも含めて破棄すべきである。
システムによっては、寿命の終わりが近づくと特定の段階を踏む必要がある。よい設計なら、適切な構造と設定オプションを最初から用意しておくことで、この段階を正しく踏みやすくできる。

\begin{exa}[購買システムの終了段階]\label{exa:ch6-purchasing-end-of-life}
購買システムは、注文の受け付けを停止した後も、給与や記録保持の目的でさらに 1 年間データを提供し続け、その後、取引記録を長期保存用にアーカイブする必要がある、といったことがありうる。
\end{exa}

%----------------------------------------------------------------------
\subsection{トレードオフ}\label{subsec:ch6-trade-offs}
% src: Making Trade-offs

簡単な選択肢がない中でトレードオフのバランスをとるには、ほかの多くの考慮事項を比較検討する高度な工学的判断が必要である。
セキュリティ上の緩和策を増やせばリスクは減るが、それも複雑さのせいで全体としてバグが増えるところまでである。収穫が逓減する中で開発の労力ばかり増えることには、常に用心すべきである。
本書は競合する優先順位の間で妥協するよう繰り返し設計者に助言するが、言うは易く行うは難しである。
この節では、こうした重要なバランスをとるための経験則をいくつか紹介する。

最悪のシナリオを予想する。つまり、あるシステム資産の機密性・完全性・可用性を守れなかったとしたら、どれほどひどいことになるかを考える。
各シナリオについて、破局の程度を考慮する必要がある。
データのどれだけが影響を受けうるか。
どれくらいの期間利用できなくなると深刻な問題になるか。
主要な緩和策は通常、最悪のケースの被害に上限を設ける。たとえば、1 時間ごとにバックアップを取っていれば、失われるおそれのある取引データはせいぜい 1 時間分で済むはずである。

\begin{rem}[機密性の喪失は上限を設けにくい]\label{rem:ch6-confidentiality-cap}
最悪のケースにおける機密性の喪失は、とくに上限を設けにくいことに注意すべきである。
いったんデータが盗まれてしまえば、開示をなかったことにする方法は普通は考えられないからである（2017 年の Equifax の侵害はその際立った例である）。
\end{rem}

ほとんどの設計作業は、企業やプロジェクトのコミュニティの中で行われ、そこでは必要なセキュリティの水準が幅広いプロジェクトを通じてたいてい一定である。
特定の設計がその水準から外れ、より高い、あるいはより低いセキュリティ水準を必要とする場合には、その前提を設計の序文で明示しておく価値が十分にある。
この重要な点を、いくつかの例で明らかにしよう。

\begin{exa}[セキュリティ水準を上げる場合と下げる場合]\label{exa:ch6-security-level}
オンラインストアの Web サイトでは、クレジットカード処理を担うソフトウェアについて、セキュリティの基準を高めに設定することを検討すべきである。クレジットカード処理は明らかに攻撃の標的になり、莫大な金銭的責任を伴うために特別な要件の対象にもなるからである。
逆に、Web デザイン会社が自社のデザインの実例を紹介する Web サイトを丸ごと公開することもあるだろう。これは情報提供だけが目的で、実際のエンドユーザーのデータを収集することはないので、セキュアにすることの重要性がそれに応じて低いのは妥当である。
\end{exa}

設計段階は、ソフトウェアに対する競合する要求の間で正しいバランスをとる最良の機会である。
率直に言えば、スケジュールの締め切り、予算と人員の制約、レガシーとの互換性の問題、そしておなじみの長い機能リストを抱えているとき --- つまりほとんどいつも --- セキュリティが最優先事項として全面的に支持されることは、まずない。
設計者は、急進的なものも含めて多くの代替案を検討し、後からでは試みることもできないような根本的な変更を加えられる、最も有利な立場にいる。

こうした理想化された原則と、現実世界のシステムを作るうえでの実際的な要求との間で正しいバランスをとることこそ、セキュアなソフトウェア設計の核心である。
目標は決して完璧なセキュリティではなく、緩和策を追加して得られる利益にも限界がある。
最適な落としどころが正確にどこにあるかを見極めるのは決して簡単ではないが、こうしたトレードオフを明示しているソフトウェア設計のほうが、分別のある妥協点を見つけられる見込みが高い。

%----------------------------------------------------------------------
\subsection{設計の単純さ}\label{subsec:ch6-design-simplicity}
% src: Design Simplicity

\begin{quote}
「単純さは究極の洗練である。」\\
\hfill --- Leonardo da Vinci
\end{quote}

皮肉なことに、da Vinci の言葉が示すとおり、単純な設計を生み出すにはかなりの思考と労力が必要になることが多い。
ルネサンス期の天文学者たちは天体の運動を説明するためにありとあらゆる複雑な計算を編み出したが、やがて Copernicus が地球の代わりに太陽を中心の基準点に据えてモデルを単純化した。そのおかげで Newton は重力の法則を導き、計算を根本から単純にできた。
著者が気に入っている見事なソフトウェア設計の例は *nix オペレーティングシステムの心臓部であり、その多くは今日でも使われ続けている。
美しく単純な設計を追い求めることは、めったに実を結ばないとしても、多くの場合そのままセキュリティの向上につながる。

ソフトウェア設計における単純さにはさまざまな形があるが、最も単純で最もエレガントな設計を見つける簡単な公式はない。
第4章で論じたパターンのいくつか、たとえば設計の簡潔さ（Economy of Design）や最小共通メカニズム（Least Common Mechanism）は、単純さを取り入れたものである。
セキュリティが、何か複雑な判断や仕組みを寸分違わず正しく行えるかどうかにかかっているときは、常に用心すべきである。同じ目的をもっと単純に達成する方法がないか探るとよい。

\begin{exa}[Three Mile Island の事故と複雑さ]\label{exa:ch6-three-mile-island}
入り組んだ機能がセキュリティの仕組みと絡み合うと、多くの場合、複雑さが爆発的に増す。
ある研究は、1979 年の Three Mile Island 原子力施設の事故には特定の原因がなく、多数の冗長な安全対策を含むシステムの途方もない複雑さが原因だったと結論付けた。
\end{exa}

セキュリティは、やろうとしていることの邪魔になりうる。そして逆に、すべてをセキュアにすることはますます厄介になる。
この場合の解決策は、多くの場合、セキュリティを機能から切り離して階層化したモデルを作ることである。通常は、セキュリティを保護の殻として「外側」に置き、機能はすべてそれとは別に「内側」に置く。

\begin{rem}[固い殻と柔らかい内側]\label{rem:ch6-hard-shell}
ただし、固い殻と「柔らかい内側」という形で設計するなら、その分離を確実に守らせることが決定的に重要になる。
城の周りにセキュアな堀を設計するのは比較的簡単だが、ソフトウェアでは、外側の保護層を迂回して内側に通じる経路を、うっかり開けてしまいやすい。
\end{rem}

\newpage

%% =======================================================
%%  part: ch7
%% =======================================================
%% =======================================================
%%  part: ch7
%% =======================================================
\section{セキュリティ設計レビュー}
% src: 7: Security Design Reviews

この章では、\cite{Kohnfelder2021} の第7章に沿って、セキュリティ設計レビュー（SDR）を扱う。

\begin{quote}
「好意的で思いやりのあるレビューは、いつだってうれしい驚きである。」\\
\hfill --- Joyce Carol Oates
\end{quote}

ソフトウェアにセキュリティを組み込む最良の方法の一つは、\underline{セキュリティの帽子}（security hat）をかぶって、設計を別途レビューすることである。
この章では、前章で論じたセキュリティとプライバシーの設計概念を、セキュリティ設計レビューでどう生かすかを説明する。

\begin{dfn}[セキュリティ設計レビュー]\label{dfn:ch7-sdr}
ソフトウェア設計のセキュリティを構造立てて検討するレビューを、\underline{セキュリティ設計レビュー}（security design review, SDR）という。
\end{dfn}

この過程は、建築家が建物を設計し、その設計が安全で堅固かどうかをエンジニアがレビューして確かめる場面にたとえるとわかりやすい。
設計者もレビュー担当者も構造工学と建築基準を理解している必要があり、両者が協力することで、品質と信頼をより高い水準に引き上げられる。

理想を言えば、セキュリティのレビュー担当者は設計作業に関わっていない人がよい。そうすれば設計から距離を置き、客観的に見られる。
同時に、そのソフトウェアが動くシステムや文脈、そしてソフトウェアの使われ方をよく知っている人であることが望ましい。
ただし、これらは必須の条件ではない。
設計にあまり詳しくないレビュー担当者は質問がずっと多くなりがちだが、それでも立派な仕事ができる。

こうした方法を広め、より多くのソフトウェア専門家に自分で SDR を行ってもらうことは、原著者が本書を書いた中心的な目的の一つであった。
自分がふだん扱っていてよく知っているソフトウェアシステムなら、セキュリティの経験は豊富でもそのシステムを知らない人より、ほぼ確実に良い SDR ができる。
本書はそのための指針を示すものであり、それがソフトウェアセキュリティの水準を引き上げることに少しでも役立てば、と原著者は願っている。

%----------------------------------------------------------------------
\subsection{SDR の実施体制}
% src: SDR Logistics

セキュリティ設計レビューの方法論を示す前に、少し背景を説明し、実施体制の基本について論じておく。
SDR は何の役に立つのか。
実施するなら、設計過程のどの段階で行うべきか。
最後に、準備のヒントをいくつか、とくに文書化の重要性について述べる。

\subsubsection{なぜ SDR を行うのか}
% src: Why Conduct an SDR?

原著者はこれまでに数百件の SDR を行ってきたが、時間の無駄だと感じたことは一度もない。
SDR にかかる時間は設計時間全体のごく一部にすぎず、しかも結果は、セキュリティを高める重要な改善点が見つかるか、設計がセキュリティにきちんと対処しているという強い保証が得られるかのどちらかである。
単純で素直な設計ならレビューはすぐに終わるし、大きな設計では、レビューの過程が主なホットスポットを洗い出して検証するための有用な枠組みになる。
一見するとセキュリティの手をすべて打ってあるように見える設計でも、それを確認するのは当然払うべき注意（デューデリジェンス）である。
そしてもちろん、SDR で重大な問題が見つかれば、かけた労力は十分すぎるほど報われる。
そうした問題を実装中に見つけるのは難しく、後から直すには大きなコストがかかるからである。

さらに、SDR から価値ある新しい気付きが得られ、セキュリティとは関係のない設計変更につながることもある。
SDR は、多様な視点（ユーザーエクスペリエンス、カスタマーサポート、マーケティング、法務など）を取り込む絶好の機会でもある。悪用の可能性や意図しない結果のような見落とされがちな論点について、全員で考えることになるからである。

\subsubsection{SDR をいつ行うか}
% src: When to Conduct an SDR

SDR は、設計（または設計の一回の反復）が完成して安定した時点で行うよう計画するとよい。典型的には機能レビューの後で、ただし変更が必要になるかもしれないので、設計を確定する前に行う。
セキュリティを機能レビューの一部として扱おうとするのは、強くお勧めしない。
考え方も注目すべき領域も大きく異なるからである。
また、レビュー担当者だけでなく全員がセキュリティに集中することが大切だが、合同レビューの場ではどうしても設計の動作のほうに関心が向きがちで、セキュリティに集中しにくい。

複雑な設計やセキュリティ上重要な設計では、設計が固まり始めたもののまだ形になりきっていない段階で、予備的な SDR を追加で行うと役に立つことが多い。
そうすれば、主要な脅威と全体的な戦略について早めに意見をもらえる。
予備的な SDR はそれほど形式ばらなくてよい。セキュリティ上とくに気になる点（あとでさらに掘り下げることになりそうな箇所）に先に目を通し、セキュリティ上のトレードオフを大まかなレベルで議論する。

\begin{rem}[設計者とレビュー担当者の責任]\label{rem:ch7-designer-responsibility}
優れたソフトウェア設計者は、設計の全期間を通じて、セキュリティとプライバシーの問題を常に考慮し、対処すべきである。
はっきり言っておくと、設計者はセキュリティを無視して、その問題を SDR で直してもらおうと当てにしては決してならない。
設計者は常に、自分の設計のセキュリティに全責任を負うものと心得ておくべきである。セキュリティのレビュー担当者は、設計者が抜かりなく仕事をできるよう手助けする支援役である。
一方で、セキュリティのレビュー担当者も上から説教すべきではない。良し悪しの判定を交えずに、自分の発見事項を明確に、説得力をもって設計者に示すべきである。
\end{rem}

\subsubsection{文書化は不可欠である}
% src: Documentation Is Essential

SDR を効果的に行えるかどうかは、最新の文書があるかどうかにかかっている。
文書があってこそ、関係者全員がレビュー対象の設計を正確に、食い違いなく理解できるからである。
口頭だけの非公式な SDR でも何もしないよりはましだが、決定的な詳細が抜け落ちたり誤って伝わったりしやすく、書面に残さなければ、せっかくの成果も簡単に失われてしまう。
原著者自身は、会議の前に設計書に目を通しておけるほうが常に好ましいと考えている。
そうすれば事前に設計の検討を始められ、何に取り組んでいるのかを把握するために会議の時間を使わずに済む。

原著者の経験では、設計書の質が高いことは、優れた SDR を行ううえでかけがえのない助けになる。
もちろん実際には、十分な文書が手に入らないこともある。この章の後半のケーススタディでは、そうした状況への対処についても述べる。

\begin{exa}[具体性を欠く記述]\label{exa:ch7-store-securely}
例えば、「顧客データを安全に保存する」と漠然と書いてあるだけの設計書は、それが具体的に何を意味し、どう実現するのかを続けて説明していない限り、大きな危険信号と受け止めるべきである。
具体性のない包括的な記述は、ほぼ例外なく、未熟さとセキュリティへの確かな理解の欠如を露呈している。
\end{exa}

%----------------------------------------------------------------------
\subsection{SDR の過程}
% src: The SDR Process

以下で説明する SDR の過程は、形式的で必須のレビュープロセスを持つ大規模なソフトウェア企業で、原著者が実際に SDR を行っていたやり方である。
とはいえ、ソフトウェア設計の進め方は千差万別であり、同じ戦略と分析を、もっと形式ばらない組織に合わせて応用することもできる。

\begin{dfn}[SDRの6つの段階]\label{dfn:ch7-six-stages}
明確で完全な設計が書面になっていることを出発点として、SDR は次の 6 つの段階からなる。
\begin{enumerate}
\item \underline{調査}（Study）：設計と補助文書を調べ、プロジェクトの基本を理解する。
\item \underline{質問}（Inquire）：設計と基本的な脅威について、設計チームに明確化のための質問をする。
\item \underline{特定}（Identify）：設計のうちセキュリティ上最も重要な部分を特定し、重点的に検討する。
\item \underline{協働}（Collaborate）：設計者と協力してリスクを特定し、緩和策を議論する。
\item \underline{報告書作成}（Write）：発見事項と推奨事項をまとめた報告書を書く。
\item \underline{フォローアップ}（Follow up）：その後の設計変更を追跡し、問題が解決されたことを確認してから承認する。
\end{enumerate}
\end{dfn}

小さな設計なら、これらの大部分を 1 回のセッションで済ませられることが多い。
大きな設計では段階ごとに作業を分け、段階によっては終えるまでに何回かのセッションが必要になることもある。
設計チームと打ち合わせるためのセッションを専用に設けるのが理想だが、必要ならレビュー担当者が一人で作業し、メールなどで設計チームとメモや質問をやりとりしてもよい。

進め方は人それぞれである。
一気に取りかかって「マラソン」のように片付けるのを好むレビュー担当者もいる。
原著者は、数日にわたって少しずつ進め、「一晩寝かせる」時間をとるやり方を好み、人にも勧めている。
原著者の場合、最良のアイデアはそこから生まれることが多い。

以下では、SDR の過程の各段階を説明し、役に立つ手法を箇条書きでまとめる。
SDR を行う際には、過程を進めながら各段階の箇条書きを参照するとよい。

\subsubsection{調査}
% src: 1. Study

レビューの準備として、設計と補助文書を調べ、そのソフトウェアの基本を理解する。
レビュー担当者には、セキュリティの知識に加えて、できればその分野固有の専門知識も求められる。
専門知識がない場合は、身につけられるだけ身につけ、過程全体を通じて好奇心を持って学び続けるようにする。
セキュリティ上の判断のほとんどにはトレードオフがつきものなので、ひたすらセキュリティの強化だけを推し進めると、やり過ぎに終わりやすく、その過程で設計を台無しにしかねない。

\begin{exa}[火災対策だけを考えた家]\label{exa:ch7-fireproof-house}
セキュリティが過剰だとなぜ困るのかは、火災のリスクを減らすことだけを目的に設計した家を考えてみるとわかる。
全体がコンクリート造りで、厚い鋼鉄の扉が一つあるだけで窓もない家は、費用がかさむうえに見た目も悪く、誰も住みたがらないだろう。
\end{exa}

この準備段階では、次のことを行う。
\begin{itemize}
\item まず、文書を読んで設計の大まかな姿をつかむ。
\item 次に、「セキュリティの帽子」をかぶり、脅威を意識しながらもう一度読み通す。
\item メモをとり、考えたことや気付いたことを後で参照できるよう記録しておく。
\item 問題になりそうな箇所には後のために印を付けておく。ただし、この段階でセキュリティ分析に深入りするのは時期尚早である。
\end{itemize}

\subsubsection{質問}
% src: 2. Inquire

システムに対する基本的な脅威を理解するために、設計者に明確化のための質問をする。
理解しやすい単純な設計の場合や、設計者がしっかりした文書を用意している場合には、この段階を省けることもある。
この段階は、先に進む前に自分の理解が正しいかを確かめ、曖昧な点や残っている疑問を解消する機会と考えるとよい。
効果的なレビューを行うのに、レビュー担当者が設計を隅々まで知っている必要はまったくない（それは設計者の仕事である）。
ただし、設計の大まかな輪郭と、主要なコンポーネントどうしがどう連携するかは、しっかり把握しておかなければならない。

この段階は、掘り下げに入る前に穴を埋めておく機会である。
いくつか指針を挙げる。
\begin{itemize}
\item 設計書が明確で完全であることを確かめる。
\item 抜けている点や訂正が必要な点があれば、文書が修正されるよう手助けする。
\item 設計について会話ができる程度には理解するが、専門家レベルである必要はない。
\item 何を一番心配しているかをチームのメンバーに尋ねる。セキュリティ上の懸念はないと言われたら、なぜないのかを知るために追加の質問をする。
\end{itemize}

セキュリティのレビュー担当者として尋ねる質問を、設計書に書かれていることだけに限る必要はない。
関連する周辺のシステムを理解しておくと、それらが設計のセキュリティにどう影響するかを見積もるのに大いに役立つ。
書かれていない詳細ほど見つけにくいものである。
例えば、データの扱いについて何も詳しく書かずに、設計が暗黙のうちにデータを保存しているなら、そのストレージとそのセキュリティについて尋ねるとよい。

\subsubsection{特定}
% src: 3. Identify

設計のうちセキュリティ上重要な部分を特定し、そこに絞って綿密に分析する。
基本原則から出発してセキュリティのレンズを通して見るようにし、C-I-A、ゴールドスタンダード、資産、攻撃対象領域、信頼境界といった観点で考える。
設計のこうした部分には特別な注意を払う価値があるが、ほかの部分をまったく見落とすことのないよう、当面はレビューの対象を設計全体に広げておく。
とはいえ、セキュリティにほとんど、あるいはまったく関係のない部分は飛ばしてかまわない。

この探索的な段階では、次のことを行う。
\begin{itemize}
\item インターフェース、ストレージ、通信を検討する。ふつうはこれらが注目の中心になる。
\item 決意を固めた攻撃者と同じように、最も露出した攻撃対象領域から最も価値の高い資産へと、内側に向かって進んでいく。
\item 設計者がどの程度セキュリティを意識しているか、設計がどの程度はっきりとセキュリティに対処しているかを評価する。
\item 必要なら、重要な保護策を指摘し、それを重要な機能として設計の中に明記してもらう。
\end{itemize}

\subsubsection{協働}
% src: 4. Collaborate

設計者と協力し、発見事項を伝え、代替案を議論する。
理想的には、設計者とレビュー担当者が集まって話し合い、問題を一つずつ検討していく。
これは全員にとって学びの過程である。
設計者はセキュリティについて学びながら自分の設計を新鮮な目で見直せるし、レビュー担当者は設計とその意図についての洞察を得て、セキュリティ上の課題と最善の緩和策の選択肢への理解を深められる。
共通の目標は設計を全体として良くすることであり、セキュリティはレビューの焦点ではあるが、考慮すべき唯一の事項ではない。
その場で変更を最終決定する必要はないが、どのような設計変更を検討すべきかについて、最終的には合意しておくことが重要である。

効果的に協働するための指針をいくつか挙げる。
\begin{itemize}
\item レビュー担当者として、必要に応じてリスクと緩和策についてセキュリティの観点から意見を述べる。設計がすでに安全な場合でも、優れたセキュリティ実践を後押しすることになるので価値がある。
\item 緩和策の必要性を設計者に納得してもらうために、セキュリティ上の変更が将来どう報われるかを示すシナリオを描いてみるのも一つの手である。
\item 可能なら、一つの問題に対して複数の解決策を示し、それぞれの代替案の長所と短所を設計者が理解できるよう手助けする。
\item 設計の最終的な責任を負うのは設計者なので、最終的な判断は設計者が下すものと受け入れる。
\item 設計に取り入れるもの、取り入れないものも含めて、意見のやりとりを文書に残す。
\end{itemize}

\begin{rem}[「最終的な判断」の実際]\label{rem:ch7-last-word}
「最終的な判断」について補足しておくと、実際にはこのバランスは、組織とその文化、適用される業界標準、規制上の要件の有無、その他の要因によって変わる。
大規模な組織や統制の厳しい組織では、最終的な判断を下すのに、アーキテクチャ委員会、標準遵守の担当者、ユーザビリティの評価者、経営陣の利害関係者など、複数の当事者の承認が必要になることもある。
承認が複数必要な場合、設計者は競合する利害のバランスをとらなければならない。セキュリティのレビュー担当者はこうした力関係にとくに配慮し、できる限り柔軟に対応すべきである。
\end{rem}

\subsubsection{報告書作成}
% src: 5. Write

レビューで得た発見事項と推奨事項について、\underline{評価報告書}（assessment report）を書く。

\begin{dfn}[評価報告書]\label{dfn:ch7-assessment-report}
セキュリティ設計レビューの結果を書面にまとめたもので、発見事項と推奨事項を優先順に並べた要約（具体的な推奨設計変更を含む）からなるものを、\underline{評価報告書}（assessment report）という。
\underline{発見事項}（findings）とは、設計のセキュリティについてセキュリティレビュー担当者が下した評価である。
\end{dfn}

報告書では、検討すべき設計変更の候補と、現状の設計のセキュリティについての分析に焦点を当てる。
設計者がすでに同意した変更は、そうとわかるように目立たせて示し、後で検証の対象にする。
提案する変更には、例えば次のような簡単な 3 段階の方式で優先度を付けるとよい。

\begin{dfn}[推奨変更の優先度]\label{dfn:ch7-priority-rankings}
推奨する変更の優先度として、次の 3 段階を用いる。
\begin{itemize}
\item \underline{Must}：最も強い段階で、選択の余地がないことを示し、多くの場合は緊急であることも意味する。
\item \underline{Ought}：中間の段階。原著者は、レビュー担当者としては「Must」寄りだが議論の余地はある、ということを表すのに使う。
\item \underline{Should}：最も弱い段階で、任意の推奨変更に使う。
\end{itemize}
\end{dfn}

もっと細かく順位を付けたい場合は、第13章に、セキュリティバグに体系的に細かな順位を付ける方法の指針があり、この目的にもそのまま流用できる。

SDR は一件一件が大きく異なるので、原著者は評価報告書に決まったテンプレートを使ったことがなく、発見事項は文章で書いている。
原著者は、レビューの過程でとった大まかなメモをもとに作業し、報告書が自然に最終的な形になっていくのを好む。
すべての詳細を頭の中に確実にとどめておけるなら、レビュー会議の後に報告書を一気に書き上げてもよい。

次のヒントは、報告書を書く際の枠組みとしても使える。
\begin{itemize}
\item セキュリティリスクに対処する具体的な設計変更を軸に、報告書を組み立てる。
\item 労力と紙幅の大部分は最も優先度の高い問題に割き、優先度の低いものにはそれに見合う分だけ割く。
\item 設計者の仕事を肩代わりしようとはせず、代替案と戦略を提案する。
\item 優先度の段階を使って、発見事項と推奨事項に優先順位を付ける。
\item セキュリティに焦点を当てるが、設計者に検討してもらいたいほかの所見を書き添えてもかまわない。SDR の範囲外のことについては控えめにし、重箱の隅をつつかず、セキュリティについてのメッセージがぼやけないようにする。
\end{itemize}

設計者とレビュー担当者の役割を分けることは重要だが、実際にどう分けるかは、それぞれの責任と協力し合える度合いによって大きく変わる。
評価報告書では、設計作業そのものには踏み込まずに、必要な変更について明確な方向性を示し、設計者が何をすべきかわかるようにする。
このレビューの結果として大幅な再設計を行う場合は、再設計をレビューしてコメントすることを申し出る。
経験則として、優れたレビュー担当者は、実際の設計変更を細かく指図するのではなく、設計者がセキュリティ上の脅威とそれが招きうる結果を理解できるよう手助けし、緩和策の戦略を提案する。
要求が厳しすぎるレビュー担当者は、たとえ言っていることが正しくても助言が受け入れられないことが多く、設計者が十分に理解していない、あるいは必要だと感じていない変更を押し付けてしまうおそれもある。

ここまで厳密にやるのは細かすぎると感じるなら、報告書は手短に済ませてもよい。ただし後になって、自分自身かそのソフトウェアに携わる誰かが、詳細を記録しておけばよかったと思う可能性は高い。
少なくとも、記録のためにチームに要約をメールで送る時間はとることを勧める。
最小限の報告書であっても、単に「良さそうです！」と言うだけで済ませず、実質のある要約で裏付けるべきである。
設計がセキュリティの要所をすべて押さえていたなら、セキュリティを支えている最も重要な設計上の特徴をいくつか挙げ、その重要性を強調しておく。
セキュリティが問題にならない設計の場合（例えば原著者は以前、個人情報を一切集めない情報提供用の Web サイトをレビューしたことがある）には、その結論に至った理由を簡単に述べておく。

こうした報告書のスタイル、長さ、詳しさは、組織の文化、使える時間、利害関係者の数など、多くの要因によって大きく異なる。
レビュー担当者がソフトウェア設計者と緊密に協力している場合は、変更が必要な問題を報告書に並べる代わりに、二人で協力して必要な規定を設計書に直接書き込めることもある。
小規模で非公式なプロジェクトであっても、設計者とレビュー担当者の役割を別々の人に割り当てることには価値がある。
そうすれば複数の目で作業を見ることになり、セキュリティがきちんと考慮されることが保証される。
とはいえ、一人で設計する場合でも、設計者がセキュリティの帽子をかぶって自分の仕事を見直し、新鮮な目で見ることは役に立つ。

\subsubsection{フォローアップ}
% src: 6. Follow Up

セキュリティレビューの結果として合意した設計変更をフォローアップし、正しく対応されたことを確認する。
協働がうまくいった場合、原著者はふつう、実装までは見ずに文書が更新されたことだけを確認する（そしてこのやり方が裏目に出たことは一度もない）。
そうでない状況では、場合によっては、レビュー担当者がもっと念入りに確認する必要があるかもしれない。
必要な変更をすべて検証し終えてレビューが完了したら、レビューを承認する。
プロジェクトのバグトラッカーに SDR を登録して担当を割り当てておくと、進捗を確実に追跡できる。
そうでなければ、好みに応じて、多少なりとも形式的な手順を使えばよい。
この最終段階についての指針をいくつか挙げる。
\begin{itemize}
\item セキュリティ設計を大きく変更する場合は、変更が正しく行われるよう、設計者と協力するとよいことがある。
\item 意見が分かれる場合、レビュー担当者は双方の立場と、採用されなかった具体的な推奨事項を記載し、未解決の問題として示すべきである（この話題は後の「意見の不一致への対処」の項で詳しく述べる）。
\item 最もうまくいった場合には、設計者がレビュー担当者をセキュリティの相談相手とみなし、必要に応じて長く関わり続けてもらえるようになる。
\end{itemize}

%----------------------------------------------------------------------
\subsection{設計のセキュリティの評価}
% src: Assessing Design Security

SDR の過程を説明したので、この節では、レビューを行うときの考え方に踏み込む。
SDR を行うのに必要な概念と道具は、本書のここまでの内容でそろっている。
基礎となる原則、脅威モデリング、設計手法、パターン、緩和策、暗号の道具は、どれも安全な設計を形づくるものである。

\subsubsection{4つの質問を指針として使う}
% src: Using the Four Questions as Guidance

第2章で脅威モデリングに使った 4つの質問（Four Questions）は、効果的な SDR を行うための優れた指針にもなる。
時間があって労力をかけられるなら明示的な脅威モデリングを行うのがよいが、そうでない場合は、4つの質問を試金石として使うと、レビューに脅威の視点をうまく取り込める。
詳しくは以下の各項で説明するが、ごく大まかに言えば、これらの質問は SDR と次のように対応する。

\begin{dfn}[SDRにおける4つの質問]\label{dfn:ch7-four-questions}
\begin{enumerate}
\item \underline{何に取り組んでいるのか？}（What are we working on?）：レビュー担当者は、レビューの前提として、設計の大まかな目標を理解しておくべきである。その目標を達成する最も安全な方法は何か。
\item \underline{何がうまくいかない可能性があるか？}（What can go wrong?）：ここで「セキュリティの帽子」をかぶった考え方の出番となり、脅威モデリングを当てはめる。設計は重大な脅威を見落としたり、過小評価したりしていないか。
\item \underline{それに対して何をするのか？}（What are we going to do about it?）：設計に組み込まれている保護策と緩和策をレビューする。重要な脅威に対して、もっと良い対応はできないか。
\item \underline{うまくやれたか？}（Did we do a good job?）：設計中の緩和策が十分か、さらに手を加えるべきものはないか、欠けているものはないかを評価する。設計はどの程度安全か。不十分なら、どうすれば水準まで引き上げられるか。
\end{enumerate}
\end{dfn}

SDR に取り組んでいる間は、4つの質問を考えるきっかけとして使える。
設計書を読んで気になる領域にメモを付けたものの、まだ何を探せばよいのかはっきりしないときは、4つの質問、とくに 2 番目と 3 番目を一通りたどり、設計の具体的な部分にどう当てはまるかを考えてみるとよい。
そこから評価は自然に 4 番目の質問へと移っていく。
答えが「まったく問題ない」でなければ、おそらくそこは議論するとよい題材か、評価報告書に必ず入れるべき項目である。

% src: What Are We Working On?
\underline{何に取り組んでいるのか？：}
この質問は、いくつかの具体的な点で、レビューが脇道にそれないようにするのに役立つ。
第一に、設計の目的を知っておくことが重要である。
目的がわかっていれば、リスクを生むのに実際には必要ない部分を削るよう、自信を持って提案できる。
逆に、変更を提案するときに、本当に必要な機能を壊してしまうことも避けたい。
おそらく最も重要なのは、リスクのある機能について、新しい方向性の代替案を提案できる場合があることである。

\begin{exa}[給与システムにおける健康に関する質問]\label{exa:ch7-payroll-health}
例えばプライバシーの領域で、全従業員から個人情報を集める給与システムをレビューしているとしよう。健康に関する質問をとくに機微なものとして特定するかもしれない。
そのデータ項目が本当に不要なら、設計から削るのが正しい対応である。
しかし、そのデータが設計の役立つ業務機能にとって重要なら、代わりに、このデータの漏えいを厳重に防ぐ方法（早い段階での暗号化や、短期間での削除など）を提案できる。
\end{exa}

% src: What Can Go Wrong?
\underline{何がうまくいかない可能性があるか？：}
レビューでは、システムが直面する重要な脅威を設計者が想定していることを確認すべきである。
しかも、設計者がこうした脅威を認識しているだけでは不十分であり、脅威に耐えられるだけの設計を実際に作り上げていなければならない。

脅威の中には、許容できるものとして緩和しないままにしておくものもある。その場合、レビュー担当者の仕事はその判断を評価することである。
ただし、設計者がその脅威を認識したうえで、あえて緩和しないことを選んだのだと確かめることが重要である。
設計書にそのことがはっきり書かれていない場合は、意図的な判断であることを再確認するために、SDR でその点を記録しておく。
また、受け入れるリスクを記録し、なぜそのリスクが許容できるのかを説明しておく。

\begin{exa}[リスク受容の記録]\label{exa:ch7-accept-risk}
例えば次のように書ける。
「通信路上の暗号化されていないデータには盗聴の脅威がある。しかし、データセンターは物理的に保護されており、PII や業務上の機密データが露出するおそれはないため、このリスクは許容できると判断した。」

このリスク受容の判断を覆しうる将来の変更も見越しておくようにする。
上の例であれば、「システムを第三者のデータセンターに移す場合は、物理ネットワークへのアクセスのリスクについてのこの判断を見直すべきである」と書き加えられる。
\end{exa}

% src: What Are We Going to Do About It?
\underline{それに対して何をするのか？：}
レビュー担当者が設計を検討していくにつれて、セキュリティの保護メカニズムと緩和策が設計の中で見えてくるはずである。
レビュー担当者はふつう、最後の二つの質問、つまり設計を安全にしているものを特定し、それがどの程度安全かを評価することに、時間の大半を費やす。
この作業の進め方の一つは、脅威と緩和策を突き合わせて、すべての要所が押さえられているかを確かめることである。
この質問から浮かび上がる問題を指摘し、設計に問題がないことを確認するのは、SDR の最も重要な貢献の一つである。

設計がセキュリティリスクを十分に緩和できていない場合は、何が欠けているかを項目ごとに挙げるべきである。
このフィードバックを役立つものにするには、対処されていない具体的な脅威と、それがなぜ重要なのかを説明し、場合によってはそれぞれに対処するための大まかな選択肢も示す必要がある。
いくつかの理由から、原著者は SDR で具体的な解決策まで提案することは勧めない。
とはいえ、非公式に手助けを申し出たり、求められれば設計者と一緒に代替案を検討したり、設計変更を詳しく詰めたりするのは大いに結構である。

\begin{exa}[監視APIについてのフィードバック]\label{exa:ch7-monitoring-api}
例えば、フィードバックは次のようになるだろう。
「監視 API は公開すべきではない。Web サイトの利用状況が明らかになり、競合他社を利することになりかねないからである。RESTful API へのリクエストを認証するために、アクセスキーを要求することを推奨する。」
\end{exa}

ある脅威に対する緩和策が設計に備わっている場合は、その有効性を評価し、もっと良い代替案がないかを検討する。
設計者はときどき、セキュリティメカニズムを一から作って「車輪の再発明」をしてしまうことがある。
その場合は、代わりに標準ライブラリを使うよう提案するのが良いフィードバックになる。
設計は安全だが、その安全性が大きな性能上のコストと引き換えになっている場合は、できれば別の方法を提案する。

\begin{exa}[冗長なセキュリティメカニズム]\label{exa:ch7-redundant-encryption}
例えば、すでに暗号化されている HTTPS 接続で送るデータをさらに暗号化するといった、冗長なセキュリティメカニズムを指摘し、設計をどう簡素化できるかを説明するのがこれにあたる。
\end{exa}

% src: Did We Do a Good Job?
\underline{うまくやれたか？：}
この最後の質問こそが核心である。
設計は安全だと考えるか。
有能な設計者ならすでにセキュリティに対処しているはずなので、SDR の価値の多くは、設計者が全体像を見渡し、主要な脅威を想定できていたことを保証する点にある。
原著者の経験では、SDR では問題点と改善の機会がすぐに見つかるか、少なくとも、いま検討しておく価値のある興味深いトレードオフの判断が浮かび上がる（後になってからでは、これほど簡単には変更できないからである）。

報告書の冒頭には、設計全体についての総合的な評価を一文でまとめておくことを勧める。
例をいくつか挙げる。
\begin{itemize}
\item この設計は現状のままで安全であり、提案する変更はない。
\item この設計は安全であるが、さらに安全にするための変更をいくつか提案する。
\item 現在の設計には懸念があり、より安全にするための一連の推奨事項を示す。
\end{itemize}

要約の後には、水準に達していない領域が複数あれば、それらを分けて一つずつ説明する。
弱点を設計の特定の部分に結び付けて示せば、設計者は問題の所在を正確に突き止めてはっきり理解でき、必要な修正もしやすくなる。

もちろん完璧な設計などないので、設計が不十分だと判断する際には、何を基準にしているのかを明確にすることが重要である。
基準を抽象的に言い表すのは難しいので、具体的な脅威、脆弱性、結果を挙げて論拠にするのがよい。
同等の製品のセキュリティと比べる形で評価を述べるのが一番よい場合もある。

\begin{exa}[競合製品との比較による評価]\label{exa:ch7-ransomware-comparison}
例えば次のように述べる。
「主要な競合他社はランサムウェアへの耐性を大きなセールスポイントに掲げているが、この設計では、従業員が Web 閲覧にも使うコンピューター上に在庫データベースをローカルに置いているため、そうした攻撃にとくに弱い。」
\end{exa}

\subsubsection{どこを掘り下げるか}
% src: Where to Dig

大きな設計の隅々まで掘り下げるのは現実的でないので、レビュー担当者はできるだけ早く、セキュリティ上重要な主要領域に的を絞る必要がある。
設計のどこに労力を向けるかを決める際には、自分の直感に従うよう、原著者はセキュリティのレビュー担当者に勧めている。
まず設計を通読し、直感に従って気になる領域にメモを付ける。
次に、最も懸念の大きい領域に戻ってもっと注意深く調べ、潜在的な脅威と 4つの質問を手がかりに、尋ねるべき質問を集める。
こうした手がかりの中には、ほかより実りの多いものもある。
実りのない方向に進んでしまっても、たいていはほどなく気付くので、労力をほかに向け直せばよい。

セキュリティとプライバシーに関係のない設計の部分は、すべての構成要素について基本を理解できる程度に目を通しながら、ざっと読み流してかまわない。

\begin{exa}[家から締め出されたとき]\label{exa:ch7-locked-out}
家から締め出されてしまったら、開いている窓や鍵のかかっていない扉がないかを確かめればよいことは、誰でもわかるだろう。
外壁全体を 1 インチずつ調べて時間を費やす人はいない。
同じように、設計の中で弱点がありそうな箇所に狙いを定めるか、設計が最も価値の高い資産をどう守っているかに注目するのが、最も効果的である。
\end{exa}

攻撃対象領域に目を光らせ、それに見合う注意を払う。
攻撃対象領域が利用しやすいほど（匿名でインターネットに露出しているのが典型的な最悪の場合である）、攻撃の発生源になる可能性は高くなる。
価値ある資源を守る信頼境界、とくに攻撃対象領域から到達できるものは、設計の主要な一般的特徴であり、レビュー担当者は分析の中で必ず重視すべきである。
価値ある資産を外部に面したコンポーネントからもっとうまく隔離できることもあるが、多くの場合、露出は避けられない。
こうした要因こそ、レビュー担当者が過程全体を通じて探し出し、評価しなければならないものである。

\subsubsection{プライバシーのレビュー}
% src: Privacy Reviews

自分のスキルと組織上の責任に応じて、情報プライバシーは SDR の範囲内で扱ってもよいし、別に扱ってもよい。
SDR でプライバシーについてフィードバックするときは、適用されるプライバシーポリシーと、それが設計の範囲内でのデータの収集、利用、保存、共有にどう関わるかを中心にすべきである。

良いやり方は、プライバシーポリシーを一通り読んで設計に関係する箇所にメモを付け、それから違反を防ぐ方法を探すことである。
前章で述べたように、技術面で焦点となるのは、設計がポリシーに確実に準拠するようにすることである。
もっと専門的な知識を要する問題については、プライバシーの専門家と法務の承認を得る。

\subsubsection{更新のレビュー}
% src: Reviewing Updates

ソフトウェアはいったんリリースされると、まるで独自の生命を持つかのように、時間とともに変化していくのは避けられない。
アジャイルやその他の反復的な開発手法では設計の変更が絶え間なく続くので、とくにそうである。
その過程で設計書は放置されやすく、何年も経つと失われたり、意味をなさなくなったりする。
しかし、ソフトウェアの設計を変更すればそのセキュリティ特性にも影響しうるので、設計が安全であり続けることを確かめるために、SDR を差分で更新するのが賢明である。

設計書は、ソフトウェアのアーキテクチャの変遷を追う、生きた文書であるべきである。
版管理された文書は、設計がどう成熟してきたか、場合によってはどう入り組んでしまったかを示す重要な記録である。
この文書を手がかりにすれば、前回の SDR 以降の変更の正確な集合（\underline{設計の差分}、design delta）に絞って差分レビューを行い、SDR を更新できる。
設計のセキュリティ上重要な領域（またはその近く）に変更がある場合は、影響の大きい小さな、しかし重要な詳細が設計書から漏れていないかを確かめるために、レビュー担当者がフォローアップするのが賢明なことが多い。
差分レビューで何か重要なことが見つかったら、それを既存の評価報告書に追加し、報告書を読めば全体の経緯がわかるようにする。
何も見つからなければ、報告書を更新して、どの版の設計を対象にしたかを書き加えるだけでよい。

「単純な変更」の影響を甘く見ることは、セキュリティ上の惨事を招くよくある原因であり、設計を再レビューすれば、そうした影響を前もって効果的に評価できる。
レビューが要らないほど軽微な設計変更なら、セキュリティへの影響がないことをレビュー担当者がすぐに確認できるのもまた事実である。
ごく些細なもの以外の設計変更については、この重要な安全策を見逃すリスクを考えれば、SDR の更新を省こうと考えても得るものはほとんどない、と原著者は考えている。

%----------------------------------------------------------------------
\subsection{意見の不一致への対処}
% src: Managing Disagreement

\begin{quote}
「人生で何をするにせよ、自分に異を唱えてくれる賢い人たちをそばに置きなさい。」\\
\hfill --- John Wooden
\end{quote}

原著者が長年セキュリティの普及に努めてきた中で得た重要な教訓（苦い経験から学んだが、後から思えば当たり前のこと）は、SDR を成功させるには良好な対人コミュニケーションが欠かせない、ということである。
分析そのものはもちろん技術的なものだが、設計を批評するには良好なコミュニケーションと協働が必要であり、人的要因もまた鍵になる。
セキュリティの専門家は、社内の人であれ外部委託であれ、（それが当たっていようといまいと）決して満足しない口うるさい部外者という評判を立てられてしまうことがあまりにも多い。
そうした見られ方はやりとりをじわじわと損ない、仕事をやりにくくするだけでなく、全員の努力の成果にも悪影響を及ぼす。
もっとうまくやるには、この要因があることを認めなければならない。

\subsubsection{如才なくコミュニケーションする}
% src: Communicate Tactfully

SDR は本質的に対立をはらむものである。
SDR の大部分は、関係者がたいてい心血を注いできた設計について、リスクや潜在的な欠陥を指摘することだからである。
設計上の弱点は、いったん見つかってしまえば、後から見ると痛いほど明らかに思えることが多く、レビュー担当者はつい、それを不注意や、ひいては無能さのせいにしてしまいがちである。
しかし、そのような伝え方が生産的であることは決してない。
そうではなく、見つかった問題を教える機会として扱う。
設計者が問題を理解すれば、レビュー担当者が見落としていたかもしれない別の実りある領域へと、設計者自身が議論を導いてくれることも多い。
自分の設計の脆弱性を誰かに指摘してもらうことほど、セキュリティを学ぶのに良い方法はない。

弱い設計を容赦なく切り刻み、何よりもセキュリティを最大化することが大切だと一方的に講釈して SDR の時間を費やしても、まず生産的にはならない（自分が言われる側に立ってみれば、理由は明らかだろう）。
残念ながらこうしたことはときどき起こるが、原著者は、必ずしもレビュー担当者が意地悪だからではなく、必要な技術的変更に気をとられて、敬意ある口調を保つことを忘れやすいからだと考えている。
善意を保ち、全員が同じチームの一員として、多様な視点を持ち寄り、適切なバランスを見つけるという共通の目標に向かっているのだと確認し合うためには、最大限の努力を払う価値が十分にある。
スポーツのコーチもしばしば同じような微妙な綱渡りをしている。
コーチは、チームが最高の試合をするのに必要な取り組みができるよう、自分が見つけた弱点（相手がそこを突いてくるとわかっているもの）を、求めすぎることなく指摘する。
Mark Cuban が言うように、「親切は意地悪よりもずっと遠くまで届く」。

歓迎されないかもしれないメッセージを伝えながら人とうまくやっていけるに越したことはないが、言うは易く行うは難しである。
本書は技術的なソフトウェアの本なので、開発者に好かれて影響を与える方法といった自己啓発的な助言はしない。
しかし人的要因は十分に重要であり、より正確に言えば、それを無視すると仕事が損なわれるおそれが十分に大きいので、ここで目立つように触れておく価値がある。
原著者の基本的な指針は単純である。
自分がメッセージをどう伝えているかを意識し、相手がそれをどう受け取り、どう反応しそうかを考えることである。
これが SDR でどう生きるかを示すために、実際にあった話と、原著者が頼りにするようになった一連のヒントを紹介する。

\subsubsection{ケーススタディ：難しいレビュー}
% src: Case Study: A Difficult Review

原著者にとって最も記憶に残る SDR の一つは、ソフトスキルの重要性を示す格好の教訓になっている。
発端は、文書をもらっていくつか基本的な質問をするためだけに原著者が始めた、つらいメールのやりとりであった。
そのやりとりから、チームリーダーが SDR をまったくの時間の無駄だとみなしていることがすぐにわかった。
そのうえチームは、SDR がこの製品のリリース要件であることを知らなかったので、懸命に取り組んできたリリースの前に、歓迎されない新たな障害が突然立ちはだかった格好になっていた。
この話から得られる第一の重要な教訓は、それが正しいか間違っているかにかかわらず、他の参加者がこの過程をどう見ているかを認識し、それに合わせて対応を変えることの大切さである。

最終的に原著者が手に入れた文書は、雑で不完全で、かなり古くなっていた。
これをそのまま率直に指摘していたら、生産的でないうえに、関係をさらにこじらせていただろう。
第二の重要な点は、改善を促し、問題をうまく回避し、SDR を効果的に進めるには、次のような戦略を使うほうが生産的だということである。
\begin{itemize}
\item 修正や追加を提案し、それぞれの提案の背後にあるセキュリティ上の理由も添える。
\item 可能なら、文書のレビューや編集の提案など、過程を円滑に進めるためにできることは何でも手伝うと申し出る（ただし、相手の仕事を肩代わりするところまではしない）。
\item 予備的な SDR のフィードバックは、要求としてではなく「私の見方」として伝える。
\item 「サンドイッチ」方式を使う。まず肯定的な所見から入り、必要な改善点を指摘し、最後に肯定的な点（変更がどう役立つかなど）で締めくくる。
\item フィードバックが大量になる場合は、まずどう伝えるのが一番よいかを尋ねる（97 項目の箇条書きのメールや、いきなり大量のバグを登録して相手を驚かせない）。
\item 気付いた手がかりはすべて調べるが、フィードバックは最も重要な点に絞る（完璧主義にならない）。
\end{itemize}

良い経験則として、欠けている情報が多くの読者にとって広く役立つものなら文書化する価値があるが、自分の必要に限ったものなら、もっと気軽に質問すればよい（必要なら、その問題の詳細を評価報告書に含めればよい）。
文書の質について不満を言ったり良し悪しを論じたりする代わりに、社内のプロトタイプがあれば使ってみる、コードやコードレビューに目を通すなど、ソフトウェアについて学ぶ別の方法を工夫して見つける。
チームの定例会議を見学させてもらうよう頼むのは、誰の時間も奪わずに設計について学べる優れた方法である。

メールでは相手が失礼にふるまっているように感じたが、ようやく会ってみると、単にストレスを抱えた開発リーダーにすぎないことがわかった。
原著者はリーダーだけを頼りにするのをやめ、それほど手一杯でなく、喜んで質問に答えてくれる別のチームメンバーを見つけた。
SDR 会議の準備の時間を節約するため、前もって解決しておくことが重要な質問だけを追いかけ、それ以外の質問は、相手が確実に聞いてくれる会議の場のためにとっておいた。

SDR 会議の準備は、バランスをとるのが難しい。
まったく準備をせずにいきなり臨むべきではない。
とくに文書を渡した後で何もかも説明させられるのは、チームにとって気持ちのよいものではないからである。
自分がよく知らない主要なコンポーネントや依存関係を事前に洗い出し、少なくとも会議で質問できる程度には理解しておくようにする。
準備の間は、問題点と質問を書き留め、それを次の分類に仕分けるのが良い習慣である。
\begin{itemize}
\item 会ったときにすぐセキュリティの掘り下げに入れるよう、事前に尋ねておく質問
\item 自分で答えを見つけられる質問
\item 会議の場で検討するのが一番よい話題
\item 議論の必要はなく、評価報告書に含める所見
\end{itemize}

ようやく会議を開いたときには、主任エンジニアは、SDR が今や製品発売の最大の障害になっていることに露骨に不満を示していた。
最初の会議は少し波乱含みだったが、全員が集中を保ち、話は大きく進んだ。
さらに何回か会議を重ねた後（回を追うごとに楽になり、短くなっていった）、原著者は設計を承認した。
いくつかの変更には最初の会議で合意していたが、詳細を確認し、変更を確定させるために会ったことは、全員にとって重要な安心材料になった。
設計に必要な変更が行われたことを確認する時間をとらないと、行き違いが見過ごされやすい。

相手の時間を使わせることがその人のためになるのだと、忙しい人たちに納得してもらうのは決して簡単ではなく、そう言い聞かせてもまずうまくいかない。
しかし、セキュリティを改善する小さな機会でも指摘し、それが最終製品にどう役立つかを示せば、互いに満足のいく結果にたどり着きやすい。

SDR が終わるころには、製品チームはセキュリティについて、そしてその延長として自分たちの製品について、はるかに深く理解するようになっていた。
最終的にチームはレビューの価値を認め、その結果として製品が良くなったことも認めた。
さらに良いことに、バージョン 2 のときにはチームのほうから進んで原著者に連絡してきて、更新の SDR は実に順調に進んだ。

\subsubsection{意見の不一致のエスカレーション}
% src: Escalating Disagreements

設計者とレビュー担当者が合意に至らない場合は、意見が違うことを互いに認め合うべきである。
問題が軽微なら、レビュー担当者は評価報告書に不一致の点を記録するだけにして、設計者の判断に従えばよい。
その場合は、例えば「採用されなかった推奨事項」（Recommendations Declined）という節を設けて不一致をはっきり示し、提案した設計変更とそれを推奨した理由、そして変更しなかった場合に起こりうる結果を説明する。
しかし、重要な決定をめぐって深刻な対立がある場合は、レビュー担当者は問題をエスカレーションすべきである。

この場合は、設計者とレビュー担当者の双方が自分の立場を書面にまとめる。
まず両者が合意できている共通の出発点を見つけることから始め、草稿を交換して、全員が双方の視点を知るようにする。
双方の立場を合わせて、リスクを説明し、提案される結論とそのコストを示すメモを作る。
このメモは評価報告書を補うもので、会議のたたき台として、あるいは経営陣がどう進めるかを決めるための指針として使える。
最終決定の結果は、エスカレーションのメモと一緒に評価報告書に収めるべきである。

原著者は長年セキュリティレビューを行ってきたが、問題をエスカレーションしたことは一度もない。ただ、それに近いことは何度かあった。
強い意見の対立は、ほぼ例外なく基本的な前提の大きな食い違いから生じており、その食い違いがいったん特定されれば、たいていは解決に向かう。
そうした食い違いは、ソフトウェアの使われ方や、ソフトウェアが処理するデータについての暗黙の前提から生じることが多い。
実際には、ソフトウェアの使われ方を制御するのはきわめて難しく、時間とともにユースケースもたいてい変わっていくので、ふつうは安全側に倒しておくのが最善である。

食い違いのもう一つの大きな原因は、データの機密性や完全性が重要であることを設計者が理解していない場合である。
たいていは、必要なエンドユーザーの視点が欠けているか、考えられるユースケースを全体まで見渡していないことが原因である。
考慮すべき重要な要因がもう一つある。
仮にリリース後に考えを変えたとしたら、その段階で変更するのはどれほど難しくなるだろうか。
後になって「だから言ったのに」と言いたい人はいないが、対立する条件を書面にしておくことが、たいていは正しい選択をするための最善の方法である。

%----------------------------------------------------------------------
\subsection{練習あるのみ}
% src: Practice, Practice, Practice

この章で学んだことを確かなものにし、本当に自分のものにするために、思い切ってソフトウェア設計を一つ見つけ、SDR を行ってみるよう、原著者は読者に強く勧めている。
いま自分の関心の範囲に進行中のソフトウェア設計がなければ、手に入る既存の設計を何でもよいので選び、練習としてレビューすればよい。
選んだソフトウェアに正式な設計書がない場合は、まず設計の大まかな表現を自分で作り（完全で洗練された文書である必要はなく、ブロック図でも十分である）、それをレビューする。
一般に、手に負えなくならないよう、手頃な大きさの設計から始めるか、大きなシステムから一つのコンポーネントを切り出してその部分だけをレビューするのがよい。
ここまで読んできたなら、始める準備はできているはずである。
評価報告書を人に見せる自信がまだなければ、自分用に手早くレビューすることから始めてもよい。

SDR に欠かせないスキルが身につくにつれて、出会うあらゆるソフトウェアにそれを応用できるようになる。
たくさんの設計を研究することは、名人のやり方を見ることと、他人の誤りを見つけることの両面から、ソフトウェア設計の技を学ぶ優れた方法である。そしてこのようにスキルを応用する練習は、スキルを伸ばすための優れた訓練になる。

手始めにとくに取り組みやすいのは、付録Aのサンプル設計書をレビューすることである。
そこではセキュリティに関する規定が強調表示されており、設計の中で何を探せばよいかの現実的な例になっている。
設計を読み、強調された部分に注意しながら、もしそうしたセキュリティ関連の詳細が欠けていたら自分ならどう見つけて補うかを想像してみるとよい。
もっと難しい課題に挑むなら、設計をさらに安全にする方法をほかにも探してみるとよい（原著者は、サンプル設計書が欠点のない理想だなどとは決して主張しないし、そう期待してもいない）。

SDR を行うたびに、腕は上がっていく。
重大な脆弱性が見つからなくても、設計についての知識とセキュリティのスキルは高まる。
セキュリティ上の手当てを必要とするソフトウェアには事欠かないので、ぜひ始めてほしい、と原著者は呼びかけている。
この価値あるスキルがどれほど速く身につくかに、きっと驚くだろう、と原著者は考えている。

\newpage

%% =======================================================
%%  part: ch8
%% =======================================================
%% =======================================================
%%  part: ch8
%% =======================================================
\section{セキュアプログラミング}
% src: 8: Secure Programming

この章では、\cite{Kohnfelder2021} の第8章に沿って、セキュアプログラミングを扱う。

\begin{quote}
「第一の原則は、自分をだまさないことである。そして、いちばんだましやすいのは自分自身なのだ。」 ---Richard P. Feynman
\end{quote}

セキュリティを念頭に置いて作成・レビューしたソフトウェア設計が完成しても、それは製品の旅路の始まりにすぎない。
その後には、実装、テスト、デプロイ、運用、監視、保守、そして最後に製品寿命を迎えたときの廃止という作業が続く。
具体的な中身は OS や言語によって大きく異なるが、セキュリティに関する大まかなテーマはほぼどこでも共通している。

開発者は、良い設計に明記されたセキュリティ上の規定を忠実に実装するだけでなく、その過程で欠陥のあるコードを書いて新たな脆弱性をうっかり持ち込まないよう注意しなければならない。
建築家の図面をもとに家を建てる大工にたとえるとわかりやすい。
質の悪い材料でずさんに施工すれば、完成した家にはあらゆる問題が出てくる。
大工が釘を打ち損じて曲げてしまえば、その失敗は目に見えるので簡単に直せる。
これに対して、欠陥のあるコードは見過ごされやすいうえに、悪用されれば悲惨な結果を招く脆弱性を生むことがある。
この章の目的はコーディングの仕方を教えることではなく（読者はすでに心得ているものとする）、コードがどのようにして脆弱になるのか、どうすればより安全にできるのかを伝えることである。
続く各章では、今もソフトウェアプロジェクトを悩ませ続けている、ありふれた実装上の脆弱性の多くを取り上げる。

設計と実装の境界は必ずしも明確ではないし、明確であるべきでもない。
思慮深い設計者は、プログラミング上の問題を予見したり、セキュリティが決定的に重要になる箇所について助言したりと、多くのことができる。
実装を担当するプログラマーは、設計に肉付けをして曖昧な点をすべて解消し、正確に定義されたインターフェースを持つ、実際に動くコードに仕上げなければならない。
プログラマーは設計を安全にコードへ落とし込む（それだけでも手ごわい作業である）だけでなく、必要なコードを細部まで書き上げる過程で新たな脆弱性を持ち込まないようにしなければならない。

理想的には、設計の段階で能動的なセキュリティ対策、すなわちシステムとその資産、ユーザーを守るために組み込むソフトウェアの機能を定めておくべきである。
一方、開発におけるセキュリティとは、ソフトウェアが陥りがちな落とし穴、いわばコンポーネントやツールの「とがった角」を避けることである。
実装の途中で新たなリスクが生じたら、そのリスクに合わせた緩和策が必要になる。
設計者がそうしたリスクまで予見できたはずだと考える理由はないからである。

この章では、バグの一部がどのようにして脆弱性になるのか、つまり脆弱性がどう生まれ、さまざまな落とし穴をどう避ければよいのかに焦点を当てる。
これらの問題を一般的な観点から扱い、歴史的にセキュリティ上の問題を多く抱えてきた主要な領域を掘り下げる以降の章への導入とする。
まず、セキュアコーディングという課題の本質を探る。
攻撃者がどのように隙を突き、その影響をコードの奥へと広げていくのかもここに含まれる。
続いてバグを取り上げる。
脆弱性がバグからどのように生まれるのか、軽微なバグがどのように脆弱性チェーンをなしてより大きな問題を引き起こしうるのか、そしてエントロピーという観点から見るとコードがどう見えるのかを扱う。

コードの脆弱性を避けるには警戒が欠かせないが、そのためにはコードがどのようにセキュリティを損なうのかを知っておく必要がある。
コーディング上の脆弱性という概念を具体的に理解するために、たった 1 行の編集ミスがインターネット全体のセキュリティを破った、壊滅的な実在の脆弱性のコードを簡略化して順に見ていく。
その後、悪用されると深刻な結果を招くバグの例として、よくある脆弱性のクラスをいくつか紹介する。

第3部を通して、コード例は Python と C で示す。
どちらも広く使われている言語であり、高水準から低水準まで幅広い抽象度をカバーしている。
コードは各言語の細部まで踏まえた本物であるが、本書の概念は言語を問わず当てはまる。
% 注: 原文はここに "Even if you are unfamiliar with Python or C, the code snippets." という文があるが、文の途中で切れているため本文には入れなかった。

%----------------------------------------------------------------------
\subsection{課題}
% src: The Challenge

この章の題名を「セキュアプログラミング」（secure programming）とするのは当然の選択であったが、誤解を招くおそれもある。
目標をより正確に言い表すなら（章題には向かないが）「安全でないコーディングを避けること」になる。
つまり、セキュアコーディングという課題の大部分は、悪用可能な脆弱性となる欠陥を持ち込まないことに尽きる。
もちろんプログラマーは、セキュリティを積極的に高める保護機構も作るが、そうした機構はたいてい設計や API の機能として明示されている。
ここで主に取り上げたいのは、意図せずに陥ってしまう落とし穴である。
こうした落とし穴は目に見えにくく、セキュリティ上の失敗の大半の根本原因となっているからである。
セキュアコーディングは、道路のどこに穴があるかを覚えておき、運転中は注意を怠らず、穴を一つ残らず避けて走ることに似ていると考えるとよい。

多くのプログラマーは、おそらくもっともな理由から、ソフトウェアセキュリティに良い印象を持っていない（もっと感情的になって、自分たちを煩わせる「セキュリティ警察」、あるいはもっとひどい呼び名で呼ばれる人々に反感を抱いている場合もある）。
筆者の考えでは、その理由は、実装について「しくじるな」というメッセージを繰り返し聞かされることにある。
希少なダイヤモンドをカットしようとしている宝石職人に「しくじるな！」と言っても何の役にも立たないのと同じ理由で、この言葉はプログラマーにも役に立たない。
職人は最善を尽くすつもりでいるのであり、余計なプレッシャーをかけても、集中して正しく仕事をするのがかえって難しくなるだけである。
善意の「警察」は必要な助言をしているのだが、親切で建設的な言い方で伝えていないことが多い。
筆者自身もこの過ちを何度も犯してきたので、ここではその微妙な一線を越えないよう努めている。読者にもその点を汲んでいただきたい。

とはいえ、注意が必要なのは確かである。
（この章の後半で GotoFail 脆弱性を見ればわかるように）プログラマーのたった 1 回のミスが、簡単に悲惨な結果を招きうるからである。
問題の根源は、現代の大規模ソフトウェアシステムがきわめて脆く複雑であることにあり、その傾向は今後さらに強まると予想されている。
プロのソフトウェア開発者はコードのテストやデバッグのやり方を知っているが、セキュリティは別問題である。
脆弱なコードは、しつこく攻撃されない限り、何の問題もなく動いてしまうからである。

ソフトウェア設計者は理想化された構想を描く。
設計はまだ実現されていないので、理論上は完全に安全でさえありうる。
しかし、実際に動くソフトウェアを作れば新たな水準の複雑さが加わり、設計に書かれていない細部まで肉付けしなければならない。
そうした作業はどれも、セキュリティ問題のリスクを避けがたく伴う。
良い知らせもある。完璧を目指す必要はないし、よくある脆弱性の大半を占めるコーディング上の失敗パターンはよく知られていて、正しく書くのもそれほど難しくない。
コツは、常に警戒を怠らず、コードに潜む危険な欠陥を見抜く「目」を養うことである。
この章では、安全なコードと脆弱なコードがそれぞれどう見えるのかをしっかりつかめるよう、いくつかの概念を例とともに紹介する。

\subsubsection{悪意ある影響}
% src: Malicious Influence

セキュアコーディングを考えるうえで重要なのは、攻撃者が実行中のコードにどのように影響を及ぼしうるかを理解することである。
大きく複雑な機械が滑らかに動いているところに、いたずら者が棒を持ってきて、機構の中に突っ込み始める場面を想像してほしい。
ガソリンエンジンのシリンダーのようにエンジンブロックの中で完全に守られている部品もあれば、ファンベルトのようにむき出しで、何かを挟み込ませて故障させやすい部品もある。
攻撃者がシステムに侵入しようとしてシステムをつつき回すやり方は、これによく似ている。
攻撃者は攻撃対象領域を起点に、巧妙に作った予期しない入力で機構を狂わせようとし、さらにシステム内部のコードをだまして思いどおりに動かそうとする。

信頼できない入力がコードに影響を及ぼす経路には、直接的なものと間接的なものの 2 つがある。
攻撃者は、信頼できない入力（たとえば文字列「BOO!」）を注入できる場所ならどこからでも始め、そのデータが拒否されずにシステムの奥まで伝わっていくことを期待して試行錯誤する。
I/O やさまざまなインターフェースの層を下っていくうちに、文字列「BOO!」はたいていいくつものコードパスに入り込み、その影響はシステムの奥深くまで浸透していく。
ときには、信頼できないデータとコードのやり取りがバグを引き起こしたり、設計どおりの機能ではあっても困った副作用のある動作を引き起こしたりする。
「BOO!」を Web 検索すると、データセンター内の何百台ものコンピューターが関わり、それぞれが検索結果の一部を担うかもしれない。
そのため、この文字列はメモリ上の何千もの場所に書き込まれることになる。
これは影響がそれだけ広く広がるということであり、害を及ぼす可能性がほんのわずかでもあれば危険になりうる。

\begin{dfn}[テインティング]\label{dfn:ch8-tainting}
データがコードに及ぼすこうした影響を指す専門用語を \underline{テインティング}（tainting）という。
言語によっては、テインティングを追跡する機能を備えているものもある。
\end{dfn}

\begin{exa}[テインティングを追跡する言語]\label{exa:ch8-taint-languages}
Perl のインタープリターは、インジェクション攻撃（第10章で扱う）を緩和するためにテインティングを追跡できる。
初期の JavaScript にも同じ理由でテイントチェックがあったが、使われなかったため、ずいぶん前に削除された。
\end{exa}

それでも、信頼できない入力元から来たデータがコードに及ぼす影響という概念は、脆弱性を防ぐうえで必ず理解しておくべきものである。

入力データがコードに影響を及ぼす経路はほかにもあり、データが保存されないまま間接的に影響することもある。
文字列「BOO!」が入力として与えられたとき、コードがそのコピーをそれ以上どこにも保存しないようにしたとしよう。
これでシステムはその影響から切り離されるだろうか。
とんでもない。

\begin{exa}[保存されない入力による間接的な影響]\label{exa:ch8-indirect-influence}
\texttt{input = "BOO!"} が与えられたときの、次のコードを考える。
\begin{verbatim}
if "!" in input:
    PlanB()
else:
    PlanA()
\end{verbatim}
入力に感嘆符が含まれているので、コードは PlanA ではなく PlanB を実行する。
入力文字列そのものは保存されず、後続の処理に渡されることもないにもかかわらず、である。
\end{exa}

この単純な例からわかるように、データ（ここでは「BOO!」）そのものは遠くまで伝わらなくても、信頼できない入力の影響はコードの奥深くまで伝わりうる。
大規模なシステムで、攻撃対象領域を起点とする推移的閉包（すべてのパスを合わせた範囲）を考えれば、影響が多くのコードに浸透しうることがわかる。
多くの層を貫いて影響が広がりうるという点が重要なのは、攻撃者が予想以上に多くのコードに手を伸ばせ、コードの動作を操る機会を得られることを意味するからである。
信頼できない入力の扱いについては、第10章でさらに述べる。

\subsubsection{脆弱性はバグである}
% src: Vulnerabilities Are Bugs

\begin{quote}
「デバッグがバグを取り除く作業だとすれば、プログラミングとはバグを埋め込む作業に違いない。」 ---Edsger Dijkstra
\end{quote}

すべてのソフトウェアにバグがあることは広く認められているので、今さらこの主張を裏付ける必要はほとんどない。
もちろん例外はある。
自明なコード、正しさを証明できるコード、そして航空や医療などの重要な機器を動かす、念入りに設計されたソフトウェアである。
しかしそれ以外のすべてについては、バグはどこにでもあると自覚することが、セキュアコーディングに取り組む良い出発点になる。
そうしたバグの一部は、攻撃者にとって使えるものになるからである。
そこで、ここではバグに焦点を当てる。

\begin{dfn}[脆弱性]\label{dfn:ch8-vulnerability}
\underline{脆弱性}（vulnerability）とは、ソフトウェアのバグのうち、攻撃者が害を及ぼすのに役立つものの部分集合である。
\end{dfn}

脆弱性をほかのバグと正確に区別することはほぼ不可能なので、まずは明らかに脆弱性ではないバグ、つまりまったく無害なバグを見分けることから始めるのがいちばん簡単である。
オンラインショッピングの Web サイトのバグを例にいくつか考えてみる。

\begin{exa}[無害なバグと有害になりうる類似のバグ]\label{exa:ch8-harmless-bug}
無害なバグの良い例は、Web ページのレイアウトが設計どおりにならないという問題である。
少し崩れてはいるが、重要なコンテンツはすべて完全に表示され、正常に機能している。
ブランドイメージや使い勝手の面からは修正が重要かもしれないが、このバグにセキュリティ上のリスクがないことは明らかである。
しかし、脆弱性を見つけるのがいかに難しいかを強調しておくと、同じくレイアウトを崩すバグでも有害なものはありうる。
たとえば、ユーザーがセキュリティ上の判断を正しく下すために見る必要のある重要な情報が、そのバグによって隠れてしまう場合である。
\end{exa}

\begin{exa}[管理インターフェースの露出]\label{exa:ch8-exposed-admin}
有害な側の極端な例として、悪夢のような脆弱性を考えてみる。
管理用インターフェースが、誤って何の保護もないままインターネットに公開されてしまったとする。
すると、Web サイトを訪れた誰もが、ボタンをクリックするだけで、管理者が価格を変更したり、機密のビジネスデータや財務データを閲覧したりするのに使うコンソールに入れてしまう。
これは認可の完全な失敗であり、誰が見てもわかる明白なセキュリティ上の脅威である。
\end{exa}

もちろん、この両極端の間は連続的につながっており、その中間には、バグが害を及ぼす可能性を主観的に判断せざるをえない、広い灰色の領域がある。
しかも次の項で見るように、複数のバグが重なったときの影響はしばしば予見できず、害の可能性を判断するのはとりわけ難しくなる。
当然ながら、セキュリティのためには安全側に倒し、脆弱性である可能性が少しでもあるバグは、できるだけ修正する方向に傾くことを勧める。

筆者がこれまで関わったプロジェクトには、どれも大量のバグで埋め尽くされた追跡データベースがあったが、既知のバグの数（実際のバグの数とはまったく別物である）をゼロにしようという組織的な取り組みすらなかった。
したがって一般に、私たちは皆、未知のバグはもちろん、既知のバグの山も抱えたままプログラミングをしていると言ってよい。
まだ取り組んでいないのであれば、既知のバグを順に確認し、脆弱性になりうるものに修正の印を付けることを検討してほしい。
また、バグが無害であることを調べて証明するよりも、単にバグを直してしまうほうがほとんど常に簡単である、ということも言い添えておきたい。
第13章では、脆弱性に優先順位を付けるのに役立つよう、セキュリティバグを評価して順位付けするための指針を示す。

\subsubsection{脆弱性チェーン}
% src: Vulnerability Chains

\begin{dfn}[脆弱性チェーン]\label{dfn:ch8-vulnerability-chain}
一見無害なバグが組み合わさって深刻なセキュリティバグを生みうる、という考え方を \underline{脆弱性チェーン}（vulnerability chain）という。
いわば、悪者にとってのバグの相乗効果である。
\end{dfn}

散歩の途中で、渡りたい小川に行き当たったところを想像してほしい。
跳び越えるには幅が広すぎるが、水面からいくつか石が突き出しているのに気付く。
石から石へと跳び移れば、靴を濡らさずに簡単に渡れる。
この石が軽微なバグにあたる。一つひとつは脆弱性ではないが、組み合わさると小川を渡る新たな道ができ、攻撃者はシステムの奥深くまでたどり着けるようになる。
こうした飛び石となるバグが組み合わさって、悪用可能な脆弱性ができあがるのである。

\begin{exa}[オンラインショッピングにおける脆弱性チェーン]\label{exa:ch8-warehouse-chain}
オンラインショッピングの Web アプリで、このような脆弱性チェーンがどのように生じうるかを、簡単な例で示す。

最近のコード変更で、注文フォームに、発送を担当する倉庫を示すコードがあらかじめ入力された新しいフィールドが追加された。
以前は、顧客が注文した後にバックエンドのビジネスロジックが倉庫を割り当てていた。
今は、顧客が編集できるフィールドによって、注文を処理する倉庫が決まる。
これを Bug \#1 と呼ぶ。
この変更を担当した開発者は、こう主張する。このフィールドが増えたことには誰も気付かないだろうし、仮に誰かがシステムが既定で入れる倉庫の指定を書き換えたとしても、別の倉庫には注文された商品の在庫がないので、そこで検知されて修正される、と。
「実害がなければ問題なし」というわけである。
この分析をもとに、ただしテストは一切せずに、チームは Bug \#1 の修正を次のリリースサイクルに回す。
緊急対応やスケジュールの遅れを避けられたことを喜びつつ、チームはバグのあるコード変更を本番環境に投入する。

一方、バグデータベースの中では、Bug \#2 が Priority-3（「いつか直す」、つまりおそらく永遠に直さないという意味）の順位のまま放置され、長らく忘れられていた。
Bug \#2 を登録したのは、何年も前のあるテスターである。
間違った倉庫を指定して注文すると、その倉庫では注文に応じられないためシステムが直ちに返金するが、その後の別の処理段階で注文が正しい倉庫に割り当て直され、その倉庫が注文を処理して発送してしまう、ということをテスターは発見したのである。
テスターはこれを深刻な問題（会社が商品をただで配ってしまう）とみなし、Priority-1 として登録した。
ところがトリアージ会議で、プログラマーたちは、（Bug \#1 が持ち込まれる前は）バックエンドが在庫を確認してから倉庫を割り当てていたのだから、テスターのやり方は「ずる」だと主張した。
つまり、見つかった時点の Bug \#2 は純粋に仮説上のもので、本番環境で起こることはありえなかった。
ビジネスロジックの各段階がどう絡み合っているかを解きほぐすのは難しいため、チームはこのバグを手つかずのまま Priority-3 とすることに決め、バグはすぐに忘れられた。

この「寝た子を起こすな」式の話をここまで追ってきたなら、結末が不幸なものになることはもう察しがつくだろう。
Bug \#1 が持ち込まれたことで、Bug \#2 と組み合わさった完全な脆弱性チェーンが、ほぼ確実に誰にも知られないまま成立してしまった。
倉庫指定フィールドを顧客が書き換えられるようになった今、Bug \#2 を引き起こす「間違った倉庫」の状況を作り出すのはたやすい。
必要なのは、ずる賢い、あるいは単に好奇心の強い顧客が 1 人、倉庫フィールドを書き換えてみることだけである。
全額返金されたうえに商品をただで受け取れて嬉しい驚きを味わった顧客は、次はもっと大量に注文しに戻ってくるかもしれないし、その秘密を人に教えるかもしれない。
\end{exa}

バグのトリアージのどこが誤っていたのかを見てみよう。
（先に見つかった）Bug \#2 は深刻な脆さであり、そもそも最初に直しておくべきであった。
放置してよいとする理屈は、倉庫側が「ほかのバックエンドロジックが自分に正しく指示してくれる」と信頼していることを前提にしており、注文の倉庫割り当てフィールドはどの攻撃対象領域からも完全に切り離されているという仮定（当時は正しかった）に基づいていた。
それでも、この脆さは明らかに悪い結果を招く憂慮すべきものであり、ビジネスロジックの修正が難しいという事実は、むしろ書き直したほうがよいことを示唆している。

後から持ち込まれた Bug \#1 は、新たな攻撃対象領域を開き、倉庫指定フィールドを改ざんの危険にさらした。
このバグを直さないという不運な判断は、システムが改ざんに対して堅牢だという誤った仮定に頼っていた。
振り返ってみれば、誰かが少しでもテスト（もちろんテスト環境で行うのであって、本番環境では決して行わない）をしていれば、自分たちの推論の欠陥に簡単に気付き、Bug \#1 をリリースする前に正しく対処できたはずである。
さらに理想を言えば、Bug \#2 を見つけたテスター、あるいはそのバグを知っている誰かがその場にいれば、二つの点を結び付けて、両方のバグを Priority-1 として修正予定に入れられたかもしれない。

実際には、バグが脆弱性チェーンをなしていることに気付くのは、この作り話の例よりずっと難しいのが普通である。
しかし、いったんこの概念を理解すれば、できる限りバグを先回りして直しておくことがいかに賢明かはすぐにわかる。
ただし注意しておくと、脆弱性チェーンがあるかもしれないと疑ったときでさえ、漠然とした仮説にしか見えないものの修正に時間を割くよう周りを説得するのは、実際には難しいことが多い。
問題のバグを直すのに大きな作業が必要な場合はなおさらである。
大規模なシステムの大半には、まだ見つかっていない脆弱性チェーンがあふれており、そのせいで私たちのシステムは弱くなっている可能性が高い。

この例は、2 つのバグが因果の連鎖として並びうることを示している。
ちょうど、手球が別の球に当たり、その球が的球をポケットに落とすという、ビリヤードの難しいショットのようなものである。

\begin{exa}[Pwn2Own における 6 つのバグの連鎖]\label{exa:ch8-pwn2own}
信じがたいかもしれないが、脆弱性チェーンはもっとずっと複雑になりうる。
ハッキング競技大会 Pwn2Own では、あるチームが 6 つのバグを連鎖させて、難度の高いエクスプロイトを成功させた。
\end{exa}

脆弱性チェーンを理解すると、コードの品質とセキュリティの関係がよりよくわかる。
特に重要な資産の周辺で脆さを生むバグは、積極的に直すべきである。
「そんなことは決して起こらない」という理由で（先の例の Bug \#2 のように）バグの修正を先送りするのは危険である。
そして、誰か 1 人が大丈夫だと言っても、それはあくまで意見であって証明ではないことを心に留めておくべきである。
こうした考え方は、隠蔽によるセキュリティ（Security by Obscurity）というアンチパターンに似ていて、せいぜい一時しのぎにすぎず、トリアージの最終判断としては良いものではない。

\subsubsection{バグとエントロピー}
% src: Bugs and Entropy

脆弱性と脆弱性チェーンをひととおり見たところで、次に、ソフトウェアはもっと輪郭のはっきりしない、損害をもたらしうる一連の事象にもさらされていることを考えてみる。
バグの中には予測できない形で物事を壊しがちなものがあり、そのため（脆弱性チェーンの場合と同じく）悪用できるかどうかを分析するのが難しい。
この現象の表れとして、私たちは、無数のバグのせいで時間とともにたまっていくエントロピーを一掃するために、スマートフォンやコンピューターをよく再起動する。
（ここでは \underline{エントロピー}（entropy）という言葉を、無秩序や比喩的な腐食を思い起こさせるために、ゆるい意味で使っている。）
攻撃者はこうしたバグやその後遺症を利用できることがあるので、対策を講じればセキュリティを高められる。

この種の問題を起こしやすいバグのクラスの一つが、実行スレッドどうしの予期しない相互作用から生じるバグである。
こうしたバグはたいてい、一見ランダムに、さまざまな形で現れるからである。
メモリ破壊のバグも、スタックやヒープの内容が絶えず変化しているため、同じようなクラスに属する。
システムを予測できない形でかき乱すこの種のバグは、可能性が事実上無限にあるため、攻撃の標的としてはかえって魅力的でさえありうる。
攻撃者はこうした厄介なバグを悪用するのにかなり長けていることがあり、しかも自動化すれば、成功率の低い試行を運よく当たるまで繰り返すのもたやすい。
その裏返しとして、たいていのプログラマーは、突き止めにくく、しばしば再現性が低すぎて気にするまでもないとみなされるこうした捉えどころのないバグに取り組むのを嫌うので、こうしたバグは手つかずのまま残りがちである。

はっきりした因果の連鎖を突き止められなくても、エントロピーを生むバグは危険になりうるので、直す価値は十分にある。
どのバグも、正しい動作からわずかに外れているという意味で、エントロピーのようなものをいくらかシステムに持ち込む。
そして、その小さな乱れは、特にずる賢い攻撃者に後押しされると、あっという間に積み重なる。
熱力学第二法則になぞらえるなら、閉じた系ではエントロピーは避けがたく増大し、この種のバグがいずれ悪用可能になって害をもたらすリスクが高まっていく。

\subsubsection{警戒}
% src: Vigilance

筆者はハイキングが大好きであるが、住んでいる地域のトレイルはぬかるんで滑りやすいことが多く、木の根や岩もむき出しなので、滑って転ぶ危険が常にある。
練習と経験を積んだおかげで滑ることはめったになくなったが、不思議なことに、特に危ない場所では、集中しているので決して滑らない。
それでもときどき転ぶことはあるが、障害物のせいというより、たいていはトレイルの楽な区間で、単に注意を怠っていたからである。
要するに、意識していれば困難な課題も乗り越えられるが、逆に不注意だと、たとえ道が楽でも簡単に足をすくわれる、ということである。

ソフトウェア開発者が直面しているのは、まさにこの種の課題である。
潜在的なセキュリティ上の落とし穴を意識し、集中を保ち続けなければ、知らず知らずのうちに落とし穴にはまりやすい。
開発者は本能的に、通常のユースケースで動くようにコードを書くが、攻撃者はエクスプロイトにつながる欠陥が見つかることを期待して、しばしば予想外のことを試す。
脆弱性チェーンとエントロピーの話で先に述べたとおり、ありうる入力や事象の組み合わせをすべて想定するよう警戒し続けることが、安全なコードを届けるうえで決定的に重要である。

次の節と以降の章では、現代のソフトウェアを悩ませている脆弱性を幅広く代表的に取り上げ、実装上の脆弱性がどのようなものかを示すために「おもちゃ」のコード例を使う。
MIT の人工知能研究の伝説的人物の一人であり、筆者が MIT に在籍していた頃に幸運にも会うことができた Marvin Minsky は、「科学においては、最も少ないものを研究することで最も多くを学べる」と指摘している。
この文脈でいえば、単純化したコード例を使うと決定的な欠陥に集中しやすく、説明がわかりやすくなる、ということである。
現実には、脆弱性は膨大なコードの中に織り込まれており、処理には重要でもセキュリティとは無関係な多くの事柄と混ざり合っているので、それほど簡単には見分けられない。
現実世界のコード例を見たければ、どれでもよいのでオープンソースソフトウェアプロジェクトのバグデータベースを眺めてみるとよい。
どのプロジェクトにも、セキュリティバグが必ずあるはずである。

警戒するには最初は自分を律する必要があるが、何に気を付けるべきかがわかっていれば、練習を重ねるうちに自然にできるようになる。
覚えておいてほしいのは、警戒が功を奏して攻撃者になりかけた者を撃退できたとしても、そのことを知る機会はおそらくない、ということである。
だから、修正するたびに将来の攻撃を一つ未然に防いだのだと考えて、小さな勝利の一つひとつを祝うとよい。

%----------------------------------------------------------------------
\subsection{ケーススタディ：GotoFail}
% src: Case Study: GotoFail

脆弱性の中には、どのパターンにも当てはまらず、なぜかテストをすり抜けてリリースされてしまう厄介なバグがある。
こうしたことが思いのほか起こりやすいのは、脆弱なコードが典型的な使い方では正しく動くことが多く、意図的な攻撃で揺さぶられたときにだけ有害な振る舞いを見せる、という性質があるからである。
2014 年、Apple は自社製品の大半に向けて一連の重大なセキュリティパッチをひっそりとリリースし、「お客様の保護のため」として問題の説明を拒んだ。
ほどなく世間は、その脆弱性が、重要なセキュリティ保護を事実上無効にしてしまった明らかな編集ミスによるものだと知ることになった。
実際のコードの短い抜粋を調べれば、何が起きたのかはすぐにわかる。
見てみよう。

\subsubsection{1行の脆弱性}
% src: One-Line Vulnerability

まず状況を説明しておくと、問題のコードは安全な接続を確立する途中で実行される。
このコードは、以後の通信を安全にするために、すべてが正しく機能していることを確認するものである。
Secure Sockets Layer（SSL）プロトコルの安全性は、サーバーのデジタル証明書に基づいて認証された形で、交渉した鍵にサーバーが署名していることを確認できるかどうかにかかっている。
より正確には、サーバーは、一時鍵の導出元となるいくつかのデータのハッシュに署名する。
SSL の基本は第11章で扱うが、そうした詳細を知らなくても、この脆弱性の背後にあるコードは追える。
問題のコードは次のとおりである。

% 注: 原典の Web 版では、コードの最初の行（1 つ目の SSLHashSHA1.update を含む if 文）が欠落しており、
%   "goto fail;" から始まっている。本文は「3 つの SSLHashSHA1.update 呼び出し」と述べているが、
%   コードには 2 つしか見えない。ここではソースのとおりに載せた。
\begin{exa}[GotoFail の脆弱なコード]\label{exa:ch8-gotofail-code}
\begin{verbatim}
  goto fail;
if ((err = SSLHashSHA1.update(&hashCtx, &serverRandom)) != 0)
  goto fail;
  goto fail;
if ((err = SSLHashSHA1.update(&hashCtx, &signedParams)) != 0)
  goto fail;

*--snip--*

fail:
  SSLFreeBuffer(&signedHashes);
  SSLFreeBuffer(&hashCtx);
  return err;
\end{verbatim}
\end{exa}

\texttt{SSLHashSHA1.update} の 3 つの呼び出しは、それぞれ一塊のデータをハッシュ関数に入力し、戻り値が 0 でない（エラーである）かどうかを検査する。
ハッシュ計算の詳細はここでの話に関係しないので示さない。
この計算がセキュリティにとって決定的に重要であることだけ知っておけばよい。
通信を認証するには、計算結果が期待値と一致しなければならないからである。

関数の末尾では、2 つのバッファを解放して \texttt{err} の値を返す。
成功なら 0、そうでなければ 0 でないエラーコードである。

このコードが意図しているパターンは明らかである。
エラーを示す 0 でない戻り値がないかを順に検査し、すべて正常なら 0 のまま最後まで通り抜けて、その値を返す。
読者はもう誤りに気付いたかもしれない。
\texttt{goto fail} の行が重複しているのである。
インデントからはそう見えないが、この行は無条件に \texttt{fail} ラベルへジャンプし、残りのハッシュ計算も、ハッシュの検査もまるごと飛ばしてしまう。
余分なジャンプの直前に \texttt{err} へ代入された値は 0 だったので、この関数は突如として、何でも無条件に承認するようになった。
このバグが見つからなかったのは、おそらく、正当な安全な接続はそれまでどおり機能していたからである。
コードはハッシュを検査していなかったが、仮に検査していたとしても、正当な接続ではすべて合格していたはずである。

\subsubsection{フットガンに注意}
% src: Beware of Footguns

GotoFail は、Python などの言語のようにインデントでコードを構造化するやり方がいかに賢明かを示す、格好の論拠である。

\begin{dfn}[フットガン]\label{dfn:ch8-footgun}
自分の足を撃ち抜くこと（自滅すること）を容易にしてしまう機能を \underline{フットガン}（footgun）という。
\end{dfn}

C 言語は、プログラムの構造をインデントではなく構文で決めるので、一種のフットガンを抱えている。
そのため、標準的なコードスタイルの慣習に照らすと誤解を招きかねないインデントを書けてしまう。
コンパイラーはインデントを完全に無視するのに、インデントは別の意味を暗示してしまうからである。

\begin{exa}[誤解を招くインデント]\label{exa:ch8-misleading-indentation}
次のコードを見たとき、
\begin{verbatim}
if ((err = SSLHashSHA1.update(&hashCtx, &serverRandom)) != 0)
  goto fail;
  goto fail;
\end{verbatim}
プログラマーは（注意深く頭の中でコンパイルしてみない限り）次のように読んでしまいやすい。
\begin{verbatim}
if ((err = SSLHashSHA1.update(&hashCtx, &serverRandom)) != 0) {
  goto fail;
  goto fail;
}
\end{verbatim}
一方、コンパイラーは曖昧さなく次のように解釈する。
\begin{verbatim}
if ((err = SSLHashSHA1.update(&hashCtx, &serverRandom)) != 0) {
  goto fail;
}
goto fail;
\end{verbatim}
\end{exa}

単純な編集ミスが、たまたま見落とされやすいうえに、重要なセキュリティ検査のまさに中心部でコードの意味をがらりと変えてしまった。
深刻な脆弱性とは、まさにこういうものである。

言語、API、その他のプログラミングツールやデータ形式に潜む、ほかのフットガンにも注意してほしい。
以降の章で多くの例を見ることになるが、ここで C の構文からもう一つ挙げておく。

\begin{exa}[代入と比較の取り違え]\label{exa:ch8-assignment-vs-comparison}
\texttt{if (x == 8)} と書くべきところを \texttt{if (x = 8)} と書いてしまう誤りである。
前者は \texttt{x} を 8 と比較し、真の場合にだけ then 節を実行する。
後者は \texttt{x} に 8 を代入し、その値は 0 でないので then 節を無条件に実行する。
両者の意味はまったく違う。
スタイルの面から反対する人もいるが、筆者はこうした C の文を \texttt{if (8 == x)} と書くのが好きである。
等号を 2 つ重ね忘れると構文エラーになり、コンパイラーが見つけてくれるからである。
\end{exa}

コンパイラーの警告は、たとえ無害に見えるものでも、この種のミスを見つける助けになる。
GCC コンパイラーの \texttt{-Wmisleading-indentation} 警告オプションは、まさに GotoFail 脆弱性を引き起こしたたぐいの問題を検出するためのものである。
警告の中には、もっと遠回しな形で潜在的な問題を知らせるものもある。
未使用変数の警告はいかにも無害に見えるが、似た名前の変数が 2 つあり、重要なアクセス検査でうっかり違うほうを打ち込んでしまったとしよう。
すると、警告が出ると同時に、決定的に重要な検査に間違ったデータが使われることになる。
警告があらゆる脆弱性を確実に示してくれるわけでは決してないが、確認するのは簡単であり、窮地を救ってくれるかもしれない。

\subsubsection{GotoFail からの教訓}
% src: Lessons from GotoFail

GotoFail からは、重要な教訓がいくつか得られる。
\begin{itemize}
\item 重要なコードでの小さなミスが、セキュリティに壊滅的な影響を与えうる。
\item 脆弱なコードでも、想定どおりの場合には正しく動く。
\item セキュリティにとっては、このようなコードが通常の正当な使い方を通すことより、不正な場合を拒否することをテストするほうが、間違いなく重要である。
\item コードレビューは、見落としから持ち込まれるバグを防ぐ重要なチェックである。注意深いレビュアーがコードの差分を見て、この誤りを見落とすとは考えにくい。
\end{itemize}

この脆弱性は、発生を防げたはずの対策をいくつも示唆している。
中にはこのバグに固有のものもあるが、そうした対策からも、欠陥のあるコードを作ってしまう苦痛を避けるためにほかの場面で使える予防策のヒントが得られるはずである。
有用な対策には次のものがある。
\begin{itemize}
\item もちろん、テストを充実させること。最低限、必要な検査がすべて機能することを確かめるために、\texttt{if} ごとにテストケースを用意しておくべきであった。
\item 到達不能なコードに注意する（多くのコンパイラーには、それを指摘するオプションがある）。GotoFail の場合、これでプログラマーが脆弱性の混入に気付けたかもしれない。
\item 省略できる場合でも丸括弧や波括弧を惜しまず使うなどして、コードをできるだけ明示的に書く。
\item 「リンター」（linter）などのソースコード解析ツールを使う。コードの品質が上がるうえに、その過程で潜在的な脆弱性の一部を指摘してくれて、先回りして直せるかもしれない。
\item 疑わしいパターン（この場合は重複したソースコード行）や、繰り返し起こるその他の誤りを検出するための、その場限りのソースコードフィルターを検討する。
\item 特にセキュリティ上重要なコードについては、完全なコードカバレッジを測定し、必須とする。この例では、2 つ目の \texttt{goto} の後から \texttt{fail} ラベルまでのコードは到達不能であり、それが危険信号になるべきだった。多くのコンパイラーには、それを指摘するオプションがある。
\end{itemize}

これらは、セキュリティを損ないうるバグを見つけるための基本的な手法のほんの一部にすぎない。
新しいクラスのバグに出会ったら、ツールをどう使えば同じ種類のバグの再発を体系的に防げるかを検討してほしい。
そうすれば、長い目で見て脆弱性は減っていくはずである。

%----------------------------------------------------------------------
\subsection{コーディング上の脆弱性}
% src: Coding Vulnerabilities

\begin{quote}
「幸福な家庭はどれも似たものだが、不幸な家庭はそれぞれに不幸である。」 ---Leo Tolstoy
\end{quote}

残念ながら、Leo Tolstoy の小説『アンナ・カレーニナ』の有名な書き出しは、ソフトウェアにもあまりによく当てはまる。
新しい種類のバグが生まれる可能性は無限にあり、潜在的なソフトウェアの脆弱性をすべて網羅したリストを作ろうとするのは愚かな試みである。
分類は役に立つので本書でもその多くを扱うが、それを、ありうるすべてを網羅した完全な分類体系だと取り違えてはならない。

本書は潜在的な欠陥をすべて網羅したリストを示すものでは決してないが、最もよくある分類の多くを代表的に幅広く取り上げる。
この基本的な概観は良い出発点になるはずであり、経験を積むにつれて、ほかの問題も直感的に察知し、安全に避ける方法が身についていくだろう。

\subsubsection{アトミック性}
% src: Atomicity

筆者がこれまでに聞いたコーディングの「武勇伝」のうち最悪の部類の多くは、事象が予期しない順序で起きたために、マルチスレッドや分散プロセスがときおり奇妙な形で干渉し合う、というものであった。
脆弱性も、しばしばこれと同じ条件から生まれる。
唯一の救いは、微妙なタイミングが求められるため、攻撃者にとってエクスプロイトの成功率が低くなりすぎる場合があることだが、だからといって攻撃者がすんなり試すのをあきらめると期待してはならない。

自分のコードがシングルスレッドで行儀よく動くものであっても、ほぼ例外なく、ほかの多くのプロセスが動いているマシンの上で実行される。
そのため、ファイルシステムやその他の共有資源を扱うときには、自分がまったく知らないコードとの競合状態（race condition）に向き合っている可能性がある。

\begin{dfn}[アトミック性]\label{dfn:ch8-atomicity}
ソフトウェアにおける \underline{アトミック性}（atomicity）とは、実質的に単一のステップとして完了することが保証された操作の性質をいう。
\end{dfn}

アトミック性は、このような状況で、脆弱性につながりうる予期しない事態を防ぐための重要な防御手段になる。

\begin{exa}[一時ファイルの作成と競合状態]\label{exa:ch8-tempfile}
何が起こりうるかを説明するために、機密データを一時ファイルにコピーする簡単な例を考える。
非推奨になった Python の \texttt{tempfile.mktemp} 関数は、まだ存在しないことが保証された一時ファイルの名前を返す。
この名前は、アプリケーションが作成して使うファイルの名前として使うことを想定している。
この関数は使ってはならない。
代わりに新しい \texttt{tempfile.NamedTemporaryFile} を使うべきである。
理由は次のとおりである。
\texttt{tempfile.mktemp} が一時ファイルのパスを返してから、コードが実際にそのファイルを開くまでの間に、別のプロセスが割り込む隙があるかもしれない。
そのプロセスが次に生成される名前を推測できれば、先回りしてそのファイルを作成し、（数ある可能性のうちの一つとして）一時ファイルに悪意あるデータを注入できる。
新しい関数は、一時ファイルの作成とオープンを一つのアトミックな操作として行い、途中に何かが割り込む余地をなくすという、すっきりした解決策を提供している。
\end{exa}

\subsubsection{タイミング攻撃}
% src: Timing Attacks

\begin{dfn}[タイミング攻撃]\label{dfn:ch8-timing-attack}
\underline{タイミング攻撃}（timing attack）とは、操作にかかる時間から情報を推測し、本来は非公開であるべきシステムの状態を間接的に知る、サイドチャネル攻撃である。
\end{dfn}

タイミングの違いは、ときにヒントになる。つまり、保護された情報をわずかに漏らし、攻撃者を利することがある。

\begin{exa}[秘密の数の推測]\label{exa:ch8-guess-number}
簡単な例として、1 から 100 までの秘密の数を当てる課題を考える。
「いいえ」と答えが返ってくるまでの時間が、推測がどれだけ外れているかに比例するとわかっていれば、この癖を手がかりに、推測する側はずっと速く正解にたどり着ける。
\end{exa}

\begin{exa}[Meltdown と Spectre]\label{exa:ch8-meltdown-spectre}
Meltdown と Spectre は、ソフトウェアより下の層で動く現代のプロセッサーに対するタイミング攻撃であるが、その原理はそのまま当てはまる。
これらの攻撃は \underline{投機的実行}（speculative execution）の癖を悪用する。
投機的実行では、プロセッサーは速度を稼ぐためにさまざまな検査をいったん緩めながら、先回りして結果を計算しておく。
その中に通常は許されない操作が含まれていれば、プロセッサーはやがてそれを検知し、結果が確定する前に取り消す。
この複雑な投機はすべてプロセッサーの設計どおりの動作であり、私たちが享受している驚異的な速度を実現するのに欠かせない。
しかし、規則が一時停止されたこの投機的実行の間に計算がメモリにアクセスすると、そのメモリがキャッシュに載るという副作用が生じる。
投機的実行が取り消されてもキャッシュはそのまま残るので、この副作用が潜在的なヒントになる。
これらの攻撃はこのヒントを使って、投機的実行の間に何が起きたかを推測する。
具体的には、攻撃コードはキャッシュの状態を調べることで、取り消された投機的実行の間に何が起きたかを推測できる。
メモリキャッシュは実行を速くするが、ソフトウェアから直接見えるわけではない。
しかし、キャッシュされたメモリへのアクセスはずっと速いので、コードはメモリアクセスにかかる時間を測れば、そのメモリ位置の内容がキャッシュにあったかどうかを知ることができる。
複雑なプロセッサーアーキテクチャに対する複雑な攻撃であるが、ここでの要点は、タイミングが保護された情報の状態と相関していれば、それが情報の漏えいとして悪用されうるということである。
\end{exa}

\begin{exa}[パスワードを忘れた場合の機能に対するタイミング攻撃]\label{exa:ch8-forgot-password-timing}
もっと単純な、純粋にソフトウェアだけで成り立つタイミング攻撃の例として、友人（あるいは友人を装った敵？）が特定のオンラインサービスにアカウントを持っているかどうかを知りたいが、そのアカウント名は知らない、という場合を考える。
「パスワードを忘れた場合」の機能は、「リマインダー」を送るために、ユーザーにアカウント名と電話番号を尋ねる。
ところがその実装は、まずデータベースで電話番号を検索し、見つかった場合にだけ、ひも付いたアカウント名を検索して入力と一致するかを確認するとしよう。
検索にはそれぞれ数秒かかり、その遅れはユーザーにもわかるとする。

まず、実在のユーザーとは一致しそうにないランダムなアカウント名（たとえばキーボードをでたらめに叩いて作る）と電話番号をいくつか試し、「そのようなアカウントはありません」という応答が返るまでに、決まって約 3 秒かかることを確かめる。
次に、自分の電話番号でサインアップし、自分の番号とランダムなアカウント名の一つで「パスワードを忘れた場合」の機能を試す。
すると今度は、応答が返るまでに 5 秒、つまりほぼ 2 倍の時間がかかることがわかる。

ここまでわかれば、同じ未使用のアカウント名で友人の電話番号を試せばよい。
応答に 5 秒かかれば友人の電話番号はデータベースにあり、3 秒ならないとわかる。
タイミングを観察するだけで、ある電話番号がデータベースにあるかどうかを推測できるのである。
特定の病気を持つ患者のためのフォーラムのように、会員であること自体が機密性の高い私的な情報を明かしてしまう場合には、このようなタイミング攻撃によって有害な情報漏えいが起こりうる。
\end{exa}

ソフトウェアでは、遅い操作が順に並んでいて（\texttt{if}…\texttt{if}…\texttt{if}…\texttt{if}… を思い浮かべるとよい）、実行がその一連の処理のどこまで進んだかがわかれば価値ある情報を推測できる、という場合に、タイミングの違いが自然に生じる。
情報が漏れるのにどれだけのタイミングの差が必要かは、多くの要因で決まる。
オンラインのアカウント確認の例では、Web のアクセスにつきものの通常の遅延を考えると、はっきりした信号になるには数秒の差が必要である。
これに対して、同じマシン上で動くコードで Meltdown や Spectre を悪用する場合には、1 ミリ秒に満たない時間差でも測定でき、しかも意味を持つことがある。

最善の緩和策は、時間差を許容できる程度、つまり気付かれない水準まで小さくすることである。
電話番号がデータベースにあることが漏れないようにするには、どちらの場合も 1 回のデータベース検索で処理するようにコードを変えれば十分だろう。
タイミングの違いが本質的に避けられず、しかもタイミングのサイドチャネルが深刻な情報漏えいにつながりうる場合には、リスクを緩和するためにできるのは、人為的な遅延を入れてタイミングの信号をぼかすことくらいしかない。

\subsubsection{シリアライゼーション}
% src: Serialization

\begin{dfn}[シリアライゼーション]\label{dfn:ch8-serialization}
\underline{シリアライゼーション}（serialization）とは、データオブジェクトをバイトストリームに変換する、よく使われる手法である。
得られたバイト列を保存または送信しておけば、後で \underline{デシリアライゼーション}（deserialization）によって同等のデータオブジェクトを再構成できる。
\end{dfn}

スタートレックの転送装置に少し似ていて、オブジェクトをバイトストリームにしてから、時間と空間を越えて「転送」するのである。
オブジェクトを「乾燥」させてから「水で戻す」この機能はオブジェクト指向プログラミングでは便利であるが、途中で改ざんされる可能性が少しでもあれば、この手法は本質的にセキュリティリスクになる。
攻撃者は重要なデータの値を変えられるだけでなく、不正なバイト列を組み立てて、デシリアライゼーションのコードに有害な操作を実行させることさえできる。
デシリアライゼーションは信頼できるシリアライズ済みデータに対して使う場合にだけ安全なので、これは信頼できない入力の問題の一例である。

問題は、これらのライブラリの作りが悪いことではない。任意のオブジェクトを組み立てるという仕事をこなすのに必要な操作をライブラリが実行するには、データを信頼することが前提になる、という点が問題なのである。
デシリアライゼーションは事実上、入力のシリアライズされたバイト列が指示することを何でも実行するインタープリターであり、したがって信頼できないデータに対して使うのは決して良い考えではない。

\begin{exa}[Python の unpickling]\label{exa:ch8-unpickling}
Python のデシリアライゼーション操作（「unpickling」と呼ばれる）は、unpickle するデータに悪意あるバイト列を埋め込むだけで、簡単にだまされて任意のコードを実行してしまう。
\end{exa}

シリアライズしたバイトデータを、改ざんの心配なく安全に保存・送信できる場合（たとえば第5章で述べた MAC やデジタル署名を使う場合）を除けば、シリアライゼーションはいっさい使わないのが最善である。

%----------------------------------------------------------------------
\subsection{常連の容疑者たち}
% src: The Usual Suspects

\begin{quote}
「悪魔が仕掛けた最大のトリックは、自分は存在しないと世間に信じ込ませたことだ。」 ---Charles Baudelaire
\end{quote}

以降のいくつかの章では、脆弱性としてコードに繰り返し現れる「常連の容疑者たち」の多くを取り上げる。
この章では、GotoFail と、アトミック性、タイミング攻撃、シリアライゼーションの問題を見てきた。
これから探るトピックを予告しておくと、次のとおりである。
\begin{itemize}
\item 固定幅整数の脆弱性
\item 浮動小数点の精度の脆弱性
\item バッファオーバーフローとその他のメモリ管理の問題
\item 入力検証
\item 文字列の誤った扱い
\item インジェクション攻撃
\item Web セキュリティ
\end{itemize}

これらの問題の多くは当たり前に思えるだろうが、どれもソフトウェアの脆弱性の根本原因として、ほとんど衰えることなく繰り返し起こり続けており、終わりは見えない。
これらの脆弱性のクラスの多くは何十年も前から存在しているので、過去の失敗から学ぶことは重要である。
しかし、ありうるセキュリティバグがすべて目録に載っているかのように、過去だけを振り返る姿勢をとるのは誤りである。
ありうる落とし穴をすべて前もって警告できる本はないが、これらの例を研究すれば、その背後にあるより深いパターンと教訓をつかめる。

\newpage

%% =======================================================
%%  part: ch9
%% =======================================================
%% =======================================================
%%  part: ch9
%% =======================================================
\section{低レベルのコーディング上の欠陥}
% src: 9: Low-Level Coding Flaws

この章では、\cite{Kohnfelder2021} の第9章に沿って、機械語に近いレベルで動くコードに特有の、算術とメモリアクセスにかかわる欠陥を扱う。

\begin{quote}
「低レベルプログラミングはプログラマーの魂によい。」 --- John Carmack
\end{quote}

続くいくつかの章では、セキュリティのためにプログラマーが知っておくべきコーディング上の落とし穴を、古典的なものから順にひととおり見ていく。
この章で扱うのは、機械語に近いレベルで動くコードによく見られる基本的な欠陥である。
ここで論じる問題は、コードが固定サイズの数値の容量か、割り当てたメモリバッファの容量を超えたときに起こる。
現代的な言語の多くは、コードをこうした危険から切り離す高水準の抽象化を備えているが、そうした安全な言語を使うプログラマーにとっても、これらの欠陥を理解しておくことは役に立つ。
少なくとも、言語が自分の代わりに何をしてくれているのか、それがなぜ重要なのかがよくわかるからである。

C や C++ のようにこうした低レベルの機能をそのまま使わせる言語は、今なお多くのソフトウェア分野で主流であり、そこから生じうる脅威は決して理論上の話ではない。
Python のような現代的な言語は、たいていハードウェアを十分に抽象化しているので、この章で述べる問題は起こらない。
しかし、効率を突き詰めるためにハードウェアのレベルに近づきたいという誘惑は今も強い。
人気のある言語の中には、両方の世界をプログラマーが選べるものもある。
Java と C\# の基本型には、型安全なオブジェクトライブラリのほかに固定幅の整数があり、通常の安全装置の多くを外す「unsafe」モードもある。
Python の float 型は、後述の「浮動小数点精度の脆弱性」（\ref{subsec:ch9-floating-point-precision} 項）で説明するようにハードウェアのサポートに頼っているので、その制約もそのまま引き継いでおり、プログラマーはそれに対処しなければならない。
% 注: 原文は「on page XX」。本ノートでは該当する項への \ref に置き換えた。

低レベルの機能を使わせる言語をまったく使わない読者は、この章を読み飛ばしたくなるかもしれないし、読み飛ばしても本書全体の流れを見失うことはない。
それでも原著者は、ひととおり読むことを勧めている。
自分が使う言語やライブラリが何を守ってくれて何を守ってくれないのかを理解し、自分の代わりに行われていることをすべて把握しておくのが一番よいからである。

ハードウェアに近いレベルのプログラミングは、うまくやれば非常に強力だが、手間が増え、壊れやすくなるという代償を払うことになる。
この章では、低レベルの抽象化を使ったコーディングに特有の、最もよくある脆弱性のクラスを取り上げる。

この章が扱うのは、もっぱらハードウェアのレベルかそれに近いところで動くコードに起因するバグである。そのため、こうした演算の多くは正確な結果がプラットフォームや言語によって異なる、という点を理解しておく必要がある。
例はできるだけ具体的になるようにしてあるが、実装の違いによって結果が変わることがある。
そして、計算結果が予測できない形で変わりうるからこそ、これらの問題は見落とされやすく、セキュリティに影響しうるのである。
細部はハードウェアやコンパイラなどの要因によって異なるが、この章で紹介する概念は一般に当てはまる。

%----------------------------------------------------------------------
\subsection{算術の脆弱性}
% src: Arithmetic Vulnerabilities

プログラミング言語には、算術演算子を数学的に定義するものもあれば、プロセッサの対応する命令に合わせて定義するものもある。
後で見るように、この二つは完全に同じではない。
ここでいう\emph{低レベル}とは、機械語の命令に頼っているために、ハードウェアの癖や制約に向き合わなければならないプログラミング言語の機能のことである。

コードは整数演算だらけである。
整数演算は数値計算だけでなく、文字列の比較やデータ構造へのインデックスアクセスなどにも使われる。
ハードウェアの命令は、より広い範囲の値を扱うソフトウェア上の抽象化よりもはるかに速くて使いやすいので、使わずに済ませるのは難しい。
しかし、その便利さと速さには、オーバーフローのリスクがついて回る。

\begin{dfn}[オーバーフロー]\label{dfn:ch9-overflow}
計算の結果が固定幅の整数の容量を超えることを \underline{オーバーフロー}（overflow）という。
オーバーフローが起こると予期しない結果になり、それが脆弱性につながることがある。
\end{dfn}

浮動小数点演算は整数演算より範囲が広いが、精度が限られているので、やはり予期しない結果になることがある。
浮動小数点数にも限界はある（単精度では $10^{38}$ 程度）が、限界を超えると無限大を表す特別な値になるという好ましい性質を持っている。

算術命令の実装をハードウェアのレベルまで掘り下げて知りたい読者は、Jonathan E. Steinhart, \textit{The Secret Life of Programs}（No Starch Press, 2019）で詳しく学べる。

%----------------------------------------------------------------------
\subsubsection{固定幅整数の脆弱性}\label{subsec:ch9-fixed-width-integer}
% src: Fixed-Width Integer Vulnerabilities

原著者が正社員として最初に就いた仕事は、ミニコンピューター上でアセンブリ言語のデバイスドライバを書くことであった。
ミニコンピューターは現代の水準からすれば笑ってしまうほど非力だったが、ハードウェアの仕組みを学ぶには絶好の機会になった。
回路基板を眺めれば、すべての配線とすべてのチップ（中の論理ゲートの数もそれほど多くない）が見えたからである。
算術論理演算ユニット（加算、減算、ブール演算しかできない）と、メモリにつながったレジスタが目に見えるので、コンピューターがどう動くのかを正確に知ることができた。
現代のプロセッサは途方もなく複雑で、数十億個の論理ゲートを含んでおり、人間がざっと眺めて理解できる範囲をはるかに超えている。

今日では、ほとんどのプログラマーが、機械語や CPU アーキテクチャの複雑さから守ってくれる高水準言語を学び、使っている。
固定幅整数は Java や C/C++ をはじめ多くの言語の最も基本的な構成要素であり、計算がその限られた範囲を超えると、何であれ\emph{黙って}誤った結果になる。

現代のプロセッサは 32 ビットか 64 ビットのアーキテクチャであることが多いが、その動作はもっと小さいサイズで考えれば理解できる。

\begin{exa}[16 ビット符号なし整数のオーバーフロー]\label{exa:ch9-overflow-300x300}
符号なし 16 ビット整数でオーバーフローの例を見てみる。
16 ビット整数は 0 から 65,535（$2^{16}-1$）までの任意の値を表せる。
例えば、300 に 300 を掛ければ 90,000 になるはずだが、この数は使っている固定幅整数の範囲を超えている。
そのため、オーバーフローが起きて、実際に得られる結果は 24,464 になる（期待される結果より 65,536 小さい）。
\end{exa}

オーバーフローを数学的に、剰余演算、つまり割り算の余りとして捉える人もいる（例えば、先ほどの計算結果は 90,000 を 65,536 で割った余りになっていた）。
2 進数や 16 進数での切り捨てとして捉える人や、ハードウェアの実装の観点から捉える人もいる。
しかし、どれもしっくりこないなら、大きすぎる値になる計算は期待どおりの結果にならない、ということだけ覚えておけばよい。

\begin{rem}[オーバーフローした値そのものは重要でない]\label{rem:ch9-overflow-value-unimportant}
オーバーフローの緩和策は、そもそもオーバーフローを起こさないことを目指すものなので、結果として得られる正確な値はふつう重要ではない。
\end{rem}

\begin{rem}[16 ビットアーキテクチャによる 2 進数の算術の簡単な復習]\label{rem:ch9-binary-refresher}
2 進数の算術になじみの薄い読者のために、本文の $300 \times 300$ の計算を図で分解して示す。
10 進数が 0 から 9 の数字で書かれるのと同じように、2 進数は 0 と 1 の数字で書かれる。
そして、10 進数では左の桁に行くごとに位が 10 倍になるのと同じように、2 進数では左に行くにつれて桁の重みが 2 倍（1, 2, 4, 8, 16, 32, 64, \ldots）になっていく。
表 \ref{tab:ch9-binary-300} は、10 進数の 300 を 16 ビットの 2 進数で表したもので、2 の冪の桁位置を 0 から 15 の 10 進数で示している。

2 進表現の値は、1 が立っている桁位置に対応する 2 の冪を合計したものである。
つまり、300 は $2^8 + 2^5 + 2^3 + 2^2$（$256 + 32 + 8 + 4$）であり、2 進数では 100101100 になる。

次に、300 にそれ自身を掛ける計算を 2 進数で見てみる（表 \ref{tab:ch9-binary-multiplication}）。
紙の上で 10 進数の掛け算をするときと同じように、被乗数を、乗数の各桁に対応する位置までずらしながら繰り返し足していく。
右から見ていくと、最初の 1 の右には 2 つの位置があるので、1 つ目のコピーは 2 桁左にずらす。
以下同様に、各コピーの右端が乗数の 1 のどれかの真下にそろうように並べる。
左にはみ出した灰色の数字は 16 ビットのレジスタに収まらないので切り捨てられる。
ここでオーバーフローが起こる。
あとは各部分を（もちろん 2 進数で）足し合わせれば結果が得られる。
2 は 2 進数では 10（$2^1$）なので、位置 5 で最初の繰り上がりが起こる（$1+1+0 = 10$）。
0 を書いて 1 を繰り上げる。
固定幅整数の掛け算はこのように行われ、こうして値が黙って切り捨てられるのである。
\end{rem}

% 図 9-1: Binary number 300 example（16 ビットの 2 進表現と桁位置 0--15）
% 注: 原文の図は画像のみ。本文から内容が一意に決まるので表として再現した。
\begin{table}[H]
\centering
\small
\setlength{\tabcolsep}{3pt}
\begin{tabular}{l|*{16}{c}}
桁位置 & 15 & 14 & 13 & 12 & 11 & 10 & 9 & 8 & 7 & 6 & 5 & 4 & 3 & 2 & 1 & 0 \\ \hline
300 & 0 & 0 & 0 & 0 & 0 & 0 & 0 & 1 & 0 & 0 & 1 & 0 & 1 & 1 & 0 & 0 \\
\end{tabular}
\caption{10 進数の 300 の 16 ビット 2 進表現}\label{tab:ch9-binary-300}
\end{table}

% 図 9-2: Binary multiplication 300x300 example
% 注: 原文の図は画像のみ。本文の説明（2 進の筆算、16 ビットを超える部分は灰色で切り捨て）から
%   内容が一意に決まるので表として再現した。16 ビットを超えるのは位置 16 の 1 桁だけである。
\begin{table}[H]
\centering
\small
\setlength{\tabcolsep}{3pt}
\newcommand{\gr}[1]{\textcolor{gray}{#1}}
\begin{tabular}{l|c|*{16}{c}}
桁位置 & \gr{16} & 15 & 14 & 13 & 12 & 11 & 10 & 9 & 8 & 7 & 6 & 5 & 4 & 3 & 2 & 1 & 0 \\ \hline
被乗数 300 & & 0 & 0 & 0 & 0 & 0 & 0 & 0 & 1 & 0 & 0 & 1 & 0 & 1 & 1 & 0 & 0 \\
乗数 300 ($\times$) & & 0 & 0 & 0 & 0 & 0 & 0 & 0 & 1 & 0 & 0 & 1 & 0 & 1 & 1 & 0 & 0 \\ \hline
位置 2 の 1 による部分積 & & & & & & & 1 & 0 & 0 & 1 & 0 & 1 & 1 & 0 & 0 & 0 & 0 \\
位置 3 の 1 による部分積 & & & & & & 1 & 0 & 0 & 1 & 0 & 1 & 1 & 0 & 0 & 0 & 0 & 0 \\
位置 5 の 1 による部分積 & & & & 1 & 0 & 0 & 1 & 0 & 1 & 1 & 0 & 0 & 0 & 0 & 0 & 0 & 0 \\
位置 8 の 1 による部分積 & \gr{1} & 0 & 0 & 1 & 0 & 1 & 1 & 0 & 0 & 0 & 0 & 0 & 0 & 0 & 0 & 0 & 0 \\ \hline
和（16 ビット: 24,464） & \gr{1} & 0 & 1 & 0 & 1 & 1 & 1 & 1 & 1 & 1 & 0 & 0 & 1 & 0 & 0 & 0 & 0 \\
\end{tabular}
\caption{$300 \times 300$ の 2 進数による掛け算（灰色の桁は 16 ビットを超えるため切り捨てられる）}\label{tab:ch9-binary-multiplication}
\end{table}

大事なのは、計算でどんな値が得られるかを正確に知ることではなく、2 進数の算術の癖を見越しておくことである。
言語やコンパイラによっては、得られる値が\underline{未定義}（not defined）の場合もある。
つまり、言語仕様が特定の値を保証していないのである。

\begin{rem}[未定義の結果に頼らない]\label{rem:ch9-undefined-behavior}
言語仕様で技術的に「未定義」とされている演算は、予測できるように見えるかもしれないが、言語仕様の保証がない以上、薄氷を踏んでいることになる。
セキュリティの観点からの結論は、言語仕様を知り、未定義になりうる計算を避けることが重要だ、ということである。
未定義の結果を検出する巧妙な方法を見つけようと、小賢しい実験をしてはならない。
ハードウェアが変わったり、コンパイラが新しいバージョンになったりすると、そのコードは動かなくなるかもしれないからである。
\end{rem}

算術の結果を誤って計算すると、コードはさまざまな形で壊れうる。
その影響はしばしば雪だるま式に不具合の連鎖を引き起こし、最後にはクラッシュやブルースクリーンに至る。
整数オーバーフローによる脆弱性のよくある例としては、バッファオーバーフロー（後述の「バッファオーバーフロー」（\ref{subsec:ch9-buffer-overflow} 項）で論じる）、値の比較の誤り、販売時に代金を請求するはずが返金してしまう状況などがある。

これらの問題は、範囲外になりうる計算を行う\emph{前に}、すべての数がまだ範囲内にあるうちに緩和するのが一番よい。
確実な方法は簡単で、許容される最大値より大きい整数サイズを使い、事前にチェックを置いて不正な値が決して紛れ込まないようにすることである。
例えば、先ほどの $300 \times 300$ を計算するには、16 ビット値どうしの積なら何でも扱える 32 ビットの算術を使う。
結果を 16 ビットに戻さなければならない場合は、範囲内にあることを確かめる 32 ビットの比較で守る。

次のコードは、16 ビット符号なし整数どうしを掛けて 32 ビットの結果を得る C のコードである。
演算子の優先順位により、キャストは掛け算より先に適用されるが、原著者はわかりやすさのためにキャストの周りにもう一組括弧を付けるのを好んでいる（こうした脆弱性がどう紛れ込むのかをより現実に即して見るために、この章の後半でもっと包括的な例を示す）。

\begin{verbatim}
uint32_t simple16(uint16_t a, uint16_t b) {
  return ((uint32_t)a) * ((uint32_t)b);
}
\end{verbatim}

固定幅整数が黙ってオーバーフローすること自体は難しい話ではないが、実際にはこうした欠陥に経験豊富なプログラマーでさえ悩まされ続けている。
原因の一つは、プログラミングでは整数演算がいたるところにあることである。
ポインタ演算や配列のインデックスのような暗黙の使い方もあり、そこにも同じ緩和策を適用しなければならない。
もう一つの難しさは、すべての変数について、妥当な値の範囲だけでなく、狡猾な攻撃者による操作も考慮したうえでコードが出会いうる値の範囲を、常に念頭に置いておくという厳密さが求められることである。

プログラミングをしていると、ただ数値を操作しているだけのように感じることが多いが、こうした計算はとても壊れやすい。
その壊れやすさを忘れてはならない。

%----------------------------------------------------------------------
\subsubsection{浮動小数点精度の脆弱性}\label{subsec:ch9-floating-point-precision}
% src: Floating-Point Precision Vulnerabilities

浮動小数点数は、多くの点で固定幅整数より頑健で、癖も少ない。
ここでは、浮動小数点数を、符号ビット（正か負か）、固定精度の仮数部、そして仮数部に掛ける 2 の指数からなるものと考えておけばよい。
広く使われている IEEE 754 の倍精度の仕様では、10 進数で 15 桁（2 進数で 53 桁）の精度が得られる。
しかも、その極めて大きい上限を超えたときには、固定幅整数のように突拍子もない値に切り捨てられるのではなく、符号付きの無限大（演算によっては NaN（not a number））になる。

15 桁の精度があれば、米国の連邦予算（現在、数兆ドル）をセント単位まで集計できるので、精度が失われるリスクが問題になることはめったにない。
それでも、精度の損失は下位の桁で黙って起こるし、浮動小数点数は 10 進数ではなく 2 進数で表現されているので、驚くような結果になることがある。

\begin{exa}[0.1 + 0.2]\label{exa:ch9-point-one-plus-point-two}
10 進数の小数は 2 進数では必ずしも正確に表せないので、$0.1 + 0.2$ は 0.30000000000000004 になり、0.3 と等しくならない。
% 注: 原文はスクレイプの崩れで「0.1」と「0.2」の間の演算子が失われているが、文脈から 0.1 + 0.2 と判断した。
こうした厄介な結果になるのは、10 進数で $1/7$ が循環小数になるのと同じように、2 進数では $1/10$ が無限に循環し（$0.00011001100\ldots$ で、1100 が永遠に続く）、最下位のビットに誤差が生じるからである。
\end{exa}

\begin{dfn}[アンダーフロー]\label{dfn:ch9-underflow}
浮動小数点の計算で、下位のビットに誤差が入り込んで精度が失われることを \underline{アンダーフロー}（underflow）という。
\end{dfn}

アンダーフローによるずれは割合としてはごくわずかだが、大きさの異なる値を扱うときには直感に反する結果になることがある。

\begin{exa}[大きさの異なる値の加減算]\label{exa:ch9-javascript-underflow}
すべての数が浮動小数点数である JavaScript で書いた、次のコードを考える。
\begin{verbatim}
var a = 10000000000000000
var b = 2
var c = 1
console.log(((a+b)-c)-a)
\end{verbatim}
数学的には、最後の行の式の結果は \texttt{b-c} に等しくなるはずである。
\texttt{a} は最初に足され、最後に引かれるからである（\texttt{console.log} 関数は、式の値を出力する便利な方法である）。
しかし実際には、\texttt{a} の値が十分に大きいので、使える精度が限られている以上、ずっと小さい値を足したり引いたりしても何も変わらない。
そのため、最後に \texttt{a} を引くと、結果は 0 になる。
\end{exa}

この例のような計算が近似で構わないものなら誤差は無害だが、完全な精度が必要な場合や、桁の大きさが異なる値が計算に混じる場合には、優れたプログラマーは注意しなければならない。
こうしたずれがコード中のセキュリティ上重要な判断に影響しうるとき、脆弱性が生まれる。
アンダーフローの誤差は、チェックサムや複式簿記のように正確な結果が欠かせない計算で問題になりうる。

多くの浮動小数点の計算では、先ほどの例のような劇的なアンダーフローがなくても、値を正確に表せない場合には下位のビットに少しずつ誤差がたまっていく。

\begin{rem}[浮動小数点数の等値比較を避ける]\label{rem:ch9-float-equality}
浮動小数点の値を等値（または非等値）で比較するのは、ほとんどの場合賢明ではない。
この演算は、計算した値のごくわずかな違いも許容できないからである。
そこで、\texttt{(x == y)} の代わりに、アプリケーションに合った \texttt{delta} の値を決め、値が小さい範囲内にあるかどうかで比較する（\texttt{(x > y - delta \&\& x < y + delta)}）。
Python には、この判定をもう少し洗練させた形で行う補助関数 \texttt{math.isclose} がある。
\end{rem}

高い精度が必要な場合は、超高精度の浮動小数点表現（IEEE 754 は 128 ビットと 256 ビットの形式を定めている）を使うことを検討する。
計算の要件によっては、任意精度の 10 進数表現や有理数表現が最善の選択になることもある。
言語がネイティブにサポートしていない場合でも、ライブラリでこの機能が提供されていることが多い。

%----------------------------------------------------------------------
\subsubsection{例：浮動小数点のアンダーフロー}\label{subsec:ch9-example-float-underflow}
% src: Example: Floating-Point Underflow

浮動小数点のアンダーフローは軽く見られがちだが、精度の損失は壊滅的な結果を招きうる。
ここでは、浮動小数点の値を使うオンライン注文システムのビジネスロジックの簡単な例を Python で示す。
次のコードは、購入注文が全額支払われているかを確認し、支払われていれば製品の出荷を承認するためのものである。

\begin{verbatim}
from collections import namedtuple
PurchaseOrder = namedtuple('PurchaseOrder', 'id, date, items')
LineItem = namedtuple('LineItem',
                      ['kind', 'detail', 'amount', 'quantity'],
                      defaults=(1,))
def validorder(po):
    """Returns an error text if the purchase order (po) is invalid,
    or list of products to ship if valid [(quantity, SKU), ...].
    """
    products = []
    net = 0
    for item in po.items:
        if item.kind == 'payment':
            net += item.amount
        elif item.kind == 'product':
            products.append(item)
            net -= item.amount * item.quantity
        else:
            return "Invalid LineItem type: %s" % item.kind
    if net != 0:
        return "Payment imbalance: $%0.2f." % net
    return products
\end{verbatim}

購入注文は、製品か支払いのどちらかを表す明細行（line item）からなる。
支払いの合計からクレジット（返金）を差し引き、さらに注文された製品の費用の合計を引くと、0 になるはずである。
支払いは事前に検証済みだが、その処理の細部を一つはっきりさせておく。
顧客が請求をすぐに全額取り消した場合は、手数料のかかるクレジットカード処理業者への問い合わせを行わずに、クレジットとデビットの両方が明細行として現れる。
また、品目に記載された価格は正しいものとする。

浮動小数点の計算に注目すると、支払いの明細行では金額が \texttt{net} に足され、製品では金額と数量の積が引かれている。
次の呼び出しは Python の doctest として書かれており、\verb|>>>| の行が実行するコード、その後に続くのが期待される戻り値である。

\begin{verbatim}
>>> tv = LineItem(kind='product', detail='BigTV', amount=10000.00)
>>> paid = LineItem(kind='payment', detail='CC#12345', amount=10000.00)
>>> goodPO = PurchaseOrder(id='777', date='6/16/2019', items=[tv, paid])
>>> validorder(goodPO)
[LineItem(kind='product', detail='BigTV', amount=10000.0, quantity=1)]
>>> unpaidPO = PurchaseOrder(id='888', date='6/16/2019', items=[tv])
>>> validorder(unpaidPO)
'Payment imbalance: $-10000.00.'
\end{verbatim}

コードは期待どおりに動き、1 つ目の、全額支払い済みのテレビの取引を承認し、支払いの記載がない注文を拒否する。

では、このコードを破ってテレビを「盗んで」みよう。
すでに脆弱性に気づいているなら、自分でこの関数をだましてみるとよい演習になる。
原著者は次のようにしてテレビ 1,000 台をただで手に入れた。説明はコードの後に続く。

\begin{exa}[アンダーフローを悪用した不正な注文]\label{exa:ch9-fraud-po}
\begin{verbatim}
>>> fake1 = LineItem(kind='payment', detail='FAKE', amount=1e30)
>>> fake2 = LineItem(kind='payment', detail='FAKE', amount=-1e30)
>>> tv = LineItem(kind='product', detail='BigTV', amount=10000.00, quantity = 1000)
>>> nonpayment = [fake1, tv, fake2]
>>> fraudPO = PurchaseOrder(id='999', date='6/16/2019', items=nonpayment)
>>> validorder(fraudPO)
[LineItem(kind='product', detail='BigTV', amount=10000.0, quantity=1000)]
\end{verbatim}
ここでのトリックは、\texttt{1e30}、つまり $10^{30}$ という法外な金額の偽の支払いと、その直後の請求の取り消しにある。
このでたらめな数どうしは合計が 0（$10^{30} - 10^{30}$）になるので、会計のチェックを通ってしまう。
取り消しのデビットとクレジットの間に、テレビ 1,000 台を注文する明細行が挟まっている点に注意してほしい。
最初の数があまりに巨大なので、テレビの費用を引いても完全にアンダーフローして消えてしまい、その後にクレジット（負の数）を足すと結果は 0 になる。
もしクレジットが支払いの直後に続き、その後にテレビの明細行が来ていれば、結果は違っていて、エラーが正しく検出されていたはずである。
\end{exa}

アンダーフローの感覚をもっと正確につかむために、そして何より、コードを安全にするための安全な値の範囲の見積もり方を示すために、もう少し掘り下げてみる。
この攻撃で $10^{30}$ を選んだことに特に意味はなく、このトリックは約 $10^{24}$ というもっと小さい数でもうまくいくが、$10^{23}$ ではうまくいかない。
1 台 \$10,000 のテレビ 1,000 台の費用は \$10,000,000、つまり $10^7$ である。
したがって、$10^{23}$ の偽の請求では、$10^7$ という値が計算にわずかに影響し始める。
これは約 16 桁（$23-7$）の精度に当たる。
先に述べた 15 桁の精度は、安全側に倒した経験則としての近似値であり（2 進数の精度は 10 進数の 15.95 桁に相当する）、私たちの多くは自然に 10 進数で考えるので便利である。
しかし、浮動小数点の表現は実際には 2 進数なので、数ビットずれることがある。

この考察を踏まえて、脆弱性を修正しよう。
浮動小数点のまま扱いたいなら、数の範囲を制限する必要がある。
製品の最低価格を \$0.01（$10^{-2}$）とし、精度を 15 桁とすると、支払い額の上限を \$$10^{13}$（$15-2$）、つまり \$10 兆に設定できる。
この上限があればアンダーフローは避けられるが、実際には、現実的な最大注文額に見合ったもっと小さい上限にするのが一番よい。

任意精度の数値型を使えばアンダーフローは避けられる。
Python では、ネイティブの整数型や \texttt{fractions.Fraction} がそれに当たる。
もっと高精度の浮動小数点計算を使えばこの特定の攻撃は防げるが、さらに極端な値に対しては依然としてアンダーフローが起こりうる。
Python は動的型付けなので、これらの型の値でコードを呼び出せば攻撃は失敗する。
しかし、このコードを任意精度の型で書いたから安全だと考えていても、攻撃者が何らかの方法で float を紛れ込ませることができれば、脆弱性は再び現れる。
だからこそ、範囲チェックを行うこと、あるいは呼び出し側が期待どおりの型を渡すと信頼できない場合には、計算の前に入力値を安全な型に変換することが重要なのである。

%----------------------------------------------------------------------
\subsubsection{例：整数オーバーフロー}\label{subsec:ch9-example-integer-overflow}
% src: Example: Integer Overflow

固定幅整数のオーバーフローによる脆弱性は、後から見ればまったく明白なことが多く、このクラスのバグは何年も前からよく知られている。
それでも経験豊富なプログラマーが何度もこの罠にはまる。
オーバーフローが起こりうるとは思っていないこともあれば、無害だと見誤ることもあるし、まったく考えていないこともある。
こうしたバグがいかに簡単に紛れ込むかを感じてもらうために、次の例では、もっと大きな計算の中に潜む脆弱性を示す。
実際には、脆弱な計算はもっと込み入っていて、変数の値も予測しにくいことが多いが、ここでは説明のために、何が起きているかが見やすい簡単なコードを使う。

単純な給与計算の式を考える。
労働時間に時給を掛けると、支払額の合計（ドル）が出る。
この単純な計算を時間とドルの端数まで含めて行えば、完全な精度が得られる。
その反面、丸めのせいで細部が少し複雑になり、後で見るように、整数オーバーフローが簡単に起こる。

正確な精度を得るために 32 ビット整数を使い、ドルの値はセント（\$0.01 単位）、時間は 1000 分の 1 時間（0.001 時間単位）で計算するので、数はかなり大きくなる。
しかし、32 ビット整数の最大値 \texttt{UINT32\_MAX} は 40 億を超える（$2^{32}-1$）ので、次のように考えて安全だと見なす。
会社の方針で、有給の労働は週 100 時間（1000 分の 1 時間単位で 100,000）までに制限されているので、時給の上限を \$400（40,000 セント）とすれば、給与の最大値は 4,000,000,000 になる（\$40,000 は 1 週間の給与としては悪くない）。

次は C による給与の計算で、変数と定数はすべて \texttt{uint32\_t} の値として定義されている。

\begin{verbatim}
if (millihours > max_millihours      // 100 hours max
   || hourlycents > max_hourlycents) // $400/hour rate max
  return 0;
return (millihours * hourlycents + 500) / 1000; // Round to $.01
\end{verbatim}

範囲外の引数に対してエラーを示す値を返す \texttt{if} 文は、後に続く計算でのオーバーフローを防ぐ欠かせないガードである。

\texttt{return} 文の計算には説明が要る。
時間を 1000 分の 1 単位で表しているので、実際の支払額を得るには結果を 1,000 で割らなければならない。
そこで、丸めのために先に 500（除数の半分）を足しておく。
簡単な例で確かめてみる。
10 時間（10,000）に時給 \$10.00（1,000）を掛けると 10,000,000 になり、丸めのために 500 を足すと 10,000,500、これを 1,000 で割ると 10,000、つまり \$100.00 となって、正しい値が得られる。
この時点ですでに、固定幅整数の制約による切り捨てと隣り合わせだという意味で、このコードは壊れやすいと考えるべきである。

ここまでのところ、コードはどんな入力に対しても正しく動くが、経営陣が新しい残業の方針を発表したとしよう。
残業時間（最初の 40 時間を超えて働いた時間）にはすべて時給を 50 パーセント上乗せするように、コードを変更しなければならない。
さらに、経営陣が後で簡単に変更できるように、このパーセンテージはパラメータにする必要がある。

残業時間の割増賃金を加えるために、\texttt{overtime\_percentage} を導入する。
そのコードは示さないが、値は 150 で、残業時間には通常の給与の 150 パーセントを支払うという意味である。
給与が増えるので、時給 \$400 の上限では整数オーバーフローを防ぐには高すぎ、もう使えない。
しかし、その時給はもともと実用上の上限としては非現実的だったので、念のために半分にして、最高時給を \$200 としよう。

\begin{verbatim}
if (millihours > max_millihours      // 100 hours max
    || hourlycents > max_hourlycents) // $200/hour rate max
  return 0;
if (millihours > overtime_millihours) {
  overage_millihours = millihours - overtime_millihours;
  overtimepay = (overage_millihours * hourlycents * overtime_percentage
                   + 50000) / 100000;
  basepay = (overtime_millihours * hourlycents + 500) / 1000;
  return basepay + overtimepay;
}
else
  return (millihours * hourlycents + 500) / 1000;
\end{verbatim}

ここでは、労働時間が残業代の閾値（40 時間）を超えているかを確認し、超えていなければ前と同じ計算を行う。
残業がある場合は、まず 40.000 時間を超えた時間（1000 分の 1 単位）を \texttt{overage\_millihours} として求める。
その時間の分については、計算した給与に \texttt{overtime\_percentage}（150）を掛ける。
パーセンテージ（10 進数で小数 2 桁）と 1000 分の 1 時間（小数 3 桁）を扱っているので、丸めのためにその半分を足してから 100,000（0 が 5 つ）で割らなければならない。
最初の 40 時間に対する、割増なしの基本給を計算した後、コードは両者を合計して支払総額を出す。
効率を考えればこれらの似た計算をまとめることもできるが、ここではわかりやすさのために、コードの構造を計算の構造に合わせている。

このコードはたいてい動くが、いつも動くわけではない。
おかしな結果の一例として、時給 \$50.00 で 60.000 時間働くと、給与は \$2,211.51 になる（本来は \$3,500.00 になるはずである）。
問題は \texttt{overtime\_percentage}（150）を掛ける掛け算にある。
それなりの時給で残業時間がそこそこあれば、この掛け算は簡単にオーバーフローする。
整数演算では 150/100 を分数として前もって計算しておくことはできない（整数では単に 1 になってしまう）ので、掛け算を先に行うしかない。

このコードを直すには、\texttt{(X*150)/100} を \texttt{(X*3)/2} に置き換える手もある。
しかし、それでは残業のパーセンテージをパラメータにした意味がなくなり、割合がもっと扱いにくい値に変わったときには動かない。
パラメータ化を保ったまま解決する方法の一つは、計算を分け、掛け算と割り算には 64 ビットの算術を使って、結果を 32 ビットにダウンキャストすることである。

\begin{verbatim}
if (millihours > max_millihours      // 100 hours max
   || hourlycents > max_hourlycents) // $200/hour rate max
  return 0;
if (millihours > overtime_millihours) {
  overage_millihours = millihours - overtime_millihours;
  product64 = overage_millihours * hourlycents;
  adjusted64 = (product64 * overtime_percentage + 50000) / 100000;
  overtimepay = ((uint32_t)adjusted64 + 500) / 1000;
  return basepay + overtimepay;
}
else
  return (millihours * hourlycents + 500) / 1000;
\end{verbatim}

説明のために、64 ビットの変数は名前でそれとわかるようにしてある。
明示的なキャストを多用してこれらの式を書くこともできるが、長くなって読みにくい。

三つの値の掛け算は、オーバーフローが起こる前に、まず二つを掛けて 64 ビットの変数に入れるように分けてある。
いったんアップキャストすれば、パーセンテージとの掛け算は 64 ビットで行われ、正しく動く。
出来上がったコードは確かにごちゃごちゃしているので、こう書いた理由を説明するコメントがあるとよいだろう。
最もすっきりした解決策は、わずかな効率の低下と引き換えに、関係する変数をすべて 64 ビットに格上げすることである。
固定幅整数で計算する場合には、こうしたトレードオフがつきものである。
% 注: 原文の修正版コードでは basepay の計算行が省かれている（原文のまま転記した）。

%----------------------------------------------------------------------
\subsubsection{安全な算術}\label{subsec:ch9-safe-arithmetic}
% src: Safe Arithmetic

整数オーバーフローは、浮動小数点のアンダーフローより問題になることが多い。
結果が劇的に変わりうるからである。
しかし、浮動小数点のアンダーフローも決して無視してよいわけではない。
コンパイラはもともと数学的な正しさから外れうるやり方で算術を行うので、その結果に対処する責任は開発者にある。
これらの問題を意識していれば、脆弱性を避けるのに役立つ緩和戦略をいくつか取り入れられる。

オーバーフローの可能性に対処するために小細工を弄したコードを書くのは避ける。
誤りがあってもテストで見つけにくく、悪用可能な脆弱性になりうるからである。
さらに、小細工は自分のマシンでは動いても、ほかの CPU アーキテクチャや別のコンパイラにはそのまま持っていけないかもしれない。
こうした計算を安全に行う方法をまとめると、次のようになる。

\begin{itemize}
\item 型変換も、計算と同じように、結果を切り捨てたり歪めたりすることがある。
\item できれば計算への入力を制限して、ありうる値がすべて表現できることを保証する。
\item オーバーフローを避けるために、より大きな固定サイズの整数を使う。より小さいサイズの整数に戻す前には、結果が範囲内にあることを確認する。
\item 最終的な結果が常に範囲内に収まっていても、計算の途中の値がオーバーフローして問題を起こしうることを忘れない。
\item セキュリティ上重要なコードの中やその周りでは、算術の正しさを確かめるときにとくに注意する。
\end{itemize}

固定幅整数や浮動小数点の計算の微妙な点がまだ難しく感じられるなら、それらに注意深く目を光らせ、初歩的に見える計算でも驚くような結果になりうると考えておくこと。
こうした計算が厄介になりうると知ったうえで、好きな言語で手早くコードを書いて少しテストしてみると、コンピューターの算術の基本的な構成要素の限界を感覚としてつかめる。

この種のバグのリスクがあるコードを見つけたら、すべての入力に極端な値を与えて計算するテストケースを作り、結果を確認する。
テストケースをうまく選べばオーバーフローの問題を見つけられるが、限られた数のテストでは、コードがオーバーフローの影響を受けないことの証明にはならない。

幸いなことに、Python のようなもっと現代的な言語では任意精度の整数を使うことが増えており、一般にこうした問題は起こらない。
算術計算を正しく行うには、まず自分が使う言語がどう動くのかを細部まで正確に理解することである。
いくつかの人気のある言語の詳細を含む優れた参考資料が、覚えやすい URL \url{floating-point-gui.de} にあり、詳しい説明とベストプラクティスのコーディング例が載っている。

%----------------------------------------------------------------------
\subsection{メモリアクセスの脆弱性}
% src: Memory Access Vulnerabilities

ここで論じるもう一つの脆弱性のクラスは、不適切なメモリアクセスにかかわるものである。
メモリを直接管理するのは強力で、非常に効率的になりうるが、コードが何かを間違えれば、どこまでもひどい結果を招くリスクがある。

ほとんどのプログラミング言語は、メモリの割り当てを完全に管理し、アクセスを適切な境界内に制限している。
しかし、効率や柔軟性のため、あるいは時には過去からの惰性で、メモリ管理の仕事をプログラマーに任せている言語もある（主に C と C++）。
プログラマーがこの仕事を引き受けると、経験豊富なプログラマーであっても、とくにコードが複雑になるにつれて簡単に間違え、深刻な脆弱性を生んでしまう。
そして先に述べた算術の欠陥と同じく、大きな危険は、メモリ管理の約束事への違反が検出されず、黙って起こり続けることである。

この節では、組み込みの安全装置なしにメモリを直接管理し、アクセスするコードのセキュリティ面を取り上げる。
コード例ではオリジナルの C 標準ライブラリにある古典的な動的メモリ関数を使うが、ここでの教訓は、同様の機能を持つ多くの変種にも一般に当てはまる。

%----------------------------------------------------------------------
\subsubsection{メモリ管理}
% src: Memory Management

ポインタはアドレスを使ってメモリに直接アクセスできるようにするもので、おそらく C 言語で最も強力な機能である。
しかし、どんな電動工具を扱うときもそうであるように、リスクを管理するために責任ある安全対策をとることが重要である。
ソフトウェアは必要なときにメモリを割り当て、使ってよい境界の中で作業し、要らなくなったら解放する。
この空間と時間の取り決めから外れたアクセスはすべて意図しない結果をもたらし、そこに脆弱性が生まれる。

C 標準ライブラリには、大きなデータ構造のためや、データ構造のサイズがコンパイル時に決まらない場合のために、動的メモリ割り当ての機能がある。

\begin{dfn}[ヒープ]\label{dfn:ch9-heap}
動的に割り当てられるメモリの供給元で、作業用メモリを提供するために使われる、プロセス内の大きなアドレス空間の塊を \underline{ヒープ}（heap）という。
\end{dfn}

C のプログラムは、メモリを割り当てるのに \texttt{malloc(3)} を使い、要らなくなったら \texttt{free(3)} を呼び出して、割り当てたものをそれぞれ再利用のために解放する。
これらの割り当て・解放関数には多くの変種がある。
ここでは簡単のためにこの二つに絞るが、考え方はコードがメモリを直接管理するあらゆる場面に当てはまるはずである。

解放後のメモリへのアクセスは簡単に起こる。多くのコードが、最後には解放されるデータ構造を共有しているのに、ポインタのコピーが残っていて誤って使われる場合である。
メモリが再利用された後に、そうした古いポインタを使えば、どんな使い方であれメモリアクセスの完全性を破ることになる。
反対に、使い終わったメモリを解放し忘れると、時間とともにヒープを使い果たし、メモリ不足に陥るリスクがある。
次のコードの抜粋は、ヒープメモリの基本的な正しい使い方を示している。

\begin{verbatim}
uint8_t *p;
// Don't use the pointer before allocating memory for it.
p = malloc(100); // Allocate 100 bytes before first use.
p[0] = 1;
p[99] = 123 + p[0];
free(p);          // Release the memory after last use.
// Don't use the pointer anymore.
\end{verbatim}

このコードは、割り当てと解放の呼び出しの間で、割り当てたメモリの境界内だけにアクセスしている。

実際には、割り当て、メモリアクセス、解放がコードのあちこちに散らばっていることがあり、いつでも寸分違わず正しく行うのは難しい。

%----------------------------------------------------------------------
\subsubsection{バッファオーバーフロー}\label{subsec:ch9-buffer-overflow}
% src: Buffer Overflow

\begin{dfn}[バッファオーバーフロー]\label{dfn:ch9-buffer-overflow}
コードが、意図した対象のバッファの外側にあるメモリ位置にアクセスすることを、\underline{バッファオーバーフロー}（buffer overflow）またはバッファオーバーラン（buffer overrun）という。
ここで \underline{バッファ}（buffer）は、メモリ上のあらゆるデータ（データ構造、文字列、配列、オブジェクト、任意の型の変数）を指す一般的な用語であり、\underline{アクセス}（access）はメモリの読み取りと書き込みの両方をまとめて指す用語である。
\end{dfn}

用語がまぎらわしいので、意味をはっきりさせておくことが大切である。
定義 \ref{dfn:ch9-buffer-overflow} のとおり、バッファオーバーフローには、意図したメモリ領域の外側での読み取り\emph{または}書き込みが含まれる。
「overflow（あふれる）」という語からは、書き込みのほうが自然に思い浮かぶにもかかわらず、である。
読み取りと書き込みでは影響が根本的に異なるが、問題を理解するには両者をまとめて考えるのが役に立つ。

バッファオーバーフローはヒープメモリに限らず、静的に割り当てた変数やスタック上のローカル変数を含め、あらゆる種類の変数で起こりうる。
どの場合も、メモリ上のほかのデータを任意に書き換えてしまう可能性がある。
意図しない範囲外への書き込みはメモリ上のほとんど何でも書き換えうるので、賢い攻撃者は、損害が最大になるようにこうした攻撃に磨きをかける。
さらに、バッファオーバーフローのバグは予期しない形でメモリを読み取ることがあり、攻撃者に情報を漏らしたり、コードを誤動作させたりするおそれがある。

明示的なメモリ割り当て、境界内でのアクセス、使われなくなった記憶領域の解放を寸分違わず正しく行うことの難しさと重要性を、甘く見てはならない。
割り当て、使用、解放の単純なパターンが一番よく、解放が決して飛ばされないことを保証する例外処理も含めるべきである。
あるコンポーネントが割り当てたメモリの参照を別のコードに引き渡す場合には、その後のメモリ解放の責任をインターフェースのどちらか一方の側に決めておくことが欠かせない。

\begin{rem}[範囲チェックのある言語でも起こりうる]\label{rem:ch9-managed-languages}
完全に範囲チェックされ、ガベージコレクションのある言語であっても、問題に陥ることはありうると認識しておく。
メモリ上のデータ構造を直接操作するコードはどれも、バッファオーバーフローと同じような誤りを犯しうる。
例えば、Python のバイト配列に入った TCP/IP パケットのようなバイナリのデータ構造を操作する場合を考える。
内容を読み取って書き換えるにはデータへのオフセットを計算する必要があり、配列の外側へのアクセスが起こらなくても、そこにバグが入り込みうる。
\end{rem}

%----------------------------------------------------------------------
\subsubsection{例：メモリ割り当ての脆弱性}\label{subsec:ch9-example-memory-allocation}
% src: Example: Memory Allocation Vulnerabilities

動的メモリ割り当てを誤ったときの危険を示す例を見てみよう。
この例はわかりやすくしてあるが、実際のアプリケーションでは、重要なコードの部分どうしが離れていることが多く、こうした欠陥ははるかに見つけにくい。

% src: A Simple Data Structure
\underline{単純なデータ構造：}
この例では、ユーザーアカウントを表す単純な C のデータ構造を使う。
この構造体は、ユーザーが管理者であれば立てるフラグ、ユーザー ID、ユーザー名、設定の集まりからなる。
\texttt{isAdmin} フィールドが 0 でない場合に無制限の認可を与える（そのため、このフィールドは攻撃者にとって魅力的な標的になる）という点を除けば、各フィールドの意味はここでは重要ではない。

\begin{verbatim}
#define MAX_USERNAME_LEN 39
#define SETTINGS_COUNT 10
typedef struct {
  bool isAdmin;
  long userid;
  char username[MAX_USERNAME_LEN + 1];
  long setting[SETTINGS_COUNT];
} user_account;
\end{verbatim}

次は、こうしたユーザーアカウントのレコードを作成する関数である。

\begin{verbatim}
user_account* create_user_account(bool isAdmin,
                                  const char* username) {
  user_account* ua;
  if (strlen(username) > MAX_USERNAME_LEN)
    return NULL;
  ua = malloc(sizeof (user_account));
  if (NULL == ua) {
    fprintf(stderr, "malloc failed to allocate memory.");
    return NULL;
  }
  ua->isAdmin = isAdmin;
  ua->userid = userid_next++;
  strcpy(ua->username, username);
  memset(&ua->setting, 0, sizeof ua->setting);
  return ua;
}
\end{verbatim}

1 つ目の引数は、ユーザーが管理者かどうかを指定する。
2 つ目の引数はユーザー名で、指定された最大長を超えてはならない。
一意な ID は、グローバルなカウンタ（\texttt{userid\_next}、宣言は示していない）から連番で振られる。
設定の値はすべて最初に 0 になり、エラーで \texttt{NULL} を返す場合を除いて、コードは新しいレコードへのポインタを返す。
コードは割り当ての\emph{前に}ユーザー名の文字列の長さを確認しているので、実際にメモリを使う場合にだけ割り当てが行われる点に注意してほしい。

% src: Writing an Indexed Field
\underline{インデックス付きフィールドへの書き込み：}
レコードを作成した後は、次の関数でどの設定の値でも設定できる。

\begin{verbatim}
bool update_setting(user_account* ua,
                    const char *index, const char *value) {
  char *endptr;
  long i, v;
  i = strtol(index, &endptr, 10);
  if (*endptr)
    return false; // Terminated other than at end of string.
  if (i >= SETTINGS_COUNT)
    return false;
  v = strtol(value, &endptr, 10);
  if (*endptr)
    return false; // Terminated other than at end of string.
  ua->setting[i] = v;
  return true;
}
\end{verbatim}

この関数は、設定のインデックスと値を 10 進数の文字列として受け取る。
それぞれを整数に変換したうえで、値をレコードの、インデックスで指定された設定に格納する。
例えば、設定 1 に値 14 を代入するには、関数を \verb|update_setting(ua, "1", "14")| のように呼び出す。

関数 \texttt{strtol} は文字列を整数値に変換する。
\texttt{strtol} が設定するポインタ（\texttt{endptr}）によって、どこまで解析したかが呼び出し側にわかる。
このポインタが NULL 終端文字を指していなければ、文字列は有効な整数ではなかったということなので、コードはエラーを返す。
インデックス（\texttt{i}）が設定の数を超えていないことを確認した後、値（\texttt{v}）を同じ方法で解析し、設定の値をレコードに格納する。

% src: Buffer Overflow Vulnerability
\underline{バッファオーバーフローの脆弱性：}
C は冗長になりがちだが、ここまでの準備はいたって単純である。
では本題に入ろう。
このコードにはバグがある。
\emph{負の}インデックス値をチェックしていないのである。
攻撃者がこの関数を \verb|update_setting(ua, "-12", "1")| のように呼び出させることができれば、管理者になれてしまう。
各要素は 4 バイトの \texttt{long} 型なので、\texttt{setting} への代入はレコードの中を 48 バイトさかのぼった位置にアクセスするからである。
したがって、この代入は値 1 を \texttt{isAdmin} フィールドに書き込み、過剰な権限を与えてしまう。

\begin{exa}[負のインデックスによる権限昇格]\label{exa:ch9-negative-index}
この例では、データ構造で負のインデックスを許したために、メモリへの認可されていない書き込みが起こり、セキュリティの保護機構が破られた。
同じテーマの変種は数多くあり、注意が必要である。
上限チェックが抜けていることによるインデックスの誤りや、オーバーフローのような算術の誤りなどである。
あるデータ構造からの不正なアクセスが、たまたま悪い位置にあったほかのデータを書き換えてしまうこともある。
\end{exa}

修正するには、負のインデックス値を受け付けないようにして、書き込みアクセスを有効な設定の範囲に限ればよい。
\texttt{if} 文に次のように条件を追加すると \texttt{i} の負の値が拒否され、抜け穴がふさがる。

\begin{verbatim}
  if (i < 0 || i >= SETTINGS_COUNT)
\end{verbatim}

\texttt{i < 0} の条件を追加したことで、負のインデックス値はすべて拒否されるようになり、この関数による意図しない書き換えを防げる。

% src: Leaking Memory
\underline{メモリの漏えい：}
負のインデックスによる上書きの欠陥を直しても、まだ脆弱性が残っている。
\texttt{malloc(3)} のドキュメントは、下線付きで「The memory is not initialized.（メモリは初期化されない）」と警告している。
つまり、メモリには何が入っていてもおかしくなく、少し実験してみると、実際に以前のデータが残っていることがわかる。
したがって、初期化されていないメモリを再利用すると、プライベートなデータが漏えいするおそれがある。

\texttt{create\_user\_account} 関数は構造体のすべてのフィールドにデータを書き込んでいるが、それでも、再利用されたメモリとしてデータ構造の中に残っているバイトを漏らしてしまう。
コンパイラはふつう、効率よく書き込めるようにフィールドのオフセットを整列（アライン）させる。
原著者の 32 ビットコンピューターでは、フィールドのオフセットは 4 の倍数になる（8 ビットの 4 バイトで 32 ビット）。
ほかのアーキテクチャでも同じような整列が行われる。
整列が必要なのは、4 の倍数のアドレスをまたぐフィールドに書き込む（例えば、アドレス 0x1000002 に 4 バイトを書き込む）には、メモリに 2 回アクセスしなければならないからである。
そのため、この例では、オフセット 0 にある 1 バイトのブール型フィールド \texttt{isAdmin} の次に、\texttt{userid} フィールドがオフセット 4 から始まり、間の 3 バイト（オフセット 1--3）は使われないまま残る。
% 図 9-3: Data structure layout diagram（user_account 構造体のメモリレイアウト）
% 注: 原文の図は画像のみ。本文から読み取れるのは isAdmin（オフセット 0）、未使用（1--3）、userid（オフセット 4）までで、
%   図の細部（例のユーザー名の具体的な文字列など）は本文からは決まらないため描かなかった。

さらに、ユーザー名に \texttt{strcpy} を使ったことで、もう一つのメモリの塊が初期化されないまま残る。
この文字列コピー関数は NULL 終端文字でコピーを止めるので、この例の 5 バイトの文字列では先頭の 6 バイトしか書き換えられず、\texttt{malloc} がたまたま確保したときに入っていた何かが 34 バイト残される。
要するに、新しく割り当てた構造体には以前のデータが残っており、すべてのバイトを上書きしない限り、それが漏えいしうるということである。

\begin{rem}[割り当てたバッファはゼロで埋める]\label{rem:ch9-zero-buffers}
こうした不注意によるメモリの漏えいのリスクを緩和するのは難しくないが、露出しうるデータ構造のすべてのバイトを根気よく上書きしなければならない。
コンパイラがフィールドのオフセットをどう割り当てるかを正確に予想しようとしてはならない。
オフセットの割り当ては、時とともに、またプラットフォームによって変わりうるからである。
こうした問題を避ける最も簡単な方法は、すべてのバイトが書き込まれることをほかの方法で保証できない限り、割り当てた直後にバッファをゼロで埋めることである。
自分のコードが機密データそのものを扱っていなくても、このメモリ漏えいの経路を通じて、プロセス内のどこにあるデータでも露出しうることを忘れてはならない。
\end{rem}

一般に、文字列のコピーに \texttt{strcpy} を使うのは避けるべきである。
間違える余地があまりにも多いからである。
\texttt{strncpy} 関数は、コピー先の使われないバイトをゼロで埋め、バッファサイズを超える文字列によるオーバーフローも防ぐ。
しかし \texttt{strncpy} は、結果の文字列が NULL 終端文字で終わることを保証しない。
だからこそ、バッファを \texttt{MAX\_USERNAME\_LEN + 1} のサイズで割り当て、NULL 終端文字の入る場所が常にあるようにすることが欠かせないのである。
もう一つの選択肢は \texttt{strlcpy} 関数で、こちらは NULL 終端を保証するが、効率のために使われないバイトをゼロで埋めることはしない。
この例からわかるように、メモリを直接扱う場合には、注意深く対処しなければならない要因が数多くある。

メモリ割り当ての仕組みを取り上げ、作った例で脆弱性がどんなものかを見たので、次はもっと現実的な事例を考えよう。
次の例は、数年前に世界の主要な Web サービスのかなりの部分を危険にさらした、驚くべきセキュリティ上の大失態に基づいている。

%----------------------------------------------------------------------
\subsubsection{ケーススタディ：Heartbleed}\label{subsec:ch9-heartbleed}
% src: Case Study: Heartbleed

2014 年 4 月初め、新たに見つかったセキュリティ上の欠陥の詳細が公になるにつれ、各紙は、主要な OS プラットフォームや Web サイトが、露出を最小限に抑えようと秘密裏に急いで準備した修正を足並みをそろえて展開し、世界規模の災害をかろうじて回避したと報じた。
Heartbleed がニュースになったのは、「かっこいいロゴを持つ最初のセキュリティバグ」だったからだけではない。
人気のある OpenSSL の TLS ライブラリを導入したあらゆるサーバーの鎧に、簡単に悪用できる穴があることを明らかにしたからである。

以下では、この 10 年で最も恐ろしいセキュリティ脆弱性の一つを詳しく見ていく。
そうすれば、誤りがどれほど深刻になりうるかを実感できるはずである。
この詳しい議論の目的は、動的に割り当てたメモリの管理にかかわるバグが、どのように壊滅的な脆弱性になりうるかを示すことにある。
そのため、脆弱性の核心が見えるように、コードと、複雑な TLS 通信プロトコルの細部の一部を単純化している。
概念的には実際に起きたことにそのまま対応しているが、動く部品はもっと少なく、コードもずっと単純である。

\begin{dfn}[Heartbleed]\label{dfn:ch9-heartbleed}
\underline{Heartbleed} は、2012 年に RFC 6520 で提案された TLS Heartbeat 拡張の OpenSSL による実装に含まれていた欠陥である。
\end{dfn}

この拡張は、TLS 接続を維持するためのオーバーヘッドの小さい方法で、クライアントがしばらく通信しなかった後でも新しい接続を確立し直さずに済むようにする。
いわゆるハートビートそのものは、メッセージを往復させるやり取りである。
16 バイトから 16,384（$2^{14}$）バイトまでの任意のデータをペイロードとして持つハートビート要求を送ると、同じペイロードを含むハートビート応答がこだまのように返ってくる。
% 図 9-4: Heartbeat protocol diagram（プロトコルの基本的な要求メッセージと応答メッセージ）
% 注: 原文の図は画像のみ。図の細部は本文から決まらないため描かなかった。

HTTPS の Web ページをダウンロードした後、クライアントは、その接続を維持したいことをサーバーに知らせるために、後でその接続上でハートビート要求を送ることがある。
ふつうの使い方では、クライアントは要求として 16 バイトのメッセージ「Hello Heartbeat!」を送り、サーバーは同じ 16 バイトを送り返して応答する（少なくとも、そう動くことになっている）。
では、Heartbleed のバグを見てみよう。

致命的な欠陥が表に出るのは、ペイロードは小さいのに、それより大きなペイロードのバイト数を申告する、不正な形式のハートビート要求を受け取ったときである。
その仕組みを正確に見るために、まず、ピア同士がやり取りするハートビートメッセージの内部構造を単純化したものを見る。
この例のコードはすべて C である。

\begin{verbatim}
typedef struct {
  HeartbeatMessageType type;
  uint16_t payload_length;
  char bytes[0]; // Variable-length payload & padding
} hbmessage;
\end{verbatim}

データ構造の宣言 \texttt{hbmessage} は、ハートビートメッセージの三つの部分を示している。
1 つ目のフィールドはメッセージの種類で、要求か応答かを示す。
次は、\texttt{payload\_length} という、メッセージのペイロードのバイト長である。
3 つ目のフィールド \texttt{bytes} は長さ 0 として宣言されているが、必要なサイズを足して動的に割り当てて使うことを想定している。

悪意あるクライアントは、まず標的のサーバーと TLS 接続を確立し、次にバイト数を 16,000 と申告した 16 バイトのハートビート要求を送ることで、サーバーを攻撃できる。
C の宣言として書くと次のようになる。

\begin{verbatim}
typedef struct {
  HeartbeatMessageType type = heartbeat_request;
  uint16_t payload_length = 16000;
  char bytes[16] = {"Hello Heartbeat!"};
} hbmessage;
\end{verbatim}

これを送るクライアントは嘘をついている。
メッセージはペイロードの長さを 16,000 バイトと申告しているが、実際のペイロードは 16 バイトしかない。
このメッセージがどのようにサーバーをだますのかを理解するために、受信したハートビート要求メッセージを処理する C のコードを見てみる。

\begin{verbatim}
hbmessage *hb(hbmessage *request, int *message_length) {
  int response_length = request->payload_length+sizeof(hbmessage);
  hbmessage* response = malloc(response_length);
  response->type = heartbeat_response;
  response->payload_length = request->payload_length;
  memcpy(&response->bytes, &request->bytes,
         response->payload_length);
  *message_length = response_length;
  return response;
}
\end{verbatim}

関数 \texttt{hb} は二つの引数で呼び出される。
受信したハートビート要求メッセージと、関数が返す応答メッセージの長さを格納する \texttt{message\_length} という名前のポインタである。
最初の 2 行で応答のバイト長を \texttt{response\_length} として計算し、そのサイズのメモリブロックを \texttt{response} として割り当てる。
次の 2 行で、応答メッセージの最初の二つの値、つまりメッセージの種類とその \texttt{payload\_length} を埋める。

次が致命的なバグである。
サーバーは要求で受け取ったメッセージのバイトを送り返さなければならないので、要求から応答へデータをコピーする。
関数は要求メッセージが長さを正確に申告していると\emph{信頼している}ので、16,000 バイトをコピーする。
しかし要求メッセージには 16 バイトしかないので、応答には何千バイトもの内部メモリの内容が入ってしまう。
最後の 2 行で応答メッセージの長さを格納し、そのポインタを返す。

\begin{exa}[Heartbleed によるプロセスメモリの漏えい]\label{exa:ch9-heartbleed-exploit}
原著の図 9-5 は、このメッセージのやり取りを図示し、先のコードがプロセスメモリの内容をどのように漏らすかを詳しく示している。
エクスプロイトの害を具体的に示すために、図には、秘密のデータを含むバッファがいくつか、要求バッファの近くのメモリにすでに置かれている様子が描かれている。
16 バイトのペイロードしかないバッファから 16,000 バイトをコピーすると（図では大きすぎる点線の領域で示されている）、秘密のデータが応答メッセージに入り込み、サーバーはそれをクライアントに送ってしまう。
% 図 9-5: Heartbleed example exploiting Heartbeat protocol
% 注: 原文の図は画像のみ。図の細部は本文から決まらないため描かず、本文の説明だけを残した。
\end{exa}

この欠陥は、作業用メモリの何千バイトものスナップショットを撮って、呼び出してきた相手に誰彼かまわず送り出す匿名の API をサーバーに用意したのと同じことである。
つまり、インターネットに露出した、メモリ分離の完全な破綻である。
HTTPS のセキュリティを使う Web サーバーが、作業用メモリにおいしい秘密をいくらでも抱えているのは驚くに当たらない。
Heartbleed のバグの発見者によれば、彼らは自分たち自身から「X.509 証明書に使われている秘密鍵、ユーザー名とパスワード、インスタントメッセージ、電子メール、ビジネス上重要な文書や通信」を簡単に盗み出せた。
どのデータが漏えいするかはメモリ割り当ての気まぐれ次第だったが、この脆弱性を悪用する攻撃者は何度でもサーバーのメモリにアクセスできたので、いずれはあらゆる種類の機密データが手に入った。

修正は、後から見れば単純であった。
自分が渡すよりも多くのペイロードを求める「嘘つきの」ハートビート要求が来ることを想定し、RFC が明示的に定めているとおり、それを無視すればよい。
Heartbleed によって世界は、どれほど多くのサーバーが OpenSSL に頼っているか、そしてインターネットのインフラの大部分が頼っている重要なソフトウェアに取り組むボランティアがいかに少ないかを知った。
このバグは、多くのセキュリティ上の欠陥がなぜ見つけにくいのかを示す典型例である。
正しい形式の要求ならすべて完璧に動き、善意のコードならまず送ることのない不正な形式の要求でだけ問題が起こるからである。
さらに、ハートビート応答で漏れたサーバーのメモリは、サーバーそのものに直接の害を与えない。
過剰に開示されたデータを注意深く分析して初めて、潜在的な損害の大きさが見えてくる。

Heartbleed は、ここ数年で見つかった中でもまず間違いなく最も深刻なセキュリティ脆弱性の一つであり、セキュリティバグの性質と、小さな欠陥がシステムのセキュリティを大規模に損ないうることを示す貴重な例として心に留めておくべきである。
機能の観点からは、これは些細なバグだと言い張るのは簡単である。
起こりそうになく、要求が渡したよりも多くのペイロードデータを送り返すことは、一見まったく無害に思えるからである。
% 図: xkcd HeartBleed explanation（https://xkcd.com/1354/）
% 注: 原文では xkcd の漫画が引用されているが、画像のため収録しなかった。

\begin{rem}[低レベル言語の壊れやすさ]\label{rem:ch9-heartbleed-lesson}
Heartbleed は、低レベル言語の壊れやすさを示す格好の教訓である。
小さな誤りが巨大な影響を及ぼしうる。
バッファオーバーフローは、価値の高い秘密がたまたま運悪くちょうどその位置のメモリにあれば、それを露出させてしまう。
設計（プロトコル仕様）は、バイト長が正しくないハートビート要求は無視すべきだと定めることで、まさにこの誤りを予期していた。
しかし明示的なテストがなかったため、2 年以上にわたって誰もこの脆弱性に気付かなかった。
\end{rem}

これは一つのライブラリの一つのバグにすぎない。
同じようなバグが、今もどれだけ残っているだろうか。

\newpage

%% =======================================================
%%  part: ch10
%% =======================================================
%% =======================================================
%%  part: ch10
%% =======================================================
\section{信頼できない入力}
% src: 10: Untrusted Input

この章では、\cite{Kohnfelder2021} の第10章に沿って、信頼できない入力とその扱い方を扱う。

\begin{quote}
「私はエンジニアリングが好きだが、創造的なインプットのほうがもっと好きだ。」---John Dykstra
\end{quote}

\begin{dfn}[信頼できない入力]\label{dfn:ch10-untrusted-input}
\underline{信頼できない入力}（untrusted input）は、安全なコードを書こうとする開発者にとって、おそらく最大の悩みの種である。
この用語自体がまぎらわしいこともあるので、システムへの入力のうち\emph{信頼できる入力}以外のものすべて、と理解するのがよいだろう。
ここで信頼できる入力とは、正しいデータを渡してくると信頼できるコードからの入力のことである。
信頼できない入力とは、自分の管理が及ばず、操作されているかもしれない入力のことであり、システムに入ってくるデータのうち完全には信頼していないものはすべて含まれる。
つまり「信頼すべきでない入力」のことであって、「誤って信頼してしまっている入力」のことではない。
\end{dfn}

外部から入ってくるデータは、すべて信頼できないものとみなすのが最善である。
システムのユーザーは親切で信頼できる人たちかもしれないが、セキュリティの観点では信頼できないものとみなしておくのがよい。
ユーザーは何でもできてしまうし、他人の策略にはまることもあるからである。
信頼できない入力が厄介なのは、それが\underline{攻撃ベクトル}（attack vector）、つまりシステムの内部に手を伸ばして問題を起こす手段になるからである。
とくに注意すべきなのは、悪意をもって作り込まれ、信頼境界を越えてくる入力である。
こうした入力はシステムの奥深くまで入り込み、特権を持つコードでエクスプロイトを引き起こしうるので、しっかりした第一の防衛線を用意しておくことが欠かせない。
世界最大の信頼できない入力の発生源はインターネットにほかならず、ネットワークから完全に切り離されたソフトウェアはきわめてまれなので、これはほぼすべてのシステムにとって深刻な脅威である。

\begin{dfn}[入力検証]\label{dfn:ch10-input-validation}
\underline{入力検証}（input validation）とは、入力に制約を課し、決められた規則に従わせる防御的なコーディングである。
\end{dfn}

入力が特定の制約を満たしていることを検証し、有効な入力であればどれに対してもコードが正しく動くようにしておけば、こうした攻撃をうまく防げる。
この章では、入力検証によって信頼できない入力を管理する方法と、それがセキュリティにとってなぜ重要なのかを中心に扱う。
地味な話題に見えるかもしれず、技術的にも難しくはない。しかし入力検証が必要になる場面はきわめてありふれているので、入力検証をもっとうまく行うことは、開発者が脆弱性を減らすためにすぐ手を付けられる、おそらく最も効果の大きい「手の届く果実」（low-hanging fruit）である。
習得する価値が十分にあるので、この章で詳しく扱う。
文字列の入力には特有の難しさがあり、Unicode がセキュリティに及ぼす影響はあまり知られていないので、それにまつわる基本的な問題も概観する。
続いて、信頼できないデータを使ったインジェクション攻撃の例を、さまざまな技術（SQL、パストラバーサル、正規表現、XML 外部実体（XXE））について見ていく。
最後に、この広範な脆弱性群に対して使える緩和の手法をまとめる。

%----------------------------------------------------------------------
\subsection{入力検証}\label{subsec:ch10-input-validation}
% src: Input Validation

\begin{quote}
「他人に承認（validation）を求める前に、まず自分の中に見つけようとしなさい。」---Greg Behrendt
\end{quote}

信頼できない入力が何かがわかったので、次は、信頼できない入力がシステムの中で及ぼしうる影響と、その害から身を守る方法を考えよう。
信頼できない入力はいつもシステムの中を流れていて、何層も下にある信頼されたコンポーネントにまで届くことも多い。
したがって、自分のコードが信頼されたコードから直接呼び出されているからといって、渡される入力が信頼できるとは限らない。
問題は、コンポーネントがどこから来たかわからないデータをそのまま受け渡しているかもしれないことである。
攻撃者がデータを操作できる経路が多いほど、そのデータは信頼できなくなる。
この点は、このあとの例ではっきりするだろう。

入力検証は、信頼できない入力を、アプリケーションが安全に処理できる範囲の値にまで絞り込むので、優れた防御策である。
入力検証の本質的な役割は、信頼できない入力が設計仕様に合っていることを保証し、検証より下流のコードが整形式（well-formed）のデータだけを扱えば済むようにすることである。

\begin{exa}[ログイン認証サービスのユーザー名の制限]\label{exa:ch10-username-restriction}
ユーザー名とパスワードを受け取り、クレデンシャルが正しければ認証トークンを発行するユーザーログイン認証サービスを書いているとする。
ユーザー名を 8 文字以上 40 文字以下とし、明確に定めた Unicode コードポイントの部分集合だけで構成するよう求めれば、入力の中身があらかじめわかっている状態になるので、入力の扱いがずっと単純になる。
後続のコードはユーザー名のコピーを固定サイズのバッファに保持できるし、見慣れない文字がどんな影響を及ぼすかを心配する必要もない。
この保証を前提にして、ほかの面でも処理を単純にできるだろう。
\end{exa}

入力検証で低レベルの脆弱性を修正する例は、前章（第9章）ですでに見た。
給与の整数計算のコードには、大きすぎる入力値を防ぐために、if 文 1 つだけから成る入力検証があった。
\begin{verbatim}
if (millihours > max_millihours       // 100 hours max
    || hourlycents > max_hourlycents) // $200/hour rate
return 0;
\end{verbatim}
説明を繰り返す必要はないが、基本的な入力検証のよい例である。
書くコードのほとんどは、何らかの制限の範囲内でしか正しく動作しない。
巨大なメモリサイズのような極端な値や、場合によっては別の言語のテキストに対しては動作しない。
制限が何であれ、設計で想定していない入力にコードをさらすべきではない。
意図しない結果を招き、脆弱性を生むおそれがあるからである。
この危険を緩和する簡単な方法の一つは、問題を起こしうる入力をすべて弾くような人為的な制限を入力に課すことである。

ただし、注意しておきたい細かな点がいくつかある。
まず当然ながら、本来正しく処理されるべき入力を制限によって拒否してはならない。
例えば給与の例では、週 40 時間の労働を無効として拒否するわけにはいかない。
有効な入力の一部をコードが処理できないなら、もっと広い範囲の入力を処理できるようにコードを直す必要がある。
また、入力検証の戦略では、複数の入力の組み合わせを考慮しなければならない場合もある。
給与の例では、第9章で見たように、時給と労働時間の積が固定幅整数に収まらなくなることがある。そこで検証では、この 2 つの入力の積に上限を設けるか、それぞれに別々の上限を設けることになる。
前者のほうが許容範囲は広いが、呼び出し側が合わせにくい場合もあるので、どちらが正しいかはアプリケーション次第である。

一般に、信頼できない入力はできるだけ早い段階で検証すべきである。
そうすれば、制約のない入力が、それをうまく扱えないかもしれない下流のコードへ流れ込む危険を最小にできる。
いったん検証してしまえば、後続のコードは行儀のよいデータだけを相手にすればよい。
開発者は入力の範囲を正確に把握できるので、安全なコードを書きやすくなる。
大事なのは一貫性である。入ってくるデータを最初に扱う層のコードに入力検証を置き、検証済みの入力を、すべての入力が有効だと安心して仮定できる深い層のビジネスロジックへ渡す、というのがよいパターンである。

入力検証は主に、信頼できない入力、とくに攻撃対象領域にある入力に対する防御と考えられているが、だからといってそれ以外の入力をすべて気楽に無視してよいわけではない。
あるデータの提供元をどれほど信頼していても、ミスによって予期しない入力が生じることはありうるし、攻撃によってシステムの一部が侵害され、実質的に攻撃対象領域が広がってしまうこともありうる。
こうした理由から、防御的な入力検証は頼れる味方である。
微妙な脆弱性を生む危険を冒すくらいなら、冗長なチェックをしておくほうが安全である。
入ってくるデータが確実に検証済みだと断言できないなら、念のため自分で検証しておくべきだろう。

\subsubsection{有効性の決定}
% src: Determining Validity

入力検証は、何を有効とするかを決めることから始まる。
これは思ったほど簡単ではない。
将来ありうる有効な入力値をすべて予測し、それ以外の値を正当な理由をもってどう禁止するかを考え出すことにほかならないからである。
この判断はふつう開発者が下し、ユーザーが求めるかもしれないものと、もっと広い範囲を許すために余分にかかるコーディングの手間とを天秤にかけなければならない。
理想を言えば、何を有効な入力とするかはソフトウェアの要件で規定され、よい設計がその指針を与えてくれるだろう。

\begin{exa}[整数入力の範囲]\label{exa:ch10-integer-range}
整数の入力なら、標準のデータ型だという理由で、32 ビット整数の全範囲を許すのが当然に思えるかもしれない。
しかし先のことを考えると、コードのどこかでこれらの値を足し合わせるならもっと大きな整数が必要になるので、32 ビットという制限は恣意的なものになってしまう。
逆に、有効な値の上限をもっと低く設定するのが妥当なら、値の合計が 32 ビットに収まることを保証できる。
\end{exa}

何を有効な入力とするかについて正しい答えを出すには、そのアプリケーションに固有の文脈を調べる必要がある。
ドメイン知識がセキュリティにとって重要であることを示す格好の例である。
有効とみなす値の範囲が決まれば、どのデータ型を使えばよいかは簡単に決まる。

ふつううまくいくのは、入力に明示的な上限を設けたうえで、実装には十分な余裕（headroom）を残し、有効な入力ならどれでも確実に正しく処理できるようにすることである。
ここでいう余裕とは、例えばテキスト文字列を 4,096 バイトのバッファにコピーするなら、有効な最大長を 4,000 バイトにして少しゆとりを持たせておく、という意味である。
（C では、終端のヌル文字の分だけバッファを 1 バイトあふれさせてしまうのが、ありがちな典型的な誤りである。）
難しい課題が好きなプログラマもいるが、気前よくしすぎる（できるだけ広い範囲の入力を許す）と、実装に必要以上に大きく難しい仕事を背負わせることになり、コードが複雑になってテストの負担も増える。
オンラインショッピングのアプリケーションが 10 億個の商品が入ったカートを扱えるとしても、そんな非現実的な取引を処理しようとするのは逆効果である。
そのような入力は拒否するのがいちばん親切だろう（誰かの猫がキーボードの上に座っただけかもしれない）。

\subsubsection{検証の基準}
% src: Validation Criteria

入力検証のチェックの多くは、いくつかの基準を組み合わせたものである。
具体的には、入力が最大サイズを超えていないこと、データが適切な形式で届いていること、そして値が許容範囲内にあることを確かめる。

値のサイズのチェックは、主にコードに対するサービス拒否の脅威を避けるための手早いテストである。
ここでいうサービス拒否の脅威とは、メガバイト単位の信頼できない入力を抱え込んで、アプリケーションの動きが鈍くなったりクラッシュしたりすることである。
データ形式としては、数値なら数字の並び、文字列なら決められた文字だけから成ること、あるいは XML や JSON のようなもっと複雑な形式が考えられる。
通常は次の順にチェックするのが賢明である。
まずサイズを制限して、過大な入力の処理に時間を無駄にしないようにする。次に、パースする前に入力が整形式であることを確かめる。最後に、得られた値が許容範囲内にあるかをチェックする。

有効な値の範囲を決めるのはいちばん主観的な判断になりがちだが、具体的な上限・下限を持つことが重要である。
範囲をどう定義するかはデータ型によって決まる。
整数なら、範囲は最小値以上かつ最大値以下である。
浮動小数点数なら、精度（小数点以下の桁数）にも制限を設けることがある。
文字列なら最大長であり、たいていは正規表現などで決まる許容される形式や構文もある。
文字列の最大長はバイト数ではなく文字数で指定することを勧める。
そうしておけば、少なくともプログラマでない人にも、この制約が何を意味するのかわかる見込みがあるからである。

入力を抽象的に考えるのではなく、\emph{ある目的にとって}有効かどうかで考えると役に立つ。

\begin{exa}[言語翻訳システムの段階的な検証]\label{exa:ch10-translation-validation}
言語翻訳システムなら、まず、サポートする文字集合に収まり、サポートするすべての言語に共通の最大長を超えないことを検証してから入力を受け付けるかもしれない。
次の処理段階でテキストを解析して何語かを判定するなら、言語が決まった時点で、テキストをその言語に合った文字集合にさらに絞り込める。
\end{exa}

\begin{exa}[注文数量の上限]\label{exa:ch10-order-quantity}
購入請求書の、注文された商品の数量を表す整数入力の検証を考える。
顧客が実際に注文しうる最大数量を決めるのは簡単ではないが、前もって考えておく価値のあるよい問いである。
過去のデータが手元にあれば、簡単な SQL クエリを投げるだけで、参考になる興味深い実例が見つかるかもしれない。
32 ビット整数の最大値なら制限が最も緩いので最良の選択だ、と主張することもできるが、実際にはそれが理にかなうことはめったにない。
ある商品を 4,294,967,295 個も注文するのを、何かの間違い以外だと考える人がいるだろうか。
2 進数に由来するそんな奇妙な数をプログラマでない人が覚えていることはまずないので、1,000,000 のように、ユーザーにわかりやすい上限を選ぶほうが理にかなっている。
正当な理由でこの上限に突き当たる人がいたら、おそらく知っておく価値のあることであり、上限は簡単に調整できるはずである。
しかもその過程で、開発者はそれまで想像もしなかった実際のユースケースを知ることになる。
\end{exa}

入力検証の第一の目的は、無効な入力を一つも通さないことである。
いちばん単純な方法は、ここまでの議論で暗黙のうちにそうしてきたように、無効な入力をただ拒否することである。
もっと寛容な方法としては、無効な入力を検出して有効な形に直すこともできる。
これらの手法と、それぞれをどんなときに使うべきかを見ていこう。

\subsubsection{無効な入力の拒否}
% src: Rejecting Invalid Input

決められた規則に合わない入力を拒否するのは、最も単純で、おそらく最も安全な手法である。
受け入れるか拒否するかのどちらかしかない方式はいちばんすっきりしていて明快であり、たいていは正しく実装するのもいちばん簡単である。
海で泳いでも安全かどうかを判断するときの常識的な助言「迷ったら、海に出るな」（When in doubt, don't go out）と同じようなものである。
拒否の仕方には、どれか一つの欄の記入が不適切なら Web フォームを処理しない、といった単純なものから、あるレコードにたった一つ違反があっただけで入ってきたデータのバッチ全体を拒否する、といった極端なものまである。

Web フォームのように人が直接入力するときは、何が悪いのかわかるエラーメッセージを示して、誤りを直して再送信しやすくするのがいちばん親切である。
ユーザーが無効な入力を送信するのは、おそらく単なる誤りか、検証規則を知らないかのどちらかであり、どちらも望ましいことではない。
処理を止めてデータの提供元に有効な入力を求めるのは入力検証の保守的なやり方であり、繰り返しデータを提供する相手にとっては、学んで慣れるよい機会にもなる。

人からの不正な入力を入力検証で拒否するときのベストプラクティスには、次のものがある。
\begin{itemize}
\item 何が有効な入力かをユーザーインターフェースの中で説明し、少なくともそれを読んだ人が推測と再試行を繰り返さずに済むようにする。（市外局番を括弧でくくるのではなくハイフンでつなぐべきだなんて、どうやって知ればいいのか。）
\item 複数のエラーをまとめて指摘し、1 回の手順で修正して再送信できるようにする。
\item 人が直接入力する場合は、規則を単純で明快なものにしておく。
\item 複雑なフォームはいくつかの部分に分け、部分ごとに別のフォームにして、入力が先に進んでいることがわかるようにする。
\end{itemize}

入力が人から直接ではなく他のコンピュータから来る場合は、もっと厳格な入力検証にするのが賢明かもしれない。
その要件を実装する最善の方法は、期待する入力形式やその他の制約を正確に記述した文書を書くことである。
専門的に運用されているシステムからの入力なら、有効なデータだけを部分的に処理しようとするより、入力のバッチ全体をまるごと拒否するのがいちばん理にかなっているかもしれない。
無効な入力があるということは、どこかが仕様から外れていることを示しているからである。
バッチ全体を拒否すれば、何が処理済みで何が未処理かを整理しなくても、エラーを直してデータセット全体を再送信できる。

\subsubsection{無効な入力の修正}
% src: Correcting Invalid Input

完全に有効な入力だけを受け付け、それ以外をすべて拒否するやり方は安全で単純かもしれないが、いつも最善というわけでは決してない。
何としても顧客を獲得したいオンラインストアにとって、チェックアウトの途中で入力を拒否すれば、恐ろしい「カゴ落ち」（abandoned cart）が増え、売上を失うことになりかねない。
対話的なユーザー入力では厳格な規則はいらだちのもとになるので、ソフトウェアがユーザーの有効な入力を手助けできるなら、そうすべきである。

ささいな誤りで処理を止めたくないなら、入力検証のコードで、無効な入力を拒否する代わりに有効な値に変換して修正してもよい。
簡単な例としては、長すぎる文字列を最大長で切り詰めたり、余分な先頭・末尾の空白を取り除いたりすることがある。
もっと複雑な修正の例もある。

\begin{exa}[郵送先住所の入力]\label{exa:ch10-mailing-address}
郵便サービスが認める形式どおりに郵送先住所を入力する、というよくある例を考える。
空白の入れ方、通りの名前の綴り、略語の形が厳密に決まっているので、これはかなりの難題である。
ほぼ唯一の方法は、似た住所を公式の形式で書いた最も近そうな候補を提示し、回答者に選んでもらうことである。
\end{exa}

扱いにくい検証要件への最善の処方箋は、入力をできるだけ単純に設計することである。
例えば、市外局番を括弧でくくったり、決まった位置にダッシュを入れたりしなければならない電話番号の入力に苦労したことのある人は多いだろう。
そうはせずに、電話番号はただの数字の列とし、そもそも構文規則を設けないようにすればよい。

修正すれば時間の節約にはなるかもしれないが、どんな修正にも、（ユーザーから見て）意図しない形で入力を変えてしまうおそれがつきまとう。

\begin{exa}[電話番号の修正の危険]\label{exa:ch10-phone-correction}
10 桁の入力を期待する電話番号のフォーム欄を例にとる。
ハイフンのようなよくある文字を取り除き、その結果が有効な 10 桁になれば受け付けるのは安全なはずである。
しかし桁数が多すぎる場合、ユーザーは国際電話番号を入力するつもりだったのかもしれないし、打ち間違えたのかもしれない。
いずれにせよ、切り詰めてしまうのはおそらく安全ではない。
\end{exa}

適切な入力検証には慎重な判断が必要だが、入力検証によってソフトウェアシステムはずっと信頼性が高くなり、その結果より安全になる。
問題空間が小さくなり、不要で厄介なエッジケースがなくなり、テストしやすくなり、システム全体がよりはっきり定義された安定したものになる。

%----------------------------------------------------------------------
\subsection{文字列の脆弱性}\label{subsec:ch10-string-vulnerabilities}
% src: Character String Vulnerabilities

\begin{quote}
「2006 年にもなってプログラマとして働いていながら、文字、文字集合、エンコーディング、Unicode の基本を知らない者を見つけたら、罰として潜水艦の中で 6 か月間タマネギの皮をむかせてやる。」---Joel Spolsky
\end{quote}

ほぼすべてのソフトウェアコンポーネントは、少なくともコマンドラインのパラメータを受け取るときや、出力を読める形で表示するときに、文字列を処理している。
文字列を大量に処理するアプリケーションもある。
ワードプロセッサ、コンパイラ、Web サーバーやブラウザなど、ほかにもたくさんある。
文字列処理はあらゆるところにあるので、文字列処理にまつわるよくあるセキュリティ上の落とし穴を知っておくことが重要である。
以下では、うっかり脆弱性を作り込まないために知っておくべき数多くの問題のうち、いくつかを取り上げる。

\subsubsection{長さの問題}
% src: Length Issues

文字列の長さには上限がないこともあるので、最初の課題は長さである。
きわめて長い文字列を固定長の記憶領域にコピーすれば、バッファオーバーフローを招く。
正しく扱えたとしても、巨大な文字列が CPU サイクルやメモリを過度に消費すれば性能の問題が起き、可用性が脅かされかねない。
したがって、第一の防衛線は、入ってくる信頼できない文字列の長さを妥当なサイズに制限することである。
言うまでもないことかもしれないが、バッファを確保するときに文字数とバイト長を混同してはならない。

\subsubsection{Unicode の問題}
% src: Unicode Issues

現代のソフトウェアはたいてい、世界中の書き言葉を網羅する豊かな文字集合である \underline{Unicode} に頼っている。しかし、その豊かさと引き換えに隠れた複雑さが多くあり、それがエクスプロイトの温床になりうる。
世界のテキストをバイト列として表す文字エンコーディングは数多くあるが、ソフトウェアはほとんどの場合、一種の共通語として Unicode を使う。
最新の Unicode 標準（執筆時点でバージョン 13.0）は 1,000 ページをわずかに超え、140,000 を超える文字、正規化のアルゴリズム、旧来の文字コード標準との互換性、右から左に書く言語のサポートを規定している。
世界のほぼすべての書き言葉をサポートし、100 万を超えるコードポイントを符号化している。

Unicode テキストには、知っておくべきエンコーディングがいくつかある。
最も一般的なのは UTF-8 だが、UTF-7、UTF-16、UTF-32 のエンコーディングもある。
バイトと文字の間を正確に変換することはセキュリティにとって重要である。
そうしないと、変換の過程でテキストの内容がいつの間にか変わってしまうおそれがある。
照合順序（collation、ソート順）はエンコーディングと言語によって変わるので、それを知らないと意図しない結果になることがある。
別の国や言語向けに設定されたコンピュータで実行した場合のように、\underline{ロケール}（locale）が違うと動作が変わる操作もあるので、そのすべての場合について正しく動くかをテストすることが重要である。
異なるロケールをサポートする必要がないなら、システムの設定からたまたま設定されているロケールを引き継ぐのではなく、ロケールを明示的に指定することを検討すべきである。

Unicode には意外な機能が多いので、セキュリティの観点からの結論は、バイトを直接いじろうとせず、信頼できるライブラリを使って文字列を扱え、ということになる。
重い仕事は専門家に任せるのがいちばんだという点で、Unicode は暗号と似ているといえる。
自分が何をしているのかわかっていないと、聞いたこともない珍しい文字や言語の癖が脆弱性を生むかもしれない。
この節では知っておく価値の大きい主な問題を詳しく述べるが、Unicode の複雑さを徹底的に掘り下げるには本が丸ごと 1 冊必要だろう。
細かな点まで理解する必要のある開発者向けに、セキュリティ上の考慮事項についての詳しい指針が Unicode コンソーシアムから提供されている。
手始めには UTR\#36: Unicode Security Considerations がよい。

% src: Encodings and Glyphs
\underline{エンコーディングとグリフ：}
Unicode が符号化するのは文字であって、\underline{グリフ}（glyph、文字を描画したときの視覚的な形）ではない。
この単純な原則からはさまざまなことが帰結するが、いちばんわかりやすい説明は、おそらく次の例だろう。大文字の I（U+0049）とローマ数字の 1（U+2160）は別々の文字なのに、まったく同じグリフで表示されることがある（このような文字を\underline{同形文字}（homomorph）という）。
Web の URL は各国の言語をサポートしており、見た目がよく似た文字を使うのは、攻撃者がユーザーをだますためによく使う手口である。

\begin{exa}[PayPal の P に見えるキリル文字]\label{exa:ch10-paypal-homomorph}
有名な事例として、PayPal の P とまったく同じに見えるキリル文字（U+0420）を使って正当なサーバー証明書を取得し、完璧なフィッシングの仕掛けを作り上げた者がいた。
\end{exa}

Unicode には\underline{結合文字}（combining character）があり、同じ文字を異なる形で表せる。

\begin{exa}[Ç の 2 通りの表現]\label{exa:ch10-combining-cedilla}
ラテン文字の Ç（U+00C7）は、大文字の C（U+0043）の後に「Combining Cedilla」（U+0327）を続けた 2 文字でも表せる。
1 文字の形も 2 文字の形も同じグリフで表示され、意味の違いもないので、コードでは一般に両者を同等なものとして扱うべきである。
\end{exa}

典型的なコーディング戦略は、まず入力文字列を正準形に\underline{正規化}（normalize）することである。ただ、あいにく Unicode には正規化が何種類もあるので、細部まで正しく扱うにはさらに学ぶ必要がある。

% src: Case Change
\underline{大文字・小文字の変換：}
文字列を大文字や小文字に変換するのは、test、TEST、tEsT などをコードが同じものとして扱えるようにテキストを正準化する、よくある方法である。
ところが、英語の A から Z 以外には、大文字・小文字の変換で意外な振る舞いをする文字がある。

\begin{exa}[ドットのない i による検査のすり抜け]\label{exa:ch10-dotless-i}
例えば、次の 2 つの文字列は異なるが、ぱっと見ではほとんど区別がつかない：「This ıs a test.」と「This is a test.」（最初の文字列では、小文字の i の一つに上のドットがない）。
大文字に変換すると、ドットのない小文字 ı（U+0131）もおなじみの小文字 i（U+0069）も大文字の I（U+0049）になるので、どちらもまったく同じ「THIS IS A TEST.」になる。

これがどう脆弱性につながるかを見るために、入力文字列に \verb|<script>| が含まれているかをチェックする場合を考える。
コードは文字列を小文字に変換してその部分文字列を探し、そのあと出力用に大文字に変換するかもしれない。
すると、文字列 \texttt{<scrıpt>} はチェックをすり抜けるのに、出力では \verb|<SCRIPT>| になり、Web ページ上でスクリプトを実行できてしまう。
これこそ、コードが防ごうとしていたことにほかならない。
\end{exa}

%----------------------------------------------------------------------
\subsection{インジェクションの脆弱性}\label{subsec:ch10-injection}
% src: Injection Vulnerabilities

\begin{quote}
「政治に真実を注入（inject）したら、政治などなくなってしまうだろう。」---Will Rogers
\end{quote}

\begin{exa}[クレジットカード契約書の書き換え]\label{exa:ch10-agarkov-contract}
郵便制度を詰まらせるほど大量のダイレクトメールのかなりの部分を、頼んでもいないクレジットカードの勧誘が占めている。ところが、ある賢い受取人は銀行を逆手に取ることに成功した。
Dmitry Agarkov は、気に入らない条件のカードへの申し込みを勧める販促の案内を捨てずに、同封の契約書をスキャンし、自分にきわめて有利な条件になるよう本文を注意深く書き換えた。
その条件には、金利 0\%、無制限の与信枠、そして銀行がカードを解約した場合には彼が気前のよい支払いを受け取れる、といったものが含まれていた。
彼は書き換えた契約書に署名して銀行に返送し、ほどなく新しいクレジットカードが届いた。
Dmitry はしばらくの間、類を見ないほど有利な契約の気前のよい条件を享受したが、銀行がようやく気付くと事態はこじれた。
長い法廷闘争では、書き換えた契約の有効性を認める彼に有利な判決も出たが、最終的に彼は示談で和解した。
\end{exa}

これは\underline{インジェクション攻撃}（injection attack）の現実世界の例である。
契約書はコードと同じではないが、プログラムの動作とよく似たやり方で、署名者に決められた行動をとらせる。
契約の条件を変えたことで、Dmitry は銀行に意に反する行動をとらせることができた。
クレジットカードの口座を管理するソフトウェアを自分に有利なように書き換えたも同然である。
ソフトウェアもこの種の攻撃を受けやすい。
信頼できない入力によってソフトウェアがだまされ、予期しないことをしてしまうことがあり、これは実際かなりよくある脆弱性である。

\begin{dfn}[インジェクション攻撃の考え方]\label{dfn:ch10-injection-attack}
実行すべき操作を表す文字列やデータ構造を組み立て、それを実行して目的の仕事を果たす、という手法はソフトウェアでよく使われる。
（銀行がクレジットカードサービスの運用方法を定めた契約書を書き、その条件がそのまま受け入れられることを期待するのに似ている。）
そこに信頼できない情報源からのデータが関わると、そのデータが実行時の動作に影響を与えられる場合がある。
攻撃者が操作の意図した効果を変えられるなら、その影響が信頼境界を越え、より高い権限を持つソフトウェアによって実行されてしまうかもしれない。
これが、抽象的に言ったときのインジェクション攻撃の考え方である。
\end{dfn}

よくあるインジェクション攻撃をいくつか具体的に説明する前に、信頼できないデータの影響がいかに人を欺きうるかを示す単純な例を考えよう。

\begin{exa}[「No Game Scheduled」という名前のチーム]\label{exa:ch10-no-game-scheduled}
真偽の怪しい話だが、ある学内ソフトボールチームが「No Game Scheduled」（試合予定なし）という巧妙な名前を付け、まさにこの種の混乱をうまく利用したという。
対戦相手のチームは何度か予定表でこの名前を見て、その日は試合がないのだと思い込み、試合に来なかったために不戦敗になった。
これがインジェクション攻撃の例だといえるのは、チーム名は予定管理システムへの入力なのに、「No Game Scheduled」が予定管理システムからのメッセージだと誤解されたからである。
\end{exa}

同じインジェクション攻撃の原理は、さまざまな技術（つまり、操作を表すために組み立てた文字列のさまざまな形式）に当てはまる。
例えば次のものがあるが、これらに限られない。
\begin{itemize}
\item SQL 文
\item ファイルパスのトラバーサル
\item 正規表現（サービス拒否の脅威として）
\item XML データ（とくに XXE 宣言）
\item シェルコマンド
\item 文字列をコードとして解釈すること（例えば JavaScript の \texttt{eval} 関数）
\item HTML と HTTP ヘッダー（第11章で扱う）
\end{itemize}

以下の節では、最初の 4 種類のインジェクション攻撃を詳しく説明する。
シェルコマンドやコードのインジェクションも SQL インジェクションと同じ仕組みで、ずさんな文字列の組み立て方が信頼できない入力に悪用される。
Web のインジェクション攻撃は次の章で扱う。

\subsubsection{SQL インジェクション}
% src: SQL Injection

% 図 10-1: Exploits of a Mom (courtesy of Randall Munroe, xkcd.com/327)
\begin{exa}[xkcd \#327「Exploits of a Mom」]\label{exa:ch10-bobby-tables}
古典的な xkcd のコミック \#327（図 10-1）は、大胆な \underline{SQL インジェクション}（SQL injection）攻撃を描いている。
両親が子どもに、特殊文字を含む、ありそうもない発音もできない名前を付ける。
その名前が地元の学区のデータベースに入力されると、学校の記録が壊されてしまう。
\end{exa}

この仕組みを理解するために、学校の登録システムが SQL データベースを使っていて、次のような SQL 文で生徒のレコードを追加するとしよう。
\begin{verbatim}
INSERT INTO Students (name) VALUES ('Robert');
\end{verbatim}
この単純化した例では、この文によって名前「Robert」がデータベースに追加される。
（実際には、2 組の括弧の中のリストには name 以外にも多くの列が並ぶが、ここでは簡単のため省略している。）

次に、\verb|Robert'); DROP TABLE students;--| というばかげた名前の生徒を想像してほしい。
% 注: 原文はこの名前の直後に「.``」という崩れた記号があり、スクレイプ時の書式の残骸と思われるので除いた。
組み立てられる SQL コマンドは次のようになる（原典では生徒の名前にあたる部分が強調表示されている）。
\begin{verbatim}
INSERT INTO Students (name) VALUES ('Robert'); DROP TABLE Students;--');
\end{verbatim}
SQL コマンドの構文規則に従うと、この文字列には実は 2 つの文が含まれている。
\begin{verbatim}
INSERT INTO Students (name) VALUES ('Robert');
DROP TABLE Students; --');
\end{verbatim}
2 つの SQL コマンドのうち最初のものは、意図どおり「Robert」のレコードを挿入する。
しかし生徒の名前に SQL の構文が含まれているので、意図しない 2 つ目のコマンド \texttt{DROP TABLE} まで注入され、テーブル全体が削除されてしまう。
二重ダッシュはコメントを表すので、SQL エンジンはそれ以降のテキストを無視する。
この手口で後ろに残る構文（一重引用符と閉じ括弧）を読み捨て、実行を妨げる構文エラーを避けることで、エクスプロイトが成立する。

では、SQL インジェクションの脆弱性が実際にどう見えるのか、どうすれば防げるのかを知るために、コードをもう少し詳しく見てみよう。
架空の学校登録システムのコードは、最初に取り上げた基本的な例のように SQL コマンドをテキスト文字列として組み立て、それを実行する。
入力データからは、生徒のレコードに記入する名前などの情報が渡される。
理屈の上では、（法的な名前に ASCII の特殊文字が含まれうるというかなり疑わしい前提に立てば）職員がこの入力を公式の記録と照らし合わせ、正確さを確認したと考えることさえできる。

プログラマの致命的な誤りは、珍しい名前が一重引用符から「抜け出す」可能性を考えずに、次のような文字列連結の文を書いたことである。
\begin{verbatim}
sql_stmt = "INSERT INTO Students (name) VALUES ('" + student_name + "');";
\end{verbatim}
インジェクション攻撃の緩和は難しくないが、気を抜いてこのようなコードを書いてしまわないよう注意を怠ってはならない。
脆弱性の根本原因は、信頼できない入力とコマンド文字列を混ぜることである。
そうした入力は引用符から抜け出し、意図しない有害な結果を引き起こしうるからである。

どんな文字列を有効な名前とするかは要件として重要な問題だが、ここではこの SQL 文で一重引用符として使われているアポストロフィ文字だけに注目しよう。
アポストロフィは SQL コマンドの構文をこじ開ける鍵となる文字だが、アポストロフィを含む名前（O'Brien など）も実在するので、アプリケーションは入力検証でこの文字を禁止するわけにはいかない。
この名前は引用符で囲んだ文字列 \verb|'O''Brien'| として正しく書けるが、脆弱性を完全に取り除くには、ほかにも特別な扱いが必要な特殊文字がたくさんあるかもしれない。

さらなる防御として、生徒を登録するソフトウェアにはテーブルを削除できる管理者権限を与えないよう、SQL データベースを設定すべきである。
生徒の登録という仕事にその権限は必要ないからである。
（これは第4章の最小権限（Least Privilege）パターンの一例である。）

独自の SQL サニタイズのコードを書いて「車輪の再発明」をするよりも、こうした問題に対処するために作られた、SQL コマンドを組み立てるライブラリを使うのがベストプラクティスである。
信頼できるライブラリが使えない場合は、インジェクション攻撃の試みが拒否されるか安全に処理されること、そして O'Brien のような名前の生徒でもすべて正しく動作することを確かめるテストケースを用意すべきである。

\begin{exa}[誤った SQL の組み立て方（Python）]\label{exa:ch10-sql-wrong}
以下に、誤ったやり方と正しいやり方を示す簡単な Python のコード片を挙げる。
まずは、Bobby Tables 攻撃を模した誤ったやり方である。
\begin{verbatim}
import sqlite3
con = sqlite3.connect('school.db')
student_name = "Robert'); DROP TABLE Students;--"
# The WRONG way to query the database follows:
sql_stmt = "INSERT INTO Students (name) VALUES ('" + student_name + "');"
con.executescript(sql_stmt)
\end{verbatim}
コードは SQL データベースへの接続（\texttt{con}）を作成したあと、生徒の名前を変数 \texttt{student\_name} に代入する。
次に、文字列 \texttt{student\_name} を VALUES のリストに差し込んで SQL の INSERT 文を組み立て、\texttt{sql\_stmt} に代入する。
最後に、その文字列を SQL スクリプトとして実行する。
\end{exa}

\begin{exa}[正しい SQL の組み立て方（Python）]\label{exa:ch10-sql-right}
正しくは、次のコード片のように、信頼できないデータを含むパラメータの挿入をライブラリに任せる。
\begin{verbatim}
import sqlite3
con = sqlite3.connect('school.db')
student_name = "Robert'); DROP TABLE Students;--"
# The RIGHT way to query the database follows:
con.execute("INSERT INTO Students (name) VALUES (?)", (student_name,))
\end{verbatim}
この実装では、プレースホルダ \texttt{?} に、後に続くタプルのパラメータ（\texttt{student\_name} の文字列）が埋め込まれる。
INSERT 文の文字列の中に引用符が要らないことに注意しよう。
引用符の処理はすべて自動的に行われる。
この書き方ならインジェクションは起きず、Bobby の奇妙な名前も安全にデータベースに入力される。
\end{exa}

\begin{rem}[複数文の実行を禁じても防御にはならない]\label{rem:ch10-execute-single-statement}
この例には、補足しておくべき細かな点がある。
元のエクスプロイトを成立させるには、\texttt{executescript} というライブラリ関数が必要である。
\texttt{execute} は単一の文しか受け付けないので、この特定の攻撃に対しては一種の防御になるからである。
しかし、インジェクション攻撃はどれも追加のコマンドを伴うと考え、この制限で十分に守られると思うのは誤りである。
例えば、その学校には、\verb|Robert', 'A+');--| という、これまた発音できない名前の生徒がもう一人いるとする。
この生徒もただの Robert もどちらも落第点なのに、別の SQL テーブルに成績が記録されると、この生徒の評価は A+ に引き上げられてしまう。
なぜだろうか。

ただの Robert の成績を送信すると、コマンドは意図どおり F の成績を次のように入力する。
\begin{verbatim}
INSERT INTO Grades (name, grade) VALUES ('Robert', 'F');
\end{verbatim}
しかし名前が \verb|Robert', 'A+');--| だと、コマンドは次のようになる。
\begin{verbatim}
INSERT INTO Grades (name, grade) VALUES ('Robert', 'A+');--', 'F');
\end{verbatim}
\end{rem}

最後に、注意深い読者なら気付いたかもしれない点について、xkcd の「Little Bobby Tables」の例に一言触れておく。
前提のばかばかしさはさておき、データベースのテーブルに勝手に付けられた具体的な名前（Students）を Bobby の両親が見越していたのは、驚くべき偶然である。
これは芸術上の自由（artistic license）として片付けるのがいちばんだろう。

\subsubsection{パストラバーサル}
% src: Path Traversal

\underline{ファイルパストラバーサル}（file path traversal）は、インジェクション攻撃と密接に関係するよくある脆弱性である。
前節の例のように引用符から抜け出すのではなく、親ディレクトリへと抜け出して、ファイルシステムの別の部分に予期しないアクセスをする攻撃である。

\begin{exa}[画像ストアに対するパストラバーサル]\label{exa:ch10-image-store-traversal}
画像のコレクションを提供するために、実装では画像ファイルを \verb|/server/data/image_store| というディレクトリに集めておき、X という名前の画像を要求されたら、（信頼できない）入力名 X から組み立てたパス \verb|/server/data/image_store/X| から画像データを取り出して返すとする。

すぐに思いつく攻撃は \verb|../../secret/key| という名前を要求することで、そうすると、非公開であるべきファイル \verb|/server/secret/key| が返されてしまう。
\texttt{.}（ドット）はカレントディレクトリを表す特別な名前であり、\texttt{..}（ドットドット）は親ディレクトリを表し、ファイルシステムのルートに向かってたどっていけることを思い出そう。
次に挙げるパス名は、どれも同じ場所を指している。
\begin{verbatim}
/server/data/image_store/../../secret/key
/server/data/../secret/key
/server/secret/key
\end{verbatim}
\end{exa}

この種の攻撃を防ぐ最善の方法は、入力（この例では X）に使える文字の集合を制限することである。
多くの場合、入力が英数字だけの文字列であることを入力検証で保証すれば、それだけで完全に扉を閉ざせる。
この方法がうまくいくのは、ファイルシステムの意図した範囲から抜け出すのに必要な、厄介なファイル区切り文字や親ディレクトリの表記を締め出せるからである。

しかし、この方法では制限が強すぎることもある。
任意のファイル名を扱う必要がある場合、この単純な方法では制限が強すぎるので、さらに手を打たなければならない。しかもファイルシステムは複雑なので、その作業も複雑になりがちである。
さらに、コードが複数のプラットフォームで動くなら、ファイルシステムの違いにも気を付ける必要がある（例えば、*nix のパス区切り文字はスラッシュだが、Microsoft Windows ではバックスラッシュである）。

\begin{exa}[不十分なパスのガード関数（Python）]\label{exa:ch10-safe-path-insufficient}
次の関数は、この Python コードが置かれているディレクトリ（\texttt{\_\_file\_\_} で表される）の中のファイルにアクセスするためのサブパスとして入力文字列を使う前に、その文字列を検査する簡単な例である。
狙いは、あるディレクトリとそのサブディレクトリの中のファイルにだけアクセスを許し、それ以外の場所にある任意のファイルには決してアクセスさせないことである。
ここに示す版では、ガード関数 \texttt{safe\_path} が、入力の先頭にスラッシュ（ファイルシステムのルートへ行く）や親ディレクトリを表すドットドットがないかを調べ、それらを含む入力を拒否する。
これを正しく行うには、場当たり的な文字列操作ではなく、Python の \texttt{os.path} の一連の機能のような標準ライブラリを使ってパスを扱うべきである。
しかし、それだけでは意図したディレクトリから抜け出されるのを防ぐには足りない。
\begin{verbatim}
def safe_path(path):
    """Checks that argument path is a safe file path. If not, returns None.
    If safe, returns the normalized absolute file path.
    """
    if path.startswith('/') or path.startswith('..'):
        return None
    base_dir = os.path.dirname(os.path.abspath(__file__))
    filepath = os.path.normpath(os.path.join(base_dir, path))
    return filepath
\end{verbatim}
この保護に残っている穴は、パスでいったん有効なディレクトリを指定してから親ディレクトリへ上がり、さらに上がっていけば抜け出せてしまうことである。
例えば、このサンプルコードを実行するカレントディレクトリがルートから 5 階層下にあるなら、パス \verb|./../../../../../etc/passwd|（ドットドットが 5 つ）は \verb|/etc/passwd| ファイルを指すことになる。
\end{exa}

ドットドットを含むパスをすべて拒否するようにして、文字列ベースの無効なパスの判定を改善することもできるが、そうしたやり方は危険をはらむ。
ありうる手口をすべて予期し、完全に塞いだと確信するのは難しいからである。
代わりに、パス文字列を自分のコードで組み立てるのではなく、\texttt{os.path} ライブラリに任せる素直な解決策がある。

\begin{exa}[安全なパスのガード関数（Python）]\label{exa:ch10-safe-path}
\begin{verbatim}
def safe_path(path):
    """Checks that argument path is a safe file path. If not, returns None.
    If safe, returns the normalized absolute file path.
    """
    base_dir = os.path.dirname(os.path.abspath(__file__))
    filepath = os.path.normpath(os.path.join(base_dir, path))
    if base_dir != os.path.commonpath([base_dir, filepath]):
        return None
    return filepath
\end{verbatim}
この保護は全面的に信頼してよい。理由は次のとおりである。
基準ディレクトリには信頼できない入力がまったく関わっていないので、信頼できるパスである。
完全にプログラマの管理下にある値だけから導かれているからである。
入力のパス文字列と結合したあとでパスを正規化すると、ドットドットによる親ディレクトリへの参照がすべて解決され、絶対パス \texttt{filepath} が得られる。
あとは、両者の最長の共通サブパスが、アクセスを制限したい意図どおりのディレクトリになっているかをチェックすればよい。
\end{exa}

\subsubsection{正規表現}
% src: Regular Expressions

効率がよく、柔軟で、使いやすい \underline{regex}（\underline{正規表現}、regular expression）は、驚くほど幅広い機能を備えており、テキスト文字列をパースするための道具としてはおそらく最も汎用的である。
正規表現は一般に、場当たり的なコードよりも（書くのも実行するのも）速く、信頼性も高い。
regex ライブラリは状態表をコンパイルし、それをインタプリタ（有限状態機械や同様のオートマトン）が実行して、文字列との照合を行う。

regex は、正しく書かれていてもセキュリティ上の問題を引き起こすことがある。
正規表現の中には実行時間が極端に長くなりやすいものがあり、攻撃者がそれを引き起こせれば深刻なサービス拒否を起こせるからである。
具体的には、regex で\underline{バックトラッキング}（backtracking）が起きる場合、つまり、ずっと先まで走査したあとで、一致を探すために何度も戻って走査し直す必要がある場合に、実行時間が膨れ上がることがある。
セキュリティ上の危険は一般に、信頼できない入力に regex を指定させてしまうことから生じる。あるいは、コードにすでにバックトラッキングする regex が含まれている場合には、計算量が最大になる長い最悪ケースの文字列を信頼できない入力として与えられることから生じる。

\begin{exa}[バックトラッキングする正規表現]\label{exa:ch10-backtracking-regex}
次の例が示すように、バックトラッキングする regex は一見無害に見えることがある。
次の Python コードは、著者のささやかな Raspberry Pi Model 4B では実行に 3 秒以上かかる。
読者のプロセッサはおそらくずっと速いだろうが、例の 26 個の D に 1 つ足すごとに実行時間が 2 倍になるので、文字列を少し長くするだけで、どんなプロセッサでも簡単に動けなくできる。
\begin{verbatim}
import re
print(re.match(r'(D+)+$', 'DDDDDDDDDDDDDDDDDDDDDDDD!'))
\end{verbatim}
\end{exa}
% 注: 原文は「26 個の D」と書いているが、掲載コードの文字列中の D は 24 個である。原文どおりとした。

実行時間が極端に長くなる危険は、バックトラッキングなどの非線形な計算が爆発しうる場合には、信頼できない入力をパースする処理なら何にでもある。
次の節ではこの種の XML 実体の例を見るが、ほかにも例はたくさんある。

こうした問題を緩和する最善の方法は計算の中身によって異なるが、この種の攻撃に対抗するための一般的な手法がいくつかある。
爆発する可能性のある計算に、信頼できない入力が影響を与えないようにすること。
正規表現でいえば、信頼できない入力に regex を定義させない、できればバックトラッキングを避ける、regex で照合する文字列の長さを制限する、ということである。
最悪の場合にどれだけの計算になりうるかを見極め、テストして極端に遅くならないことを確かめるべきである。

\subsubsection{XML の危険性}
% src: Dangers of XML

XML は、強力であると同時に人間にも読めるので、構造化データを表現する方法として最もよく使われているものの一つである。
しかし、XML の強力さは武器にもなりうることを知っておくべきである。
信頼できない XML が XML 実体を使って害を及ぼす主な方法は 2 つある。

\begin{exa}[XML 実体の指数的な展開]\label{exa:ch10-xml-entity-expansion}
XML の実体宣言はあまり知られていない機能だが、あいにく攻撃者はその悪用方法を次々と編み出してきた。
次の例では、名前付き実体 \texttt{big1} を 4 文字の文字列として定義している。
別の名前付き実体 \texttt{big2} は \texttt{big1} を 8 個並べたもの（計 32 文字）として定義し、\texttt{big3} はさらにそれを 8 個並べたもの、というように続く。
\texttt{big7} まで来ると扱うデータは 1 メガバイトになり、そこからさらに大きくするのも簡単である。
この例では 8 メガバイトの XML の塊ができあがる。
見てのとおり、あと数行足すだけでギガバイト単位に達する。
\begin{verbatim}
<!DOCTYPE dtd[
  <!ENTITY big1 "big!">
  <!ENTITY big2 "&big1;&big1;&big1;&big1;&big1;&big1;&big1;&big1;">
  <!ENTITY big3 "&big2;&big2;&big2;&big2;&big2;&big2;&big2;&big2;">
  <!ENTITY big4 "&big3;&big3;&big3;&big3;&big3;&big3;&big3;&big3;">
  <!ENTITY big5 "&big4;&big4;&big4;&big4;&big4;&big4;&big4;&big4;">
  <!ENTITY big6 "&big5;&big5;&big5;&big5;&big5;&big5;&big5;&big5;">
  <!ENTITY big7 "&big6;&big6;&big6;&big6;&big6;&big6;&big6;&big6;">
]>
<mega>&big7;&big7;&big7;&big7;&big7;&big7;&big7;&big7;</mega>
\end{verbatim}
\end{exa}

\begin{exa}[XML 外部実体（XXE）によるファイルの読み出し]\label{exa:ch10-xxe}
外部実体宣言を使えば、さらに多くの手口が可能になる。
次の例を考えよう。
\begin{verbatim}
  <!ENTITY snoop SYSTEM "file:///etc/passwd>" >
\end{verbatim}
この宣言は、見たとおりのことをする。
パスワードファイルを読み込み、以降の XML で \verb|&snoop;| が現れる箇所にはどこでもその内容を展開できるようにする。
攻撃者がこれを XML として送り込み、実体を展開した結果を見ることができれば、名前さえ指定できればどんなファイルの内容でも漏えいさせられる。
\end{exa}

この種の問題に対する第一の防衛線は、コードが処理する XML に信頼できない入力を一切入れないことである。
XML の外部実体が必要ないなら、入力から外部実体を取り除くか、そうした宣言の処理を無効にして、この種の攻撃を防ぐべきである。

\subsubsection{インジェクション攻撃の緩和}
% src: Mitigating Injection Attacks

さまざまな種類のインジェクション攻撃は、信頼できない入力を使って、アプリケーションの文脈で実行される文やコマンドに影響を与える、という共通の手口に頼っている。同じように、それらへの緩和策にも、細部は異なるものの共通する考え方がある。
入力検証はいつでもよい第一の防衛線だが、許容する入力がどんなものかによっては、それだけで十分とは限らない。

実行のために組み立てる文字列に、信頼できないデータを（例えばコマンドとして）差し込もうとするのは避けるべきである。
SQL をはじめ、インジェクション攻撃を受けやすい機能を扱う現代的なライブラリなら、データをコマンドとは別に渡せるヘルパー関数が用意されているはずである。
こうした関数は、引用符付けやエスケープなど、どんな入力に対しても意図した操作を安全に実行するために必要なことをすべて引き受けてくれる。
ライブラリの文書に、セキュリティについての具体的な記述があるかを確認することを勧める。
API の体裁の裏で文字列を貼り合わせているだけの、インジェクション攻撃を受けやすいずさんな実装も実際にあるからである。
迷ったときは、セキュリティテストケース（第12章を参照）で確かめるのがよい。

安全なライブラリを使えない、あるいは使う気がない場合は---繰り返しになるが、「何がまずいことになるというのか」と考え始めると滑りやすい坂を転げ落ちることになるので注意してほしい---まず、インジェクションの危険を避けられる別の方法がないかを検討すべきである。

\begin{exa}[シェルコマンドの代わりにシステムコールを使う]\label{exa:ch10-readdir}
ディレクトリの内容を一覧するなら、*nix の \texttt{ls} コマンドを組み立てる代わりに、システムコールを使うべきである。
理由は明快である。
\texttt{readdir(3)} にできるのはディレクトリエントリの情報を返すことだけだが、シェルコマンドを呼び出せば、ほとんど何でもできてしまう可能性がある。
\end{exa}

ファイルシステムを自家製のデータストアとして使いたくなる誘惑はときに抗いがたく、場合によってはそれがいちばん手早い解決策かもしれないが、安全なやり方としてはとても勧められない。
どうしても危険なやり方をとるというのなら、完全に安全にするために、ありうる攻撃をすべて予期して塞ぐのに必要な手間を甘く見てはならない。
ここでは入力検証が頼りになる。
文字列を安全な文字集合（例えば ASCII の英数字だけから成る名前）に制限できるなら、おそらく問題ないだろう。
さらなる防御層として、組み立てるコマンドや文の構文をよく調べ、何もまずいことが起きないよう、必要な引用符付けやエスケープを漏れなく施すべきである。
自分の知らないあまり知られていない書き方があるかもしれないので、関係する仕様をよく読んでおく価値がある。

よい知らせは、インジェクションの危険がある操作は、たいていソースコードを走査すれば簡単に見つかることである。
SQL コマンドが場当たり的な文字列ではなく、パラメータを使って安全に組み立てられているかを確認すべきである。
シェルコマンドのインジェクションについては、\texttt{exec(3)} とその変種の使用箇所に注意し、コマンドの引数を適切に引用符で囲むようにすべきである（Python にはまさにこの目的のための \texttt{shlex.quote} がある）。
JavaScript では \texttt{eval} の使用箇所をレビューし、組み立てる式に信頼できない入力が影響を与えうる場合には、安全な形に制限するか、使わないことを検討すべきである。

この章ではインジェクション攻撃とそれに関連するよくある脆弱性を数多く取り上げたが、インジェクションはきわめて応用の利く手法であり、さまざまな姿で現れる。
次の章では、Web の脆弱性の文脈で、インジェクションに再び（2 回）出会うことになる。

\newpage

%% =======================================================
%%  part: ch11
%% =======================================================
\section{Web セキュリティ}
% src: 11: Web Security

この章では、\cite{Kohnfelder2021} の第11章に沿って、Web セキュリティモデルの成り立ちと、その正しい使い方・誤った使い方を扱う。

\begin{quote}
「言葉が現れたとき、誰もがそれを奇跡だと言った。けれども、Web（クモの巣）そのものが奇跡だと指摘した者は一人もいなかった。」\\
\hfill --- E. B. White（『シャーロットのおくりもの』（Charlotte's Web）より）
\end{quote}

World Wide Web がこれほど成功したのは、何百万、何千万という人々が、仕組みをまったく理解しないまま日常的に使っているという驚くべき事実（今ではすっかり当たり前になっている）によるところが大きい。
これほど複雑な技術の寄せ集めでこうした比類ない成果を上げられたことは、恵みであると同時に呪いでもある。
Web が使いやすいからこそ広く成長し続けてきたことは疑いない。
その裏返しとして、端点で何百万、何千万という無頓着な人々が使う、独立したデジタルサービスの世界規模のネットワークをセキュアにするのは、きわめて困難な課題である。
そして、この大きく難しい問題の中でも、おそらく最も難しいのがセキュリティである。

セキュリティをとりわけ難しくしている要因の一つは、初期の Web がセキュリティをあまり考慮せず、かなり素朴に設計されていたことである。
その結果、現代の Web は、標準が長い時間をかけて進化してきた産物となり、その進化は激しい「ブラウザ戦争」と後方互換性の制約によって複雑になった。
要するに、Web は、セキュリティを後から付け足した「セキュリティのボルト留め」（bolting on of security）の、史上最も極端な例である。ただし、発明から四半世紀をゆうに超えた今では、それなりにまともなものになりつつある。

現代の Web はセキュアにできるものの、複雑な歴史を背負っているために非常に壊れやすい。Web クッキーの仕様 RFC 6265 の著者たちの実に味わいのある言い回しを借りれば、多くの「セキュリティとプライバシーの不適切さ」（security and privacy infelicities）に満ちてもいる。
ソフトウェア専門家は、Web 向けに開発するときにこうした問題に引っかからないよう、これらすべてを理解しておく必要がある。
ほんの小さな踏み外しが、たやすく脆弱性を生む。
インターネットは「無法の西部」（Wild West）のような場所なので、悪意ある者は Web サイトの仕組みを自由に、簡単に探ることができ、匿名のまま攻撃の糸口を探してうろつくこともできる。

この章では、Web セキュリティモデルがどのように進化してきたかという基本と、その正しい使い方・誤った使い方を中心に扱う。
脆弱性は細部から生じるものであり、セキュアな Web サイトが寸分違わず正しく行わなければならないことは非常に多い。
まず、複雑な細部を代わりに処理してくれるセキュアなフレームワークの上に構築してほしいという訴えから始めて、Web セキュリティの基本をひととおり扱う。
次に、セキュアな通信（HTTPS）、HTTP プロトコルの適切な使い方（クッキーを含む）、そして同一オリジンポリシーが組み合わさって、Web サイトの安全をどのように保っているかを見る。
最後に、Web に特有の主要な脆弱性 2 つ（XSS と CSRF）を取り上げ、組み合わせれば現代の Web サーバーのセキュリティを大きく高められる、その他の多くの緩和策を論じる。
それでも、この章は Web セキュリティを網羅した解説ではまったくない。Web セキュリティの詳細は膨大で、しかも急速に変化している。

ここでの目標は、よくある主要な落とし穴を大まかにつかみ、それらを見分けて対処できるようになってもらうことである。
Web アプリケーションは、本書のほかの箇所で扱った多くの脆弱性の影響も受ける。この章で取り上げたものだけがセキュリティ上の懸念だと受け取らないでほしい。

\begin{rem}[前提とする知識]\label{rem:ch11-prerequisites}
以下の議論では、読者が Web の基本、すなわちクライアント／サーバーモデル、クッキーを含む HTTP と HTML の基本、少しの CSS、「入門」レベルの JavaScript、そして Document Object Model（DOM）とは何かを、最低限知っていることを前提とする。
Web にあまりなじみのない読者でも、ところどころで補足的な資料を読んで穴を埋めれば、おおむね話についてこられるはずである。
\end{rem}

%----------------------------------------------------------------------
\subsection{フレームワークの上に構築する}\label{subsec:ch11-framework}
% src: Build on a Framework

\begin{quote}
「混沌から秩序を引き出すための枠組み（framework）として、デザインを使え。」\hfill --- Nita Leland
\end{quote}

現代の Web 開発ツールのおかげで、Web サイトを作るのは、Web サイトを使うのとほとんど変わらないくらい簡単になった。
セキュアな Web サイトを作るために筆者が何より勧めるのは、高品質なフレームワークに頼ること、フレームワークが備えている安全装置を決して無効にしないこと、そして面倒な細部はすべて有能な専門家に任せることである。

しっかりしたフレームワークに頼れば、以下の節で扱うような脆弱性からは守られるはずである。それでも、フレームワークを効果的に使うには、フレームワークが何をして何をしないのかを正確に理解しておくことに価値がある。
また、最初からセキュアなフレームワークを選ぶことも決定的に重要である。コードはフレームワークに大きく依存することになるので、期待を裏切られたときに後から乗り換えるのは大変だからである。
では、Web フレームワークが本当にセキュアかどうかは、どうすればわかるのか。
結局のところ、それは信頼の問題であり、作り手の善意と専門性の両方を信頼できるかどうかに行き着く。

Web フレームワークの人気や話題性はパリのファッションと同じくらい目まぐるしく移り変わり、どれを選ぶかは多くの要因で決まるので、筆者は特定のものを推奨しない。
ただし、自分で評価するときに考慮すべき一般的な指針なら挙げられる。
\begin{itemize}
\item 絶えず変化する Web の技術と実践についていけるよう、フレームワークを積極的に開発・保守している、信頼できる組織が作ったものを選ぶ。
\item ドキュメントにセキュリティについての明示的な宣言があるかを探す。見つからなければ、そのフレームワークは候補から外す。
\item 過去の実績を調べる。完璧な記録である必要はないが、対応が遅い、問題が繰り返し起きている、といった傾向があれば危険信号である。
\item 小さなプロトタイプを作り、（この章の例にあるような入力を使って）生成された HTML が適切にエスケープ・クォートされているかを確認する。
\item この章の後半で説明するように、基本的な XSS 攻撃と CSRF 攻撃を試せる簡単なテスト環境を作る。
\end{itemize}

%----------------------------------------------------------------------
\subsection{Web セキュリティモデル}\label{subsec:ch11-web-security-model}
% src: The Web Security Model

\begin{quote}
「Web がある意味まったくの無政府状態であることを、私はちょっと嬉しく思っている。私はそれで構わない。」\hfill --- Roger Ebert
\end{quote}

Web はクライアント／サーバー型の技術であり、そのセキュリティモデルを理解するには、クライアントとサーバーの両方の視点を同時に考える必要がある。
そうするとすぐに話は面白くなる。2 者のセキュリティ上の利害はしばしばぶつかり合い、しかもインターネット経由で入り込んでくる攻撃者の脅威まであるからである。

典型的なオンラインショッピングの Web サイトを考えてみよう。
ここで働くセキュリティの原則は、程度の差こそあれ、Web 上のあらゆる活動に当てはまる。
商売が成り立つには、販売者と顧客がある程度互いを信頼しなければならず、ほとんどの場合、実際にそうなっている。
それでも、悪意ある者がどうしても少しは紛れ込むので、Web サイトはすべてのクライアントを完全には信頼できないし、その逆も同じである。
次に挙げる点は、販売者と顧客の間にある、条件付きの相互信頼の微妙なところをよく表している。

\begin{exa}[販売者と顧客の要件]\label{exa:ch11-merchant-customer}
販売者の基本的な要件には、たとえば次のようなものがある。
\begin{itemize}
\item ほかの Web サイトが、自分と顧客とのやり取りに干渉できないようにしたい。
\item 正当な顧客には役立つ情報を提供しつつ、競合他社に商品や在庫の詳細をスクレイピングされる余地はできるだけ小さくしたい。
\item 顧客が価格を操作したり、在庫のない商品を注文したりできないようにしたい。
\end{itemize}
顧客の要件には、たとえば次のようなものがある。
\begin{itemize}
\item アクセスしている Web サイトが本物だという保証がほしい。
\item オンライン決済がセキュアだと確信できなければ困る。
\item 自分の買い物の履歴を、販売者が秘密にしておいてくれることを期待する。
\end{itemize}
\end{exa}

Web がうまく機能するには、明らかに両者とも警戒を怠ってはならない。
とはいえ、顧客が販売者に期待することは多い。
混乱しやすい顧客やだまされやすい顧客を教育するという難しい問題は、そもそも解決できるものかどうかも怪しいが、ここでは扱わない。
その代わり、Web セキュリティでは、販売者の視点から Web サイトをセキュアにすることに焦点を当てる。
Web は、サーバーがきちんとセキュリティを提供して初めて機能する。そうなって初めて、誠実なエンドユーザーがセキュアに Web を利用できる見込みが生まれるのである。
販売者は、顧客をどこまで信頼できるかを決めるだけでなく、顧客が自分をどこまで信頼してくれそうかも直観的に見極めなければならない。

Web のセキュリティモデルにはもう一つ奇妙な点がある。クライアントのブラウザの役割である。
Web サービスの設計が難しいのは、自分ではまったく制御できないブラウザとやり取りしなければならないからである。
悪意あるクライアントなら、何でもできるように改造したブラウザを簡単に使えるかもしれない。
不注意なクライアントなら、セキュリティホールだらけの古いブラウザを使っているかもしれない。
Web サーバーがクライアントの使うブラウザを特定のバージョンに制限しようとしても、ブラウザは身元を偽ってそうした制限を簡単にすり抜けられることを忘れてはならない。
救いは、誠実なクライアントなら、自分の利益を守るためにセキュアなブラウザを使い、定期的に更新したいと考えることである。
何より重要なのは、サーバーがセキュアである限り、悪意あるクライアントが 1 つあっても、ほかのクライアントが受けるサービスに干渉することはできないという点である。

\begin{rem}[クライアントのブラウザを信頼しすぎない]\label{rem:ch11-overtrusting-browsers}
Web サーバーが、信頼できるとは限らないクライアントのブラウザを信頼しすぎることが、多くの Web セキュリティ脆弱性の根本にある。
繰り返しになるのを承知でこの点を強調するのは、この章全体で説明していくように、これがあまりにも簡単に、そして頻繁に忘れられるからである。
\end{rem}

%----------------------------------------------------------------------
\subsubsection{HTTP プロトコル}\label{subsec:ch11-http}
% src: The HTTP Protocol

\begin{quote}
「プロトコル（作法）を重要でないと考える者は、猫を相手にしたことがないのだ。」\hfill --- Robert A. Heinlein
\end{quote}

HTTP プロトコルそのものが Web の中心にあるので、Web セキュリティに踏み込む前に、その仕組みを簡単に振り返っておく価値がある。
ここでのごく単純化した説明は、以降のセキュリティの議論の概念的な枠組みとなるもので、セキュリティが関わってくる部分に焦点を当てる。
多くの人にとって Web の閲覧は日常生活のありふれた一部になっているので、一歩引いて、その過程を一段階ずつ考えてみる価値がある。現代のプロセッサーとネットワークはいつも猛烈な速さで応答を返すので、その段階の多くは、ふだんほとんど意識されることがない。

Web の閲覧は、必ず \underline{URL}（uniform resource locator）から始まる。
次の例で、その各部分を示す。
\begin{verbatim}
http://www.example.com/page.html?query=value#fragment
\end{verbatim}
\underline{スキーム}（scheme）はコロンの前の部分で、要求するリソースを得るためにブラウザが使うべきプロトコル（ここでは http）を指定する。
IP ベースのプロトコルは \texttt{//} で始まり、その後に\underline{ホスト名}（hostname）が続く。Web ページの場合、ホスト名は Web サーバーのドメイン名（この例では \texttt{www.example.com}）である。
残りはすべて省略可能で、\texttt{/} に続く\underline{パス}（path）、\texttt{?} に続く\underline{クエリ}（query）、\texttt{\#} に続く\underline{フラグメント}（fragment）からなる。
パスは、ブラウザがどの Web ページを要求しているかを指定する。
クエリを使うと、Web ページの内容をパラメーターで変えられる。
例えば、Web で何かを検索すると、結果の URL のパスは \texttt{/search?q=something} のようになる。
フラグメントはページ内の副次的なリソースを指すもので、多くの場合はリンクの行き先となるアンカーである。
まとめると、URL は、内容をどのように、どこへ要求するか、サイト上のどのページか、ページをカスタマイズするクエリパラメーター、そしてページの特定の部分を指す方法を指定する。

URL を渡されてから Web ページを表示するまでに、ブラウザがしなければならない作業はたくさんある。
まず、要求の送り先を知るために、ホスト名の IP アドレスを Domain Name System（DNS）に問い合わせる。
要求には、URL のパスと、要求ヘッダー（request header）としてエンコードされたその他のパラメーター（クッキー、ユーザーの優先言語など）が含まれ、Web サーバーのホストに送られる。
サーバーは、ステータスコードと応答ヘッダー（クッキーの設定をはじめ、さまざまなことを行う）を含む応答を返し、続いて Web ページの HTML からなる内容本体（content body）を返す。
スクリプトや画像などの埋め込みリソースのそれぞれについても、内容がすべて読み込まれて表示されるまで、同じ要求／応答の過程が繰り返される。

では、Web サーバーがセキュアであり続けるために、何を正しく行わなければならないかを見ていこう。
まだ触れていない重要な細部の一つに、要求が \underline{HTTP 動詞}（HTTP verb）を指定するという点がある。
ここでは、そうした動詞のうち最もよく使われる 2 つだけを取り上げる。
\texttt{GET} はサーバーに内容を要求する。
一方、クライアントはフォームの送信やファイルのアップロードには \texttt{POST} 動詞を使う。
\texttt{GET} 要求は明示的に状態を変更しないものとされ、\texttt{POST} 要求はサーバーの状態を変更することを意図している。
CSRF 攻撃を扱うときにわかるように、この意味上の区別を守ることは重要である。
今のところは、要求に使う動詞を指定するのはクライアントでも、その要求をどう扱うかを決めるのはサーバーだということを覚えておいてほしい。
さらに、サーバーは自分のページにハイパーリンクやフォームを置くことで、事実上、クライアントが次に行う \texttt{GET} 要求や \texttt{POST} 要求を誘導している。

杓子定規な人なら、\texttt{GET} 動詞の要求に応じて状態を変更し、あまのじゃくにもフォームの \texttt{POST} 送信では状態の変更を拒むサーバーを動かすことも確かにできる、と指摘するだろう。
しかし、標準の規則に厳密に従ったほうが、サーバーをセキュアにするのはずっと簡単になる。
こう考えるとよい。展望台の崖で「立入禁止！」と書かれた柵を乗り越え、落ちずに断崖の縁を歩くことは確かにできるが、それは必要もなく自分の身を危険にさらす行為である。

\begin{rem}[機密データを URL に埋め込まない]\label{rem:ch11-no-secrets-in-url}
これに関連するセキュリティ上のタブーが、機密データを URL に埋め込むことである。データをサーバーに送るには、代わりにフォームの \texttt{POST} 要求を使うべきである。
そうしないと、\texttt{REFERER} ヘッダーによって、要求の発端となった Web ページの URL が開示され、データが露出するおそれがある。
例えば、URL が \texttt{https://example.com?param=SECRET} の Web ページ上のリンクをクリックすると、SECRET を含むその URL が入った \texttt{REFERER} ヘッダー付きの \texttt{GET} 要求でリンク先に移動するので、秘密のデータが漏えいする。
さらに、ログや診断メッセージから URL に含まれるデータが開示されるリスクもある。
サーバーは \texttt{Referrer-Policy} ヘッダーを使ってこれを防げるが、クライアントがこのヘッダーに従ってくれることに頼るしかなく、完璧な解決策とはとても言えない。
（\texttt{REFERER} ヘッダーは仕様の中で実際に綴りを間違えており、もはやそのまま使うしかないが、ポリシーの名前のほうは正しく綴られている。）
\end{rem}

\begin{rem}[URL にユーザー名を含めない]\label{rem:ch11-no-usernames-in-url}
犯しやすい誤りの一つに、URL にユーザー名を含めることがある。
ユーザー名のハッシュのような不透明な識別子であっても、情報は漏えいする。盗聴者が別々に観察した 2 つの URL を突き合わせ、同じユーザーを指していると推測できてしまうおそれがあるからである。
\end{rem}

%----------------------------------------------------------------------
\subsubsection{デジタル証明書と HTTPS}\label{subsec:ch11-https}
% src: Digital Certificates and HTTPS

\begin{quote}
「伝えられることが偽りであるなら、それはとうてい伝達（communication）とは呼べない。」\hfill --- Benjamin Mays
\end{quote}

セキュアに Web を閲覧するための最初の課題は、正しいサーバーと確実に通信することである。
そのためには、正しい URL を知り、正しい IP アドレスを返してくれる DNS サービスに問い合わせなければならない。
そのうえでネットワークが要求を正しく経路制御して送信すれば、要求は意図したサーバーに届くはずである。
正しく働かなければならない要素がこれだけ多いということは、攻撃対象領域も大きいということである。悪意ある者は、DNS の問い合わせにも、経路制御にも、経路上のどの地点でも回線上のデータにも干渉できる。
要求が悪意あるサーバーにすり替えられても、ユーザーはまず気付かないかもしれない。ほとんど誰でも簡単にだませるそっくりな Web サイトを立ち上げるのは、難しいことではない。

\underline{HTTPS} プロトコル（HTTP over TLS/SSL とも呼ばれる）は、まさにこうした脅威を緩和するために作られたものである。
HTTPS は、第5章で扱った技法の多くを使って Web をセキュアにする。
HTTPS は、端から端まで改ざんから守られた、暗号化されたセキュアなチャネルを提供し、さらに、そのチャネルの向こう側に意図したサーバーが本当にいることをクライアントに保証する。
セキュアなチャネルは、サーバーのアイデンティティを確認したうえでデータを流す、改ざんされればすぐにわかるパイプラインだと考えるとよい。
盗聴する攻撃者は暗号化されたデータを目にできるかもしれないが、秘密鍵がなければ、そのデータはランダムなビットと区別がつかない。
攻撃者は保護されていないネットワーク上でデータを改ざんできるかもしれないが、HTTPS を使っていれば、どんな改ざんも必ず検出される。
攻撃者は、例えばケーブルを物理的に切断して通信を妨げることはできるかもしれないが、偽のデータが通り抜けることは決してないと保証される。

Web 上の金融取引をセキュアにするのに HTTPS が必要だという点に異論を唱えた者はいなかったが、主要なサイトはあまりにも長いあいだ完全な HTTPS 化を先延ばしにした。
（例えば、Facebook がそれを果たしたのは 2013 年になってからである。）
最初に実装された当時、このプロトコルには微妙な欠陥があり、必要な計算も当時のハードウェアには重すぎて、広く採用するだけの理由がなかった。
幸い、時がたつにつれて開発者がバグを修正し、プロトコルを最適化していった。
より効率的な暗号アルゴリズムとより高速なプロセッサーのおかげで、今日の HTTPS は高速かつ堅牢になり、急速にどこでも使われるものになりつつある。
HTTPS はプライベートなデータ通信の保護に広く使われているが、公開情報だけを提供する Web サイトであっても、真正性と強い完全性を確保するうえで HTTPS は重要である。
言い換えれば、HTTPS は、クライアントが要求 URL で指定された正真正銘のサーバーと通信していること、そして両者の間でやり取りされるデータがのぞき見も改ざんもされていないことを保証する。
今日では、Web サイトを HTTPS 専用に設定しない正当な理由は、なかなか思いつかない。
とはいえ、セキュアでない HTTP の Web サイトはまだたくさんある。
そうしたサイトを使う場合は、HTTPS の優れたセキュリティ特性が当てはまらないことを心に留め、適切な予防措置をとるべきである。

HTTPS がクライアントとサーバーのやり取りをセキュアにするために何をするのか（そして何をしないのか）を正確に理解することは、その価値、つまりどう役立ち、何を変えられて何を変えられないのかを把握するうえで欠かせない。
セキュアなチャネルは、サーバーの真正性と、Web の要求と応答の内容の機密性・完全性を保証するだけでなく、URL のパス（要求ヘッダーの最初の行にあるもの。例えば \texttt{GET /path/page.html?query=secret\#fragment}）も保護するので、のぞき見している者には、クライアントが Web サイトのどのページを要求したかがわからない。
（HTTPS では、オプションでクライアントをサーバーに対して認証することもできる。）
ただし、HTTPS のトラフィックそのものはネットワーク上で依然として観察でき、端点の IP アドレスも保護されないので、盗聴者はサーバーのアイデンティティを推測できることが多い。

表 \ref{tab:ch11-http-vs-https} は、クライアント／サーバー通信の 2 つの端点の間に潜む攻撃者に何ができるかという観点から、HTTP と HTTPS のセキュリティ上の性質を比べたものである。

\begin{table}[H]
\centering
\caption{HTTP と HTTPS のセキュリティ上の性質}\label{tab:ch11-http-vs-https}
\small
\begin{tabularx}{\textwidth}{|X|c|l|}
\hline
攻撃者は\ldots & HTTP & HTTPS \\
\hline
クライアント／サーバーの端点間の Web トラフィックを見られるか？ & はい & はい \\
\hline
クライアントとサーバーの両方の IP アドレスを特定できるか？ & はい & はい \\
\hline
Web サーバーのアイデンティティを推測できるか？ & はい & 場合による（下の注を参照） \\
\hline
サイト内のどのページが要求されたかを見られるか？ & はい & いいえ（暗号化されたヘッダー内） \\
\hline
Web ページの内容と POST の本体を見られるか？ & はい & いいえ（暗号化） \\
\hline
ヘッダー（クッキーを含む）と URL（クエリ部分を含む）を見られるか？ & はい & いいえ \\
\hline
URL、ヘッダー、内容を改ざんできるか？ & はい & いいえ \\
\hline
\end{tabularx}
\end{table}

\begin{rem}[サーバーのアイデンティティの推測]\label{rem:ch11-server-identity}
Web サーバーの IP アドレスを DNS で逆引きすれば、そのドメイン名がわかる。
複数の Web サーバーが 1 つの IP アドレスを共有している場合、SNI（Server Name Indication）は見えるが、ESNI（Encrypted SNI）なら保護される。
\end{rem}

HTTPS と技術環境が成熟してくると、広く普及するうえで最後に残った障害は、サーバー証明書を取得する手間になった。
大企業なら信頼された認証局が課す料金を払う余裕があり、更新作業を管理するスタッフもいたが、小規模な Web サイトの所有者は、余分な費用と管理の手間に二の足を踏んだ。
2015 年には HTTPS は成熟し、インターネットに接続されたハードウェアのほとんどはそれを処理できるだけの速さで動くようになっていた。Web のプライバシーが重要だという認識も急速に広がる中、インターネットコミュニティは、Web トラフィックの大部分をセキュアにする必要があるという合意に近づいていた。
最大の障害として残ったのは、無料で簡単に手に入るサーバー証明書がないことだった。

すばらしい Electronic Frontier Foundation の強力な推進と、幅広い業界の企業の後援のおかげで、今日では、非営利団体 Internet Security Research Group の製品である \underline{Let's Encrypt} が、無料で自動化されたオープンな認証局を世界に提供している。
Let's Encrypt は、\underline{ドメイン認証}（Domain Validation, DV）のデジタル証明書を、どんな Web サイトの所有者にも無料で発行する。

\begin{exa}[Let's Encrypt の仕組み]\label{exa:ch11-lets-encrypt}
Let's Encrypt の仕組みを単純化して説明すると次のようになる。
なお、実際には、この過程はふつう完全に自動化されている。
\begin{enumerate}
\item 鍵ペアを生成して公開鍵を送り、Let's Encrypt に対して自分を識別させる。
\item ドメインを自分が管理していることを証明するには何をすればよいかを、Let's Encrypt に問い合わせる。
\item Let's Encrypt は、ドメインに指定の DNS レコードを用意せよ、といったチャレンジを出す。
\item 要求された DNS レコードを作成してチャレンジに応え、自分が行ったことを検証するよう Let's Encrypt に依頼する。
\item 検証が済むと、生成した鍵ペアの秘密鍵が、Let's Encrypt によってそのドメインについて認可される。
\item 以後は、認可された秘密鍵で署名した要求を Let's Encrypt に送れば、新しい証明書を要求できる。
\end{enumerate}
\end{exa}

Let's Encrypt は有効期間 90 日の DV 証明書を発行し、自動更新を担う「certbot」を提供している。
自動更新できる証明書が無料のサービスとして手に入るようになったことで、今日では、セキュアな Web の提供は追加費用なしで導入できる手軽な仕組みとして広く普及した。
2020 年には HTTPS が Web トラフィックの 85\% 以上を占めるようになり、Let's Encrypt が始まった 2016 年の 40\% から 2 倍以上に増えている。

Web サイトのアイデンティティを証明するには、たいてい DV 証明書があれば十分である。
DV のデジタル証明書は、認証された Web サーバーのドメイン名を主張するだけで、それ以上のことは何もしない。
つまり、\texttt{example.com} のデジタル証明書は、\texttt{example.com} の Web サーバーの所有者にしか発行されない。
これに対して、\underline{組織認証}（Organization Validation, OV）証明書や\underline{拡張認証}（Extended Validation, EV）証明書のように、より高い水準の信頼を提供するデジタル証明書は、Web サイトのアイデンティティだけでなく、所有者のアイデンティティと評判もある程度まで認証する。
しかし、無料の DV 証明書が普及したことで、ほかの種類の証明書が今後も生き残れるかどうかはますます怪しくなっている。
そうした信頼の違いを気にするユーザーはほとんどおらず、OV 証明書と EV 証明書の技術的・法的な違いも微妙である。
それぞれの利点を正確に理解するのは、弁護士でもない限り（いや、弁護士であっても）難しい。

デジタル証明書を使って Web サーバーを HTTPS プロトコルで動くように設定したら、常に HTTPS が使われるようにしなければならない。
そのためには、\underline{ダウングレード攻撃}（downgrade attack）を拒否しなければならない。ダウングレード攻撃とは、通信を弱い暗号化で、あるいは暗号化なしで行わせようとする攻撃である。

\begin{exa}[ダウングレード攻撃の 2 つの方法]\label{exa:ch11-downgrade}
ダウングレード攻撃には 2 つの方法がある。
\begin{itemize}
\item 最も単純なのは、攻撃者が HTTPS の要求を（のぞき見も改ざんもできる）HTTP に変えようとするもので、設定の甘い Web サーバーはだまされてこれに応じてしまうかもしれない。
\item もう一つは、暗号化されたチャネルで使う暗号スイート（cipher suite）を 2 者の間で交渉できるようにする、HTTPS プロトコルのオプションを悪用する方法である。
例えば、サーバーはある暗号の「方言」の組を「話せ」、クライアントは別の組を「話す」かもしれないので、両者はまず、双方のレパートリーに共通する方言について合意する必要がある。
この過程が攻撃者に付け入る隙を与える。攻撃者は両者をだまして安全でない選択肢を選ばせ、セキュリティを損なうことができる。
\end{itemize}
\end{exa}

最善の防御は、HTTPS の設定を、安全な現代的暗号アルゴリズムでしか動かないようにしておくことである。
どの暗号スイートが安全かを正確に判断するのは非常に専門的な作業であり、暗号学者に任せるのが一番よい。
また、古くて非力なクライアントを締め出したり、その使い勝手を損なったりしないよう、バランスもとらなければならない。
信頼できる専門家の助言が得られない場合は、信頼できる主要な Web サイトのやり方を調べて、それにならえばよい。
デフォルトの設定がいつまでもセキュアだと安易に思い込むのは、失敗のもとである。

こうした攻撃は、HTTP を常に HTTPS にリダイレクトし、Web クッキーを HTTPS 限定にすることで緩和する。
応答の HTTP ヘッダーに \texttt{Strict-Transport-Security} ディレクティブを含め、この Web サイトでは常に HTTPS を使うことをブラウザに知らせる。
HTTPS の Web ページが完全にセキュアであるためには、純粋に HTTPS だけでできていなければならない。
つまり、サーバー上のすべての内容が HTTPS を使うべきであり、スクリプト、画像、フォント、CSS、その他の参照するリソースもすべて同様である。
必要な予防措置を一つでも怠れば、セキュリティ上の保護は弱まる。

%----------------------------------------------------------------------
\subsubsection{同一オリジンポリシー}\label{subsec:ch11-sop}
% src: The Same Origin Policy

\begin{quote}
「疑いは知恵の源（origin）である。」\hfill --- Rene Descartes
\end{quote}

\begin{dfn}[同一オリジンポリシー]\label{dfn:ch11-same-origin-policy}
ブラウザは、異なる Web サイトのリソース（典型的にはウィンドウやタブ）を、互いに干渉できないよう隔離する。
この規則は\underline{同一オリジンポリシー}（Same Origin Policy）と呼ばれ、ホストのドメイン名とポート番号が一致する場合にだけ、リソース間のやり取りを許可する。
\end{dfn}

同一オリジンポリシーの起源は Web の初期にさかのぼり、JavaScript の登場とともに必要になった。
Web のスクリプトは、\underline{Document Object Model}（DOM）を介して Web ページとやり取りする。これは、ブラウザのウィンドウとその内容に対応するオブジェクトを木構造にまとめたものである。
どんな Web ページでも、スクリプトでほかの任意のサイトを \texttt{window.open} し、その内容をプログラムで好き勝手に操作できるとしたら、数え切れない問題が起きることは、セキュリティの専門家でなくてもわかった。
最初に実装された制限は、長年のあいだに見つかったさまざまな巧妙な抜け道への修正も取り込みながら、今日の同一オリジンポリシーへと進化してきた。

同一オリジンポリシーは、スクリプトとクッキー（クッキーにはいくつか追加のひねりがある）に適用される。どちらも、独立した Web サイトの間でデータを漏えいさせるおそれがあるからである。
一方、Web ページには、Web 広告などの画像やその他の内容を、ほかの Web サイトから取り込むことができる。
こうした内容は、自分が表示されているウィンドウの内容にアクセスできないので、取り込んでも安全だとして許可されている。

同一オリジンポリシーは、ほかの Web サイトのページのスクリプトが入り込んでくるのは防ぐが、Web ページのほうからは、望めばいつでも別の Web サイトに手を伸ばし、その内容をウィンドウに取り込める。
Web ページがほかの Web サイトの内容を取り込んで、画像を表示したり、スクリプトや CSS を読み込んだりするのはごく一般的なことである。
しかし、ほかの Web サイトの内容を取り込むことは、重要な信頼の判断である。取り込んだ先から悪意ある内容がやって来れば、その Web ページは無防備になるからである。

%----------------------------------------------------------------------
\subsubsection{Web クッキー}\label{subsec:ch11-cookies}
% src: Web Cookies

\begin{quote}
「事態が厳しくなると、たくましい者はクッキーを焼く。」\hfill --- Erma Bombeck
\end{quote}

\begin{dfn}[クッキー]\label{dfn:ch11-cookie}
サーバーがクライアントに預けておき、以後の要求のたびにサーバーへ送り返してもらう小さなデータ文字列を、\underline{クッキー}（cookie）という。
\end{dfn}

この巧妙な発明のおかげで、開発者は Web ページを特定のクライアント向けに簡単にカスタマイズできる。
サーバーは応答の中で、名前付きのクッキーに何らかの値を設定できる。
以後、クッキーの期限が切れるまで、クライアントのブラウザは、各要求でそのページに該当するクッキーを送る。
% 注: 原文 "applicable to a give page" は "a given page" の誤植とみなした。
クッキーはクライアントが自分で保持するので、サーバーはクッキーの値をクライアントに結び付けるためにクライアントを識別する必要が必ずしもない。その意味で、この仕組みはプライバシーを守るものにもなりうる。

\begin{exa}[来店回数を数える紙片]\label{exa:ch11-cookie-analogy}
簡単なたとえで説明する。
私が店を経営していて、顧客ごとの来店回数を数えたいとすると、手軽なのは、顧客一人ひとりに「1」と書いた紙片を渡し、次に来るときに持ってきてもらう方法である。
顧客が再び来るたびに、私はその紙を受け取り、書かれた数に 1 を足して返す。
顧客が協力してくれる限り、帳簿を付けなくても、顧客の名前を覚えていなくても、正確な回数を数え続けられる。
\end{exa}

Web ではクッキーをあらゆることに使っているが、中でも最も議論を呼んでいるのがユーザーの追跡である。
クッキーはセキュアなセッションを確立するためによく使われ、これによってサーバーはすべてのクライアントを確実に区別できる。
新しいクライアントごとに一意のセッションクッキーを生成しておけば、サーバーは要求に含まれるクッキーからクライアントを識別できる。

どのクライアントも自分のクッキーを改ざんして別のセッションになりすますことはできるが、セッションクッキーが適切に設計されていれば、有効なセッションクッキーを偽造することはできないはずである。
また、クライアントは自分のクッキーのコピーを第三者に渡すこともできるが、そうして損なわれるのは自分自身のプライバシーだけである。
そうした振る舞いは無実のユーザーを脅かすものではなく、自分のパスワードを人に教えるのと変わらない。

\begin{exa}[クッキーに保存したショッピングカート]\label{exa:ch11-cart-cookie}
顧客のショッピングカートの現在の中身を、商品のリストと合計金額としてクッキーに保存する、架空のオンラインショッピング Web サイトを考えよう。
悪知恵の働く不誠実な買い物客がローカルのクッキーストアを書き換えようとしても、止めるものは何もない。
例えば、高価な商品一式の価格を、ほんのわずかな金額に書き換えられてしまう。
だからといってクッキーが役に立たないわけではない。顧客の好みやお気に入りの商品などを覚えておくのにクッキーを使うことはでき、それらが改ざんされても販売者は困らない。
要するに、クライアント側の保存領域は、常に「信頼せよ、されど検証せよ」（trust but verify）の原則で使うべきだということである。
役に立つなら商品の価格やカートの合計をクライアントに保存してもかまわないが、取引を受け付ける前には、必ずサーバー側で各商品の価格を検証し、改ざんされたデータはすべて拒否すること。
この例では問題は一目瞭然である。
しかし、同じ信頼の誤りでも、ほかの形をとるものはもっと見えにくく、攻撃者はこの種の脆弱性を頻繁に悪用する。
\end{exa}

では、同じ例をクライアントの視点から見てみよう。
2 人がオンラインショッピングの Web サイトで同じ \texttt{/mycart} という URL を開くと、それぞれが別のセッションを持っているので、表示されるショッピングカートも異なる。
通常は、一意のクッキーによって独立した匿名のセッションが確立されるか、ログインしたユーザーの場合はクッキーによって特定のアカウントが識別される。

サーバーはセッションクッキーに有効期限を設定するが、クライアントがその指定を必ず守ってくれるとは限らないので、更新が必要なセッションクッキーについては、有効性の制限をサーバー側でも強制しなければならない。
（ユーザーから見ると、この期限切れは、しばらく操作しないでいると再びログインを求められる、という形で現れる。）

クッキーは同一オリジンポリシーの対象であり、サブドメイン間での共有については明示的な規定がある。
つまり、\texttt{example.com} が設定したクッキーはサブドメインの \texttt{cat.example.com} と \texttt{dog.example.com} から見えるが、それぞれのサブドメインで設定されたクッキーは互いに隔離されている。
また、サブドメインは親ドメインが設定したクッキーを見ることはできても、変更することはできない。
たとえて言えば、州政府はパスポートのような国が発行する資格証明を利用するが、自分で発行することはできない、というのと同じである。
ドメイン内では、クッキーの範囲をパスでさらに絞ることもできる（ただし、これは強力なセキュリティの仕組みではない）。
表 \ref{tab:ch11-cookie-sharing} に、これらの規則を詳しく示す。
さらに、クッキーで \texttt{Domain} 属性を指定して、明示的に制御することもできる。

% 注: 原文の表 11-2 は Markdown の表が崩れて 1 行につながっているが、セルの内容はすべて読み取れるので表として再構成した。
\begin{table}[H]
\centering
\caption{サブドメインがある場合の同一オリジンポリシー（SOP）のもとでのクッキーの共有}\label{tab:ch11-cookie-sharing}
\small
\begin{tabular}{|l|c|c|c|c|}
\hline
\multirow{2}{*}{\shortstack[l]{下のホストが提供する\\Web ページは\ldots}} & \multicolumn{4}{c|}{\ldots これらのホストに設定されたクッキーを見られるか？} \\
\cline{2-5}
 & \texttt{example.com} & \texttt{dog.example.com} & \texttt{cat.example.com} & \texttt{example.org} \\
\hline
\texttt{example.com} & \shortstack{はい\\（同じドメイン）} & \shortstack{いいえ\\（サブドメイン）} & \shortstack{いいえ\\（サブドメイン）} & \shortstack{いいえ\\（SOP）} \\
\hline
\texttt{dog.example.com} & \shortstack{はい\\（親ドメイン）} & \shortstack{はい\\（同じドメイン）} & \shortstack{いいえ\\（兄弟ドメイン）} & \shortstack{いいえ\\（SOP）} \\
\hline
\texttt{cat.example.com} & \shortstack{はい\\（親ドメイン）} & \shortstack{いいえ\\（兄弟ドメイン）} & \shortstack{はい\\（同じドメイン）} & \shortstack{いいえ\\（SOP）} \\
\hline
\texttt{example.org} & \shortstack{いいえ\\（SOP）} & \shortstack{いいえ\\（SOP）} & \shortstack{いいえ\\（SOP）} & \shortstack{はい\\（同じドメイン）} \\
\hline
\end{tabular}
\end{table}

\begin{rem}[httponly 属性と secure 属性]\label{rem:ch11-cookie-attributes}
スクリプトからは建前上 DOM を介してクッキーにアクセスできるが、この便利さは、Web ページの中でまんまと実行された悪意あるスクリプトにクッキーを盗む隙を与えてしまう。そこで、\texttt{httponly} クッキー属性を指定して、スクリプトからのアクセスを阻止するのが最善である。
HTTPS の Web サイトでは、さらに \texttt{secure} 属性も付けて、クッキーをセキュアなチャネルでのみ送るようクライアントに指示すべきである。
残念ながら、ここで扱うには込み入りすぎたレガシーの制約があるため、両方の属性を使っても完全性と可用性の問題は残る（おぞましい詳細は RFC 6265 を参照）。
筆者がこの点に触れるのは、注意を促すためだけでなく、Web セキュリティで繰り返し現れるパターンの好例だからでもある。後方互換性と現代のセキュアな使い方とがせめぎ合うと、妥協の解決策が生まれる。そうした解決策を見ると、セキュリティは最初から組み込んでおかなければ、後から実現するのが難しいことが多いとよくわかる。
\end{rem}

HTML5 では、セキュリティモデルが数多く拡張された。
代表的なものが \underline{Cross-Origin Resource Sharing}（CORS）で、ほかの信頼できる Web サイトがデータにアクセスできるよう、同一オリジンポリシーの制限を選択的に緩められる。
さらにブラウザは、Web アプリ向けのより現代的なクライアント側の保存機能として Web Storage API も提供しており、これも同一オリジンポリシーの対象となる。
こうした新しい機能はセキュリティの観点からずっとよく設計されているが、それでもクッキーを完全に置き換えられるわけではない。

%----------------------------------------------------------------------
\subsection{よくある Web の脆弱性}\label{subsec:ch11-common-vulns}
% src: Common Web Vulnerabilities

\begin{quote}
「Web サイトは内から見ても外から見ても見栄えがよくあるべきだ。」\hfill --- Paul Cookson
\end{quote}

ここまでで、Web サイトを作り、使ううえでのセキュリティ上の主な要点をひととおり見てきたので、いよいよ、よく起きる具体的な脆弱性の話に移ろう。
Web サーバーは、本書のほかの箇所で扱った多くの脆弱性を含め、あらゆる種類のセキュリティ脆弱性を抱えうるが、この章では Web に特有のセキュリティ問題に絞る。
ここまでの節では Web セキュリティモデルを説明し、セキュリティを弱めないための多くの方法や、Web 上の自分の拠点をよりセキュアにするのに役立つ機能を紹介した。
この節では、そうしたことをすべて正しく行ったとしても、Web サーバーが誤りを犯して脆弱になる道筋がまだあることを見ていく。

Web の脆弱性の最初の分類であり、おそらく最もよく見られるのが、\underline{クロスサイトスクリプティング}（cross-site scripting, XSS）である。
ここで扱うもう一つの脆弱性は、筆者のお気に入りといってよい\underline{クロスサイトリクエストフォージェリ}（cross-site request forgery, CSRF）である。

%----------------------------------------------------------------------
\subsubsection{クロスサイトスクリプティング}\label{subsec:ch11-xss}
% src: Cross-Site Scripting

\begin{quote}
「私は Web で『サーフィン』しないようにしている。さもないと、たぶん溺れてしまうから。」\hfill --- Aubrey Plaza
\end{quote}

同一オリジンポリシーによる隔離は、セキュアな Web サイトを作るうえでの土台であるが、必要な予防措置をとらなければ、この保護はたやすく崩れる。

\begin{dfn}[クロスサイトスクリプティング]\label{dfn:ch11-xss}
\underline{クロスサイトスクリプティング}（XSS）は、悪意ある入力によって Web サイトの振る舞いを変える、Web に特有のインジェクション攻撃であり、典型的には、認可されていないスクリプトが実行される結果になる。
\end{dfn}
% 注: 原文は "*Cross-site scripting (*XSS)" と強調記号が崩れている。

この攻撃がどのように機能し、なぜ防ぐことが欠かせないのかを、簡単な例で見てみよう。
攻撃はたいてい、無実のユーザーが信頼できる Web サイトにすでにログインしている状態から始まる。
ユーザーはその後、別のウィンドウやタブを開いてネットサーフィンをしたり、不用意にメールのリンクをクリックしたりして、攻撃サイトにたどり着く。
攻撃者の典型的な狙いは、標的サイトでのユーザーの認証済みの状態を乗っ取ることである。
クッキーさえ残っていれば、被害者サイトのタブが開いていなくても乗っ取りは可能である（だからこそ、銀行の Web サイトを使い終えたらログアウトする習慣が大切なのである）。
被害者サイトの XSS 脆弱性がどのようなものか、それを具体的にどう悪用するのか、そして最後にどう修正するのかを見ていこう。

\underline{脆弱な例：}
被害者の Web サイト（\texttt{www.example.com}）のあるページで、何らかの理由から、1 行のテキストをいくつかの色で表示し分けたいとする。
開発者は、その行の色だけが違う同じページを何枚も作る代わりに、表示したい色を URL のクエリパラメーターで指定することにする。
例えば、テキストの行が緑色の版の Web ページの URL は次のようになる。
\begin{verbatim}
https://www.example.com/page?color=green
\end{verbatim}
サーバーは、このクエリパラメーター（\texttt{green}）を次の HTML 断片に挿入する。
\begin{verbatim}
<h1 style="color:green">This is colorful text.</h1>
\end{verbatim}
% 注: 原文では挿入されるパラメーターが強調表示されているが、スクレイプで強調は失われている。
正しく使われている限り、これは問題なく動く。だからこそ、この種の欠陥は見過ごされやすいのである。
問題の根本を見抜くには、この処理を担当するサーバー側の Python コードを見る必要がある（ついでに、少々ひねくれた発想も必要である）。

脆弱なコード：
\begin{verbatim}
query_params = urllib.parse.parse_qs(self.parts.query)
color = query_params.get('color', ['black'])[0]
h = '<h1 style="color:%s">This is colorful text.</h1>' % color
\end{verbatim}
1 行目は URL のクエリ文字列（疑問符より後の部分）を解析する。
次の行は \texttt{color} パラメーターを取り出し、指定がなければ黒（\texttt{black}）をデフォルトにする。
最後の行は、見出しレベル 1 のタグ（\texttt{<h1>}）にインラインのスタイル指定を付け、指定された色でテキストを表示する HTML 断片を組み立てる。
変数 \texttt{h} は、その後、Web ページを構成する HTML 応答の一部になる。

XSS 脆弱性は、この最後の行にある。
ここでプログラマーは、URL の内容（インターネット上では誰でもサーバーに送れる）から、クライアントに返す HTML の内容へと直結する経路を作ってしまっている。
これは第10章で見たおなじみのインジェクション攻撃のパターンであり、パラメーターの入力文字列が Web ページの HTML の内容の中に入り込むので、保護されていない信頼境界の越境にあたる。
この状態だけでも危険信号としては十分だが、この XSS 脆弱性の全貌をつかむために、実際に悪用してみよう。

\underline{悪用：}
攻撃には少し想像力が要る。
\texttt{<h1>} HTML タグに戻り、強調表示された色の名前の代わりに、ほかに何を入れられるかを考えてみよう。
枠にとらわれずに、この場合は二重引用符で囲まれた文字列 \texttt{style="color:green"} の外側まで考えてみることである。
あるいは、\texttt{<h1>} タグから完全に抜け出せないだろうか。
「抜け出す」とは、例えば次のようなことである。
{\footnotesize
\begin{verbatim}
https://www.example.com/page?color=orange"><SCRIPT>alert("Gotcha!")</SCRIPT><span%20id="dummy
\end{verbatim}
}
強調表示されたこの部分がすべて、前と同じように \texttt{<h1>} HTML タグに律儀に挿入され、まったく異なる結果になる。

実際の HTML ではこのコードは 1 行に並ぶが、ここでは読みやすさのために、どのように解析されるかがわかるようインデントしてある。
\begin{verbatim}
<h1 style="color:orange">
<SCRIPT>alert("Gotcha!")</SCRIPT>
<span id="dummy">This is colorful text.
</h1>
\end{verbatim}
% 注: 原文のインデントはスクレイプで失われている。

新しい \texttt{<h1>} タグは構文的に正しく、オレンジ色を指定している。
しかし、閉じ山括弧を供給しているのが攻撃者の URL パラメーターの値であることに注意してほしい。
これは親切心からではない。整形式の \texttt{<SCRIPT>} タグを作って HTML に注入し、スクリプトを確実に実行させるには、攻撃者は \texttt{<h1>} タグを閉じる必要があったのである。
この例では、スクリプトは警告ダイアログを開く。無害ではあるが、悪用に成功したことを示す紛れもない証拠である。
閉じタグ \texttt{</SCRIPT>} より後のインジェクション部分は、改ざんの痕跡を目立たなくするための詰め物にすぎない。
新しい \texttt{<span>} タグに \texttt{id} 属性があるのは、後ろに続く二重引用符と閉じ山括弧を \texttt{<span>} タグの一部として収めるためだけである。
ブラウザは欠けている閉じタグ \texttt{</span>} をふだんから補うので、悪用されたページは整形式の HTML になり、（HTML のソースを調べない限り）ユーザーには変更が見えない。

実際に被害者を遠隔から攻撃するには、攻撃者は、人々を悪意ある URL に誘導するために、さらに手間をかけなければならない。
この種の攻撃がうまくいくのは、一般に、ユーザーが標的の Web サイトですでに認証されているとき、つまり有効なログインセッションのクッキーがあるときだけである。
そうでなければ、攻撃者が自分のブラウザにその URL を打ち込んでも同じことである。
攻撃者が狙っているのは、銀行の残高やプライベートな文書が表示される、ユーザーの Web サイトのセッションである。
本気の攻撃者が仕込むスクリプトなら、ただちに追加のスクリプトを読み込み、データを持ち出したり、ユーザーになりかわって認可されていない取引を行ったりするだろう。

攻撃者にとって、XSS 脆弱性を見つけるのは難しくない。Web ページの内容を見れば、HTML の内部の仕組みが簡単にわかるからである。
（正確に言えば、攻撃者にはサーバー上のコードは見えないが、いろいろな URL を試して返ってくる Web ページを観察すれば、仕組みについて役に立つ推測をするのは難しくない。）
URL の内容が Web ページに注入されていることに気付いたら、ここで示した例のような簡単なテストで、サーバーが XSS に対して脆弱かどうかを確かめられる。
さらに、山括弧や引用符のような HTML のメタ文字が、URL のクエリパラメーター（あるいは別の攻撃対象領域）から結果の Web ページへ流れ込むことを確認できれば、あとはページのソースコードを見ながら、大当たりを引くまで試行錯誤すればよい。

\begin{dfn}[XSS 攻撃の種類]\label{dfn:ch11-xss-kinds}
XSS 攻撃にはいくつかの種類がある。
\begin{itemize}
\item \underline{反射型 XSS 攻撃}（reflected XSS attack）：HTTP 要求によって始まり、サーバーが直後に返す応答の中に現れるもの。この章の例はこれにあたる。
\item \underline{格納型 XSS 攻撃}（stored XSS attack）：2 つの要求を伴うもの。まず、攻撃者は何らかの方法で、サーバー上かクライアント側の保存領域に悪意あるデータを保存する。その準備が整うと、続く要求で Web サーバーをだまし、保存しておいたデータをその後の要求に注入させて、攻撃を完成させる。格納型 XSS 攻撃は、異なるクライアントをまたいで成立しうる。
\item \underline{DOM ベースの XSS}（DOM-based XSS）：HTML の Document Object Model を悪意あるインジェクションの源として使うもの。それ以外の点は、ほかの種類とほぼ同じように機能する。
\end{itemize}
\end{dfn}

\begin{exa}[ブログのコメントによる格納型 XSS]\label{exa:ch11-stored-xss-blog}
例えばブログで、コメントの表示時に XSS を引き起こすようなコメントを攻撃者が投稿できるなら、その後にその Web ページを閲覧したユーザーは悪意あるスクリプトを受け取ってしまう。
\end{exa}

分類はともかく、要するに、これらの脆弱性はすべて、Web サーバーが信頼できないデータを Web ページに流れ込ませてしまい、悪意あるスクリプトやその他の有害な内容が持ち込まれることから生じる。

\underline{防御：}
セキュアな Web フレームワークには XSS 保護が組み込まれているはずであり、その場合、フレームワークの範囲内で作業している限りは安全なはずである。
ほかのインジェクション脆弱性と同様に、防御策は、信頼できない入力が Web ページに流れ込んで抜け出す可能性を一切なくすか、入力検証を行って入力が安全に扱われるようにするかのどちらかである。
色付きテキストの例なら、前者の技法は、扱いにくいクエリパラメーターを使わず、名前の付いた別々の Web ページ（例えば \texttt{/green-page} と \texttt{/blue-page}）を提供するだけで実現できる。
URL に \texttt{color} パラメーターを持たせたいなら、クエリパラメーターの値を許可リストにあるものだけに制限すればよい。

%----------------------------------------------------------------------
\subsubsection{クロスサイトリクエストフォージェリ}\label{subsec:ch11-csrf}
% src: Cross-Site Request Forgery

\begin{quote}
「クモの巣（web）の形を、それが生まれた過程から切り離すことはできない。」\hfill --- Neri Oxman
\end{quote}

\begin{dfn}[クロスサイトリクエストフォージェリ]\label{dfn:ch11-csrf}
\underline{クロスサイトリクエストフォージェリ}（CSRF、ときに XSRF）は、同一オリジンポリシーが抱える根本的な限界を突く攻撃である。
\end{dfn}

この種の攻撃が悪用する脆弱性は、概念としては単純だが非常にとらえにくく、問題がどこにあるのか、どう直せばよいのかは、最初はなかなかわからないことがある。
Web フレームワークは CSRF 保護を提供しているはずだが、それがきちんと機能していることを確かめ、その仕組みを邪魔しないようにするために、根底にある問題をしっかり理解しておくことにはやはり価値がある。

Web サイトが、HTTP の \texttt{GET} で取得したほかの Web サイトの画像などを取り込むことはもちろんできるし、実際よく行われている。
同一オリジンポリシーは、そうした要求を許可しつつ内容を隔離するので、ドメインの異なる Web サイトの間で画像データが漏えいすることはない。
例えば、サイト X は自分のページにサイト Y の画像を表示できる。ユーザーには埋め込まれた画像がページの一部として見えるが、DOM を介してスクリプトが画像データにアクセスすることをブラウザが阻止するので、サイト X 自身はその画像を「見る」ことができない。

しかし、同一オリジンポリシーは \texttt{POST} に対しても \texttt{GET} と同じように働き、しかも \texttt{POST} 要求はサイトの状態を変更しうる。
具体的には次のようなことが起こる。ブラウザは、サイト X がサイト Y にフォームを送信することを許し、その際に Y のクッキーも付けて送る。
ブラウザは、サイト Y からの応答をサイト X から完全に隔離することは保証する。
問題は、\texttt{POST} によって Y のサーバー上のデータが変更されうることである。X にそんなことができてはならないが、設計上、どの Web サイトもほかのどの Web サイトにも \texttt{POST} できてしまう。
ブラウザがこうした認可されていない要求に手を貸してしまう以上、Web 開発者は、サーバー上のデータを変更しようとするこうした試みに対して、明示的に防御しなければならない。

\begin{exa}[投票を盗む CSRF 攻撃]\label{exa:ch11-csrf-voting}
CSRF 脆弱性がどのようなもので、どう悪用され、それに対してどう防御するかを、簡単な攻撃シナリオで示す。
多くのユーザーがそれぞれアカウントを持っているソーシャル Web サイト Y を考えよう。
サイト Y では投票を実施しており、各ユーザーは 1 票ずつ投じられる。
サイトは、投票ページで認証済みのユーザーそれぞれに一意のクッキーを発行し、ユーザー 1 人につき 1 票だけを受け付ける。

投票ページに「投票する前にこれを見て！」というコメントが投稿され、投票の仕方を指南する別の Web サイト X のページへのリンクが張られている。
多くのユーザーがリンクをクリックして、そのページを読む。
同一オリジンポリシーが守ってくれているのだから、何も問題は起きないはずだと思うかもしれない。

まだ問題がわからないなら、大きなヒントを出そう。サイト X のウィンドウの中で何が起きているかもしれないかを考えてほしい。
サイト X を運営しているのが卑劣でずる賢い不正者で、票を盗もうとしているとする。
ユーザーが X を閲覧するたびに、そのページのスクリプトが、そのユーザーのブラウザのコンテキストで（そのユーザーの Y のクッキーを使って）、サイト所有者が望む票をソーシャル Web サイトに送信する。

サイト X は各ユーザーの Y のクッキーを使ってフォームを送信できるので、票を盗むにはそれで十分である。
攻撃者はサーバー上の状態を変更させたいだけであり、ユーザーの投票を確認する応答ページを見る必要はない。同一オリジンポリシーが阻止するのは、その応答だけなのである。
\end{exa}

\underline{防御：}
CSRF を防ぐには、有効な状態変更の要求を推測できないものにする。
言い換えれば、有効な \texttt{POST} 要求の一つひとつを、本来の用途のコンテキストで一度しか通用しない、かけがえのないもの（special snowflake）として扱う。
簡単なのは、すべてのフォームに秘密のトークンを隠しフィールドとして入れておき、各要求に、その Web セッションに対応する秘密が含まれているかを確認する方法である。
CSRF 保護のための秘密のトークンの作成と確認には、微妙な細部がたくさん詰まっているので、掘り下げて見ておく価値がある。
まともな Web フレームワークなら代わりに処理してくれるはずだが、詳細を見てみよう。

次は、CSRF 対策の秘密のトークン（\texttt{csrf\_token} の行）を含む投票フォームの例である。
\begin{verbatim}
<form action="/ballot" method="post">
<label for="name">Voting for</label>
<input type="text" id="name" name="name" value=""/>
<input type="hidden" name="csrf_token"
value="mGEyoi1wE6NBWCyhBN9IZdEmaJLQtrYxi0J23XuXR4o="/>
<input type="submit" value="Vote"/>
</form>
\end{verbatim}
% 注: 原文ではトークン部分が強調表示されているが、スクレイプで強調は失われている。
隠しフィールド \texttt{csrf\_token} は画面には表示されないが、\texttt{POST} 要求には含まれる。
このフィールドの値は、セッションクッキーの内容の SHA-256 ハッシュを base-64 でエンコードしたものだが、クライアントごとの秘密であれば何でもよい。
次は、セッション用の CSRF 対策トークンを作成する Python コードである。
\begin{verbatim}
def csrf_token(self):
    digest = hashlib.sha256(self.session_id.encode('utf-8')).digest()
    return base64.b64encode(digest).decode('utf8')
\end{verbatim}
このコードはトークンをセッションクッキー（文字列値 \texttt{self.session\_id}）から導き出すので、トークンはクライアントごとに異なる。
同一オリジンポリシーのため、サイト X は被害者のサイト Y のクッキーを知ることができない。したがって、Y の作り手が、これらの条件を満たす本物のフォームをでっち上げて \texttt{POST} し、勝利を盗み取ることはできない。
% 注: 原文は "it's impossible for Y's creators to concoct an authentic form" だが、文脈からは X の作り手（攻撃者）のはずである。原文どおりに訳した。

Y のサーバー上の検証コードは、期待されるトークンの値を計算し、受け取ったフォームの対応するフィールドがその値と一致するかを確認するだけである。
次のコードは、トークンが一致しなければ、フォームを実際に処理する前にエラーメッセージを返すことで、CSRF の試みを防ぐ。
\begin{verbatim}
token = fields.get('csrf_token')
if token != self.csrf_token():
    return 'Invalid request: Cross-site request forgery detected.'
\end{verbatim}
% 注: 原文のコードは "csrf\_token" と Markdown のエスケープが残っているので、バックスラッシュを除いた。

CSRF 攻撃を緩和する方法はたくさんあるが、トークンをセッションクッキーから導き出すのはよい解決策である。確認に必要な情報がすべて \texttt{POST} 要求の中に含まれて届くからである。
別の緩和策として、ナンス（推測できない一度限りのトークン）を使う方法もあるが、CSRF 攻撃を防ぐには、やはりナンスを本来のクライアントのセッションに結び付けなければならない。
この方法では、フォームの CSRF トークンとしてランダムなナンスを生成し、セッションをキーとする表にそのトークンを保存しておき、フォームを検証するときにはそのセッションのナンスを引いて一致するかを確認する。

\begin{rem}[SameSite 属性はクライアント側の追加の緩和策]\label{rem:ch11-samesite}
現代のブラウザは、CSRF 攻撃を緩和するためのクッキーの \texttt{SameSite} 属性をサポートしている。
\texttt{SameSite=Strict} は、ページ上のサードパーティへの要求（ほかのドメインへの要求）ではクッキーを一切送らないようにする。これで CSRF は防げるが、自分のクッキーを前提とする別のサイトへ移動するときの便利な振る舞いが壊れることがある。
ほかの設定値もあるが、ブラウザの種類や古いバージョンによってサポート状況がまちまちなことがある。
これはクライアント側の CSRF 防御なので、サーバーが全面的に頼るのは危険であり、唯一の防御ではなく、追加の緩和策と位置付けるべきである。
\end{rem}
% 注: 原文 "considered at additional mitigation" は "an additional mitigation" の誤植とみなした。

%----------------------------------------------------------------------
\subsection{その他の脆弱性と緩和策}\label{subsec:ch11-more}
% src: More Vulnerabilities and Mitigations

\begin{quote}
「線がどこにあるかを知る唯一の方法は、それを越えてみることだ。」\hfill --- Dave Chappelle
\end{quote}

まとめると、セキュアであるためには、品質の高いフレームワークを使い、純粋な HTTPS で Web サイトを構築すべきである。
自分が何をしているかを本当にわかっているのでない限り、フレームワークが備える保護機能を無効にしてはならない。ここで「わかっている」とは、XSS や CSRF のような脆弱性がどのように生じるかを理解しているということである。
現代の Web サイトは外部のスクリプト、画像、スタイル指定などを取り込むことが多いが、そうしたリソースには自分の Web ページへ内容を注入させることになるので、信頼できる提供元のものだけに頼るべきである。

もちろん、話はそれで終わりではなく、サーバーを Web に公開すれば、トラブルに陥る道はまだいくらでもある。
Web サイトは公衆のインターネットに大きな攻撃対象領域をさらしており、そこから入る信頼できない入力は、サーバーのコードで SQL インジェクション（Web サーバーはデータの保存にデータベースを使うことが多い）をはじめ、ありとあらゆる脆弱性をたやすく引き起こしうる。

ほかにも、触れておくべき Web 特有の落とし穴がいくつかある。
特に注意すべき、よくある問題を次に挙げる（ただし、このリストはとても網羅的とはいえない）。
\begin{itemize}
\item 攻撃者が HTTP ヘッダーに信頼できない入力を注入できないようにする（XSS に似た問題である）。
\item ブラウザが応答を正しく処理できるよう、正確な MIME コンテンツタイプを指定する。
\item \underline{オープンリダイレクト}（open redirect）は問題になりうる。任意の URL へのリダイレクトを許さない。
\item \texttt{<IFRAME>} には信頼できる Web サイトだけを埋め込む。（多くのブラウザは、緩和策として \texttt{X-Frame-Options} ヘッダーをサポートしている。）
\item 信頼できない XML データを扱うときは、\underline{XML 外部エンティティ}（XML external entity, XXE）攻撃に注意する。
\item CSS の \texttt{:visited} セレクターによって、ある URL がブラウザの履歴にあるかどうかが知られてしまうおそれがある。
\end{itemize}

さらに、Web サイトは、すばらしい新機能である HTTP の \texttt{Content-Security-Policy} 応答ヘッダーを使って、XSS にさらされる危険を減らすべきである。
このヘッダーでスクリプトや画像（をはじめ、同様の多くの機能）の認可された提供元を指定しておくと、インラインのスクリプトや、ほかのドメインからの悪意ある内容を注入しようとする試みを、ブラウザが阻止できるようになる。
世の中にはたくさんのブラウザがあり、この機能のブラウザ互換性はまだまちまちなので、このヘッダーを使ったからといって、脆弱性が完全に修正されたとみなすことはできない。
追加の防御線と考えるべきものであり、クライアント側の仕組みで自分の制御の外にある以上、XSS に対する完全な免疫を与えてくれる免罪符だとは思わないことである。

信頼できないサードパーティの Web サイトへのリンクは危険になりうる。この章で先に述べたようにブラウザが \texttt{REFERER} ヘッダーを送るかもしれないうえ、DOM の \texttt{window.opener} オブジェクトをリンク先のページに渡してしまうからである。
それらが役に立ち、リンク先が信頼できる場合を除いて、\texttt{rel="noreferrer"} 属性と \texttt{rel="noopener"} 属性をそれぞれ使って阻止すべきである。

既存の大規模な Web サイトにとって、後から新しいセキュリティ機能を追加するのは気が重いかもしれないが、比較的簡単に正しい方向へ進む方法がある。
テスト環境で、すべての Web ページに制限の厳しいセキュリティポリシーを追加し、それから Web サイトをテストして、阻止されたものを一つずつ突き止めていくのである。
安全だとわかっていて使うつもりだったサイトからのスクリプトの読み込みまで禁止してしまった場合は、スクリプトのポリシーを少しずつ緩めていけば、正しいポリシーの例外にすぐたどり着く。
サイト全体を漏れなくテストできるよう、ブラウザ上での自動テストも組み合わせれば、ささやかな労力でセキュリティを大きく前進させられるはずである。

ブラウザに何を許可させ、何を許可させないかを指定するのに役立つ HTTP 応答ヘッダーは数多くあり、\texttt{Content-Security-Policy}、\texttt{Referrer-Policy}、\texttt{Strict-Transport-Security}、\texttt{X-Content-Type-Options}、\texttt{X-Frame-Options} ヘッダーなどがある。
仕様はまだ進化の途中で、サポートもブラウザごとに異なることがあるので、扱いの難しい、変化し続ける状況である。
理想は、まずサーバー側で Web サイトをセキュアにし、そのうえでこれらのセキュリティ機能を第 2 の防御層として使うことである。その際、クライアント側の仕組みだけに頼るのは危険だということを心に留めておいてほしい。

うまくいかなくなりうる道筋の多さ、Web がどんなところから進化してきたか、そしてそこを流れる重要なデータの量を考えると、Web が実際にこれほどセキュアなのは驚くべきことである。
今にして思えば、Web が世界中に広く普及していく中で、セキュリティ技術が時間をかけてゆっくり成熟してきたのは、むしろ最善だったのかもしれない。
もし初期の革新者たちが当時、完全にセキュアなシステムを設計しようとしていたら、その課題はとてつもなく困難なものになっていただろうし、失敗していたら、この取り組みそのものが何も生まずに終わっていたかもしれない。

\newpage

%% =======================================================
%%  part: ch12
%% =======================================================
%% =======================================================
%%  part: ch12
%% =======================================================
\section{セキュリティテスト}
% src: Security Testing

この章では、\cite{Kohnfelder2021} の第12章に沿って、セキュリティテストを扱う。

\begin{quote}
「テストは失敗につながり、失敗は理解につながる。」\hfill ---Burt Rutan
\end{quote}

この章では、信頼性が高くセキュアなコードを開発するうえで欠かせないものとして、セキュリティテストを紹介する。
セキュリティ脆弱性を見つけるために先回りしてテストする手法はよく知られており、難しくもない。
それにもかかわらず実際にはほとんど活用されていないので、セキュリティの保証を高める大きな余地が残っている。

章の構成は次のとおりである。
まずセキュリティテストの用途をざっと見たうえで、セキュリティテストさえあれば世界中のユーザーを重大な脆弱性から守れたはずの事例を順にたどる。
次に、脆弱性やその前兆を見つけて捕まえるためのセキュリティテストケースについて、書き方の基本を見る。
ファズテストは、より深いところに潜む問題をあぶり出すのに役立つ強力な補助手段である。
また、既存の脆弱性への対応として作成し、同じ誤りを二度と繰り返さないようにするセキュリティリグレッションテストも扱う。
最後に、サービス拒否やそれに関連する攻撃を防ぐためのテストを論じ、セキュリティテストのベストプラクティスをまとめる（セキュリティテストに関するアイデアを幅広く扱うが、決して網羅的ではない）。

\subsection{セキュリティテストとは}
% src: What Is Security Testing?

はじめに、ここでいうセキュリティテストが何を指すのかをはっきりさせておく。

\begin{dfn}[セキュリティテスト]\label{dfn:ch12-security-testing}
ほとんどのテストは、コードを実行して、機能が意図どおりに動くことを確かめるものである。
\underline{セキュリティテスト}（security testing）はこれを逆にしただけのもので、\emph{許されるべきでない}操作が実際に許されないことを確かめる（この違いは、後でコードの例を見ればすぐにはっきりする）。
\end{dfn}

緩和策がきちんと働いていることを保証するには、セキュリティテストが欠かせない。
コードを書く人は当然、ふつうの使い方で意図した機能が動くことに集中するので、攻撃者が仕掛けてくる予想外の操作を完全に見通すのは難しい。
これまでの章の内容から、セキュリティテストのネタはいくらでも思いつくはずである。
以前に扱った主な脆弱性のクラスに対応する、基本的なセキュリティテストケースの種類を次に挙げる。

\begin{itemize}
\item \underline{整数オーバーフロー}：許容する値の範囲を決め、範囲外の値を検出して拒否する処理が働くことを確かめる。
\item \underline{メモリ管理の問題}：極端に大きなデータ値をコードが正しく扱い、大きすぎるものは拒否することをテストする。
\item \underline{信頼できない入力}：さまざまな不正な入力でテストし、それらが拒否されるか、安全に処理できる有効な形式に変換されることを確かめる。
\item \underline{Web}：HTTP のダウングレード攻撃、不正な認証トークンや CSRF トークン、XSS 攻撃が失敗することを確かめる（詳細は前の章を参照）。
\item \underline{例外処理の欠陥}：コードにさまざまな例外処理の経路を強制的に通らせ（まれな経路には依存性注入を使う）、適切に回復することを確かめる。
\end{itemize}

これらのテストに共通するのは、どれもふつうの使い方から外れたところにあるという点で、だからこそ忘れられやすい。
しかも、これらの領域はどれも攻撃の格好の的なので、徹底してテストするかどうかで大きな差が出る。
セキュリティテストは、こうしたケースを見越して、必要な保護メカニズムが常に働くことを確かめることで、コードをよりセキュアにする。
加えて、セキュリティ上重要なコードについては、品質をできる限り高めるために、コードカバレッジを徹底することを勧める。
そうした領域のバグは、壊滅的な結果を招きがちだからである。

セキュリティテストは、アプリケーションのセキュリティを実際に高め始めるための、おそらく最良の方法である。しかも難しくはない。
ソフトウェア業界でセキュリティテストがどの程度行われているか（あるいは行われていないか）を示す公的な統計はないが、同じような脆弱性が繰り返し起きている状況は、セキュリティテストという大きな機会がみすみす逃されていることを強くうかがわせる。

\subsection{GotoFail 脆弱性のセキュリティテスト}\label{subsec:ch12-gotofail}
% src: Security Testing the GotoFail Vulnerability

\begin{quote}
「逆境こそが、人の真価を試す。」\hfill ---Harry Emerson Fosdick
\end{quote}

第8章で調べた GotoFail 脆弱性を思い出してほしい。セキュアな接続のチェックが素通りされてしまった脆弱性である。
そこで示した簡略版の例を広げて、セキュリティテストがあればこのような問題をいかに簡単に見つけられたかを見ていく。

GotoFail 脆弱性は、1 行のコードがうっかり二重になったことで生じた（次のコード片で 2 つ続いている \texttt{goto fail;} の 2 行目）。
その行が goto 文だったため、一連の重要なチェックが飛ばされ、検証関数は無条件に「合格」の戻り値を返すようになった。
以前は（筆者の簡略版で）重要な行だけを見たが、セキュリティテストを行うには関数全体を調べる必要がある。

% 脆弱なコード（vulnerable code）。長いため itembox に入れず地の文に置く（itembox はページをまたげない）。
\noindent\underline{脆弱なコード：}
% 注: 原文ではハイライトされた行が重複した goto fail;（下の 3 番目の if の直前の行）。
%   行頭の 1〜4 は本文で参照される位置の目印。原文のまま転記した
%   （fail: ラベル後の余分な閉じ括弧、空行も原文のまま）。
\begin{verbatim}
/*
* Copyright (c) 1999-2001,2005-2012 Apple Inc. All Rights Reserved.
*
* @APPLE_LICENSE_HEADER_START@
*
* This file contains Original Code and/or Modifications of Original Code
* as defined in and that are subject to the Apple Public Source License
* Version 2.0 (the 'License'). You may not use this file except in
* compliance with the License. Please obtain a copy of the License at
* http://www.opensource.apple.com/apsl/ and read it before using this
* file.
*
* The Original Code and all software distributed under the License are
* distributed on an 'AS IS' basis, WITHOUT WARRANTY OF ANY KIND, EITHER
* EXPRESS OR IMPLIED, AND APPLE HEREBY DISCLAIMS ALL SUCH WARRANTIES,
* INCLUDING WITHOUT LIMITATION, ANY WARRANTIES OF MERCHANTABILITY,
* FITNESS FOR A PARTICULAR PURPOSE, QUIET ENJOYMENT OR NON-INFRINGEMENT.
* Please see the License for the specific language governing rights and
* limitations under the License.
*
* @APPLE_LICENSE_HEADER_END@
*/
int VerifyServerKeyExchange(ExchangeParams params,
                            uint8_t *expected_hash, size_t expected_hash_len) {
  int err;
  HashCtx ctx = 0;
  uint8_t *hash = 0;
  size_t hash_len;
  if ((err = ReadyHash(&ctx)) != 0)
    goto fail;
1 if ((err = SSLHashSHA1_update(ctx, params.clientRandom, PARAM_LEN)) != 0)

    goto fail;
2 if ((err = SSLHashSHA1_update(ctx, params.serverRandom, PARAM_LEN)) != 0)
    goto fail;
goto fail;
3 if ((err = SSLHashSHA1_update(ctx, params.signedParams, PARAM_LEN)) != 0)
    goto fail;
  if ((err = SSLHashSHA1_final(ctx, &hash, &hash_len)) != 0)
    goto fail;
  if (hash_len != expected_hash_len) {
    err = -106;
    goto fail;
  }
4 if ((err = memcmp(hash, expected_hash, hash_len)) != 0) {
    err = -100; // Error code for mismatch
}
  SSLFreeBuffer(hash);

fail:
  if (ctx)
    SSLFreeBuffer(ctx);
  }
  return err;
}
\end{verbatim}

\begin{rem}[コードの出典]\label{rem:ch12-gotofail-source}
このコードは、バグを含んでいたオリジナルの \texttt{sslKeyExchange.c} をもとにしている。
致命的な脆弱性に直接関係しない部分は簡略化し、簡潔にするため一部の名前を変えてある。
たとえば、実際の関数名は \texttt{SSLVerifySignedServerKeyExchange} である。
\end{rem}

\texttt{VerifyServerKeyExchange} 関数は 3 つのフィールドからなる引数 \texttt{params} を受け取り、その内容からメッセージダイジェストのハッシュを計算して、データを認証するための \texttt{expected\_hash} の値と比べる。
戻り値 0 はハッシュが一致したことを表し、正当なリクエストであるための条件になる。
0 以外の戻り値は何か問題があったことを意味する：ハッシュ値の不一致（$-100$）、ハッシュ長の不一致（$-106$）、あるいは何らかのエラーによりハッシュ計算ライブラリが 0 以外のエラーコードを返した場合のいずれかである。
セキュリティはこの仕組みに頼っている：ハッシュ値やデータが少しでも改ざんされればハッシュが一致しなくなり、何かがおかしいとわかるのである。

まず、重複した goto 文が紛れ込む前の、正しい版のコードを順に見ていく。
\texttt{HashCtx ctx} というコンテキスト変数を準備したら、\texttt{params} の 3 つのデータフィールドを順にハッシュする（1、2、3 の位置）。
途中で何かエラーが起きれば、ラベル \texttt{fail} に飛んで変数 \texttt{err} のエラーコードを返す。
エラーがなければ処理を続け、ハッシュの結果をバッファにコピーして、期待されるハッシュ値と比べる（4 の位置）。
比較関数 \texttt{memcmp} は、両者が等しければ 0 を返す。ハッシュが異なる場合、コードは \texttt{err} にエラーコード $-100$ を代入し、そのまま下へ進んでその値を返す。

\subsubsection{機能テスト}
% src: Functional Testing

セキュリティテストを考える前に、まず \texttt{VerifyServerKeyExchange} 関数の機能テストから始めよう。

\begin{dfn}[機能テスト]\label{dfn:ch12-functional-testing}
コードが期待どおりに動くことを確かめるテストを\underline{機能テスト}（functional testing）という。
\end{dfn}

次の簡単な例は、決して完全なものではない。
この例では、C 用のテストフレームワーク MinUnit を使う。
理解するのに必要な知識は、\texttt{mu\_assert(test, message)} が式 \texttt{test} が真であることを確かめ、真でなければテストを失敗させて、指定した \texttt{message} を表示する、ということだけである。

\begin{verbatim}
mu_assert(0 == VerifyServerKeyExchange(test0, expected_hash, SIG_LEN),
    "Expected correct hash check to succeed.");
\end{verbatim}

このテストは正しいとわかっているパラメーターで関数を呼び出すので、合格するには戻り値 0 が返ってくる必要がある。
関数の中では、3 つのフィールドがハッシュされ（1、2、3 の位置）、4 の位置でハッシュが等しいと判定される。
ここには示していないが、3 つのデータフィールドのテスト値（\texttt{test0} という名前の \texttt{ExchangeParams} 構造体）と、サーバーが署名するはずの正しいハッシュをあらかじめ計算したもの（\texttt{expected\_hash}）を用意してある。

\subsubsection{脆弱性がある状態での機能テスト}
% src: Functional Testing with the Vulnerability

次に、GotoFail 脆弱性（ハイライトされた行）を入れたらどうなるかを見てみよう。
余分な goto がある状態で機能テストを実行し直しても、テストは\emph{そのまま合格してしまう}。
コードは重複した goto までは正常に動くが、そこで 3 番目のデータフィールドのハッシュ（3 の位置）とハッシュの比較（4 の位置）を飛び越えてしまう。
関数は正しい入力をこれまでどおり検証に通すが、本来拒否すべき不正な入力の一部まで通すようになっている。
ところが、この時点ではまだそれに気付けない。
セキュリティテストが非常に重要なのはまさにこのためであり、見過ごされやすいのもこのためである。

もっと徹底した機能テストなら、検証の失敗（0 以外の戻り値）を確かめるテストケースなども用意していたかもしれない。
しかし機能テストは、セキュリティのために検証関数が入力を拒否しなければならないケースを、すべて網羅するところまではいかないことが多い。
そこでセキュリティテストの出番になる。次にそれを見ていく。

\subsubsection{セキュリティテストケース}
% src: Security Test Cases

では、セキュリティテストケースをいくつか書いてみよう。
ハッシュするデータの塊は 3 つあるので、それぞれに対応する 3 つのテストを書くのがよさそうである。
どのテストもデータの値をどこか変えてあり、そのためハッシュが期待値と一致しなくなる。
値が変わっているのはデータが改ざんされたことを表しているのかもしれず、ハッシュの比較はまさにそれを防ぐためのものなので、検証関数はこれらの入力を拒否しなければならない。
実際の値（\texttt{test1}, \texttt{test2}, \texttt{test3}）は、正しい \texttt{test0} をコピーし、3 つのデータフィールドのうち 1 つだけを少し変えたものである。値そのものは重要ではないので示さない。
3 つのテストケースは次のとおりである。

\begin{verbatim}
mu_assert(-100 == VerifyServerKeyExchange(test1, expected_hash, SIG_LEN),
    "Expected to fail hash check: wrong client random.");
mu_assert(-100 == VerifyServerKeyExchange(test2, expected_hash, SIG_LEN),
    "Expected to fail hash check: wrong server random.");
mu_assert(-100 == VerifyServerKeyExchange(test3, expected_hash, SIG_LEN),
    "Expected to fail hash check: wrong signed parameters.");
\end{verbatim}

バグがあると、この 3 つはすべて失敗する。
検証関数は問題の goto までは正常に動くが、そこで無条件にラベル \texttt{fail} へ飛んでしまい、ハッシュの処理は中途半端なまま、4 の位置のハッシュ値の比較は一度も行われない。
これらのテストは検証が失敗するのを正しい結果として期待するように書いてあるので、戻り値 0 が返るとテストは失敗する。
これで、リリース前にこの脆弱性を捕まえ、その後に起きた大失態を避けられたはずの、テストというセーフティネットが手に入ったことになる。
% 注: 原文 "go``to" は "goto" の誤植とみなした。

網羅性を考えると、もう一つのセキュリティテストケースが思い浮かぶ。
3 つの値はどれも \texttt{test0} と同じく正しいのに、署名されたハッシュのほうが異なる（\texttt{wrong\_hash}）場合はどうだろうか。
そのテストケースは次のとおりである。

\begin{verbatim}
mu_assert(-100 == VerifyServerKeyExchange(test0, wrong_hash, SIG_LEN),
    "Expected check against the wrong hash value to fail.");
\end{verbatim}

予想どおり、誤った goto があればこのテストも失敗する。
この脆弱性に限っていえば、これらのテストのどれか 1 つで捕まえられたはずだが、セキュリティテストの目的は、潜在的な脆弱性をできるだけ広く網羅することにある。

\subsubsection{セキュリティテストの限界}
% src: The Limits of Security Tests

セキュリティテストはコードが故障しうる主な箇所を見つけることを目指すが、コードがおかしくなりうる無数の道筋をすべて網羅することは決してできない。
今書いたテストでは検出できない脆弱性を作り込むことはできるが、それが\emph{うっかり}起きる可能性は低い。
テストカバレッジがよほど徹底していない限り、テストをすり抜けるバグをわざと作り込む余地は残る。
しかし、ここで主に警戒すべきなのはうっかり入るバグなので、ささやかなセキュリティテストケースの集まりでも、十分に効果を発揮しうる。

セキュリティテストケースをどこまで徹底すべきかは判断が要るところだが、目安ははっきりしている。

\begin{itemize}
\item セキュリティにとって決定的に重要なコードほど、セキュリティテストは重要になる。
\item 最も重要なセキュリティテストは、多くの場合、（成功ではなく）アクセスの拒否、入力の拒否、失敗といった動作を確かめるものである。
\item セキュリティテストケースでは、主要なステップ（この例では 3 つのハッシュとハッシュの比較）がそれぞれ正しく動くことを確かめるべきである。
\end{itemize}

ここまで、単純な（しかし思いがけない）原因による現実のセキュリティ脆弱性と、そうした事態に備えるセキュリティテストの方法を詳しく見てきた。
次は一般的な場合を考え、この種の問題をどうすれば予測し、先回りして避けられたかを見ていく。

\subsection{セキュリティテストケースを書く}
% src: Writing Security Test Cases

\begin{quote}
「よいテストケースとは、まだ見つかっていない誤りを見つける確率が高いテストケースである。」\hfill ---Glenford Myers
\end{quote}

\begin{dfn}[セキュリティテストケース]\label{dfn:ch12-security-test-case}
特定のセキュリティ上の失敗が起こらないことを確かめるテストケースを\underline{セキュリティテストケース}（security test case）という。
\end{dfn}

この種のテストの出発点は、4つの質問の 2 番目「何がうまくいかなくなりうるか」である。
セキュリティテストはペネトレーションテスト（善意の人々が倫理的にソフトウェアを攻めて脆弱性を探し、悪意のある者に見つかる前に修正できるようにするもの）とは違い、考えうるエクスプロイトをすべて洗い出そうとはしない。
また、将来脆弱性が入り込むのを防ぐ守りになるという点でも、セキュリティテストはペネトレーションテストと異なる。

セキュリティテストケースは保護メカニズムが正しく動くことを確かめるものであり、多くの場合、不正な入力や許されない操作を拒否したり無害化したりすることを確かめる。
GotoFail のバグそのものを予測できた人はいなかっただろうが、\texttt{VerifyServerKeyExchange} 関数の if 文がどれもセキュリティにとって決定的に重要であることは、すぐにわかる。
一般に、このようなコードでは、セキュリティチェックを強制する条件の一つ一つにテストカバレッジが必要である。
そこまでのテストが揃っていれば、余分な goto が脆弱性を生んだ時点で、どれかのテストケースが失敗して問題に気付かせてくれる。

セキュリティテストケースは、脆弱性が見つかってから慌てて作るのではなく、ほかの単体テストを書くときに一緒に作っておくべきである。
セキュアなシステムは、不適切な操作をブロックし、悪意のある入力を拒否し、アクセスを拒否するといった方法で、価値ある資源を守る。
そうしたセキュリティメカニズムがある箇所には必ずセキュリティテストケースを作り、許可されていない操作が確かに失敗することを確かめる。

\begin{exa}[ありふれたセキュリティテストケース]\label{exa:ch12-common-test-cases}
ありふれたセキュリティテストケースの例としては、誤ったパスワードでのログインが失敗すること、ユーザー空間からカーネル資源への許可されていないアクセスが失敗すること、さまざまな形で無効な、あるいは不正な形式のデジタル証明書が必ず拒否されること、を確かめるものがある。
\end{exa}

コードを読むのは、よいセキュリティテストケースのアイデアを得るための優れた方法である。

\subsubsection{入力検証のテスト}
% src: Testing Input Validation

入力検証のためのセキュリティテストケースを考えよう。
簡単な例として、長さが 10 文字以上 20 文字以下で、英数字の ASCII 文字だけからなる文字列を要求する入力検証のコードをテストする。

この種の標準的な入力検証を行うヘルパー関数を作っておけば、検証を一様に、漏れなく行える。
そのうえで、入力検証とそれに対応するテストケースを組にして、検証のチェックが働くこと、許容される限界ぎりぎりまでコードが正しく動くことを確かめる。
実際、off-by-one エラーはプログラミングにつきものなので、限界ちょうどの値と、限界をわずかに超えた値の両方を確かめる習慣をつけるとよい。
次の単体テストで、この例の入力検証のテストケースを網羅できる。

\begin{exa}[長さと文字種の入力検証のテストケース]\label{exa:ch12-input-validation-tests}
\begin{itemize}
\item 長さ 10 の有効な入力は通り、長さ 9 以下の入力は失敗することを確かめる。
\item 長さ 20 の有効な入力は通り、長さ 21 以上の入力は失敗することを確かめる。
\item 不正な文字を 1 つでも含む入力は必ず失敗することを確かめる。
\end{itemize}
\end{exa}

もちろん、すべての制約を満たすサンプル入力が正しく動くことは、機能テストですでに確かめてあるはずである。

\begin{exa}[固定長バッファ]\label{exa:ch12-fixed-buffer-tests}
似た例をもう一つ挙げる。テスト対象のコードが、バイト配列のパラメーターを $N$ バイトの固定長バッファに格納するとしよう。
この場合のセキュリティテストケースでは、サイズが $N$ 以下の入力ならコードが期待どおりに動き、サイズ $N+1$ の入力は安全に拒否されることを確かめるべきである。
\end{exa}

\subsubsection{XSS 脆弱性のテスト}
% src: Testing for XSS Vulnerabilities

次に、もっと手ごわいセキュリティテストケースと、そこで使えるいくつかのテスト戦略を見ていく。
第11章の XSS 脆弱性を思い出してほしい。
信頼できない入力が Web サーバー上で生成される HTML に入り込み、ページの中へ抜け出して、たとえば攻撃を仕掛けるスクリプトを実行させる、という脆弱性であった。
この脆弱性の根本原因はエスケープの不備だったので、セキュリティテストもそこに焦点を当てる。

テスト対象のコードは、次の Python 関数だとしよう。内容を記述した文字列をもとに、HTML の断片を組み立てる関数である。

\begin{exa}[HTML タグを組み立てる関数（脆弱なコード）]\label{exa:ch12-html-tag}
\begin{verbatim}
def html_tag(name, attrs):
    """Build and return an HTML fragment with attribute values.
    >>> html_tag('meta', {'name': 'test', 'content': 'example'})
    '<meta name="test" content="example">'
    """
    result = '<%s' % name
    for attr in attrs:
        result += ' %s="%s"' % (attr, html.escape(attrs[attr]))
    return result + ">"
\end{verbatim}
\end{exa}

コメント（\texttt{"""} で囲まれた部分）にある doctest（\texttt{>>>} で始まる部分）の例は、この関数で \texttt{<meta>} タグの HTML テキストを生成する方法を示している。
最初の行は、テキスト文字列 \texttt{result} の先頭部分を組み立てる：あらゆる HTML タグの始まりを示す山括弧（\texttt{<}）の後ろにタグ名を置く。
次のループでは属性（\texttt{attrs}）を一つずつたどり、属性ごとに空白とその宣言（\texttt{X="Y"} の形）を付け足していく。

このコードは属性の文字列値ごとに \texttt{html.escape} 関数を正しく適用しているが、それでもテストはしておくべきである。
（ここでは、エスケープが必要な信頼できない入力が入りうるのは属性値だけだと仮定する。
実際にはふつうはこれで足りるが、何が起こるかはわからないので、アプリケーションによってはさらにエスケープや入力検証が必要になるかもしれない。）

Python の \texttt{unittest} ライブラリでテストケースを書いてみよう。

\begin{verbatim}
class SecurityTestCases(unittest.TestCase):
    def test_basic(self):
        self.assertEqual(html_tag('meta', {'name': 'test', 'content': '123'}),
                         '<meta name="test" content="123">')

    def test_special_char(self):
        self.assertEqual(html_tag('meta', {'name': 'test', 'content': 'x"'}),
                         '<meta name="test" content="x&quot;">')

if __name__ == '__main__':
    unittest.main()
\end{verbatim}

最初のテストケースは基本的な機能テストで、こうした単体テストがどう動くかを示している。
コマンドラインから実行すると、このモジュールは最後の行で単体テストフレームワークの \texttt{main} を呼び出す。
この関数は、単体テストを含む \texttt{unittest.TestCase} のすべてのサブクラスについて、各メソッドを自動的に呼び出す。
\texttt{assertEqual} メソッドは引数どうしを比べ、等しくなければテストは失敗する。

次に、\texttt{test\_special\_char} という名前のセキュリティテストケースを見てみよう。
信頼できない入力を囲む二重引用符の外へ抜け出せば XSS でこのコードを悪用できることがわかっているので、二重引用符を含む文字列でエスケープをテストする。
HTML のエスケープが正しく行われれば、assert 文の期待値の文字列にあるとおり、二重引用符は HTML エンティティ \texttt{\&quot;} に変換されるはずである。
対象のメソッドから \texttt{html.escape} 関数を取り除くと、狙いどおり、このテストはきちんと失敗する。

ここまでは順調である。
ただし、このテストを書くには、どんな入力が問題になりうるか（二重引用符）を前もって知っている必要があったことに注意してほしい。
HTML の仕様はかなり込み入っているので、もっと重要なテストケースがほかにないと、どうして言い切れるだろうか。
ほかの特殊文字を片っ端から試すこともできる。その多くは、エスケープ関数によってさまざまな HTML エンティティに変換される（たとえば、大なり記号は \texttt{\&gt;} になる）。
しかし、こうしてあらゆる可能性を網羅するようにテストケースを揃えていくのは、大変な手間である。

扱っているのは HTML なので、仕様を隅々まで知り尽くしたライブラリに面倒な作業を任せることができる。
次のテストケースは、先ほどと同じ 2 つのテスト値、つまり通常の文字列と二重引用符を含む文字列を、変数 \texttt{content} に順に入れて HTML タグを組み立て、その結果を確かめる。

\begin{verbatim}
    def test_parsed_html(self):
        for content in ['x', 'x"']:
            result = html_tag('meta', {'name': 'test', 'content': content})
            soup = BeautifulSoup(result, 'html.parser')
            node = soup.find('meta')
            self.assertEqual(node.get('name'), 'test')
            self.assertEqual(node.get('content'), content)
\end{verbatim}

ループの中は両方のケースに共通のテストコードで、まず対象の関数を呼び出して HTML の \texttt{<meta>} タグの文字列を組み立てる。

明示的な期待値と比べる代わりに、ここでは BeautifulSoup パーサーを呼び出す。
このパーサーは、パースした HTML の構造を論理的に表すオブジェクトの木を生成する（このオブジェクト群は、しゃれて「スープ（soup）」と呼ばれる）。
変数 \texttt{soup} は HTML のノード構造の根で、これを起点にオブジェクトモデルをたどって内容を調べられる。

\texttt{find} メソッドはスープの中から最初の \texttt{<meta>} タグを探し出し、変数 \texttt{node} に代入する。
\texttt{node} オブジェクトには、属性の値を名前で調べる \texttt{get} メソッドがある。
このコードでは、\texttt{<meta>} タグの \texttt{name} 属性と \texttt{content} 属性がどちらも期待どおりの値になっていることをテストする。
パーサーを使う大きな利点は、HTML テキスト中の空白や改行の扱い、エスケープとその解除、エンティティ表現の変換など、HTML のパースにまつわる処理をすべてパーサーが引き受けてくれることである。

パーサーライブラリを使っているので、このセキュリティテストケースはパース済みのオブジェクトを相手に動き、HTML の癖に振り回されない。
XSS によって二重引用符の外へ抜け出す悪意のある入力を注入できたなら、パースされた HTML の \texttt{<meta>} タグの \texttt{node} オブジェクトには、元と同じ値が入っていないはずである。
したがって、二重引用符が一部の XSS 攻撃で問題になることをまったく知らなかったとしても、さまざまな特殊文字を手軽に試し、どれが正しく動くか（あるいは動かないか）の判定はパーサーに任せられる。
次の話題は、このようにテストケースのバリエーションを数多く試すという考え方を、大規模に自動化するものである。

\subsection{ファズテスト}
% src: Fuzz Testing

\begin{quote}
“Rock and roll to the beat of the funk fuzz.”\hfill ---A Tribe Called Quest
\end{quote}

\begin{dfn}[ファズテスト]\label{dfn:ch12-fuzz-testing}
テストケースを自動的に生成し、大量のテスト入力を対象のコードに浴びせる手法を\underline{ファズテスト}（fuzz testing）という。
\end{dfn}

ファズテストは、特定の入力がコードを失敗させたりプロセスをクラッシュさせたりするかどうかを調べるのに役立つ。

\begin{exa}[食器洗い機のたとえ]\label{exa:ch12-dishwasher}
食器洗い機は、回転するアームからさまざまな角度で水を噴射して洗う。
食器がどう並んでいるかも、どの角度で水を当てれば効果的かも知らずにランダムに噴射しているのに、それでもすべてをきれいに洗い上げる。
\end{exa}

特定の狙いをもって書くセキュリティテストケースとは対照的に、ファズテストの乱れ撃ち方式は、より幅広いバグを見つけるのにかなり効果を発揮しうる。そのうちのいくつかは脆弱性である。

セキュリティテストケースでの典型的な使い方は、信頼できない入力を「ファズ」する（つまり、さまざまな値を大量に試す）ことで、異常な結果やクラッシュを探すというものである。
実際にセキュリティ脆弱性を突き止めるには、ファズテストの結果から得た手がかりを調べる必要がある。

前節のテストケース \texttt{test\_parsed\_html} は、二重引用符だけでなく多数の特殊文字を確かめるようにすれば、簡単にファズテストに作り替えられる。

\begin{exa}[HTML 組み立てのファズテスト]\label{exa:ch12-fuzzy-html}
\begin{verbatim}
    def test_fuzzy_html(self):
        for fuzz in string.punctuation:
            content = 'q' + fuzz
            result = html_tag('meta', {'name': 'test', 'content': content})
            soup = BeautifulSoup(result, 'html.parser')
            node = soup.find('meta')
            self.assertEqual(node.get('name'), 'test')
            self.assertEqual(node.get('content'), content)
\end{verbatim}
\end{exa}

このコードは、選んだテストケースのリストを試す代わりに、標準の \texttt{string} ライブラリの定数で定義されている ASCII の句読点文字すべてについてループする。
ループを回るたびに、変数 \texttt{fuzz} に句読点文字が入り、その前に文字 \texttt{q} を付けて 2 文字の \texttt{content} の値を作る。
残りのコードは元の例とまったく同じだが、実行するテストケースはずっと多い。

この例は単純すぎて、ファズテストと呼ぶには少し無理があるほどだが、テストケースを一つずつ手で選んで書く代わりに、プログラムで 32 ケースを力ずくでテストできることの威力はわかるだろう。
このコードをもっと凝ったものにすれば、HTML の厄介な引用符やエスケープ文字を並べた長い文字列を使って、はるかに多くのケースを作れるだろう。

さまざまなファジング機能を提供するライブラリは数多くあり、ランダムなファジングから、HTML、XML、JSON といった特定の形式の知識に基づくバリエーションの生成まで、内容は多岐にわたる。
特定のテスト戦略を思い描いているなら、もちろん自分でテストケースを書いて試してもよい。
要するに、テストケースは安上がりであり、それを大量に生成するのは、よいテストカバレッジを得るための手軽な方法なのである。

\subsection{セキュリティリグレッションテスト}
% src: Security Regression Tests

\begin{quote}
「後ずさりする者は、決して前に進めない。」\hfill ---Umar ibn al-Khattâb
\end{quote}

いったん見つかって修正されたセキュリティ脆弱性こそ、二度とぶり返して痛い目にあわされたくないバグの筆頭である。
それでもぶり返しは実際に起きており、しかも本来あるべき頻度よりずっと多い。そしてそれが起きたなら、セキュリティテストが足りていないことの明らかな兆候である。

\begin{dfn}[セキュリティリグレッションテスト]\label{dfn:ch12-security-regression-test}
以前に修正された既知のセキュリティバグが再発することを\underline{セキュリティリグレッション}（security regression）という。
新たに見つかったセキュリティ脆弱性に対応するときに作成し、その原因となったバグ（複数のこともある）を検出するテストを\underline{セキュリティリグレッションテスト}（security regression test）という。
\end{dfn}

セキュリティリグレッションテストを作ることは、重要なベストプラクティスである。
このテストは手軽な \underline{repro}（バグを再現するテストケース）として役立つうえ、修正によって脆弱性が本当に取り除かれたことの確認にもなる。

少なくとも考え方としてはそうなのだが、最大手で最も洗練されたソフトウェアメーカーでさえ、この慣行を十分きちんと守っているとはいえないようである。

\begin{exa}[iOS 12.4 でのリグレッション]\label{exa:ch12-ios-regression}
Apple は 2019 年に iOS 12.4 をリリースした際、iOS 12.3 ですでに見つかって修正されていたのとまったく同じバグを再び持ち込み、しっかり閉じておくべき扉となる脆弱性をたちまち復活させてしまった。
元の修正にセキュリティリグレッションテストケースが含まれていれば、こんなことは決して起きなかったはずである。
\end{exa}

場合によっては、セキュリティリグレッションのほうが新しい脆弱性よりはるかに深刻になりうる点は注目に値する。
この iOS のリグレッションが特に痛手だったのは、そのバグがセキュリティ研究者の間ですでによく知られていたため、iOS 12.3 向けに作られていた既存の脱獄（jailbreak）ツールがすぐに手直しされ、iOS 12.4 でも使えるようになってしまったからである
（\underline{脱獄}とは、ユーザーがデバイスでできることを制限するためにメーカーが課した制約をかいくぐる権限昇格のことである）。

筆者は、実際の修正に取りかかる前に、まずテストケースを書くことを勧める。
緊急時に修正方法がはっきりしているなら修正を先にしてもよいが、一人で作業しているのでない限り、並行して誰かにリグレッションテストを作ってもらうのがよい習慣である。
効果的なリグレッションテストを作る過程で、その問題についての理解が深まり、関連する潜在的な脆弱性の手がかりまで得られることもある。

よいセキュリティリグレッションテストは、既知の攻撃とまったく同じテストケースを 1 つだけ試すのではなく、もっと一般化したものであるべきである。

\begin{exa}[SQL インジェクションのリグレッションテスト]\label{exa:ch12-sql-injection-regression}
第10章で述べた SQL インジェクション攻撃の場合、既知の「Bobby Tables」攻撃 1 つが失敗するようになったことをテストするだけでは足りない。
極端に長い名前も試すこと。そうすれば、入力検証で名前の入力文字列の長さもチェックしなければならないとわかるかもしれない。
攻撃のバリエーション、たとえば単一引用符の代わりに二重引用符を使うものや、名前の末尾にバックスラッシュ（SQL の文字列のエスケープ文字）を置くものも試すこと。
同じテーブルの別の列や、別のテーブルに対する同様の攻撃も試すこと。
SQL インジェクションのバグを、\texttt{Robert');} で始まる名前だけを狭く拒否するという方法では修正しないのと同じく（それでもその攻撃自体は止められるが）、リグレッションテストもそのような書き方をすべきではない。
\end{exa}

新たに見つかった脆弱性に対処していると、調べるうちに、システムのほかの場所にも悪用されうる似た脆弱性があるとわかることはよくある。
システム内部についての深い知識とソースコードへの精通を活かして、潜在的な敵の先回りをすること。
できれば似たバグがないかをすぐに調べ、元の脆弱性を塞ぐ更新に、それらの修正もあわせて含められるようにする。
攻撃者も間違いなく同じ方向で考えているし、修正をリリースすれば、システムを狙う新たな方法について攻撃者に大きな手がかりを与えることになるので、この点は重要になりうる。
すべての手がかりを調べる時間がなければ、詳細を記録しておき、時間ができたときに改めて調べる。

\begin{exa}[Heartbleed のリグレッションテスト]\label{exa:ch12-heartbleed-regression}
例として、Heartbleed 脆弱性に対するセキュリティリグレッションテストの書き方を考えよう。
このエクスプロイトは、任意のバイト列からなるペイロードに、実際よりはるかに大きなバイト数を付けたパケットを送ることで成立していたことを思い出してほしい。
サーバーはそのバイト数に従って余分なメモリの内容まで応答に含めて送り返し、多くの場合、内部データの深刻な漏えいを招いた。

正しい振る舞いは、そのような不正なリクエストを無視することである。
よいセキュリティリグレッションテストケースとしては、次のようなものがある。
\begin{itemize}
\item 既知のエクスプロイトのリクエストに、もう応答が返らないことをテストする。
\item リクエストのバイト数が 16,384（最大値）を超える場合をテストする。
\item ペイロードが 0 バイトのリクエストと、最大バイトサイズのリクエストをテストする。
\item TLS プロトコルのほかの種類のパケットにも同様の問題がありうるかを調べ、あるならそれらもテストする。
\end{itemize}
\end{exa}

\subsection{可用性のテスト}
% src: Availability Testing

\begin{quote}
「利用できなくなることを心配しなさい。不在であることや、偽物であることを心配しなさい。」\hfill ---Anne Lamott
\end{quote}

サービス拒否攻撃は特殊な潜在的脅威である。システムがどこまでの負荷に耐えられるべきかを明確にするのが難しいからである。
とくに、この文の「負荷」という言葉には多くの意味が詰まっている：処理能力、メモリ消費、オペレーティングシステムの資源、ネットワーク帯域幅、ディスク容量、その他のボトルネックになりうるもの（第5章の CSPRNG のエントロピープールを思い出してほしい）である。
運用スタッフはふつう、本番環境での利用状況に合わせてこれらの要素を監視するが、性能上の脆弱性を意図的に突く攻撃のなかには、セキュリティテストで防げるものもある。

セキュリティテストには、性能が非線形に劣化するおそれのあるコードを見つけるためのテストケースを含めるべきである。
この種の脆弱性の例は、第10章でバックトラッキングする正規表現や XML エンティティ展開の爆発を取り上げたときに見た。
これらは性能を指数関数的に悪化させうるので、特に強力な脆弱性である。
もちろん、これらはより大きな現象の 2 つの例にすぎず、同じ問題はあらゆる種類のコードで起こりうる。

以下の節では、この種の問題をテストするための 2 つの基本戦略を説明する：特定の機能の性能を測ることと、さまざまな負荷のもとで全体の性能を監視することである。

\subsubsection{資源の消費}
% src: Resource Consumption

可用性への攻撃を受けやすいとわかっている機能については、入力の妥当な上限を測って決め、爆発が起きないよう守るセキュリティテストケースを追加する。
そのうえでさらにテストを行い、それより大きな入力が入力検証で弾かれ、システムが過負荷にならないことを確かめる。

たとえば、バックトラッキングする正規表現なら、長さ $N$ と $N+1$ の文字列でテストして、計算時間が幾何級数的に増える割合を見積もれる。
その係数を使って最長の有効な入力にかかる時間を外挿し、それが最大のしきい値を下回っていればテストは合格とする。

\begin{exa}[計算時間の外挿]\label{exa:ch12-extrapolate-runtime}
仮に、$N = 20$ で 1 秒、$N = 21$ で 2 秒かかるとしよう。つまり 1 文字増えるごとに実行時間が 2 倍になる。
入力の最大長が 30 文字なら、処理には 1,024（$2^{10}$）秒かかると見積もれるので、それが許容できるかどうかを判断できる。
$N = 30$ の場合を実際に実行する代わりに処理時間を数学的に外挿すれば、極端に遅いテストケースを避けられる。
\end{exa}

\begin{rem}[測定の回数]\label{rem:ch12-more-measurements}
実際の処理時間はほかの要因にも左右されることがあるので、適切なモデルを検証するには 2 回より多く測定する必要がある場合もある点に注意すること。
\end{rem}

この種の的を絞ったテストに加えて、システム全体の性能指標を測り、余裕を持たせた上限を設定しておく。そうすれば、ある反復で性能が大きく劣化したときに、テストがそれを検査すべき問題として知らせてくれる。
多くの場合、これらの測定は、スモークテスト、負荷テスト、互換性テストなど、既存の大規模なテストに簡単に組み込める。

コードの変更によってメモリ消費が急増するのを防ぐ簡単な手法の一つは、資源を人為的に制約した条件でテストを実行することである。
ここでいうメモリとは、スタック領域やヒープ領域、スワップ領域、ディスク上のファイルやデータベースなどを指す。
単体テストは、使えるメモリを少なくして実行すべきである。テストスイートがその限界に達することがあれば、調べてみる価値がある。
大規模な統合テストには本番環境並みの資源が必要になるが、余裕を最小限にして実行すれば「炭鉱のカナリア」として役立てられる。
たとえば、本番環境で使えるメモリの 80 パーセントでシステムのテストに成功すれば、20 パーセントの余裕（予備の容量）があるという、ある程度の保証が得られる。

\subsubsection{しきい値のテスト}
% src: Threshold Testing

システムの可用性を守るうえで重要なのに見落とされやすいのが、根本的な限界に達する前に警告が出るようにしておくことである。

\begin{exa}[32 ビットカウンターのラップアラウンド]\label{exa:ch12-counter-wrap}
こうした限界を超えてしまった古典的な例が、少し前にある有名なソフトウェア企業で起きた。
システムが管理するオブジェクトに一意の ID を割り当てる 32 ビットのカウンターが 2,147,483,647 から 0 に一周し、その結果、番号の小さいオブジェクトの ID が重複してしまったのである。
復旧には何時間もかかった。カウンターが限界に近づいていないかを監視し、たとえば \texttt{0.99*INT\_MAX} に達した時点で警告を出していれば、たやすく避けられたはずの惨事である。
製品の初期には、カウンターが最大値に達するなど想像しにくかったに違いないが、会社が成長してそれが現実の問題になりうる段階になっても、誰もその可能性を考えなかった。
\end{exa}

このようなしきい値の警告は、セキュリティテストではなく運用監視の仕事だと考えられることが多い。
しかし、あまりにも頻繁に見落とされるうえ、修正もごく簡単なので、両方の枠組みでこうした事態に備えておく価値がある場合が多い。
カウンターだけでなく、システムが壁に激突しかねないほかの限界にも必ず注意を払うこと。

ストレージ容量も、余裕をもって対応できるよう、早めに警告がほしい領域である。
限界の 99 パーセントといった恣意的なしきい値を決めるよりも、\underline{時系列}（time series、時間を追って記録した一連の測定値）を見て、限界に達するまでの時間を外挿するほうが役に立つ。

時間に関する限界にも先手を打つことを忘れてはならない。
デジタル証明書の有効期限は、ある日突然検証に失敗するようになるまで、つい見過ごされがちである。
データフィードを提供するパートナーの証明書に依存しているシステムは、その有効期限を監視し、事前に通知するようにしておくべきである。そうすれば、顧客の目には自社の問題に見える障害を避けられる。

「Y2K バグ」は、今では何も起きなかった出来事として遠い記憶になっている（年を 2 桁の値で保存していたコンピューターシステムで、1999 年から 2000 年に変わる際に起きかねなかった混乱を避けるため、当時並外れた努力が払われたおかげかもしれない）。
しかし今度は、2038 年 1 月 19 日に「Y2k38 バグ」が待ち構えている。
この日、1970 年 1 月 1 日 00:00:00 UTC（Unix エポック）から 2,147,483,647 秒が経過する。
20 年足らずのうちに、エポックからの経過秒数が 32 ビットの数の範囲からあふれる時点が来る。そうなれば、ほぼ確実にあらゆる種類の厄介なバグとなって現れるだろう。
今コードベースにその備えとしての計測や監視を仕込むのが早すぎるというなら、いったいいつがその時なのだろうか。

% 図 12-1: Bug (courtesy of Randall Munroe, xkcd.com/376)
%   （xkcd の漫画。画像のみのため収録しない）

\subsubsection{分散型サービス拒否攻撃}
% src: Distributed Denial-of-Service Attacks

\begin{dfn}[分散型サービス拒否攻撃]\label{dfn:ch12-ddos}
サービス拒否（DoS）攻撃は、可用性を損なう単一の行為である。
これに対し、\underline{分散型サービス拒否攻撃}（distributed denial-of-service, DDoS）は、協調した多数の行為の効果を積み重ねることで同じ結果をもたらす。
\end{dfn}

インターネットに接続されたシステムには、インターネットの開かれたアーキテクチャのせいで、協調したボットネットなどからの DDoS 攻撃という新たなリスクも加わる。
分散した匿名の発信源から力ずくで過負荷をかけられると、たいていは計算能力の規模の勝負になる。
これらの攻撃を緩和するには、ふつう、膨大なデータセンターの容量に支えられたネットワークの専門知識を持つ DDoS 対策ベンダーに頼る必要がある。

筆者がこれをほかの可用性の脅威とは別に取り上げるのは、運悪くサーバーが本格的な DDoS 攻撃の標的になってしまった場合、自力で簡単に緩和できるものではないからである。

\subsection{セキュリティテストのベストプラクティス}
% src: Best Practices for Security Testing

しっかりしたセキュリティテストケースを書くことは、どんなコードベースでもセキュリティを高める重要な方法である。
セキュリティテストケースで完璧なセキュリティを保証することはできないが、保護策や緩和策がきちんと働いていることは確かめられるので、正しい方向への大きな一歩になる。
充実したセキュリティテストケースのスイートとセキュリティリグレッションテストを組み合わせれば、重大なセキュリティ上の失態が起きる可能性は劇的に下がる。

\subsubsection{テスト駆動開発}
% src: Test-Driven Development

セキュリティテストケースは、決定的に重要なコードを書きながら、そのセキュリティ上の意味をじっくり考えているときに特に重要になる。

\begin{dfn}[テスト駆動開発]\label{dfn:ch12-tdd}
新しいコードと並行してテストケースを書く開発手法を\underline{テスト駆動開発}（test-driven development, TDD）という。
この手法を厳格に実践する人は、実際にはテストを先に作り、最初は失敗するそのテストを通すためだけに新しいコードを書く。
\end{dfn}

筆者は TDD を強く支持する。
最初からセキュリティテストケースを含めて TDD を行えば、セキュリティが後付けではなく、コードに組み込まれていることが保証される。
とはいえ、テストにどんな方法論を使うにしても、セキュリティテストケースは必ずテストスイートに含めなければならない。

ほかの人がテストを書く場合は、必要なセキュリティテストケースを説明した手引きを開発者が用意すべきである。
コードに対するセキュリティ上の要求をよく理解していないと、セキュリティテストケースを直感的に思いつくのは難しいからである。

\subsubsection{統合テストの活用}
% src: Leveraging Integration Testing

\underline{統合テスト}（integration testing）は、個別に単体テストを済ませたコンポーネントがすべて、想定どおりに連携して動くことを確かめるために、システムを一通り動かしてみるテストである。
品質保証のうえで重要なテストだが、いったんその手間をかけたなら、ちょっとしたセキュリティテストまで広げるのは簡単である。

\begin{exa}[ログへのパスワード流出]\label{exa:ch12-password-log-leak}
2018 年、ある大手ソーシャルメディアプラットフォームは、自ら招いたセキュリティ侵害のために、顧客にパスワードの変更を呼びかけた：バグのせいで、アカウントのパスワードが平文のまま内部ログに書き出されていたのである。
統合テストをうまく活用していれば、この脆弱性を持ち込んだコードを本番環境にリリースする前に、たやすく見つけて修正できたはずである。
このサービスの統合テストには、偽のユーザーアカウント（たとえば \texttt{USER1}）で何らかのパスワード（たとえば \texttt{/123!abc\$XYZ}）を使ってログインする手順を含めておくべきだった（偽のアカウントでもセキュアなパスワードを使うべきである）。
テストが終わったら、セキュリティテストで出力をスキャンしてその特徴的なパスワード文字列を探し、一致するものが見つかればエラーにする。
\end{exa}

このテスト手法はログファイルに限らず、漏えいが起こりうるあらゆる場所に使える：ほかに残っているファイル、誰でもアクセスできる Web ページ、クライアントのキャッシュなどである。
こうしたテストは \texttt{grep(1)} コマンドで済む程度の簡単なものでよい。

パスワードは説明に都合のよい例だが、この手法はあらゆるプライベートなデータに使える。
テストシステムには本番環境の実際のユーザーデータの代わりとなる大量の合成データが必要であり、そのプライベートな内容はどれも、まったく同じように漏えいするおそれがある。
もっと包括的な漏えいテストでは、機密として明示的に保護されていないシステムの出力をすべてスキャンし、プライベートなテスト入力データの痕跡がないかを調べる。

\subsubsection{セキュリティテストの遅れを取り戻す}
% src: Security Testing Catch-Up

セキュリティテストケースがまったくないコードベースを扱っているなら、セキュリティを優先するのであれば、やるべき重要な作業がある。
セキュリティを考慮した設計があり、脅威モデリングとレビューを経ているなら、その設計を、どのコードから手を付けるべきかを示す地図として使う。
ふつうは、作業を段階的なマイルストーンのある部分に分け、達成できそうな最初の 1、2 回の反復をこなし、進めながら残りに何が必要かを見極めるのが賢明である。

保護メカニズムと機能領域を重要なものから順に対象にし、何をテストすべきかはコードを見ながら判断する。
既存のテストケースも見直すこと。すでに何らかのセキュリティテストを行っているものもあるかもしれないし、少し手直しすればセキュリティ用に使えるものもあるかもしれない。
プロジェクトに新しく加わってコードを覚える必要がある人がいれば、セキュリティテストケースの一部を書いてもらうとよい。その人にとって優れた学習の機会になり、長く役立つ成果も残る。

\newpage

%% =======================================================
%%  part: ch13
%% =======================================================
%% =======================================================
%%  part: ch13
%% =======================================================
\section{セキュアな開発のベストプラクティス}
% src: 13: Secure Development Best Practices

この章では、\cite{Kohnfelder2021} の第13章に沿って、開発プロセスそのものとセキュリティの関係を扱う。

\begin{quote}
「誰も完璧ではないと言われる。そのくせ、練習すれば完璧になるとも言われる。どちらかにしてほしいものだ。」---Winston Churchill
\end{quote}

ここまで本書では、開発フェーズで生じるさまざまなセキュリティ脆弱性を概観してきた。
この章では、開発プロセスそのものがセキュリティとどう関わり、どこでうまくいかなくなりうるかに焦点を当てる。
まずコード品質を取り上げる。具体的には、コード衛生を保つことの価値、エラーと例外を漏れなく処理すること、セキュリティ上の性質を文書化すること、そしてセキュリティを高めるうえでのコードレビューの役割である。
次に依存関係の扱い、とくに依存関係がシステムにどのように脆弱性を持ち込むかを見る。
3つ目はバグのトリアージである。これは、セキュリティとほかの差し迫った要求との兼ね合いをとるための重要な技能である。
最後に、セキュアな開発は安全な作業環境があってこそ成り立つので、侵害されないために何をすべきかについて基本的な助言を述べる。

以下の指針は、実用上の理由から一般的なものにとどめてある。
読者はこれを自分の開発実務に当てはめられるはずである。
ほかにも有効な手法は数多くあるが、その多くはプログラミング言語やオペレーティングシステムなど、システム固有の事情に左右される。
そのため、以下の議論にある大きなパターンをつかむことが重要であると同時に、自分の仕事の中で出てくるセキュリティ関連の問題や改善の機会にも目を配る必要がある。

%----------------------------------------------------------------------
\subsection{コード品質}
% src: Code Quality

\begin{quote}
「品質は決して流行遅れにならない。」---Robert Genn
\end{quote}

第3部のこれまでの章では、脆弱性がコードに紛れ込むさまざまな経路を説明してきた。ここでは、バグ全般とセキュリティとの関係に焦点を当てたい。
コードの品質を高められれば、意識しているかどうかにかかわらず、長い目で見てコードはより安全になる。
脆弱性はすべてバグであり、バグが減れば脆弱性も脆弱性チェーンも減る。
とはいえ当然、バグをすべてなくすよりずっと手前で収穫逓減が始まるので、バランスのとれたやり方をとるのが一番である。

以下では、セキュリティのために力を入れるべき主な領域をいくつか取り上げる。

\subsubsection{コード衛生}
% src: Code Hygiene

プログラマーはたいてい、自分の扱っているコードの品質をよくわかっている。それでもさまざまな理由から、必要な改善をせずに既知の欠陥をそのままにしておきがちである。
\underline{コードの臭い}（code smell）や\underline{スパゲッティコード}（spaghetti code）、まだ作業が残っていることを示す先送りされた「TODO」コメントは、どれも脆弱性の温床になりやすい。
少なくともセキュリティがとくに気になる領域では、こうした粗い部分を見つけて整えておくことが、脆弱性を避ける最善の方法の一つになりうる。
しかもこの方法なら、バグがどう悪用されうるかを見極めるセキュリティ分析をしなくて済む。

コードの状態についての自分の感覚に加えて、ツールを使ってこうした問題を検出するとよい。
コードはすべての警告を有効にしてコンパイルし、指摘された問題がすべて解消されるようにコードを直す。
こうした自動警告のなかには、誤解を招くインデントや、実行経路がなく使われないコードについての警告がある。こうした警告があれば、第8章と第12章で扱った GotoFail 脆弱性は見つかっていたはずである。
Lint などの\underline{静的コード解析}（static code analysis）ツールを使えばコードをさらに詳しく精査でき、ときにはバグや脆弱性の発見につながるヒントが得られる。

コード解析は、セキュリティバグをセキュリティバグとして見分けてくれるとは限らない。そのため、もっと広く網を張る必要がある。
潜在的なバグの総数を減らすため、開発中はこうしたツールをこまめに使うこと。
そうしておけば、ツールの出力が大きく変わったときに気付きやすい。
新しい出力が、古いメッセージの洪水に埋もれてしまわないからである。

簡単に直せる場合や、問題が深刻になりうるとわかる場合には、警告をすべて直すこと。
たとえば到達不能なコードがあるということは、誰かが何かの理由でそのコードを書いたのに、今ではまったく使われていないということであり、そんな状態が正しいはずはない。
一方、変数の命名規則についての警告は、よい提案ではあっても、セキュリティ脆弱性につながることはまずない。

この種の整理に充てる時間を確保するのは、いつでも難しい。
少しずつ進めるとよい。
週に1〜2時間でも続ければいずれ大きな差になるし、大きなコードベースに慣れるよい方法にもなる。
警告が多すぎて対処しきれない場合は、最も効果が見込めるもの（たとえば GCC の \texttt{-Wmisleading-indentation}）から始めて、指摘された箇所を直すこと。

\subsubsection{例外とエラーの処理}
% src: Exception and Error Handling

1996年の Ariane 5 Flight 501 Failure Report（アリアン5型ロケット501便の事故報告書）には、ずさんな例外処理が招いた結果が痛々しいほど詳しく記されている。
この破滅的なバグは、悪意ある行為者が関わらない純粋な自損事故であった。しかし、そこで生じた挙動を攻撃者が悪用すればシステムを侵害しうることを示す例にもなっている。

\begin{exa}[Ariane 5 Flight 501 の事故]\label{exa:ch13-ariane5}
Ariane 5 の打ち上げ直後、ある計算で浮動小数点数を整数に変換したところ、例外が発生した。
例外処理の仕組みは働いたが、この変換エラーは想定外だったので、例外ハンドラのコードにはこの状況への備えがなかった。
コードはエンジンを停止させ、飛行開始から 36.7 秒後に壊滅的な失敗に至った。
\end{exa}

こうした問題を防ぐには、まず場当たり的な例外処理の危険を認識し、最も起こりそうにない例外についてさえ、正しい対応を考え抜くことである。
一般に、例外はできるだけ発生源の近くで処理するのが最善である。
発生源の近くには最も多くの文脈があり、さらに別の問題が起きるまでの時間も最も短いからである。

とはいえ大規模なシステムでは、上位へ伝わってくる未処理の例外をすべて受け止めるトップレベルのハンドラが必要になることもある。
これをうまく行う方法の一つは、処理の単位を決めておき、その単位全体を失敗させることである。
たとえば Web サーバーなら、HTTP リクエストの処理中に発生した例外を捕捉し、汎用の 500（サーバーエラー）応答を返せばよい。
一般に Web アプリケーションは、状態を変更するリクエストをトランザクションとして扱い、どんなエラーが起きても状態がまったく変わらないようにすべきである。
こうすれば、システムを脆い状態に置きかねない中途半端な変更を避けられる。

ずさんな例外処理が潜在的な脆弱性につながるという理屈の多くは、エラー処理全般にも当てはまる。
エラーケースも例外と同じく、めったに起きないことがある。そのため開発者に忘れられやすく、不完全なまま、あるいはテストされないまま残りやすい。
攻撃者がエクスプロイトを探すときによく使う手口は、何らかのエラーを起こそうとして、そのときのコードの振る舞いを観察し、弱点を見つけることである。
したがって、最善の防御は最初からしっかりしたエラー処理を実装しておくことである。

\begin{rem}[エラーが攻撃の目的になる]\label{rem:ch13-error-as-goal}
これは、セキュリティ脆弱性がほかのバグとどう違うかを示す典型例である。
通常の使い方ではきわめてまれにしか起きないエラーでも、組織的な攻撃の場面では、そのエラーを起こすこと自体がはっきりした目的になりうる。
\end{rem}

エラーと例外を正しく処理するには、しっかりしたテストが重要である。
すべてのコード経路、とくにあまり通らない経路がテストでカバーされているか確認すること。
本番環境の例外ログを監視して原因を突き止め、例外からの回復が正しく機能しているか確かめること。
ときどきしか発生しない例外は、積極的に調査して直すこと。
賢い攻撃者がその例外を起こす方法を知れば、そこから悪意あるエクスプロイトに仕立て上げるかもしれないからである。

\subsubsection{セキュリティの文書化}
% src: Documenting Security

セキュリティ上重要な意味を持つコードを書くとき、自分の判断をコメントでどこまで説明しておけばよいだろうか。ほかの人（あるいは数か月後、数年後の、忘れっぽくなった自分）がうっかり壊してしまわないようにするためである。

重要なコードや、セキュリティ上の影響を説明しておくべき箇所には、コメントを書くことが大切である。
コメントがあれば、そのコードを変更しようとする人は誰でも、何が懸かっているかを理解できる。
セキュリティについてコメントを書くときは、セキュリティ上の影響を具体的に説明すること。
単に \verb|// Beware: security consequences| と書くだけでは説明にならない。
はっきりと、要点を絞って書くこと。
言葉を重ねすぎると、読み流されるか、読むのをやめられてしまう。

\begin{exa}[Heartbleed に対する良いコメント]\label{exa:ch13-heartbleed-comment}
第10章と第12章で扱った Heartbleed バグでいえば、良いコメントとは次のことを説明するものである。実際に渡されたデータより大きいバイト数を指定した不正なリクエストを拒否することが決定的に重要であり、拒否しなければバッファの範囲を超えて秘密のデータが漏えいしうる。
\end{exa}

セキュリティ分析が複雑になりすぎてコメントでは説明しきれない場合は、詳細を別の文書にまとめ、コメントにはその文書への参照を書いておくこと。

とはいえ、セキュリティが依存しているコードすべてに印を付けようとすべきだ、という意味ではない。
むしろ、将来見落とされやすい、自明でない問題について読み手に注意を促すことを目指すべきである。
結局のところ、セキュリティ上の影響に常に気を配っている知識あるプログラマーの代わりを、コメントが完全に務めることはできない。
だからこそ、この作業は簡単ではないのである。

良いセキュリティテストケース（第12章で扱った）を書くことは、文書を裏付けると同時に、将来の変更でほかの人が知らずにセキュリティを壊してしまうのを防ぐ仕組みを用意する、理想的な方法である。
そうしたテストは攻撃がどのようなものかを示す動く模型であり、意図しない悪い変更を防ぐだけでなく、コードがどうおかしくなりうるかを正確に示す役割も果たす。

\subsubsection{セキュリティコードレビュー}
% src: Security Code Reviews

プロのソフトウェア開発プロセスには、同僚によるコードレビューが標準的な実践として含まれている。
著者は、そのレビューにセキュリティをはっきりと含めるべきだと主張したい。
ふつうは、コードレビューのワークフローの一つのステップとして行うのが最もよい。コードの正しさ、読みやすさ、スタイルなど、レビュアーが注意すべき問題のチェックリストと並べて扱うのである。

著者は、同じコードレビュアーがセキュリティを検討するステップを明示的に設けることを勧める。
典型的には、まずコードを一通り読み、そのあと「セキュリティの帽子」（security hat）をかぶってもう一度読み直す。
レビュアーがセキュリティを見る自信がない場合は、その部分を力のある人に任せるべきである。
もちろん、明らかにセキュリティに関係しないコード変更なら、このステップは省いてよい。

コード変更をセキュリティの観点からレビューするのは、セキュリティ設計レビュー（SDR、第7章のテーマ）とは違う。
設計全体をレビューするときのような全体像がないまま、システムのごく一部分を見ることになるからである。
コードがさまざまな信頼できない入力をどう扱うかを必ず検討し、入力検証が堅牢であることを確かめ、混乱した使節（Confused Deputy）の問題が起きそうな箇所をなくすこと。
当然、セキュリティにとって重要なコードにはとくに注意を払うべきであり、たいていはより高い品質基準を課すだけの価値がある。
コードのセキュリティをもう一組の目で見る機会は、システム全体をよくする大きな可能性を秘めている。

コードレビューは、すでに作成されているセキュリティテストケース（第12章で述べた）が十分かどうかを確かめる絶好の機会でもある。
レビュアーとして、ある入力が問題を起こすかもしれないと考えたら、推測で済ませずに、セキュリティテストケースを書いて何が起きるか確かめること。
その探索的なテストケースで脆弱性が見つかったら、問題を報告し、確実に修正されるようそのテストケースも提供すること。

%----------------------------------------------------------------------
\subsection{依存関係}
% src: Dependencies

\begin{quote}
「依存は隷属につながる。」---Thomas Jefferson
\end{quote}

現代のシステムは、外部コンポーネントを大量に積み重ねて作られることが多い。
\underline{依存関係}（dependency）には、一つどころではない問題がある。
npm をはじめ多くのプラットフォームは、追跡しにくい大量の依存関係を自動的に取り込む。
また、既知の脆弱性を抱えた古いバージョンの外部コードを使い続けることは、業界がいまだに体系的になくせていない、最大級の慢性的な脅威の一つである。
さらに、ソフトウェアの\underline{サプライチェーン}（software supply chain）を通じて、悪意あるコンポーネントを取り込んでしまう危険もある。
これにはいくつかの経路がある。
たとえば、有名なパッケージに似せた名前のパッケージを誤って選んでしまうこともあれば、ほかのコンポーネントの依存関係を通じて間接的にマルウェアを取り込むこともある。

システムにコンポーネントを追加すると、たとえセキュリティを強化するためのコンポーネントであっても、セキュリティが損なわれるおそれがある。
コンポーネントの提供元を信頼するだけでなく、その提供元が信頼しているものすべてまで信頼しなければならないからである。
コードが増えればバグも全体の複雑さも増すという避けがたい危険に加えて、コンポーネントは思いもよらない新しい形で攻撃対象領域を広げることがある。
バイナリ配布物は事実上中身が見えない。ソースコードとドキュメントがあっても、パッケージの中身をすべて注意深くレビューして理解するのはたいてい現実的でなく、結局は盲目的に信頼することになりがちである。

\begin{exa}[アンチウイルスソフトウェア]\label{exa:ch13-antivirus}
アンチウイルスソフトウェアはマルウェアを検出してブロックできる。その一方で、システムの奥深くまで入り込むフックを広く使い、スーパーユーザー権限を必要とし、攻撃対象領域を広げてしまうこともある。
たとえば、最新のマルウェアデータベースを取得したり検出結果を報告したりするために、外部へ通信（phone home）する場合がそうである。
\end{exa}

脆弱なコンポーネントを軽率に選ぶと、防御層を足すつもりだったのに、かえってセキュリティを下げてしまうことがある。

\subsubsection{安全なコンポーネントの選択}
% src: Choosing Secure Components

システム全体が安全であるためには、それを構成する各コンポーネントが安全でなければならない。
さらに、コンポーネント間のインターフェースも安全でなければならない。
安全なコンポーネントを選ぶときに考えるべき基本的な要素を、以下に挙げる。
\begin{itemize}
\item そのコンポーネントと作り手は、セキュリティに関してどのような実績を持っているか。
\item コンポーネントのインターフェースは独自仕様か、それとも互換性のある代替品があるか（選択肢が多いほど、より安全な代替品がある可能性も高い）。
\item コンポーネントにセキュリティ脆弱性が見つかったとき（見つかるかどうかではなく、いつ見つかるかの問題である）、開発者がすばやく対応して修正版をリリースしてくれると信頼できるか。
\item コンポーネントを最新に保つための運用コスト（つまり、労力、ダウンタイム、費用）はどれくらいか。
\end{itemize}

コンポーネントは、セキュリティの観点を念頭に置いて選ぶことが重要である。
秘密のデータを処理するのに使うコンポーネントは、情報漏えいを起こさないことを保証すべきである。
データ処理の副作用として内容をログに記録したり、保護されていないストレージに保存したりするなら、それだけで漏えいの可能性がある。
たとえば、プライバシーの心配がまったくない海水温を扱うために書かれたソフトウェアを、機微な医療データに転用してはならない。
また、プロトタイプのコンポーネントなど、高品質な本番リリース以外のものは避けること。

\subsubsection{インターフェースの保護}
% src: Securing Interfaces

よく文書化されたインターフェースであれば、セキュリティとプライバシーに関する性質を明記しているべきだが、実際にはそうなっていないことが多い。
プログラマーは効率を優先して入力検証を省きがちであり、検証はもう済んでいると思い込んでいる場合はとくにそうである。
かといって、すべてのインターフェースに重複した入力検証をさせるのは、たしかに無駄である。
確信が持てないときは、できればテストしてインターフェースの挙動を確かめ、それでも疑わしければ、念のためインターフェースの手前に入力検証の層を追加すること。

\underline{非推奨}（deprecated）の API は使わないこと。
非推奨の API には、潜在的なセキュリティ問題が隠れていることが多いからである。
API の提供者は、安全でない機能を含む API を完全に削除せず、非推奨にとどめるのが一般的である。
そうすれば、既存の API 利用者に対しては後方互換性を保ちつつ、ほかの人が脆弱なコードを使わないよう促せる。
もちろん、非推奨になる理由はほかにもある。しかし API の利用者としては、非推奨になった理由にセキュリティが関わっているかどうかを調べることが重要である。
攻撃者も API の非推奨化を追いかけて、攻撃の準備をしているかもしれないことを忘れてはならない。

こうした基本的な例のほかにも、インターフェースが内部構造を公開している場合は、いつでもとくに注意すること。
公開された内部構造は意図しない形で使われやすく、簡単に脆弱性を生むからである。

\begin{exa}[The Most Dangerous Code in the World]\label{exa:ch13-most-dangerous-code}
論文 ``The Most Dangerous Code in the World'' は優れたケーススタディである。
研究者たちは、広く使われている SSL ライブラリが何度も安全でない形で使われ、本来提供すべきセキュリティ上の性質がまったく得られていないことを発見した。
論文の著者らは、「これらの脆弱性の大半の根本原因は、基盤となる SSL ライブラリの API のひどい設計にある」と結論づけている。
\end{exa}

複雑な設定オプションを持つ API にも用心すること。
とくに、セキュリティがその設定次第で決まる場合はそうである。
自分で API を設計するときは、セキュアバイデフォルト（Secure by Default）のパターンを守り、システムを安全に設定する方法を文書化し、必要に応じて正しい設定を保証するヘルパーメソッドを用意すること。
安全でないかもしれない機能をどうしても公開しなければならないときは、自分が何をしているかを正確に理解しないまま使うことが事実上ありえないよう、あらゆる手を尽くすこと。

\subsubsection{セキュリティの車輪を再発明しない}
% src: Don’t Reinvent Security Wheels

基本的なセキュリティ機能には、できるかぎり標準的で高品質なライブラリを使うこと。
たとえば、クエリパラメータでの XSS 攻撃を一から緩和しようとすると、そのたびに、たとえ HTML の構文を隅々まで知っている人でも、目立たない形の攻撃を見落とす危険がある。

よい解決策が手に入らない場合は、特定の潜在的欠陥に対処するライブラリを作ってコードベース全体で使うことを検討し、そのライブラリを徹底的にテストすること。
場合によっては、脆弱性になりやすい特定の欠陥をコードから見つけるのに、自動化ツールが役立つ。
たとえば、C のコードから古い「安全でない」文字列関数（\texttt{strcpy} など）を探し出し、同じ機能を持つ新しい「安全な」版（\texttt{strlcpy}）に置き換える。

ライブラリやフレームワークを書いているなら、セキュリティ上の弱点を注意深く探し、一か所で確実に正しく処理されるようにすること。
そのうえで、どのような保護を提供し、どのような保護を提供しないのかを明記すること。
「このライブラリを使えばセキュリティの心配はすべて解決します」と宣伝するだけでは役に立たない。
利用者があなたのコードに頼っているなら、何がきちんと処理され、何が処理されないのかを、利用者はどうやって知ればよいのか。

\begin{exa}[Web フレームワークが文書化すべきこと]\label{exa:ch13-web-framework-docs}
Web フレームワークであれば、セッション管理に Cookie をどう使うか、XSS をどう防ぐか、CSRF 対策のナンスをどう提供するか、HTTPS だけを使うか、などを記述すべきである。
\end{exa}

すべての卵を一つのかごに盛るように思えるかもしれないが、潜在的なセキュリティ問題はライブラリやフレームワークで一度だけ解決するのがふつうは最善である。
そうした層を一貫して使えば自然なボトルネックができ、潜在的な問題が起きうる箇所すべてに対処できる。
あとで新しい脆弱性が見つかっても、共通コードを一か所変えるだけで済む。その変更は修正もテストも容易で、すべての使用箇所をカバーできるはずである。

セキュリティを意識したライブラリでも、完全には保護しきれない基盤機能への生のアクセスを提供しなければならないことがある。
たとえば HTML フレームワークのテンプレートでは、アプリケーションが任意の HTML を挿入できる場合がある。
これが必要なら、通常の保護が効かなくなる箇所をすべて徹底的に文書化し、API の利用者が負う責任を説明すること。
理想的には、\texttt{unsafe\_raw\_html} のように、危険をはっきりと伝える名前を API に付けるとよい。

要するに、セキュリティ脆弱性はわかりにくいことがあり、考えられる攻撃は数多く、しかもそのうち一つでも成功すれば攻撃者には十分なのである。
だから、こうした課題に自分一人で取り組むのは避けるのが賢明である。
同じ理由で、誰かがある問題をうまく解決しているなら、それを汎用の解決策として再利用するのが賢い。
人為的ミスは攻撃者の味方なので、安全なやり方を簡単にとれる解決策を使うのが一番である。

\subsubsection{レガシーなセキュリティへの対処}
% src: Contending with Legacy Security

デジタル技術は急速に進化するが、セキュリティツールはいくつかの理由で後れをとりがちである。
これは、いまも続く重要な課題である。
ことわざにある「ぬるま湯の中のカエル」のように、誰かが厳しく見直し、危険をはっきり指摘し、より安全な解決策と移行計画を提案しない限り、レガシーなセキュリティ手法はあまりにも長く使われ続けることが多い。

誤解のないように言えば、既存のセキュリティ手法が必ずしも弱いと言っているのではない。ほとんどすべてのものには「賞味期限」がある、と言っているのである。
さらに、変化し続ける脅威の状況に照らして、既存のシステムを定期的に見直す必要がある。

\begin{exa}[セキュリティの賞味期限]\label{exa:ch13-sell-by-date}
パスワードによる認証は、フィッシング攻撃を受けやすくなってきたなら、第2要素で補強する必要があるかもしれない。
暗号の実装は、その時点のハードウェアのコストと能力の見積もりに基づいているが、ムーアの法則が示すとおり、その前提は常に動き続ける標的である。
量子コンピューティングが成熟するにつれ、高いセキュリティが求められるシステムはすでに、この新技術にも耐えられると考えられている耐量子（post-quantum）アルゴリズムへの移行を進めている。
\end{exa}

弱いセキュリティが有効期限をとうに過ぎても残り続けるのには、いくつかの理由がある。
第一に、慣性は強力な力である。
システムはたいてい少しずつ進化するので、認証や認可の今のやり方を誰も疑わない。
第二に、エンタープライズのセキュリティアーキテクチャではふつう、すべてのサブシステムに互換性が求められる。そのため何かを変えようとすると、すべてのコンポーネントを新しい方式で相互運用できるよう改修しなければならない。
これはしばしば途方もない作業に思え、強い抵抗を招く。

また、古いサブコンポーネントも問題になりうる。
レガシーなハードウェアやソフトウェアは、新しいセキュリティ技術に対応していないことがあるからである。
加えて、今のセキュリティでこれまで問題なかったのだから、壊れていないものを直す必要はない、という安易な反論もある。
さらに、レガシーなセキュリティを設計した人がもういなくなっていて、ほかに誰もその仕組みを完全には理解していないかもしれない。
あるいは、元の設計者がいたとしても、自分の仕事を守ろうとする態度をとるかもしれない。

こうした懸念すべてに答えられる単純な答えはない。しかし、弱いレガシーセキュリティの具体的な問題を脅威モデリングで洗い出せれば、それがもたらすリスクははっきりするはずである。

レガシーコードを段階的に廃止すべきだと判断したら、移行の計画を立てる必要がある。
互換性のあるインターフェースを持つ新しいコンポーネントをコードベースに組み込めれば作業は楽になるが、それができない場合もある。
場合によっては、より良いセキュリティを段階的に実装するのがよいやり方である。
システムの各部分を少しずつ新しい実装に切り替えていき、レガシーコードが不要になった時点で削除する。

%----------------------------------------------------------------------
\subsection{脆弱性のトリアージ}
% src: Vulnerability Triage

\begin{quote}
「『トリアージ』という言葉はふつう、誰を最初に手当てするかを決めることを意味する。」---Bill Dedman
\end{quote}

ほとんどのセキュリティ問題は、いったん特定されれば直すのは簡単で、チームもその直し方にすんなり合意できる。
しかし、ときにはセキュリティ問題をめぐって意見が分かれることもある。
とくに、バグが悪用できるかどうかがはっきりしない、あるいは修正が難しいという、どっちつかずの領域でそうなりやすい。
急がなければならない重大な事情がない限り、少しでも悪用される可能性のあるバグは、どれでも直すのが一般に賢明である。
第8章で見たように、小さなバグがいくつか組み合わさって重大な脆弱性を生む脆弱性チェーンが生じうることを、心に留めておくこと。
そして、自分にバグの悪用方法が思いつかないからといって、決意を固めた攻撃者にもできないという証明には決してならないことを、常に忘れないこと。

\subsubsection{DREAD による評価}
% src: DREAD Assessments

バグを直すかどうかについてチームがすぐに合意できないというまれなケースでは、そのバグがもたらすリスクを構造化された方法で評価すること。

\begin{dfn}[DREAD]\label{dfn:ch13-dread}
\underline{DREAD} モデルは、特定の脅威のリスクを評価するための単純なツールである。
もともと Jason Taylor が考案し、Microsoft で彼と著者の二人が広めたものである。
DREAD は、脆弱性がもたらすリスクを 5 つの側面に分けて挙げる。
\begin{itemize}
\item \underline{被害の大きさ}（Damage potential）：悪用されたら、どれほどひどいことになるか。
\item \underline{再現性}（Reproducibility）：攻撃は毎回成功するか、ときどき成功するか、それともまれにしか成功しないか。
\item \underline{悪用のしやすさ}（Exploitability）：技術的な難しさ、労力、コストの点で、その脆弱性を悪用するのはどれほど難しいか。攻撃経路はどれくらい長いか。
\item \underline{影響を受けるユーザー}（Affected users）：ユーザー全員が影響を受けるか、一部か、それともごく少数か。特定の相手を狙って攻撃しやすいか、それとも被害者は無作為か。
\item \underline{発見のしやすさ}（Discoverability）：攻撃者がその脆弱性を見つける可能性はどれくらいか。
\end{itemize}
\end{dfn}

著者の経験では、DREAD による評価は、5 つの独立した軸として考えるのが最もうまくいく。
著者個人としては、各項目に数値のスコアを付けることは勧めない。
深刻度はあまり線形ではないからである。
著者が好む方法は、以下の例のように、主観的な大きさを T シャツのサイズ（S、M、L、XL）で表すことである。

\begin{rem}[DREAD スコアを合計しない]\label{rem:ch13-dread-no-sum}
数値のスコアを使う場合でも、5 つのスコアを足し合わせ、その合計で脅威どうしの順位を付けることは、とくに勧めない。
それは本質的にリンゴとオレンジを比べるようなものだからである。
いくつかの要素の DREAD スコアがかなり低いのでない限り、その脅威は緩和する価値のある重大なものとみなすこと。
\end{rem}

問題を決着させるためにトリアージ会議が必要な場合は、DREAD を使って自分の主張を示すこと。
脆弱性がもたらす結果をはっきりつかむために、必要なら要素ごとに議論すること。
多くの場合、ある要素のスコアが低いと、それが全体の影響にとって何を意味するかに議論が集中する。

\begin{exa}[Heartbleed の DREAD 評価]\label{exa:ch13-heartbleed-dread}
DREAD が実際にどう使われるかを見てみよう。
Heartbleed バグを発見したばかりで、その DREAD 評価をしたいとする。
この脆弱性では、匿名の攻撃者が悪意ある Heartbeat リクエストを送ると、Web サーバーのメモリの大きな塊を受け取れるのであった。
情報漏えいの脅威に対する簡単な DREAD 評価を表 \ref{tab:ch13-heartbleed-dread} に示す。
\end{exa}

\begin{table}[H]
\centering
\small
\begin{tabularx}{\linewidth}{|l|c|X|}
\hline
要素 & 評価 & 理由 \\ \hline
被害の大きさ（Damage potential） & XL & サーバーの内部メモリから秘密鍵が漏えいするおそれがある。 \\ \hline
再現性（Reproducibility） & M & 漏えいするメモリの内容は多くの要因で変わる。無害な場合もあるが、予測はできない。 \\ \hline
悪用のしやすさ（Exploitability） & L & 匿名の攻撃者が単純なリクエストパケットを送るだけでよい。役に立つ秘密を抜き出すには多少の専門知識が要る。 \\ \hline
影響を受けるユーザー（Affected users） & XL & サーバーとすべてのユーザーが危険にさらされる。 \\ \hline
発見のしやすさ（Discoverability） & L & 攻撃者がこのアイデアを思いつくかどうか次第である（公表されれば明白）。試すのも確かめるのも簡単である。 \\ \hline
\end{tabularx}
\caption{Heartbleed の情報漏えいの脅威に対する DREAD 評価}\label{tab:ch13-heartbleed-dread}
\end{table}

この DREAD 評価は主観的なものである。
ここでの想定では、バグを手早く確認しただけで、それ以上脆弱性を調べる時間がなかったからである。
サーバーの鍵が漏えいすることは確認できた（だから被害の大きさは XL）が、何度かテストするとメモリの内容は大きく変わったので、再現性は M とした、という想定である。
発見のしやすさはとくに厄介である。
誰かがこれを試してみようと思いつく可能性を、どうやって測ればよいのか。
著者は、自分が思いついたのなら、ほかの人もそう遠くないうちに思いつくと考えておくのが最善だと主張する。

DREAD スコアをめぐる議論は、こうした判断の細かな違いを引き出すのに優れた方法である。
議論になったら、注意深く耳を傾け、ほかの意見も十分に考慮すること。
Heartbleed は史上最悪の脆弱性の一つなのに、DREAD の全要素を最大に評価したわけではない。
このことは、評価を注意深く読み解かなければならない理由をよく示している。
この欠陥は何百万台もの Web サーバーで動くコードにあり、HTTPS のセキュリティを損なったのだから、被害の大きさと影響を受けるユーザーのスコアは実際には桁外れ（たとえば XXXXXXXL）であり、いくつかの中程度の評価を補って余りある、と言うこともできる。
DREAD 評価の価値は、脆弱性のさまざまな側面がそれぞれどれほど重要かを明らかにし、その脆弱性がもたらすリスクをはっきり示すところにある。

\subsubsection{動作するエクスプロイトの作成}
% src: Crafting Working Exploits

実際に動く原理実証（proof-of-principle）の攻撃を作ってみせることは、脆弱性を直すべきだと主張する最も強力な方法である。
攻撃方法が明白なバグもあり、エクスプロイトを簡単にコードにできるなら、それで決着がつく。
しかし著者の意見では、そこまで必要になることはまれであり、それにはいくつかの理由がある。
まず、デモ用のエクスプロイトを作るにはふつう多くの作業が要る。
実際に動くエクスプロイトは、根本の脆弱性を特定したあとも、さんざん手を加えなければならないことが多い。
さらに重要なのは、たとえ経験豊富なペネトレーションテスターでも、動くエクスプロイトを作れなかったからといって、その脆弱性が悪用できないという証明には決してならないことである。

これは意見の分かれる話題だが、著者の考えでは、こうした理由から、セキュリティ脆弱性に対処するために動くエクスプロイトを作る手間をかけるのは正当化しにくい。
とはいえ、本格的な攻撃でなくても、バグを直接引き起こすリグレッションテスト（第12章で扱った）はぜひ書くべきである。

\subsubsection{トリアージの意思決定}
% src: Making Triage Decisions

DREAD を使う場合でも、ほかのどんな方法で脆弱性を評価する場合でも、実際の脅威は過大評価するより過小評価するほうがはるかに簡単だということを心に留めておくこと。
潜在的な脆弱性に気付いていながら何もしないのは悲劇的な誤りになりうるし、明らかに避けるべきことである。
著者はそうした議論に何度か負けたことがあり、あとになって「だから言ったのに」と言っても何の満足も得られないことは保証できる。
重大な欠陥を直さないのは、やる価値のないロシアンルーレットである。
「とにかく直す」は、いつでも通用するすばらしい方針である。

セキュリティのトリアージでよりよい判断を下すための、一般的な経験則を以下に挙げる。
\begin{itemize}
\item 特権コードや価値ある資産にアクセスするコードのバグは直すべきであり、新たなバグを持ち込まないよう注意深くテストすべきである。
\item どの攻撃対象領域からもよく隔離されていて、無害に見えるバグは、ふつうは後回しにしても安全である。
\item バグが無害だという主張は注意深く確かめること。そのバグが及ぼしうる影響全体を正確に見積もるより、直してしまうほうが簡単なこともある。
\item 脆弱性チェーン（第8章で扱った）の一部になりうるバグは、積極的に直すこと。
\item 最後に、判断が五分五分のときは、著者は必ず問題を直すよう助言する。転ばぬ先の杖である。
\end{itemize}

さらに調査が必要な場合は、誰かにその問題の調査を割り当て、提案を添えて報告してもらうこと。
仮定の話をめぐる議論に時間を浪費してはならない。
議論では、相手の考えを変えさせようとするより、ほかの視点を理解することに力を注ぐこと。
自分の直感を信じること。
練習を積んで何に注目すべきかがわかってくれば、すぐに楽にできるようになる。

%----------------------------------------------------------------------
\subsection{安全な開発環境の維持}
% src: Maintaining a Secure Development Environment

\begin{quote}
「風景の秘密は創造ではない……維持である。」---Michael Dolan
\end{quote}

衛生管理は、よいたとえになる。
安全な食品を作るには、信頼できる仕入れ先からの新鮮な原材料、衛生的な作業環境、殺菌された道具などが必要である。
同じように、できあがる製品が安全であるためには、開発プロセス全体を通じてよいセキュリティの実践を守らなければならない。

開発中にたった一度気を緩めただけでも、悪意あるコードが製品に紛れ込むことがある。
この事実には、立ち止まって考えさせられるはずである。
開発者にとって最も避けたいのは、自分のソフトウェアがマルウェアを運ぶ媒体になってしまうことである。

\subsubsection{開発環境と本番環境の分離}
% src: Separating Development from Production

まだやっていないなら、開発環境と本番環境を厳格に分けること。
中心となる考え方は両者の間に「壁」を設けることであり、典型的には別々のサブネットワークにするか、少なくとも互いに排他的なアクセス権限の体系にする。
つまり、ソフトウェアを開発しているプログラマーは、本番データにアクセスできるべきではない。
また、本番マシンや運用スタッフは、開発環境とソースコードにアクセス（書き込みアクセス）できるべきではない。
一人が本番と開発の両方を担う小さな組織なら、ユーザーアカウントを切り替えればよい。
切り替えの手間は、たった一度のミスから製品を守れるだけで十分に元が取れる。
しかも、安心感も得られる。

\subsubsection{開発ツールの保護}
% src: Securing Development Tools

開発ツールやライブラリのコードは、インストールして使う前に注意深く吟味すること。
「どこか」からダウンロードしたちょっとしたユーティリティは、一度使うだけであっても、得られるもの以上の問題をもたらしかねない。
中核の開発プロセスに含まれない実験や雑用のために、安全に隔離されたサンドボックスを用意することを検討するとよい。
仮想マシンを使えば簡単に用意できる。

できあがるものが安全であるためには、開発に関わるコンピューターがすべて安全でなければならない。
ソースコードリポジトリなどのサービスもすべて同様である。
これらはどれも、最終製品に脆弱性が忍び込む入り口になりうるからである。
実際にはもっと深く、開発プロセスに関わるオペレーティングシステム、コンパイラ、ライブラリもすべて安全でなければならない。
気の遠くなるような課題で、ほとんど不可能に聞こえるかもしれないが、幸い完璧を目指す必要はない。
こうしたリスクをすべて認識したうえで、少しずつ改善できる機会を見つけていかなければならない。

こうしたリスクを緩和する最善の方法は、開発環境と開発プロセスを脅威モデリングすることである。
ソースコードを主要な資産として扱い、さまざまな脅威について攻撃対象領域を分析すること。
典型的な開発作業に対する基本的な緩和策には、次のものがある。
\begin{itemize}
\item 開発用コンピューターを最新の状態に保ち、できるかぎり安全に設定する。
\item 開発に使うコンピューターの私的な利用を制限する。
\item 新しいコンポーネントや依存関係を体系的にレビューする。
\item ビルドとリリースのプロセスに使うコンピューターを安全に管理する。
\item 秘密情報（コード署名鍵など）を安全に管理する。
\item 強力な認証によって、ログインのクレデンシャル管理を保護する。
\item ソース変更のコミットを定期的に監査し、異常な活動がないか調べる。
\item ソースコードとビルド環境のバックアップを安全に保管する。
\end{itemize}

\subsubsection{製品のリリース}
% src: Releasing the Product

開発から本番への受け渡しには、正式なリリースプロセスを使うこと。
たとえば、開発スタッフだけが変更でき、運用スタッフは読み取りしかできない共有リポジトリを介して受け渡せばよい。
この\underline{職務の分離}（Separation of Duty）によって、それぞれの当事者の責任が明確になるだけでなく、その責任が強制される。
その結果、承認された経路を通さずに手早くコードを変更し、新しいバージョンを本番に押し込むような、単独の「カウボーイ」的な行為は事実上できなくなる。
そうした行為は、セキュリティ上の欠陥をたやすく持ち込むものである。

\newpage

%% =======================================================
%%  part: ch14
%% =======================================================
%% =======================================================
%%  part: ch14
%% =======================================================
\section{今後の展望}
% src: Looking Ahead

この章では、\cite{Kohnfelder2021} の第14章に沿って、ソフトウェアセキュリティの今後の方向性を扱う。

\begin{quote}
「我々は未来の建築家となるべく召されているのであって、その犠牲者となるためではない。」---R. Buckminster Fuller
\end{quote}

過去 50 年にわたってコンピューティングの進化を見てきた著者は、未来を予測しようとするのは愚かなことだと学んだ。
それでも本書を締めくくるにあたり、実現の見込みが薄いものも含め、価値があると考えるセキュリティの今後の方向性について、著者の考えを述べておきたい。
以下は決して予測ではない。実現すれば大きな進歩になりうる可能性である。

黎明期のインターネットに警鐘を鳴らしたのは、1988年の \underline{Morris ワーム}（Morris worm）である。
このワームは、オンラインのマルウェアがどれほどの力を持ちうるか、そして既存の脆弱性を悪用してどのように広がりうるかを初めて示した。
それから 30 年以上が経ち、多くの面で驚くべき進歩を遂げたが、私たちはこうしたリスクを十分に理解し、緩和の取り組みに十分な優先度を与えてきたのだろうか、と著者は疑問に思う。
攻撃や個人データの漏えいはいまだに日常的に報じられており、終わりは見えない。
攻撃者がやりたい放題をしている一方で、防御側は必死に立ち泳ぎをしているだけのように見えることもある。
しかも、多くのインシデントは秘密にされていたり、気付かれないまま続いていたりする可能性すらあるので、現実は私たちが知っているよりほぼ確実に悪いことを心に留めておく必要がある。
私たちはおおむね、脆弱なソフトウェアと共存することに慣れてしまっている。

注目すべきは、不完全なシステムが侵害され続けているのに、すべてがどうにか回り続けていることである。
セキュリティ問題がなくならないのは、まさにこのためかもしれない。
つまり、今のままで十分なのである。
しかし、投資対効果という冷静な理屈はわかっていても、著者は心の底ではそれを受け入れられない。
業界として、今の状態を自分たちにできる最善だと受け入れてしまえば、本当の進歩を妨げることになると著者は考える。
セキュリティにもっと労力を割くべきだと正当化するのは、いつも難しい。
失敗に終わった攻撃のことも、どの防御線が効いたのかさえも、私たちが知ることはめったにないからである。

この最終章では、ソフトウェアセキュリティの水準を全体として引き上げるための、有望な今後の方向性を素描する。
最初の節では本書の中心テーマを振り返り、本書の手法をどう活用できるかをまとめる。
章の残りでは、さらなる革新と将来のベストプラクティスを構想するが、こちらはより思索的な内容である。
モバイル端末のデータ保護についての議論は、「ラストワンマイル」で本当に効果的なセキュリティを届けるには、まだどれほど多くのことが必要かを示す例になる。
本書の概念的・実践的なアイデアが、この重要で変化し続ける分野への読者の関心を呼び起こし、ソフトウェアを安全にするための読者自身の取り組みの足がかりになることを願っている。

%----------------------------------------------------------------------
\subsection{行動への呼びかけ}
% src: Call to Action

\begin{quote}
「教育の偉大な目的は知識ではなく行動である。」---Herbert Spencer
\end{quote}

本書は、ソフトウェアセキュリティの向上につながると著者が信じる、2 つの単純な考えを広めることを軸に構成されている。
一つは、ソフトウェアを作るすべての人をそのソフトウェアのセキュリティの推進に関わらせること、もう一つは、要件定義と設計の段階からセキュリティの観点と戦略を組み込むことである。
著者は本書の読者に、この動きの先頭に立つ手助けをしてくれるよう切に願う。

加えて、自分たちが作るソフトウェアの品質に継続して力を注ぐことも、セキュリティの向上につながる。
バグが少なければ、悪用できるバグも少なくなるからである。
高品質なソフトウェアを作るには労力がかかる。
有能な設計、注意深いコーディング、網羅的なテスト、漏れのない文書化が必要であり、しかもそれらすべてをソフトウェアの進化に合わせて最新に保たなければならない。
この取り組みが続くよう、開発者もエンドユーザーも、より高い品質と完成度を求め続けなければならない。

\subsubsection{セキュリティは全員の仕事}
% src: Security Is Everyone’s Job

セキュリティ分析は、そのソフトウェアを深く理解している人が行うのが一番である。
本書は良いセキュリティ実践の概念的な土台を示しており、ソフトウェアに携わる人なら誰でも、設計のセキュリティ面を理解したり、セキュアコーディングを学んだりできるようにしている。
セキュリティがほとんど顧みられてこなかったからといって脆弱性の発見と修正を専門家任せにするのではなく、私たちが作るすべてのソフトウェアが少なくとも最低限の基準線を満たすよう、全員で協力しよう。
そうすれば、より難解で技術的なセキュリティの作業は、その技能が最も生きる専門家に任せられる。
その理由は次のとおりである。
\begin{itemize}
\item 専門のコンサルタントがどれほどセキュリティに精通していても、部外者である以上、ソフトウェアとその要件を文脈ごと完全に理解することはできない。そこには、企業の文化や、エンドユーザーのもとでソフトウェアがどう運用されなければならないかも含まれる。
\item セキュリティはソフトウェアのライフサイクル全体に組み込まれているときに最もよく機能するが、セキュリティコンサルタントを長期間雇い続けるのは現実的ではない。
\item 腕のよいソフトウェアセキュリティの専門家は引く手あまたで、見つけるのが難しく、急な依頼には都合がつかない。雇うのにも費用がかさむ。
\end{itemize}

セキュリティの考え方は難しくはないが、抽象的なので、最初はなじみにくく感じるかもしれない。
ほとんどの脆弱性は後から見れば明白なのに、私たちは同じ過ちを何度も繰り返しているようである。
もちろん、こつは問題が現実になる前に見抜くことである。
本書では、まさにそのやり方を身につけるためのさまざまな手法を紹介している。
よい知らせは、これを完璧にこなせる人などいないので、小さな貢献から始めるだけでも何もしないよりはましだということである。
続けていけば、だんだんうまくなる。

セキュリティへの参加を広げるという話は、各人が最も得意な部分を受け持つチームの取り組みとして捉えるのが一番わかりやすい。
一人ひとりが仕事全体を一人でこなせるという考え方ではない。さまざまな技能を持つメンバーの意見を合わせることで相乗効果が生まれ、最良の結果が得られる、という考え方である。
ソフトウェア製品の制作、保守、サポートのなかで自分の担当が何であれ、それを自分の主な貢献として力を注ぐこと。
ただし、関連するコンポーネントのセキュリティを検討したり、チームメイトの仕事を見直して見落としがないか確かめたりすることにも価値がある。
自分の役割が小さくても、サッカーのゴールキーパーがときどきゴールを決めるように、重大な欠陥を見つけることがあるかもしれない。

ここではっきりさせておきたいのは、外部の専門知識には価値があるということである。ギャップ分析やペネトレーションテストなどの作業を行ったり、組織の能力の偏りを補ったり、深い経験を持つ「新鮮な目」として役立ったりする。
しかし専門のコンサルタントは、社内のしっかりしたセキュリティ理解と地に足のついた実践を補う存在であるべきで、セキュリティの重荷を一手に引き受けさせるために呼ぶべきではない。
また、専門家が全体のセキュリティ態勢に貢献してくれるとしても、結局は別の案件へと去っていく。
だから、ソフトウェアに責任を持つチームのできるだけ多くの人が、日頃からセキュリティについて考えているのが常に最善である。

\subsubsection{セキュリティを組み込む}
% src: Baking In Security

橋、道路、建物、工場、船、ダム、港、ロケットは、どれもまず設計され、品質と安全性を確保するために綿密なレビューを受け、そのうえで初めて建設される。
ほかのどの工学分野でも、あとからセキュリティ対策を付け足すより、紙の上で設計を練り上げるほうがよいと認められている。
それなのに、ほとんどのソフトウェアはまず作られ、あとからセキュリティが確保される。

本書の中心となる前提は、早い段階からセキュリティに力を入れれば時間の節約になり、大きな見返りが得られ、できあがるものの品質も上がるということであり、著者はこれが業界で何度も実証されるのを見てきた。
設計がセキュリティを十分に考慮していれば、実装者は安全な解決策をずっと届けやすくなる。
セキュリティを高めるようにコンポーネントを構成すれば、潜在的な問題を予見しやすくなる。
こうしたことを最初から正しく行えば、作業ははるかに楽になり、しかもより安全になる。

セキュリティを設計段階に前倒しする（業界の流行りの言葉でいう「左へ移る」（moving left））ことは最悪のシナリオへの備えであり、その最も説得力のある理由は、\underline{設計上のセキュリティ欠陥}（by-design security flaw）を避けることである。
コンポーネントの分割、API の構造、プロトコル設計など、アーキテクチャのどの側面であれ、設計に組み込まれたセキュリティ欠陥は壊滅的な結果を招きうる。
互換性を壊さずにあとから直すことは、ほぼ不可能だからである。
こうした問題を早期に見つけて直すことが、苦痛で時間のかかる、後手に回った再設計を避ける最善の方法である。

よいセキュリティ設計の判断には、見過ごされがちなもっと大きな利点がある。
よい設計の本質は、必要な機能を損なわない範囲での最小主義である。
これをセキュリティに当てはめると、設計によって攻撃対象領域の広さと重要なコンポーネント間のやり取りを最小限に抑えるということであり、ひいては実装者が誤りを犯す機会が減ることを意味する。

セキュリティに焦点を当てた設計レビューが重要なのは、ソフトウェア設計の機能面のレビューは視点が異なり、セキュリティを考慮しない問いを立てるからである。
「必要な要件をすべて満たしているか。運用と保守は容易か。もっとよい方法はないか。」
実際、安全でない設計が、壊滅的な攻撃に対して脆弱でありながら、こうした審査にすべて見事に合格してしまうこともある。
設計レビューをセキュリティ評価で補えば、設計が直面する脅威を理解し、設計がどう破綻しうるか、どう悪用されうるかを検討することで、設計のセキュリティを吟味できる。

ソフトウェアセキュリティの実装面では、うっかり脆弱性を作り込んでしまう数多くの経路について学び、それを用心深く避けるか、少なくともそうしたよくある落とし穴を緩和することが求められる。
安全な設計にすれば実装で脆弱性が入り込む機会は最小限になるが、それでソフトウェアが魔法のように完全無欠になるわけでは決してない。
開発者は、数多くの罠のどれかにはまってセキュリティを損なうことのないよう、気を抜かずに取り組まなければならない。

セキュリティは、ソフトウェアシステムの構想から、避けられない廃止に至るまで、ライフサイクル全体を貫くプロセスである。
デジタルシステムは複雑で脆く、ソフトウェアが「世界を飲み込む」につれて、私たちはますますソフトウェアに依存するようになっている。
私たちは不完全な人間であり、不完全な人間のために、不完全な部品を使って、十分に良いシステムを作っている。
しかし、完璧が達成できないからといって、進歩できないわけではない。
むしろ、直したバグの一つひとつ、改善した設計の一つひとつ、追加したセキュリティテストケースの一つひとつが、大小さまざまな形でシステムの信頼性を高める助けになるということである。

%----------------------------------------------------------------------
\subsection{将来のセキュリティ}
% src: Future Security

\begin{quote}
「未来は、今日あなたが何をするかにかかっている。」---Mahatma Gandhi
\end{quote}

本書は、著者自身が実践し、一貫して効果を上げるのを見てきたセキュリティ向上の手法を軸に構成されているが、その先にもやるべきことはたくさんある。
以下の各項では、著者が有望だと考えるアイデアをいくつか素描する。
どれもさらに練り上げる必要はあるが、大きな前進につながる可能性があると著者は考えている。

人工知能などの先進技術には大いに期待が持てるが、著者の直感では、必要な作業の多くは「薪を割り、水を運ぶ」ような地道なものである。
私たち全員が貢献できる方法の一つは、自分たちが作るソフトウェアの品質を確保するよう努めることである。
脆弱性はバグから生まれるからである。
第二に、システムの能力と適用範囲が広がるにつれて複雑さは必然的に増すが、それに圧倒されないよう管理しなければならない。
第三に、本書のために調べるなかで、世界のソフトウェアの状態やその安全性について確かなデータが乏しいことに、著者は失望した（が驚きはしなかった）。
透明性が高まれば、前に進むための見通しがもっとはっきりするはずである。
第四に、真正性、信頼、責任は、ソフトウェアコミュニティが安全に協働するための土台であるにもかかわらず、今の仕組みはおおむね場当たり的で当てにならない。
これらの領域で進歩があれば、状況は一変しうる。

\subsubsection{ソフトウェア品質の向上}
% src: Improving Software Quality

「プログラマーはバグを入れることで給料をもらい、バグを取り除くことでも給料をもらう。」
これは、25 年前に著者が Microsoft でプログラムマネージャーをしていたころに聞いた、最も印象に残っている言葉の一つである。
バグは避けられないというこの態度は今も根強く、近いうちに変わる気配はほとんどない。
しかしバグは脆弱性の材料なので、バグの多いソフトウェアが本当はどれだけのコストを生むかを認識しておくことが重要である。

セキュリティを高める方法の一つは、従来のバグトリアージを広げることである。つまり、各バグが攻撃チェーンの一部になりうるかどうかも検討し、その可能性が高く、しかも懸かっているものが大きいバグを優先して直す。
こうしたバグ修正のうち実際の脆弱性を塞ぐのはごく一部だとしても、その取り組みには十分な価値があると著者は主張する。

\subsubsection{複雑さの管理}
% src: Managing Complexity

\begin{quote}
「進化するシステムは、それを減らす作業がなされない限り、その複雑さを増していく。」---Meir Lehman
\end{quote}

ソフトウェアシステムが大きくなるにつれて、それに伴う複雑さの管理は難しくなり、システムは脆くなっていくおそれがある。
最も信頼性の高いシステムは、複雑さを単純なインターフェースのコンポーネントの中に閉じ込め、それらのコンポーネントを障害に強い構成で疎結合にすることで成功している。
大規模な Web サービスは、リクエストを特定の機能を担う多数のマシンに振り分け、それぞれの結果を組み合わせて応答全体を作ることで、高い回復力を実現している。
冗長性を組み込んで設計されているので、障害やタイムアウトが起きても、必要に応じて別のマシンで再試行できる。

大規模な情報システムを構成する多数のコンポーネントについて、それぞれのセキュリティモデルを区画化しておくことは、成功のための基本的な要件である。
組み合わせたコンポーネント間の微妙なやり取りがセキュリティに影響することがあり、相互依存が積み重なるほど、システムを安全にする作業は格段に難しくなる。
複雑なシステムを扱うときは、優れたテストに加えて、よく文書化されたセキュリティ要件と依存関係が、重要な第一の防御線になる。

\subsubsection{透明性の最小化から最大化へ}
% src: From Minimizing to Maximizing Transparency

ソフトウェアセキュリティの現状について最も暗い見方は、おそらく次の格言（誰の言葉かは諸説ある）から導かれる。
「測定できないものは改善できない。」
残念ながら、世界のソフトウェアの品質、とくにセキュリティについての測定はほとんどない。
セキュリティ脆弱性について一般に知られているのは、一部のケースに限られる。
すなわち、オープンソースのソフトウェア、プロプライエタリなソフトウェアの公開リリース（通常はバイナリのリバースエンジニアリングが必要）、あるいは研究者が欠陥を見つけて詳細な分析を公表した場合である。
自社のソフトウェアセキュリティの実績を詳細まですべて公表しようと考える企業は、ほとんどない。
業界として、私たちはセキュリティインシデントからほとんど学べていない。
詳細がすべて開示されることはめったになく、その大きな原因は恐れである。
この恐れには根拠がないわけではないが、もっと情報量の多い開示がほかの人にもたらしうる価値と比べて、釣り合いをとる必要がある。

セキュリティ脆弱性をすべて完全に公開するには障壁があることを受け入れたとしても、改善の余地は大いにある。
主要なオペレーティングシステムのセキュリティアップデートの開示には、たいてい役に立つ詳細が欠けており、そのしわ寄せはユーザーに来る。
ユーザーにとっては、リスクを評価して対応するうえで、追加の情報があれば役立つはずである。
著者の意見では、大手ソフトウェア企業は、提供している情報でさえ、ダブルスピーク（ごまかしの言葉）と呼べるほどぼかしていることが多い。

\begin{exa}[最近の OS セキュリティアップデートの記述]\label{exa:ch14-vague-release-notes}
最近のあるオペレーティングシステムのセキュリティアップデートから、いくつか例を挙げる。
\begin{itemize}
\item 「ロジックの問題が、制限の改善によって対処された。」（これはほぼどんなセキュリティバグにも当てはまる。）
\item 「バッファオーバーフローの問題が、メモリ処理の改善によって対処された。」（バッファオーバーフローをほかにどうやって直せるというのか。）
\item 「検証の問題が、入力のサニタイズの改善によって対処された。」（これも、どんな入力検証の脆弱性にも言えることである。）
\end{itemize}
\end{exa}

このように詳細を伏せることは、あまりにも多くの製品で反射的な習慣になっている。
これは顧客に害を与えるものであり、ソフトウェアセキュリティのコミュニティは、もっと情報量の多い開示から恩恵を受けられるはずである。
ソフトウェアの発行者は、ほとんどの場合、将来のセキュリティを損なわずに追加の情報を提供できる。
現実には、敵対者はアップデートの変更点を解析して基本的な詳細を拾い集めるので、役に立たないリリースノートは、誠実な顧客から重要な詳細を奪うだけである。
これからの責任あるソフトウェア提供者は、完全な開示を出発点とし、セキュリティを弱めないよう必要に応じて一部を伏せる（墨塗りする）ほうがよい。
さらに望ましいのは、悪用の危険が去ったあとで、保留していた詳細を追加で開示することである。
たとえ後知恵であっても、主要な商用ソフトウェア製品のセキュリティを理解するうえで価値があるはずである。

脆弱性を詳しく報告するのは、後から見れば問題がたいていあからさまに明白なので、恥ずかしいかもしれない。
しかし著者は、こうした過失に正直に向き合うことは健全で生産的だと主張する。
完全な開示から学べることは非常に大きいので、長期的にセキュリティに真剣に取り組むのであれば、もっと透明性が必要である。
著者が顧客なら、セキュリティ修正のリリースノートに次の内容を載せているソフトウェアベンダーにこそ、ずっと感心するだろう。
\begin{itemize}
\item バグが報告され、トリアージされ、修正され、テストされ、リリースされた日付と、不当な遅れがあった場合はその説明。
\item その脆弱性がいつ、どのようにして作り込まれたかの説明（たとえば、不注意な編集、セキュリティ上の影響についての無知、意思疎通の行き違い、あるいは悪意ある攻撃）。
\item 欠陥のあるコードを含むコミットがレビューされたかどうか。レビューされたなら、なぜ見逃されたのか。されなかったなら、なぜされなかったのか。
\item 同じ種類の似た欠陥を探す取り組みをしたかどうか。したなら、何が見つかったか。
\item 今後のリグレッションや似た欠陥を防ぐためにとった予防策の詳細。
\end{itemize}

脆弱性、その原因、その緩和策について、もっと率直に開示し共有する文化へと業界を変えていけば、私たち全員がこうしたインシデントから学べるようになる。
詳細や文脈がほとんどない開示は形だけのものであり、誰の役にも立たない。

\begin{exa}[国家運輸安全委員会]\label{exa:ch14-ntsb}
ベストプラクティスの優れた例が、米国の国家運輸安全委員会（National Transportation Safety Board, NTSB）である。
NTSB は詳細な報告書を公表しており、航空業界もパイロットもそれを読んで事故から学べる。
\end{exa}

さまざまな理由から、ソフトウェアの世界がこの仕組みをそのまま真似ることはできないが、目指すべき手本にはなる。
理想をいえば、主要なソフトウェアメーカーは、公開の開示を、舞台裏で何が起きたかを正確に説明し、対応における自社の能力とプロ意識を示す機会と捉えるべきである。
そうすれば、ほかの製品で似た問題を防ぐための幅広い学びにつながるだけでなく、自社製品への信頼を取り戻す助けにもなるだろう。

\subsubsection{ソフトウェアの真正性・信頼・責任の向上}
% src: Improving Software Authenticity, Trust, and Responsibility

現代の大規模なソフトウェアシステムは多くのコンポーネントから作られている。そのすべてが真正でなければならず、しかもそれぞれが、信頼できる作り手によって、安全なサブコンポーネントから、安全なツール群を使って作られていなければならない。
この連鎖はどこまでも続き、文字どおり現代のデジタルコンピューティングの夜明けにまで遡る。
私たちのシステムのセキュリティは、現代のソフトウェアスタックを築き上げてきた、こうした積み重ねすべてのセキュリティに依存している。
しかし、その正確な系譜は今ではコンピューティングの歴史の霧の中に消えており、すべての始まりとなった初期のいくつかのセルフコンパイルするコンパイラにまで遡る。

\begin{exa}[Reflections on Trusting Trust]\label{exa:ch14-trusting-trust}
Ken Thompson の古典的な論文 ``Reflections on Trusting Trust'' は、セキュリティがこの歴史全体にどう依存しているか、そしていったん深く埋め込まれたマルウェアを見つけるのがいかに難しいかを、見事に示している。
そこに何かよからぬものが潜んでいないと、私たちは本当にどうやって知ることができるのだろうか。
\end{exa}

ソフトウェアがどう作られたかの完全性を確保するのに必要なツールは、今では無料で手に入り、宣伝どおりに機能すると考えてよい。
しかし、その使われ方は残念なほど場当たり的で手作業に頼りがちであり、そのためにプロセスは、妨害工作とまではいかなくても、人為的ミスの影響を受けやすい。
単に時間を節約するために確認を省いてしまうこともあり、それも無理はない。

\begin{exa}[*nix ディストリビューションの検証]\label{exa:ch14-verify-distro}
たとえば、*nix ディストリビューションが正当なものかどうかを検証する場面を考える。
信頼できる Web サイトからイメージをダウンロードしたら、それとは別に権威ある鍵とチェックサムのファイルもダウンロードし、（信頼できる情報源から入手した）いくつかのコマンドを使ってそれらをすべて検証する。
これらの確認がすべて通って初めて、インストールに進むべきである。
しかし実際のところ、管理者はこうした追加の手順をどこまで徹底しているだろうか。
主要なディストリビューションでこうした確認が失敗した例など聞いたこともないとなれば、なおさらである。
しかも、たとえ常に実行されているとしても、それを裏付ける記録は私たちの手元にない。
\end{exa}

今日、ソフトウェアの発行者はリリースするコードに署名しているが、その署名が保証するのは、ビットが改ざんされていないという完全性だけである。
署名されたコードは信頼できると言っているように受け取られがちだが、あとで脆弱性が見つかっても署名が無効になるわけでは決してないので、そう解釈するのはまったく安全ではない。

将来、真正性の連鎖を監査できる記録などを備えたよりよいツールがあれば、もっと高い完全性の保証が得られ、システムのセキュリティが依存している信頼の判断や依存関係の手がかりにできるだろう。
たとえば新しいコンピューターには、オペレーティングシステム、ドライバー、アプリケーションなどが真正であることを記録した\underline{ソフトウェアマニフェスト}（software manifest）を付けるべきである。
コンポーネントとビルド環境の\underline{ソフトウェア部品表}（software bill of materials）を文書化して認証するには大変な労力がかかるが、その難しさにひるまず、完全な解決策の一部から始めて、時間をかけて少しずつ改善していくべきである。
ソフトウェアの来歴と真正性に真剣に取り組み始めれば、重要なソフトウェアのリリースが安全なコンポーネントから作られていることを、はるかにうまく保証できるようになり、未来は私たちに感謝するだろう。

%----------------------------------------------------------------------
\subsection{ラストワンマイルを届ける}
% src: Delivering the Last Mile

\begin{quote}
「最も長い一マイルは、家へ帰る最後の一マイルである。」---作者不詳
\end{quote}

あらゆるベストプラクティスに忠実に従い、本書で述べた手法を当てはめ、フットガンを避けるよう注意してコードを書き、レビューを行い、徹底的にテストし、システム全体を漏れなく文書化すれば、仕事は完璧に安全になる、と言えればよいのだが、もちろんそれほど単純ではない。
セキュリティの作業は決して終わらないうえに、よく設計され、よく作り込まれたシステムでさえ、現実の世界で実際に使われるなかで、意図したセキュリティ水準を届けられないことがある。

\begin{dfn}[ラストワンマイル]\label{dfn:ch14-last-mile}
\underline{ラストワンマイル}（last mile）は通信業界や運輸業界から借りた用語で、個々の顧客をネットワークにつなぐという課題を指す。
多くの場合、これがサービスを提供するうえで最も費用がかかり、最も難しい部分である。
\end{dfn}

\begin{exa}[インターネットサービスプロバイダのラストワンマイル]\label{exa:ch14-isp-last-mile}
たとえば、インターネットサービスプロバイダは、すでにあなたの近所まで高速な光ファイバーのインフラを引いているかもしれない。
それでも、新しい顧客を一人獲得するたびに、作業員の派遣や、場合によってはケーブルの敷設、モデムの設置が必要になる。
どれもスケールしにくく、その時間と費用は大きな追加の先行投資になる。
\end{exa}

それとまったく同じように、よく設計された安全なシステムを配備しても、それは実際のセキュリティを届けることの始まりにすぎない場合が多い。

セキュリティにおけるこうした「ラストワンマイル」の課題を理解するために、「スマートフォンをなくしたら、他人に中身を読まれてしまうだろうか」という単純な問いを手がかりに、モバイル端末のデータセキュリティの最先端を詳しく見てみよう。
長年にわたる集中的なエンジニアリングの努力の結果、よく作られた強力な最新の暗号技術一式が生まれた。それでも、今日のハイエンドのスマートフォンでさえ、答えは「はい、おそらくデータの大部分は読まれてしまう」となるようである。
これはおそらく近年で最大のソフトウェアセキュリティの取り組みなので、どこが、なぜ不十分なのかを理解することが重要である。

以下の議論は、Johns Hopkins 大学の 3 人のセキュリティ研究者が書いた 2021 年の論文 ``Data Security on Mobile Devices: Current State of the Art, Open Problems, and Proposed Solutions'' に基づいている。
この報告書は、堅牢なソフトウェアセキュリティを届けることがいまだに難しい場合が多い理由を、いくつかの重要な観点から述べている。
この例から得られる、セキュリティ全般にわたる大きな教訓を浮き彫りにするために、議論は大幅に単純化する。

まず、データ保護のレベルについて述べる。
モバイルアプリはあらゆる種類の便利なことをする。
その種類は一つの暗号化方式ですべてをまかなうには多すぎるので、モバイルオペレーティングシステムはさまざまな選択肢を用意している。
iOS プラットフォームには 3 つのデータ保護レベルがあり、主な違いは、保護されたデータにアクセスできるよう暗号鍵がメモリ上に置かれている時間を、どこまで積極的に短くするかという点にある。

\begin{exa}[銀行の金庫の扉]\label{exa:ch14-bank-vault-door}
これは、銀行の金庫の扉をどれくらいの頻度で開けたままにしておくかにたとえられる。
大きく重い扉を朝に開け、閉店時刻にだけ閉めるのであれば、職員は一日中手軽に出入りできるが、使っていないときにも金庫が侵入にさらされやすくなる。
これに対して、職員が金庫に入る必要があるたびに支店長を探して開けてもらわなければならないなら、利便性と引き換えにセキュリティが高まる。
金庫はほとんどの時間しっかり施錠されているからである。
モバイル端末で、保護されたデータにアクセスするために（パスワード、指紋、顔認識で）暗号鍵のロック解除をユーザーに求めることは、おおよそ支店長に金庫を開けてもらうことにあたる。
\end{exa}

最高レベルの保護では、暗号鍵はスマートフォンのロックが解除されて使われている間しか利用できない。
これは非常に安全だが、ほとんどのアプリにとっては足かせになる。
端末がロックされるとデータにアクセスできなくなるからである。
たとえば、会議の時間になると知らせてくれるカレンダーアプリを考えてみよう。
スマートフォンがロックされていると、アプリはカレンダーのデータにアクセスできない。
同期を含むバックグラウンド処理も、ロックされている間は止められる。
つまり、スマートフォンがロックされている間にカレンダーに予定が追加されると、たまたまその前にロックを解除して同期させていない限り、通知を受け取れないことになる。
最も制限の緩い保護クラスである \underline{After First Unlock}（AFU）でさえ、起動後に暗号鍵を組み立て直すにはユーザーのクレデンシャルが必要であり、大きな制約がある。
名前のとおり、再起動した直後の端末では暗号鍵が使えないので、その間もカレンダーの通知は止められてしまう。

こうした制約をかわすようにアプリを設計することは考えられる。
データを、いつ必要になるかに応じて、保護クラスの異なる別々の保存場所に分けるのである。
たとえばカレンダーなら、時刻は使えるよう保護せずにおき、通知では「午後4時に会議があります」とだけ伝え、詳細を見るには端末のロック解除を求める、というやり方がありうる。
件名のない通知はわずらわしいだろうが、ユーザーはプライバシーのためにカレンダーが暗号化されていることも期待しているので、どこかで折り合いをつける必要がある。
この情報がどれほど機微かはユーザーによって違い、会議の中身にもよるが、毎回ユーザーにはっきり判断させるのも現実的ではない。
アプリは勝手にうまく動いてくれるものだと、人々は思っているからである。
結局、ほとんどのアプリは、自分が管理するデータにアクセスしやすくするほうを選び、低いレベルのデータ保護を使うことになる。
あるいは、多くの場合、保護をまったく使わない。

ほとんどのアプリが利便性のために「保護なし」の選択肢で動いていると、攻撃者が端末を調べられる場合、そうしたデータはすべて持ち出し放題の格好の標的になる。
簡単ではないが、Johns Hopkins の報告書が詳しく述べているように、高度な手法を使えばメモリに入り込む方法が見つかることが多い。
AFU 保護の場合、攻撃者は暗号鍵を見つけさえすればよい。端末はほとんどの時間をこの状態で過ごすので、鍵はメモリ上にあることが多い。

秘密のメッセージングアプリはこの法則の主な例外で、「完全な保護」クラスを使っている。
目的が特殊なので、ユーザーも、端末がロックされているときに機能が使えないことや、使うのに余計な手間がかかることを我慢する覚悟ができている。
こうしたアプリは少数派であり、端末に保存されたユーザーデータのごく一部を占めるにすぎない。
それなのに、ほとんどのスマートフォンユーザー（セキュリティについて考えることがあればの話だが）は、おそらく自分のデータはすべて安全だと思っている。

状況がまだ暗さが足りないとでもいうように、多くのアプリにとってクラウドとの連携がどれほど重要で、それが強力なデータ保護とどれほど相容れないかを考えてみよう。
クラウドコンピューティングのモデルは現代のコンピューティングに革命をもたらした。私たちは今や、Web 検索、リアルタイム翻訳、画像や音声の保存をはじめ数え切れないサービスをすぐに使え、どこからでも、つながったデータセンターを指先で使えることに慣れている。
顔認識を使って写真コレクションから人物を探すような機能は、現代の端末のかなりの計算能力でさえはるかに及ばないので、クラウドに大きく頼っている。
クラウドのデータモデルなら、複数の端末からのアクセスも簡単になる（もう同期は要らない）。
また、端末をなくしても、データはクラウドに安全に保存されているので、新しいハードウェアを買うだけで済む。
しかし、クラウドの力を生かすには、端末上で暗号化によってデータを固く守るのではなく、データをクラウドに預けなければならない。

もちろん、こうした途切れのないデータアクセスは、どれも強力なデータ保護とは相容れない。
クラウドにつながったスマートフォンをなくした場合は、とくにそうである。
ほとんどのモバイル端末はクラウドのデータにいつでもアクセスできるので、端末を拾った人なら誰でも、保存されたデータにアクセスできるおそれがある。
しかも、そのデータはおそらく暗号化されていない。
たとえば、クラウドにエンドツーエンドで暗号化したデータを保存する写真アプリを考えてみても、それでは中身の見えないビットの塊しか保存できず、検索や写真共有を提供するクラウドの力は失われる。
また、復号鍵は端末上だけに厳格に保持しなければならないので、複数の端末からアクセスする使い方は難しくなる。
さらに、端末上の鍵に何かあれば、クラウド上のデータがすべて役に立たなくなるおそれもある。
こうした理由から、クラウドに頼るアプリは、暗号化によるデータ保護をほぼ完全にあきらめている。

ここではモバイル端末のデータ保護がどこまで有効かという技術的詳細の、ほんの表面をなぞったにすぎない。それでも私たちの目的には、より一般的な問題の輪郭が見えてきたはずである。
モバイル端末は豊かで複雑なエコシステムのなかにあり、データ保護がすべてのコンポーネントとシナリオで機能しない限り、たちまち使いものにならなくなる。
最善の助言は今も変わらない。なくして漏れても大して困らないこと以外には、スマートフォンを使わないことである。

この話から著者が強調したい教訓は、モバイル端末の暗号化の設計にとどまらず、セキュリティを届けようとするあらゆる大規模システムにおおむね当てはまる。
要するに、データ保護のための機能を豊富に備えた入念な設計であっても、ラストワンマイルでセキュリティを届けきれずに終わってしまうことは、あまりにもたやすいのである。
強力なセキュリティモデルも、開発者がそれを使い、ユーザーがその利点を理解して初めて効果を発揮する。
効果的なセキュリティを実現するには、アプリの邪魔をするのではなくアプリとうまくかみ合う、役に立つバランスのとれた機能を提供する必要がある。
保護が必要なデータはすべて保護されなければならず、インフラ（この例ではクラウド）とのやり取りやインフラへの依存が、保護の効果を損なってはならない。
最後に、こうしたことすべてを典型的な作業の流れに組み込み、エンドユーザーがセキュリティの仕組みと格闘するのではなく、むしろ仕組みに協力するようにしなければならない。

\begin{exa}[.NET Framework のコードアクセスセキュリティ]\label{exa:ch14-dotnet-cas}
何年も前、著者は .NET Framework のリリースで、ラストワンマイルで届けきれなかった事例を目の当たりにした。
セキュリティチームは、この新しいプログラミングプラットフォームに\underline{コードアクセスセキュリティ}（Code Access Security, CAS、第3章で説明した）を組み込むために懸命に働いたが、その利用を十分に広めることができなかった。
CAS では、マネージドコードが特権操作を行うための権限を付与され、必要なときにその権限を宣言（assert）しなければならなかったことを思い出してほしい。
これは最小権限（Least Privilege）のパターンにとって理想的なツールである。
残念ながら、ランタイムチーム以外の開発者はこれを負担と受け止め、この機能のセキュリティ上の利点を理解しなかった。
その結果、アプリケーションは、システムが用意したきめ細かな権限を必要な箇所でだけ使うのではなく、たいていプログラムの開始時に一度だけ全権限を宣言し、あとはまったく制限なしに動くようになった。
機能としてはそれで動いたが、アプリケーションが過剰な権限のもとで動くことになった。
いわば、銀行の金庫の扉が常に開けっ放しの状態である。
その結果、どんな脆弱性も、CAS が意図どおりに使われていた場合よりはるかに大きなリスクをもたらすことになった。
\end{exa}

こうした考察は、どのシステムも直面する課題を代表するものであり、セキュリティの作業が本当の意味で終わらない大きな理由でもある。
優れた解決策を作ったら、それが開発者にもユーザーにも理解され、実際に使われ、しかも正しく使われるようにする必要がある。
ソフトウェアは作り手がまったく予想しなかった新しい使われ方をすることがよくある。そうしたケースを知ったら、そのセキュリティ上の影響を検討し、必要なら対応することが重要である。
こうした要素すべてに加え、さらに多くの要素が、本当に機能する安全なシステムを作るのに欠かせない。

%----------------------------------------------------------------------
\subsection{結論}
% src: Conclusion

ソフトウェアには、完全にビットだけでできているという、独特で恵まれた性質がある。
ただの 0 と 1 の集まりにすぎないので、文字どおり無から生み出せる。
材料は無料で無制限に手に入るので、制約になるのは想像力と創造力だけである。
これは善をなす側にも、害をなそうとする側にも同じように当てはまるので、可能性も、手ごわい課題も、どちらも限りがない。

この章では、行動への呼びかけと、将来を見据えたいくつかのアイデアを示した。
ソフトウェアを開発するときは、プロセスの早い段階でセキュリティ上の影響を検討し、より多くの人にセキュリティについて考えてもらい、この話題についてより多様な視点を取り入れること。
セキュリティへの意識が高まれば、ソフトウェアのライフサイクル全体を通じて、健全な懐疑心と警戒心が生まれる。
手作業による確認に頼る部分を減らし、自動化された検証を増やすこと。
システムを実現するまでの重要な決定と行動をすべて監査できる記録に残し、システムのセキュリティ上の性質がはっきり定義されるようにすること。
コンポーネントは賢く選び、同時に前提やシステムの重要な性質もテストすること。
脆さを減らし、複雑さと変化を管理すること。
脆弱性が見つかったら、その根本原因を調べ、そこから学び、今後のリスクを先回りして減らすこと。
現実的なシナリオを批判的に検討し、ラストワンマイルまでセキュリティを届けるよう努めること。
自分が出会った問題とその対応からほかの人が学べるよう、責任の許す範囲で詳細を公表すること。
セキュリティを高め、プライバシーを尊重するために、小さな歩みを絶えず積み重ねること。

ソフトウェアセキュリティの山や谷をめぐるこの旅に付き合ってくれたことに感謝する。
すべてを隅々まで歩き尽くしたわけではないが、読者はもう土地の全体像をつかめているはずである。
本書のなかに役立つアイデアを見つけ、このテーマへの理解を深めたうえで、それを実践し始めてくれることを願っている。
本書は答えそのものではないが、ソフトウェアセキュリティの水準を引き上げるための答えをいくつか示している。
何より大切なのは、ときどき「セキュリティの帽子」をかぶり、今日から自分の仕事にこれらの概念と手法を当てはめることである。

\newpage

%% =======================================================
%%  part: appendix
%% =======================================================
\appendix

%%  appendix A
%% =======================================================
%%  part: appendix A
%% =======================================================
\section{サンプル設計書}
% src: Sample Design Document

この付録では、\cite{Kohnfelder2021} の付録Aに沿って、セキュリティを意識した設計書の例を扱う。

以下の文書は、設計書のセキュリティレビューの過程を説明するために用意した架空の設計である。
学習用の教材として作られたもので、実際の設計書なら含まれるはずの詳細の多くを省き、セキュリティの側面に焦点を絞っている。
したがって、実際のソフトウェア設計書の完全な例ではない。

原典では、イタリック体の文はこの設計書についてのメタな説明として使われており、文書の教育上のねらいについての所見や、原著者が省略した部分の説明にあてられている。
また太字は、セキュリティに関連する内容を強調するために使われている。
具体的には、設計における優れたセキュリティ実践の例、優れた設計者なら加える機能、あるいはセキュリティのレビュー担当者が指摘すべき点である。
% 注: スクレイプした Markdown では太字・イタリックの強調がほぼ失われているため、
%   どの文が強調されていたかはこのノートでは再現していない。イタリックの名残が明らかな
%   メタ説明（Section 1 末尾）だけを「注意」環境にした。

%----------------------------------------------------------------------
\subsection{表題：私的データロギングコンポーネント設計書}
% src: Title – Private Data Logging Component Design Document

設計書の表題は「私的データロギングコンポーネント設計書」（Private Data Logging Component Design Document）である。

% src: Table of Contents
% 注: 原典の設計書の目次（Section 1〜10 と 2.1〜2.7、6.1〜6.4 の一覧、END OF DOCUMENT）は、
%   以下の \subsection / \subsubsection の見出しそのものと同じなので省略した。

%----------------------------------------------------------------------
\subsection{セクション1：製品の説明}
% src: Section 1 – Product Description

この文書では、監査、システム監視、デバッグを支援するために、標準的なソフトウェアイベントのロギング機能を提供するロギングコンポーネント（以下 \underline{Logger} と呼ぶ）について述べる。
Logger は、意図しない情報漏えいのリスクを緩和するように設計されている。
Logger はログ内の私的データを明示的に区別して扱い、私的でないデータは日常的な用途のために自由にアクセスできるようにする。
このアクセスレベルでは足りないまれな場合には、保護された私的なログデータへの限定的なアクセスを、明示的な承認を条件に、露出をできるだけ抑える制限を付けて提供できる。
% 注: 原文の "eminimize" は "minimize" の誤記と解した。

\begin{rem}[セキュリティ中心の設計思考]\label{rem:appa-security-centric}
ロギングシステムの中で私的データを明示的に分けて扱うという考え方は、セキュリティを中心に据えた設計思考の一例である。
この機能を既存のシステムに後から追加すると、最初から設計に組み込む場合に比べて効率が悪く、コードの大幅な手直し（code churn）が必要になるだろう。
\end{rem}

%----------------------------------------------------------------------
\subsection{セクション2：概要}
% src: Section 2 – Overview

プロジェクトの設計の基本的な前提については、セクション10に挙げた文書を参照のこと。

\subsubsection{2.1 目的}
% src: 2.1 Purpose

データセンター内のアプリケーションはどれも重要なソフトウェアイベントの詳細をログに記録する必要があり、ログには私的データが含まれうるので、アクセスを注意深く制御しなければならない。
Logger は、ログの生成、ログの保存、権限のあるスタッフへの適切なアクセスの強制を担う標準的なコンポーネントを提供し、あわせて、実際に行われたアクセスについて信頼でき否認できない記録を残す。
ロギング、アクセス、保持の要件はシステムごとに異なるので、Logger はアクセスポリシーを指定する簡単なポリシー設定に基づいて動作する。

\subsubsection{2.2 範囲}
% src: 2.2 Scope

この文書では、実装言語、デプロイ、運用上の考慮事項の選択は指定せずに、Logger のソフトウェアコンポーネントの設計を説明する。

\subsubsection{2.3 概念}
% src: 2.3 Concepts

\begin{dfn}[ログのフィルタ済みビュー]\label{dfn:appa-filtered-view}
私的な詳細を取り除いたうえでログを比較的自由に閲覧できるようにしたものを、ログの\underline{フィルタ済みビュー}（filtered view）という。
フィルタ済みビューでは、私的データは隠され、つまり「ラップ」（wrap）され、機微でないデータは表示される。
これに対し、私的データを含むすべてのデータを表示するものを\underline{非フィルタビュー}（unfiltered view）という。
\end{dfn}

ログのフィルタ済みビューという考え方が、この設計の中核である。
狙いは、私的な詳細を取り除いたうえでログを比較的自由に閲覧できるようにすることであり、ほとんどの用途にはこのアクセスレベルで足りるはずである。
さらに必要なときには、追加の認可を条件に、ログに記録された機微なデータを閲覧できる。
そのアクセス自体もイベントとしてログに記録されるので、閲覧が行われたという事実を監査できる。
このように段階的にアクセスできるようにすることで、アプリケーションは重要な私的詳細をログに記録しながら、社内スタッフの正当な利用であってもそのデータが露出する範囲を最小限に抑えられる。
ログに決して現れてはならないほど機微なデータは、そもそもログに記録すべきではない。

例えば Web アプリケーションは、システム利用の記録として、またほかにもさまざまな理由から、HTTPS リクエストを日常的にログに記録する。
こうしたログには、記録しておく必要はあるがめったに使われない私的な情報（IP アドレス、Cookie など多数）が含まれていることが多い。
例えば IP アドレスは、悪意ある攻撃を調査するときには（攻撃の発生源を突き止めるのに）役立つが、それ以外の用途には関係がない。
フィルタ済みビューで割り当てる仮名を使えば、例えば「IP7」と付いたイベントはすべて同じ IP アドレスから来ていることを、実際のアドレスを明かさずに示せる。
このようなフィルタ済みビューは、監視、統計の収集、デバッグといった目的には十分な情報を提供することが多く、その場合、私的データを一切露出させずに済むという利点がある。
ログ自体には完全なデータが残っているので、保護された情報が必要になるまれな場合には、適切な認可のもとで、制御された形で非フィルタビューを使える。

\begin{exa}[ログイン試行のログ]\label{exa:appa-login-log}
Web アプリケーションが受け取ったユーザーのログイン試行が、プロセスをクラッシュさせるバグを引き起こしたとする。
ログの内容を簡略化すると、例えば次のようになる。
% 注: Markdown では password の値が ">1<}2{]3[\\4/" と書かれているが、"\\" は Markdown の
%   エスケープで、実際は "\" 1 文字と解した（後で password の長さが 12 とされることと一致する）。
{\footnotesize
\begin{verbatim}
2018/06/23 08:09:10 66.77.88.99 POST login.htm {user: "SAM", password: ">1<}2{]3[\4/"}
\end{verbatim}
}
このログに含まれる項目は、タイムスタンプ（機微でない）、IP アドレス（機微）、HTTP の動詞と URL（機微でない）、ユーザー名（機微）、パスワード（非常に機微）である。
バグを再現するには、調査でこれらの情報をすべて考慮しなければならないかもしれないが、このデータは本当に必要な場合を除いて平文で表示したくないし、表示するとしても権限のある人だけにしたい。
\end{exa}

さまざまなシステムのセキュリティ上の要求に応えられるよう、ログデータの種類ごとの機微度は設定可能にし、ロギングシステムは機密データを必要に応じて選択的にだけ明かすべきである。
例えば、ベストプラクティスでは URL に機微な情報を含めるべきではないが、この経験則に反していることがわかっている古いシステムがあって、ふつうなら不要な保護が必要になる場合がある。
その場合、デバッグの種類によってはフィルタ済みビューはあまり役に立たなくなる。
URL については、正規表現を使えば、特定の URL をほかより機微なものとして設定しやすくなる。

前の例のログで、機微なデータを省略またはラップしたフィルタ済みビューは、次のようになるだろう。
% 注: 原典ではここにフィルタ済みビューのログ行が示されているはずだが、スクレイプした Markdown では
%   欠落している。同じ内容の行がセクション5（例 \ref{exa:appa-filtered-entry}）にあり、
%   以下の説明（US1, USER1, PW1, (v4), (3), (12)）はその行と対応する。ここでは行を補わない。

IP アドレス、ユーザー名、パスワードはどれも、データを隠すために識別子としてラップされている。ただし、置き換えた識別子を使えば、文脈の中で値が一致するほかのリクエストを問い合わせることができる。
この例では、US1 は米国内の IP アドレスを表し、USER1 はイベントに関連するユーザー名を具体的に明かさずに表し、PW1 は送信されたパスワードを表す。
括弧内の接尾辞は実際のデータの形式や長さを示しており、具体的な詳細を明かさずに手がかりを添えている。
この例からは、IPv4 アドレスであること、ユーザー名が 3 文字であること、パスワードが 12 文字であることがわかる。
例えば、長すぎるパスワードが問題を引き起こしたのなら、想定外の長さだけからそのことがわかるだろう。
パスワードの長さがわかるとわずかに情報が漏れるが、実際上、それで危険にさらされることはないはずである。

目の前の作業にフィルタ済みビューでは足りない場合は、US1 のような識別子のラップを解くよう追加でリクエストできる。
こうすることで、機微なデータを見るには明示的にそれを選ぶ必要があり、データを段階的に明かせるようになる。
例えば IP アドレスだけが必要なら、ユーザー名とパスワードの値は明かされないままになる。

\subsubsection{2.4 要件}
% src: 2.4 Requirements

ログは確実に保存され、認可があればすぐにアクセスでき、必要な保持期間を過ぎたら破棄される。
大量の利用に耐えられるよう、ログ取得インターフェースは高速でなければならない。また、インターフェースが成功を返したら、ログを生成したアプリケーションはログが保存されたものとみなしてよい。

ログは私的な詳細を知らなくても監視できるので、フィルタ済みのログビューはほとんどの用途に広く提供できる。完全なデータ（私的データを含む）を見るには、本当に必要な場合に限って特別な認可を必要とする。

ログシステムを使うのは技術スタッフだけなので、利用者は全員、複雑なインターフェースや、ソフトウェアがログを生成する仕組みを理解できるものとする。

この設計の重要な目標は、非常に機微な私的データもログに記録できるようにし、セキュリティインシデントの疑いの調査や、まれに本番環境でしか起きない問題のデバッグに使えるようにすることである。
内部者による攻撃を完全に緩和するのは現実的な目標ではないが、合理的な予防策をすべて講じ、抑止力として強力な監査証跡を残しておくことが重要である。

物理媒体が盗まれた場合の漏えいを防ぐために、ログのストレージは暗号化する。

ログを生成するソフトウェアは完全に信頼されているものとする。
Logger が私的データを正しく扱うには、ログを生成するソフトウェアが私的データを正しく識別しなければならない。

\subsubsection{2.5 目標としないこと}
% src: 2.5 Non-Goals

Logger は管理者が使うことを想定しているので、洗練された UI は不要である。

コードの改ざんや、管理者による root 権限の悪用といった内部者の攻撃は、範囲外とする。

Logger が効果を発揮するには、注意深い設定と監督が欠かせない。
それをどう実施するかはシステム管理部門が定めなければならないが、レビューの過程と、相互牽制を伴う監査を含めるべきである。

\subsubsection{2.6 未解決の問題}
% src: 2.6 Outstanding Issues

ログへのアクセスの設定、ユーザー認証、非フィルタアクセスの認可の付与については、詳細がまだ決まっていない。

暗号化された私的データへの問い合わせは、どうしても遅くなる。
この設計では、ログデータの量は十分に少なく、レコードを必要に応じて復号しながら総当たりで走査する（つまり索引に頼らない）方法でも十分な性能が出るものと想定している。
将来のもっと意欲的なバージョンでは、暗号化されたデータに索引を付けて高速に問い合わせられるようにすることに取り組むかもしれない。

エラーになる場合を洗い出し、その処理を規定する必要がある。

Logger の将来のバージョンで検討すべき機能強化には、次のものがある。
\begin{itemize}
\item より詳しい、またはより粗い情報を提供するフィルタ済みビューのレベルを定義する。
\item いずれ定期的に削除されるログの一部を取り出し、長期間安全に保存できるようにする。
\end{itemize}

\subsubsection{2.7 代替設計}
% src: 2.7 Alternative Designs

最終的に選んだ設計は、ログ内の機微な情報をすべて保存することについて Logger を全面的に信頼するもので、いわば「すべての卵を一つのかごに盛る」設計である。
機微な情報を発生源ごとに区画化できる代替案も検討した。
この代替案は、重要な利用シナリオと両立しないと考えられるいくつかの理由（以下に概略を示す）から採用しなかったが、ロギングの解決策としてはこちらのほうが間違いなく安全であろうことは、書き留めておく必要がある。

% src: Alternative design
\underline{代替設計：}
ログの発生源は非対称暗号の鍵ペアを作り、それを使って、ログレコードのうち機微なデータの部分を Logger に送る前に暗号化する。
注意深く行えば、Logger は（おそらく）それでもフィルタ済みビュー用の仮名（例えば米国内のある IP アドレスに対する US1）を生成できる。
そのうえで、非フィルタビューへの権限あるアクセスには、データを復号するための秘密鍵が必要になる。
この方式の主な利点は、保存されたログデータが漏えいしても暗号化された機微なデータは漏れず、そもそも Logger は必要な鍵すら持たないことである。

% src: Reasons not chosen
\underline{採用しなかった理由：}
この設計では、暗号化と鍵管理の負担が、ログの発生源と権限のあるアクセス者の両方にかかる。
どのデータが機微で、それをどう区分するかは、ログの発生源が決めることになり、その時点で固定されてしまう。
信頼を Logger に集中させれば、この両方を必要に応じて設定し直せるし、ログの閲覧者を認証することでアクセスを細かく制御できる。

%----------------------------------------------------------------------
\subsection{セクション3：ユースケース}
% src: Section 3 – Use Cases

データセンター内のアプリケーションは、Logger を使って重要なソフトウェアイベントのログを生成する。
日常的な監視ソフトウェアと、しかるべき運用スタッフには、日常業務のためにフィルタ済みアクセス（私的データを一切明かさないデータビュー）を許可する。
トラフィック量、アクティブユーザー数、エラー率などの運用統計は、すべてフィルタ済みのログビューから作る。

サポートやデバッグのために非フィルタのログへのアクセスがまれに必要になった場合、権限のあるスタッフは、ポリシーに従って限定的なアクセスを得られる。
アクセスを要求するときは、必要なログの部分集合、その時間枠、アクセスの理由を指定する。
承認されるとアクセスを許可するトークンが発行され、そのアクセスは監査のためにログに記録される。
作業が終わったら、要求者が調査の結果を記したメモを追加し、承認者がそれをレビューして適正だったことを確かめる。

要求、承認、監査レビューの要約に加え、ログ量の傾向と、期限切れのログデータが削除されたことの確認を詳しく記した報告書を作り、経営陣に報告する。

%----------------------------------------------------------------------
\subsection{セクション4：システムアーキテクチャ}
% src: Section 4 – System Architecture

データセンター内で、Logger のサービスインスタンスは、サービスを提供する相手のアプリケーションとは独立して運用される物理的に別のマシン上で動き、標準的なパブリッシュ／サブスクライブのプロトコルを介してアプリケーションにサービスを提供する。
Logger は、次の機能を担う 3 つの新しいサービスで構成される。

\subsubsection{Logger Recorder}
% src: Logger Recorder

ログを保存するサービスである。
アプリケーションはログイベントのデータを、暗号化された通信路を通じて Logger Recorder サービスにストリーミングし、データはそこで永続ストレージに書き込まれる。
一つのインスタンスで複数のアプリケーションのログを扱うように設定してもよい。

\subsubsection{Logger Viewer}
% src: Logger Viewer

技術スタッフがフィルタ済みのログを手作業で閲覧するための Web アプリケーションであり、ポリシーに従った認可があれば、非フィルタビューも表示できる。
Logger Viewer は、アクセスポリシーに従って、フィルタ済みのログと、（認可があれば）非フィルタのログへのアクセスを提供する。

\subsubsection{Logger Root Recorder}
% src: Logger Root Recorder

Logger Recorder と Logger Viewer のイベントをログに記録する、Logger Recorder の特別なインスタンスである。
簡単のため、このログのフィルタ済みビューと非フィルタビューの詳細は省略する。

%----------------------------------------------------------------------
\subsection{セクション5：データ設計}
% src: Section 5 – Data Design

ログデータは、どのイベントをどこまで詳しく記録するかを決めるアプリケーションから、直接集める。
ログはソフトウェアイベントの追記専用の記録であり、期限切れで削除される以外には変更されない。

アプリケーションは、ログイベントの型ごとにスキーマを定義する。
各型は、あらかじめ設定した 0 個以上のデータ項目を持つ。
次の例はその様子を示したものである。
ログイベントはどれも、タイムスタンプと、少なくとも一つのほかの識別用データ項目を持たなければならない。

\begin{verbatim}
{LogTypes: [login, logout, . . .]}
{LogType: login, timestamp: time, IP: IPaddress, verb: string,
 URL: string, user: string, password: string, cookies: string}
{LogType: logout, timestamp: time, IP: IPaddress, verb: string,
 URL: string, user: string, cookies: string}
{Filters: {timestamp: minute, IP: country, verb: 0, URL: 0,
 user: private, password: pw_mask, cookies: private}}
\end{verbatim}

組み込み型や書式などについての細かな点は多くを省略している。
どう定義するかの基本的な考え方は、この部分的な例から読み取れるはずだからである。

リクエストとレスポンスは、UTF-8 でエンコードされた有効な JSON 式で、長さが 100 万文字未満でなければならない。
個々のフィールドの値は、最大 10,000 文字までとする。

最初の行（LogTypes）は、このアプリケーションが生成するログイベントの型を列挙している。
型ごとに、対応する LogType の項目を持つ JSON レコード（2 行目は LogType: login のもの）があり、その型のログに付けてよいデータ項目を列挙している。

4 番目の行（Filters）は、データ項目ごとの扱いを宣言している。
0 は機微でないデータ、private は「ラップ」すべき私的データを表し、そのほかに特別な種類のデータ向けの扱いがあり、次のものが含まれる。
\begin{itemize}
\item minute：時刻の値を最も近い分に丸める（正確な時刻をぼかす）。
\item country：フィルタ済みビューでは、IP アドレスを発信元の国に置き換える。
\end{itemize}

フィルタは差し替え可能なコンポーネントで定義し、アプリケーションごとに必要となる独自のデータ型にも簡単に対応できるよう拡張可能にすべきである。

\begin{rem}[「機微でない」データの意味とスキーマ定義の精査]\label{rem:appa-nonsensitive}
「機微でない」データは社内での限定的な閲覧にだけ使うべきであり、この指定はデータを公開してよいという意味ではないことに注意する。
すべてのデータ項目について扱い（私的かどうか）も含めて宣言するよう求めているのは、アプリケーションの文脈の中で、項目ごとに明示的に判断が下されるようにするためである。
ログ処理の完全性を確保するには、これらの定義とその更新を注意深く精査することが決定的に重要である。
\end{rem}

\begin{exa}[非フィルタビューとフィルタ済みビューのログエントリ]\label{exa:appa-filtered-entry}
このスキーマに基づく非フィルタビューのログエントリの例は次のとおりである。
{\footnotesize
\begin{verbatim}
2018/06/23 08:09:10 66.77.88.99 POST login.html {user: "SAM", password: ">1<}2{]3[\4/"}
\end{verbatim}
}
これに対応するフィルタ済みビューは次のとおりである。
\begin{verbatim}
2018/06/23 08:09 US1(v4) POST login.html {user: USER1(3), password: PW1(12)}
\end{verbatim}
\end{exa}

データは永続的に保存され、ポリシーで設定した有効期限に達するまで利用できる。
有効期限は、イベントログのタイムスタンプからの経過時間で測る。

ログは、監視とデバッグ、あるいはセキュリティ侵害が起きた場合のフォレンジックのためだけの一時的なデータなので、限られた期間だけ保持する。
データ損失のおそれは、専用のマシンにデータを保存し、冗長な永続ストレージとして冗長な RAID（または同様の）ディスクアレイを使うことで緩和する。
ログは、監査と診断のための短期間の保存を想定している。
このデータを長期間保存したい場合は、ログから派生させた別のデータとして作成し、アクセス制御を確実にするために暗号化すべきである。

%----------------------------------------------------------------------
\subsection{セクション6：API}
% src: Section 6 – API

Logger Recorder のネットワークインターフェースは、次のリモートプロシージャコールを受け付ける。
\begin{itemize}
\item Hello：セッションの最初の API 呼び出しでなければならない。アプリケーションとそのバージョンを識別する。
\item Schema：ログデータのスキーマを定義する（セクション5を参照）。
\item Log：指定したログに記録するイベントデータを送る（セクション5を参照）。
\item Goodbye：アプリケーションが終了するときに送り、セッションを終える。
\end{itemize}

各アプリケーションは、専用の通信路を介して自分用のロギングサービスに接続する。
認証されたエンドポイント間の API 呼び出しは HTTPS で保護する。
クライアントは、あらかじめ設定したサーバー名によって、正当な Logger サービスインスタンスに接続していることを（そのデジタル証明書で）確認する。
リクエストの種類は次のとおりである。

\subsubsection{6.1 Hello リクエスト}
% src: 6.1 Hello Request

Logger サービスを使うプロセスは、ロギングを始めるために次のリクエストを送る。
\begin{verbatim}
{"verb": "Hello", "source": "Sample application", "version": "1"}
\end{verbatim}
レスポンスは、次のように OK またはエラーを返し、セッション用の文字列トークンを渡す。
\begin{verbatim}
{"status": "OK", "service": "Logger", "version": "1", "token": "XYZ123"}
\end{verbatim}
トークンは以降のリクエストで、Hello を送った開始元アプリケーションの文脈を識別するのに使う。これにより、複数のアプリケーションが一つの接続でログを記録できる。
トークンは、推測できないだけの複雑さとエントロピーを持つようランダムに生成する。
トークンの大きさは最低でも 120 ビット、つまり base64 エンコーディングで約 20 文字にすることを推奨する。
ここでは簡潔にするために短いトークンを使っている。

\subsubsection{6.2 スキーマ定義リクエスト}
% src: 6.2 Schema Definition Request

このリクエストは、セクション5で述べたように、以降のロギングで使うデータスキーマを定義する。
\begin{verbatim}
{"verb": "Schema", "token": "XYZ123", ...}
\end{verbatim}
このリクエストの詳細は簡潔さのために省略する。

スキーマは、ログの内容に現れるフィールド名、型、その他の属性を定義する。
次の項に示すイベントログリクエストの例（フィールド timestamp、ipaddr、http、url、error を含む）を見れば、その様子がわかる。

\subsubsection{6.3 イベントログリクエスト}
% src: 6.3 Event Log Request

このリクエストは、Logger サービスに実際に 1 件のレコードを記録する。
\begin{verbatim}
{"verb": "Event", "token": "XYZ123", "log":
 {"timestamp": 1234567890, "ipaddr": "12.34.56.78",
  "http": "POST", "url": "example", "error": "404"}}
\end{verbatim}
log の JSON はログに記録する内容を表し、スキーマと一致していなければならない。

レスポンスは、OK またはエラーを返す。
\begin{verbatim}
{"status": "OK"}
\end{verbatim}
エラーの詳細は簡潔さのために省略する。
ロギングのエラー（例えばストレージ容量の不足）は深刻であり、すぐに対処しなければならない。
ロギングができなければ、システムの動作を監査できないからである。

\subsubsection{6.4 Goodbye リクエスト}
% src: 6.4 Goodbye Request

このリクエストで、ロギングのセッションを終える。
\begin{verbatim}
{"verb": "Goodbye", "token": "XYZ123"}
\end{verbatim}
レスポンスは、OK またはエラーを返す。
\begin{verbatim}
{"status": "OK"}
\end{verbatim}
以後そのトークンは無効になる。
ロギングを再開するには、クライアントはまず Hello リクエストを送らなければならない。

%----------------------------------------------------------------------
\subsection{セクション7：ユーザーインターフェース設計}
% src: Section 7 – User Interface Design

Logger のユーザーインターフェースは、Logger Viewer が提供する Web インターフェースであり、ログを調べるのに使う。
この Web アプリには権限のある運用スタッフだけがアクセスでき、企業のシングルサインオンで認証される。
認証されたユーザーには、アクセスを許可されたログの一覧が表示され、最新のフィルタ済みログエントリを閲覧・検索するためのリンクと、許可されていれば、承認を条件に非フィルタのログへのアクセスを要求するためのリンクが示される。

この例では簡潔にするため、Web インターフェースについては大まかな説明にとどめる。

承認の要求は、次の基本情報を入力する Web フォームから送られ、処理待ちのキューに入る。
\begin{itemize}
\item アクセスを要求する理由（顧客の問題チケット番号などの具体的な情報を含む）
\item 要求するアクセスの範囲（典型的には特定のユーザーアカウントや IP アドレス）
\end{itemize}

承認の要求が来ると、要求をレビューするための Web アプリのページへのリンクを付けた自動メールが承認者に送られる。
判断が下されると、要求者には次の内容がメールで通知される。
\begin{itemize}
\item 承認または却下
\item 却下の場合は、その理由
\item 承認されたアクセスの時間枠
\end{itemize}

フィルタ済みのログと非フィルタのログは、ログごとに用意されたページに表示される。
どのログエントリを表示するかを指定する問い合わせを入力できる。
問い合わせが空なら最新のエントリが表示され、結果をページ送りするための「次へ／前へ」のリンクが付く。

問い合わせでは、ログエントリのフィールドと値を指定し、それらをブール演算子で組み合わせて、一致するログエントリを選ぶ。
問い合わせで順序を明示しない限り、新しい順に並ぶ。
簡潔にするため、問い合わせの構文の詳細は省略する。

フィルタ済みのログでは、生のログの内容の代わりに記号的な識別子（例はセクション2.3を参照）が表示される。
問い合わせには、フィルタ済みのログの内容に現れる記号的な識別子を使ってよい。
例えば、フィルタ済みのログエントリに IP アドレス US1 が表示されていれば、\verb|[IP = US1]| という問い合わせで、アドレスそのものを明かさずに、同じ IP アドレスからのほかのログを見つけられる。

\begin{rem}[フィルタ済みフィールドに対する完全一致検索の禁止]\label{rem:appa-no-exact-search}
フィルタ済みのログへの問い合わせでは、フィルタされたフィールドを正確な値で検索できないようにしなければならない。
例えば、IP アドレスが表示されなくても、ユーザーが \verb|[IP = 1.1.1.1]| のように当て推量で検索できると、いずれ USA888 などと表示されるログエントリに行き当たり、実際の値を推測できてしまうかもしれない。
\end{rem}

非フィルタアクセスが承認された場合でも、非フィルタの閲覧と問い合わせを始めるには、ユーザーがそのためのオプションを選ばなければならない。
ベストプラクティスは、フィルタ済みのログをできるだけ活用し、フィルタされた値は必要なときにだけ明かすことであり、ユーザーインターフェースがそうした使い方を促すことが重要である。

作業が終わったら、ユーザーは非フィルタのログへのアクセス権を手放せる。
不要なアクセスのリスクを最小限にするために、ユーザーインターフェースは、一定時間操作がなかったらアクセス権の放棄を促すべきである。

ログの内容を表示する Web ページは、意図しない漏えいを防ぎ、有効期限が切れたログデータを確実に利用できなくするために、ユーザーエージェントにローカルにキャッシュさせるべきではない。

%----------------------------------------------------------------------
\subsection{セクション8：技術設計}
% src: Section 8 – Technical Design

Logger Recorder サービスは、アプリケーションがログイベントのデータをストリーミングして永続ストレージに書き込むための書き込み専用のインターフェースと、それらのログのビューを得るための問い合わせインターフェースからなる。
ストレージは、UTF-8 のテキスト行を追記していくファイルの列であり、ログイベント 1 件が 1 行に対応する。
該当するスキーマ（前述）で記述されたログデータは、テキストによる正準的な表現と相互に変換される。
書式の詳細はこの例では省略する。

フィルタリングの対象となるログデータのフィールドは、フィルタ済みの表現で保存するのに加えて、生のデータを、サービスが生成した AES 鍵で暗号化して保存すべきである。鍵は毎日新しいものを使う。
これらの鍵を安全に保護するために、ハードウェアの鍵ストレージなど適切な手段を使う。

ロギングサービスにとって、使えるストレージを使い果たすことは致命的なエラーになる。そこで、書き込み速度を空き容量と比べて測定し（\texttt{free\_storage\_MB / avg\_logging\_MB\_per\_hour}）、書き込み量が一定だと仮定して残りが 10 時間分のデータを下回ったら、優先度の高い運用アラートを出す（アラートを出す時間数は設定できる）。

性能のために、実際のログファイルを補って効率よくアクセスできるよう、フィルタ済みのログイベント情報（タイムスタンプ、ログの型、ファイル名、オフセット）を記録する SQL データベースを使うことを検討する。

フィルタ済みのログでは、私的データを記号的な識別子（例えば米国内の IP アドレスに対する US1）で隠す。
フィルタされていない私的データを保存しないで済むよう、この対応表は、フィルタされていないデータ値の安全なダイジェストから、フィルタ済みの呼び名への対応として持つ。
この対応は一時的なもので、Logger Viewer がログごと、ユーザーの文脈ごとに別々に保持する。
ユーザーは対応を消去して最初からやり直すことができ、また時間とともに無駄にたまっていくのを防ぐため、24 時間使われなかった対応は自動的に消去される。

Logger Root Recorder は、Logger 自身の動作をログに記録する Logger Recorder の特別なインスタンスである。

%----------------------------------------------------------------------
\subsection{セクション9：設定}
% src: Section 9 – Configuration

ログの保持は次のように設定する。
データは保持期間を過ぎると自動的に削除される。
保持期間を過ぎたデータは、安全かつ完全に削除する（単にゴミ箱に移すのではなく、shred(1) コマンドまたは同様のものを使う）。
\begin{verbatim}
Retention: {
  "Log1": {"days": 10},
  "Log2": {"hours": 24},
}
\end{verbatim}

ログへのアクセス権は、権限のあるユーザーのリストを設定して付与する。
\begin{verbatim}
Access: {
  "Log1": {"filtered": ["u1", "u2", "u3", . . .],
  "unfiltered": ["x1", "x2", "x3", . . .],
  "approval": ["a1", "a2", "a3", . . .]},
}
\end{verbatim}

上の例では、Log1 で表されるログへのフィルタ済みアクセスを許可されたユーザーを、角括弧の中に列挙している（例えば u1、u2、u3）。
次に、非フィルタアクセスを許可されたユーザーを同じように列挙する。
これらのユーザーには、要求が承認された後にだけアクセスが与えられる。
最後に、限定的な非フィルタアクセスを承認する権限を持つユーザーを、同じ方法で列挙する。

%----------------------------------------------------------------------
\subsection{セクション10：参考文献}
% src: Section 10 – References

この設計書を理解するには、次の文書が役に立つ。
これらは架空のものである。
\begin{itemize}
\item 企業の基本設計前提文書（セクション2で参照）
\item 企業の一般データ保護ポリシーとガイドライン
\item パブリッシュ／サブスクライブ・プロトコル設計書（セクション4で参照）
\end{itemize}

% src: END OF DOCUMENT
\begin{center}
（設計書 終わり）
\end{center}

\newpage

%%  appendix B
%% =======================================================
%%  part: appendix B
%% =======================================================
\section{用語集}
% src: Glossary

この付録では、\cite{Kohnfelder2021} の付録Bに沿って、本書に現れるソフトウェアセキュリティの用語をまとめる。
% 注: 配列は原典どおり英語のアルファベット順。「参照」「対比」は原典の (See ...) / (Cf. ...) にあたる。

\begin{description}
\item[影響を受けるユーザー（Affected users）]
特定の脆弱性が悪用されたときに影響を受けうるユーザーが、全体のどれくらいの割合かを見積もったもの。（DREAD の構成要素）

\item[評価報告書（Assessment report）]
セキュリティ設計レビュー（SDR）の結果をまとめた文書。所見と推奨事項を重要度順に並べた要約で構成され、セキュリティを高めるための具体的な設計変更や戦略を含む。（第7章を参照）

\item[資産（Asset）]
守るべき価値のあるデータや資源。とくに攻撃の標的になりやすいもの。

\item[非対称暗号（Asymmetric encryption）]
暗号化と復号に別々の鍵（暗号化には公開鍵、復号には秘密鍵）を使うデータの暗号化。（対比：対称暗号）

\item[攻撃（Attack）]
セキュリティを侵害しようとする行為。

\item[攻撃者（Attacker）]
システムのセキュリティを侵害しようとする悪意ある主体。（脅威アクター（Threat actor）ともいう）

\item[攻撃対象領域（Attack surface）]
システムを攻撃するための侵入口になりうる箇所をすべて合わせたもの。

\item[攻撃ベクトル（Attack vector）]
一つの攻撃を最初から最後まで構成する一連のステップ。攻撃対象領域から始まり、資産へのアクセスで終わる。

\item[監査（Auditing）]
プリンシパルの行為を信頼できる形で記録し、定期的に点検して、不適切な活動の兆候となる疑わしい振る舞いを検出すること。（ゴールドスタンダードの構成要素）

\item[認証（Authentication, authN）]
プリンシパルの身元を高い確度で確定すること。（ゴールドスタンダードの構成要素）

\item[真正性（Authenticity）]
データの値が改ざんされていないことの保証。言い換えれば、システムがデータの無許可の変更を許さないこと。

\item[認可（Authorization, authZ）]
特権的なアクセスを、認証された特定のプリンシパルだけに限定するためのセキュリティポリシーによる制御。（ゴールドスタンダードの構成要素）

\item[可用性（Availability）]
許可されたプリンシパルがいつでもデータにアクセスできることの保証。言い換えれば、正当なアクセスを妨げるような大きな遅延や停止をシステムが起こさないこと。

\item[バックトラッキング（Backtracking）]
正規表現のマッチングなどのアルゴリズムで、処理が前進と後退を繰り返し、指数的に反復しうる振る舞い。バックトラッキングで計算量が膨れ上がって可用性が下がると、セキュリティ上の問題になりうる。（第10章を参照）

\item[ブロック暗号（Block cipher）]
データを 1 ビットずつではなく、固定長のブロック単位で処理する対称暗号のアルゴリズム。

\item[ボトルネック（Bottleneck）]
コードの実行パス上にあって、特定の資産へのアクセスをすべて見張る単一の地点。すべてのアクセスに同じ認可チェックが確実に行われるようになるので、セキュリティにとって重要である。

\item[バッファオーバーフロー（Buffer overflow）]
割り当てられたメモリの境界の外側に不正にアクセスしてしまう脆弱性の総称。

\item[認証局（Certificate authority, CA）]
デジタル証明書の発行者。

\item[チョークポイント（Chokepoint）]
ボトルネックを参照。

\item[選択平文攻撃（Chosen plaintext attack）]
攻撃者が自分で選んだ平文に対応する暗号文を知ることができ、そこから秘密鍵を突き止める暗号の解析。（第5章を参照）

\item[C-I-A]
情報セキュリティの基本的なモデル。（機密性、完全性、可用性を参照）

\item[暗号文（Ciphertext）]
メッセージを暗号化したもの。鍵がなければ意味をなさない。（対比：平文）

\item[衝突（Collision）]
2 つの異なる入力から同じメッセージダイジェストの値が得られること。

\item[衝突攻撃（Collision attack）]
既知の衝突を利用して、暗号学的メッセージダイジェストの値が一意であることを前提にした真正性の保証を覆すこと。

\item[コマンドインジェクション（Command injection）]
悪意ある入力によって、攻撃者の思いどおりの任意のコマンドが実行されてしまう脆弱性。

\item[機密性（Confidentiality）]
データへのアクセスを許可されたものだけに限るという、情報セキュリティの基本的な性質。

\item[混乱した使節（Confused Deputy）]
許可されていない主体が、許可された主体やコードをだまして、自分の代わりに有害な行為を実行させることができてしまう、脆弱なパターン。（第4章を参照）

\item[暗号（Cryptography）]
データを隠すために、元に戻せる形で変換する数学的な技術。

\item[暗号論的に安全な擬似乱数生成器（Cryptographically secure pseudo-random number generator, CSPRNG）]
推測が事実上不可能なほど予測できないとみなされ、暗号に使える乱数の供給源。（対比：擬似乱数生成器（PRNG））

\item[被害の大きさ（Damage potential）]
特定の脆弱性を悪用されたときに、どれだけの害が生じうるかを見積もったもの。（DREAD の構成要素）

\item[復号（Decryption）]
暗号文を元の平文メッセージに戻す処理。

\item[サービス拒否（Denial of service, DoS）]
計算資源を浪費させて可用性を低下させる攻撃。（STRIDE の構成要素でもある）

\item[依存関係（Dependency）]
ソフトウェアの動作に必要な、ソフトウェアライブラリやシステムのその他の構成要素。

\item[ダイアログ疲れ（Dialog fatigue）]
何度も出てくる、あるいは中身の乏しいソフトウェアのダイアログに対する人間の反応。目的を果たすためにダイアログを片付けようとして、反射的に応答してしまうことが多い。ユーザーが自分の行為のセキュリティ上の帰結を理解しない、あるいは考えないときに、セキュリティ上の問題になる。

\item[ダイジェスト（Digest）]
任意の大きさの入力データから計算される、固定サイズの一意な数値。ダイジェストの値が異なれば入力も異なることが保証されるが、衝突は起こりうる。（ハッシュ（Hash）ともいう）

\item[デジタル証明書（Digital certificate）]
署名者が特定の主張を表明し、デジタル署名を付けた文書。よく使われるデジタル証明書の規格には、TLS/SSL の安全な通信用（サーバー側とクライアント側の両方）、コード署名用、メール署名用、認証局用（ルート、中間、リーフ）のものがある。

\item[デジタル署名（Digital signature）]
秘密鍵を知っていることを示すことで、署名者が本物であることを証明する計算。

\item[発見のしやすさ（Discoverability）]
特定の脆弱性があることを、攻撃をもくろむ者がどれだけ簡単に知りうるかを見積もったもの。（DREAD の構成要素）

\item[DREAD]
脆弱性の深刻度を見積もるために、5 つの要素で評価する方式の頭字語。（被害の大きさ、再現性、悪用のしやすさ、影響を受けるユーザー、発見のしやすさを参照）

\item[ECB（electronic code book）モード（ECB mode）]
各ブロックを独立に暗号化するブロック暗号の暗号化モード。同じブロックからは同じ出力が得られるので、多くの用途で ECB は弱く、通常は推奨されない。（第5章を参照）

\item[権限昇格（Elevation of privilege）]
主体がより大きな権限を手に入れるあらゆる手段。とくに、攻撃者が脆弱性を悪用して権限を得る場合をいう。（STRIDE の構成要素）

\item[暗号化（Encryption）]
メッセージを秘密裏に伝えるために、平文を暗号文に変換するアルゴリズム。

\item[エントロピー源（Entropy source）]
乱数生成器へのランダムな入力の供給源。

\item[エクスプロイト（Exploit）]
セキュリティを侵害して害をもたらす、実際に通用する攻撃の手順。

\item[悪用のしやすさ（Exploitability）]
特定の脆弱性をどれだけ簡単に悪用できるかを見積もったもの。未知の要素が多いため、主観的な推測になりがちである。（DREAD の構成要素）

\item[通信の事実（Fact of communication）]
2 者が情報をやり取りしたかどうかという知識。たとえば盗聴者は、解読できない暗号化されたメッセージでも、それを観察するだけでこの知識を得られる。

\item[欠陥（Flaw）]
設計または実装上のバグ。脆弱性である場合も、そうでない場合もある。

\item[フットガン（Footgun）]
バグ、とくに脆弱性を持ち込みやすくしてしまうソフトウェアの機能。

\item[ファズテスト（Fuzz testing）]
ソフトウェアの欠陥を見つけるために、任意の入力を使って自動的に行う総当たりのテスト。

\item[ゴールドスタンダード（Gold Standard）]
セキュリティを強制する 3 つの基本的な仕組みの愛称。（監査、認証（authN）、認可（authZ）を参照）

\item[ガード（Guard）]
資源へのアクセスを制御する、ソフトウェア内で認可を強制する仕組み。

\item[ハードウェア乱数生成器（Hardware random number generator, HRNG）]
ランダム性の高いデータを効率よく生成するように設計されたハードウェア装置。（暗号論的に安全な擬似乱数生成器（CSPRNG）を参照）

\item[ハッシュ（Hash）]
ダイジェストを参照。

\item[HMAC（Hash message authentication code）]
鍵の値ごとに固有のメッセージダイジェスト関数が定まる、メッセージダイジェスト関数の族。

\item[HTML インジェクション（HTML injection）]
攻撃者が悪意ある入力を細工して、HTML ページに任意のマークアップやスクリプトを注入できてしまう脆弱性。

\item[インシデント（Incident）]
セキュリティ攻撃の具体的な一件。

\item[情報漏えい（Information disclosure）]
無許可で情報が漏れること。（STRIDE の構成要素）

\item[インジェクション攻撃（Injection attack）]
入力の一部が予期しない形で解釈される脆弱性を、悪意ある入力で突くセキュリティ攻撃。よくある形態として、SQL インジェクション、HTML インジェクション、コマンドインジェクション、パストラバーサルがある。（第10章を参照）

\item[入力検証（Input validation）]
入力データが下流で正しく処理されるよう、有効な形式であることを確かめる防御的な検査。

\item[結合テスト（Integration testing）]
複数のコンポーネントを組み合わせて動かして行うソフトウェアテスト。（対比：単体テスト）

\item[完全性（Integrity）]
データを正確に保つ、すなわち許可された変更と削除だけを認めるという、情報セキュリティの基本的な性質。（C-I-A を参照）

\item[鍵（Key）]
データをどう変換するかを決める、暗号アルゴリズムのパラメータ。（秘密鍵、公開鍵を参照）

\item[鍵付きハッシュ関数（Keyed hash function）]
HMAC（Hash message authentication code）を参照。

\item[メッセージ認証コード（Message authentication code, MAC）]
メッセージが本物で改ざんされていないことの証拠として、メッセージに添えるデータ。（対比：HMAC）

\item[メッセージダイジェスト（Message digest）]
ダイジェストを参照。

\item[緩和策（Mitigation）]
起こりうる攻撃を防ぐ、あるいはその害を減らすために先回りして講じる対策。たとえば、被害を最小限に抑える、攻撃を受けても回復できるようにする、攻撃を検知しやすくする、といったものである。

\item[ナンス（Nonce）]
リプレイ攻撃を防ぐためなどに通信プロトコルで使われる、一度だけ使う任意の数。

\item[ワンタイムパッド（One-time pad）]
メッセージを暗号化するための共有秘密鍵。再利用するとセキュリティが弱まるので、一度しか使えない。

\item[オーバーフロー（Overflow）]
値が変数に収まる範囲を超えたときに算術命令が返す、誤った結果。オーバーフローが気付かれないまま起きると、予期しない結果が紛れ込んで、脆弱性につながることが多い。

\item[パストラバーサル（Path traversal）]
悪意ある入力によってファイルシステムのパスに予期しない内容が入り込み、本来のアクセス範囲の外にあるファイルを指定できてしまう、よくある脆弱性。

\item[平文（Plaintext）]
暗号化する前の、あるいは本来の受信者が復号した後の、元のメッセージ。

\item[原像攻撃（Preimage attack）]
メッセージダイジェスト関数に対する攻撃で、特定のメッセージダイジェストの値になる入力値を見つけようとするもの。（第5章を参照）

\item[プリンシパル（Principal）]
認証された主体。人、企業、組織、アプリケーション、サービス、またはデバイス。

\item[秘密鍵（Private key）]
復号に必要なパラメータで、許可された受信者が秘密にしておくもの。

\item[来歴（Provenance）]
データの出所と、誰の手を経てきたかについての信頼できる履歴。データが正当であることの裏付けになる。

\item[擬似乱数生成器（Pseudo-random number generator, PRNG）]
「そこそこ良い」乱数生成器だが、高度な解析によって予測されうるもの。この乱数はシミュレーションなど多くの目的に役立つが、ランダム性が足りないので暗号には向かない。（対比：暗号論的に安全な擬似乱数生成器（CSPRNG））

\item[公開鍵（Public key）]
特定の受信者宛てのメッセージを暗号化するのに必要な、広く公開されたパラメータ。

\item[乱数（Random number）]
でたらめに選ばれ、確実には予測できない数。

\item[レート制限（Rate limiting）]
処理の速度を落とす方法。総当たりで繰り返すことで成功を狙う攻撃を緩和するのによく使われる。

\item[リプレイ攻撃（Replay attack）]
以前の本物のメッセージを再送して、安全な通信プロトコルを攻撃すること。攻撃者が以前の本物の通信のコピーを再送し、それが元の送信者から改めて送られた同じメッセージだと誤認されれば、リプレイ攻撃は成功する。（第5章を参照）

\item[再現性（Reproducibility）]
特定の脆弱性の悪用がどれだけ確実に成功するかを見積もったもの。（DREAD の構成要素）

\item[否認（Repudiation）]
自分の行為をもっともらしく否定できること。とくに、それによって攻撃者が責任を逃れられることをいう。（STRIDE の構成要素）

\item[ルート証明書（Root certificate）]
認証局を信頼することを認める、自己署名されたデジタル証明書。

\item[サンドボックス（Sandbox）]
内部で実行されるコードが使える権限に上限を設けるように設計された、制限付きの実行環境。

\item[セキュリティ設計レビュー（Security design review, SDR）]
ソフトウェア設計のセキュリティについて、体系立てて行うレビュー。

\item[セキュリティの帽子（Security hat）]
物事がどううまくいかなくなりうるかを考えるために、セキュリティの考え方を「かぶる」ことを表す言い回し。

\item[セキュリティリグレッション（Security regression）]
以前に修正した既知のセキュリティバグが再び現れること。

\item[セキュリティテストケース（Security test case）]
セキュリティ制御が常に効いていることを確かめるソフトウェアのテストケース。

\item[セキュリティテスト（Security testing）]
セキュリティ制御が正しく機能することを確かめるためのソフトウェアテスト。

\item[サイドチャネル攻撃（Side channel attack）]
保護の仕組みを直接破るのではなく、間接的に機密情報を推測する攻撃。たとえば、計算結果が出るまでにかかる時間から、その結果について確実に推測するもの。（第8章を参照）

\item[投機的実行（Speculative execution）]
現代のプロセッサーが使う最適化手法。時間を節約できる見込みに賭けて先の命令を前倒しで実行し、不要だった場合は後でその結果を捨てるバックトラッキングの仕組みを持つ。投機的実行がキャッシュの状態に残す影響から、本来アクセスできない情報が漏れるおそれがあり、セキュリティ上の脅威となる。（第8章を参照）

\item[なりすまし（Spoofing）]
攻撃者が許可されたプリンシパルのふりをして、認証をすり抜けること。（STRIDE の構成要素）

\item[SQL インジェクション（SQL injection）]
攻撃者が悪意ある入力を細工して、任意の SQL コマンドを実行できてしまう脆弱性。

\item[STRIDE]
ソフトウェアセキュリティの脅威の基本的な 6 つの種類を表す頭字語。脅威モデリングの手引きとして役立つ。（なりすまし、改ざん、否認、情報漏えい、サービス拒否、権限昇格を参照）

\item[対称暗号（Symmetric encryption）]
暗号化と復号に同じ鍵を使う暗号化の方法。「対称」とは、暗号化できる者は誰でも復号もできるという意味である。（対比：非対称暗号）

\item[テインティング（Tainting）]
ソフトウェアの中でデータの出所を追跡する処理。信頼できない入力や、その入力の影響を受けたデータが、インジェクション攻撃などの特権的な操作に使われるのを防ぐために用いる。（第8章を参照）

\item[改ざん（Tampering）]
データを無許可で変更すること。（STRIDE の構成要素）

\item[脅威（Threat）]
起こりうる、あるいは仮定上のセキュリティ問題。

\item[脅威アクター（Threat actor）]
攻撃者を参照。

\item[脅威モデリング（Threat modeling）]
緩和が必要な脅威を洗い出すために、システムのモデルを分析すること。（第2章を参照）

\item[タイミング攻撃（Timing attack）]
操作にかかる時間を測ることで情報を推測する、サイドチャネル攻撃。

\item[アンダーフロー（Underflow）]
浮動小数点の計算結果の精度が失われること。

\item[単体テスト（Unit testing）]
個々のモジュールを他のコンポーネントから切り離して行うソフトウェアテスト。

\item[信頼できない入力（Untrusted input）]
信頼できない入力元から来る入力データ。とくに、攻撃対象領域になりうるものとして見た場合をいう。

\item[脆弱性（Vulnerability）]
セキュリティ攻撃を可能にするソフトウェアの欠陥。

\item[脆弱性チェーン（Vulnerability chain）]
組み合わせると一つのセキュリティ攻撃が成り立つ、脆弱性の集まり。

\item[弱点（Weakness）]
脆さの原因となり、そのため脆弱性であるかもしれないバグ。
\end{description}

\newpage

%%  appendix C
%% =======================================================
%%  part: appC
%% =======================================================
\section{演習}\label{sec:appc-exercises}
% src: Exercises

この付録では、\cite{Kohnfelder2021} の付録Cに沿って、さらに探究するための演習を扱う。

\begin{quote}
「探検はイノベーションを駆動するエンジンである。」\\
\hfill --- Edith Widder
\end{quote}

この付録には、本書の内容からさらに先へ進みたい読者に向けて、探究のアイデア、未解決の問い、課題を収めている。

\subsection{第1章 基礎}
% src: Chapter 1, Foundations

\noindent\textbf{[1]（物理世界と関わるソフトウェアの C-I-A）}\quad
本書は従来のコンピューターシステムの情報セキュリティを中心に扱っているが、家電や各種の機器もソフトウェアで動いており、インターネットにつながるものがますます増えている。
物理世界とやり取りするソフトウェアを守るには、C-I-A のような原則をどう拡張すればよいか。

\subsection{第2章 脅威}
% src: Chapter 2, Threats

\noindent\textbf{[1]（既存の設計の脅威モデリング）}\quad
既存のソフトウェア設計を脅威モデリングせよ。大きなシステムの 1 つのコンポーネントだけでもよい。

\noindent\textbf{[2]（映画や本の場面の脅威モデリング）}\quad
お遊びとして、敵対する者どうしが貴重な資産をめぐって争う、お気に入りの映画や本の場面を脅威モデリングしてみよ。

\subsection{第3章 緩和策}
% src: Chapter 3, Mitigations

\noindent\textbf{[1]（メモリ上の機密データの露出の制限）}\quad
第3章の「データの露出の最小化」（Minimize Data Exposure）で述べたように、メモリ上の機密データの露出を制限するヘルパー関数を書け。
% 注: 原文の参照先は "on page XX"（ページ番号未確定のまま）。

\noindent\textbf{[2]（混乱した使節の作成と修正）}\quad
混乱した使節（Confused Deputy）をわざと作り込んだコードを書き、悪用してみよ。同僚に悪用を挑んでもらってもよい。
その後、脆弱性を修正し、コードが安全になったことを確認せよ。

\noindent\textbf{[3]（拡張可能なアクセスポリシーのライブラリ）}\quad
既存のデータアクセス API に対して、拡張可能なアクセスポリシーを強制するライブラリを設計せよ。

\subsection{第4章 パターン}
% src: Chapter 4, Patterns

\noindent\textbf{[1]（パターンの適用）}\quad
既存の設計を取り上げるか新しい設計に取り組み、できるだけ安全にするためにこの章のパターンをいくつ使えるか試してみよ。

\noindent\textbf{[2]（パターンとアンチパターンの追加）}\quad
ほかにどのようなセキュリティのパターンとアンチパターンを思いつくか。
この章で示したものに書き足しながらリストを育て、同僚と共有せよ。

\noindent\textbf{[3]（許可リストと拒否リスト）}\quad
許可リストは常に拒否リストより良いのだろうか。
例外を一つ考えよ。あるいは、例外がない理由を説明せよ。

\subsection{第5章 暗号}
% src: Chapter 5, Cryptography

\noindent\textbf{[1]（OpenSSL による実験）}\quad
本物の暗号の道具を手軽に試すには、OpenSSL のコマンドライン（\url{https://wiki.openssl.org/index.php/Main_Page}）を使うとよい。
対称暗号と非対称暗号、MAC（openssl(1) ではダイジェストと呼ばれる）を試せるほか、自分で証明書を作成・検査することもできる。

\noindent\textbf{[2]（暗号ライブラリによる基本操作の実装）}\quad
高品質な暗号ライブラリを見つけ、この章で述べた基本的な操作を実装してみよ。
API は使いやすかったか。また、自分の実装が安全だとどれくらい自信を持てるか。

\noindent\textbf{[3]（API の再設計）}\quad
前の演習が難しかったなら、もっと使いやすく、誤用もしにくくするには API をどう再設計すればよいか。

\noindent\textbf{[4]（API の改善の実装）}\quad
思いついた暗号 API の改善を実際にコードにせよ。あるいは、元のライブラリをラップして、より良い API を提供せよ。

\subsection{第6章 セキュアな設計}
% src: Chapter 6, Secure Design

\noindent\textbf{[1]（Google の設計文書のガイダンス）}\quad
Google の設計文書の書き方のガイダンス（\url{https://www.industrialempathy.com/posts/design-docs-at-google/}）を調べよ。

\noindent\textbf{[2]（設計文書を書いてみる）}\quad
ソフトウェアの設計文書を書いたことがなければ、次の機会に書いてみよ（くだけた書き方でも、大まかな内容でもかまわない）。

\noindent\textbf{[3]（設計文書の事後作成）}\quad
設計文書のないコードベースで作業しているなら、後追いで設計文書を作成せよ。
大きなシステムでは、セキュリティにとって最も重要なコンポーネントや、ほかに関心のあるコンポーネントに絞って、1 つずつ設計を書き起こせばよい。

\subsection{第7章 設計レビュー}
% src: Chapter 7, Design Reviews

\noindent\textbf{[1]（既存の設計のレビュー）}\quad
学習のための演習として、既存の設計を見つけてレビューせよ。
脆弱性を探すだけで終わらせず、長所と短所の両方を幅広く評価すること。セキュリティが最も重要な箇所、設計がセキュリティを高めている点、緩和策の代替案、セキュリティを改善したり使いやすくしたりできる点なども評価に含める。

\noindent\textbf{[2]（評価の共有）}\quad
前の演習で得た知見を同僚と共有し、議論せよ。

\subsection{第8章 セキュアプログラミング}
% src: Chapter 8, Secure Programming

\noindent\textbf{[1]（既知のセキュリティバグの調査）}\quad
セキュリティ脆弱性の実例の感触をつかむために、自分のコードベースやオープンソースのソフトウェアプロジェクトで、すでに発見・修正されたセキュリティバグを探せ。
原著者のおすすめはオープンソースのプロジェクトである。脆弱性がたいてい詳しく説明されていて、コードも見られるからである。
米国国土安全保障省は、公に知られた脆弱性の大規模なデータベース（\url{https://cve.mitre.org/}）を後援している。
公開された脆弱性の情報源としては、Chromium のバグデータベース（\url{https://bugs.chromium.org/p/chromium/issues/list/}）もよい。
まずはこれらのデータベースを修正済みのセキュリティバグで絞り込み、実際のコードの変更を見てみるとよい。

\noindent\textbf{[2]（アンダーハンデッドコーディング）}\quad
\underline{アンダーハンデッドコーディング}（underhanded coding）は\underline{難読化コーディング}（obfuscated coding）とも呼ばれ、フットガンなどの策略を駆使して、コードを読んだ誰もが予想しない動き方をするコードを書く技芸である。
アンダーハンデッドコーディングのコンテストでは、プログラミング言語を限界まで使い倒す創造性をプログラマーが競い合う。
しかし、悪意のあるコードを無害に見せかける技法は、知らずに出くわせばフットガンにもなる。
手始めに次のサイトを見るか、自分で作ってみよ：\url{http://www.underhanded-c.org/}、\url{https://underhandedcrypto.com/}。

\subsection{第9章 低レベルのコーディング上の欠陥}
% src: Chapter 9, Low-Level Coding Flaws

\noindent\textbf{[1]（固定幅整数のオーバーフロー検出）}\quad
固定幅の整数型を備えた言語は、なぜオーバーフローを検出する仕組みを何も用意していないのか。
そうした仕組みがあれば役に立つだろうか。役に立つなら、その仕組みを生かすために C 言語をどう拡張するか。

\noindent\textbf{[2]（Valgrind の仕組み）}\quad
Valgrind のような解析ツールがメモリ管理の問題をどのように検出するかを調べよ（\url{https://valgrind.org/docs/manual/mc-manual.html} を参照）。

\noindent\textbf{[3]（メモリ管理の脆弱性と解析ツール）}\quad
読み取りと書き込みの両方のバッファオーバーフローなど、何種類かのメモリ管理の脆弱性を含む小さなプログラムを書け。
Valgrind のようなツールを使って、そのバグを検出するかどうか確かめよ。
さらに、ツールが解析しにくくなるようにコードを変えてみて、バグをツールの目から逃れさせられるか試してみよ。

\subsection{第10章 入力検証}
% src: Chapter 10, Input Validation

\noindent\textbf{[1]（信頼できない入力の特定）}\quad
自分が取り組んでいるシステムの主な攻撃対象領域にある信頼できない入力を特定し、入力検証がどこまで徹底して実装・テストされているかを確かめよ。

\noindent\textbf{[2]（入力検証の実装）}\quad
信頼できない入力が脆弱性となりうることがわかったら、入力検証を実装せよ。

\noindent\textbf{[3]（入力検証の共通化）}\quad
システムの入力検証は、同じような処理の繰り返しになりがちである。
共通のコードやヘルパー関数で確実に処理できる箇所を探せ。
うっかり書き忘れることがないよう、入力検証をフレームワークに組み込む方法を検討せよ。

\subsection{第11章 Web セキュリティ}
% src: Chapter 11, Web Security

\noindent\textbf{[1]（Web セッションのコンポーネントの設計）}\quad
Web セッションを作成・認証するコンポーネントのセキュリティ要件を書け。
そのコンポーネントを設計して脅威モデリングし、友人にセキュリティレビューをしてもらえ。

\noindent\textbf{[2]（Web セッションの実装と攻撃）}\quad
自分の Web セッションの実装を簡単な Web アプリに組み込め。
別のセッションになりすましたり、必要なセッション状態を盗んだりしてみよ。
友人に自分の実装を「攻撃」してもらえば、なおよい。

\noindent\textbf{[3]（CSRF 対策の追加）}\quad
コンポーネントに CSRF 対策の仕組みを追加し、自分の Web アプリでテストせよ。

\noindent\textbf{[4]（cookie を使わない Web セッション）}\quad
セキュリティ上の課題の本質を理解する実験として、cookie を使わずに Web セッションを保護する方法を探れ。

\noindent\textbf{[5]（Web フレームワークの調査）}\quad
Web フレームワークのソースコード（できれば設計文書も）を手に入れ、そのフレームワークがセッションをどう実装し、XSS と CSRF の脆弱性をどう防ぎ、すべての Web のやり取りが HTTPS で保護されることをどう保証しているかを学べ。
脅威モデリングやその他の手段によって、何らかの脆弱性を見つけられるか。
攻撃を試したい場合は、自分専用のテストサーバーを立てて行うこと。

\subsection{第12章 セキュリティテスト}
% src: Chapter 12, Security Testing

\noindent\textbf{[1]（セキュリティテストケースの追加）}\quad
好きなコードベースでセキュリティが重要な領域を見つけ、追加すべきセキュリティテストケースを探せ。
新しいセキュリティテストケースを書いて貢献せよ。

\noindent\textbf{[2]（トリガー付きの GotoFail）}\quad
本書で書いたセキュリティテストでは検出できない、GotoFail の脆弱性の別の例を考えてみよう。
余分な \texttt{goto fail;} の代わりに、次の行を挿入する。
\begin{verbatim}
if (expected_hash[0] == 0x23) goto fail;
\end{verbatim}
この種の手法は、特定のトリガーで発動する脆弱性を、一種のバックドアとしてひそかに仕込むのに使われるかもしれない。
これを検出するには、期待されるハッシュの最初のバイトが \texttt{0x23} であるテストケースが必要になる。
具体的な中身を知らないまま、この種の脆弱性を検出するテストをどう書けばよいか。

\noindent\textbf{[3]（セキュリティリグレッションテストの作成）}\quad
既知の脆弱性を持つオープンソースのソフトウェアプロジェクトの古いバージョンをチェックアウトせよ。
テストスイートを実行し、すべてのテストが通ることを確認せよ。
その脆弱性を確認するセキュリティリグレッションテストを書け。
脆弱性を修正した次のバージョンに同期し、自分のリグレッションテストをマージせよ。
これでセキュリティリグレッションテストは通るはずである。通らなければ修正せよ。
そのうえで、最新バージョンに関連する脆弱性がほかにないかを確認せよ。

\subsection{第13章 セキュアな開発のベストプラクティス}
% src: Chapter 13, Secure Development Best Practices

\noindent\textbf{[1]（コード品質の漸進的な改善）}\quad
lint やコードスキャンツールを使う、エラー処理と例外処理のテストカバレッジを確認するなど、コード品質を少しずつ高める手軽な方法を探れ。

\noindent\textbf{[2]（セキュリティの文書化）}\quad
自分のコードベースで、セキュリティにかかわる部分がどの程度きちんと文書化されているかを確認し、必要な改善を行え。

\noindent\textbf{[3]（コードレビューでのセキュリティの帽子）}\quad
コードレビューのときは毎回、必要に応じてセキュリティの帽子をかぶり、もう一度見直せ。

\noindent\textbf{[4]（バグのトリアージとセキュリティ）}\quad
バグのトリアージではセキュリティも考慮せよ。あるいは、セキュリティの観点からバグデータベースを見渡し、セキュリティに影響するバグが先送りにされていないか確かめよ。

\subsection{第14章 今後の展望}
% src: Chapter 14, Looking Ahead

\noindent\textbf{[1]（小さな改善の機会）}\quad
小さな一歩でもよいので、この章で述べた方向に沿って改善する機会を探せ。
たとえば、より多くの人をセキュリティに参加させる、セキュリティの視点と戦略をより早い段階で取り入れる、複雑さを減らすか管理する、セキュリティの実践についての透明性を高める、といったことである。

\noindent\textbf{[2]（再利用可能なコンポーネント）}\quad
独自のセキュリティ上の課題を見つけ、その課題に対処する再利用可能なコンポーネントを設計・開発せよ。

\noindent\textbf{[3]（自分のアイデアの追求）}\quad
セキュリティの水準を引き上げ、その輪を広げるための自分なりのアイデアを、ほかにも追求せよ。

\newpage

\begin{thebibliography}{9}
\bibitem{Kohnfelder2021}
Loren Kohnfelder, \textit{Designing Secure Software: A Guide for Developers}, No Starch Press, 2021.
\end{thebibliography}

\end{document}
